% =====================================================================
%  Fondamenti di Unreal Engine, lezione 12 — Il menu principale e il
%  sistema di interazione.
%  Dodicesima e ultima lezione del corso breve di introduzione a Unreal
%  Engine (focus sui Blueprint) con questo docente, online, nel
%  pomeriggio (pausa alle 17:26, ripresa "alle diciassette cinquantacinque,
%  diciotto"). Docente titolare (come nelle lezioni precedenti, nella
%  trascrizione non viene mai chiamato per nome). Da qui in poi il corso
%  prosegue con Davide e il level design.
%  Contenuti, nell'ordine della lezione: le domande d'apertura (recuperare
%  le lezioni o seguire in diretta; l'esame, i Big One dei tre anni e il
%  pezzo da portfolio; il designer che prototipa con i Blueprint; "less is
%  more" e il testo iniziale di Star Wars); gli annunci (ultima lezione,
%  le lezioni con Davide, il server Discord e i suoi canali, le
%  registrazioni fino a settembre); la scelta dell'ordine (prima il menu,
%  poi le interazioni); la dimostrazione del sistema di interazione; il
%  menu: il Widget Blueprint WB_Menu, la mappa vuota M_Menu (Empty Level,
%  il piano con il materiale a griglia), il Canvas Panel, i due Button
%  (Size 300 x 80, ancora al centro, Alignment), i nomi Button_StartGame e
%  Button_QuitGame, i Text dentro i bottoni, lo Style (Normal, Hovered,
%  Pressed), il codice esadecimale sRGB del colore, l'Image dietro ai
%  bottoni e l'ordine della Hierarchy, l'alfa a 0,2; il Level Blueprint
%  (Begin Play, Create Widget, Add to Viewport), il GameMode Override con
%  la Game Mode base e la Default Pawn, il Camera Actor, la Directional
%  Light, Set View Target with Blend (prima cercato sul Camera Manager,
%  poi trovato sul Player Controller); il Graph del menu (Event Construct,
%  Set Show Mouse Cursor, Set Input Mode UI Only, On Clicked, Quit Game,
%  Set Input Mode Game Only, Open Level by Name con ThirdPersonMap); i
%  colori di hover e pressione; il diorama del menu (il bot in Idle, la
%  porta modulare, la luce); i problemi degli studenti (il personaggio che
%  compare ancora, World Settings, il Player Start fuori campo, il target
%  dell'Add to Viewport, il cursore, la spunta del Set Show Mouse Cursor);
%  le interfacce (l'intermediario, il volo di linea) e i Blueprint
%  Component (Actor Component e Scene Component); il componente
%  BPC_RayInteract: la funzione LookForInteractions, Sphere Trace by
%  Channel con Radius 5, Get World Location, il Select con il trucco della
%  variabile Scene Component, Get Owner e Get Component by Class (Camera
%  Component), Is Valid, Get Forward Vector, DistanceTrace a 300, il
%  componente aggiunto al bot, Draw Debug Type, l'Event Tick; il Trace
%  Channel InteractionTrace (Ignore); Break Hit Result, Hit Actor, Does
%  Object Implement Interface; la Blueprint Interface BPI_InteractionRay
%  con le tre funzioni; la pausa (la stamina di Tommaso, il Call che
%  mancava; il rumore della trascrizione); dopo la pausa (la BPI aperta
%  "male"); la variabile Interactable, il diverso (!=), il Validated Get,
%  Interaction End e Interaction Start come messaggi, lo svuotamento
%  della variabile; il problema di tipo di uno studente; l'Event Tick del
%  componente; l'actor BP_InteractableTest (Cube e Widget), Class Settings
%  e Implemented Interfaces, gli eventi dell'interfaccia con tre Print
%  String colorati, la collisione Block su InteractionTrace; IA_Interact
%  (copia di IA_Run, Pressed e Released) e il tasto E nell'IMC; nel bot:
%  Do Once, Interaction (Message), Set Actor Rotation con il solo Yaw da
%  Find Look at Rotation, Play Montage, la variabile Interacting, Disable
%  Input, Stop Movement Immediately, Enable Input, il Reset del Do Once; il
%  Montage dall'animazione "Interact press button"; il widget
%  WB_InteractUI (la cartella UI degli asset, i due Border, Panel e Panel
%  Border, la E centrata); il Widget Component in Screen Space, Draw Size
%  100 x 100, Set Hidden in Game; il problema dei segni di spunta; la luce
%  accesa e spenta con Flip Flop e BP_ActivableLight; la chiusura del
%  corso.
%  Fonte: trascrizione automatica della dodicesima lezione (UE12).
%  Integrazioni dal libro "Blueprints Visual Scripting for Unreal Engine 5"
%  (Romero, Sewell; capitoli 4, 8, 11, 13, 15, 16 e 18) e nozioni tecniche
%  generali su Unreal Engine 5, sempre segnalate come approfondimenti. Il
%  libro non tratta la Default Pawn, il Widget Component in Screen Space,
%  il trucco del Select, i Trace Channel per le interazioni né Stop
%  Movement Immediately: su questi argomenti gli approfondimenti si basano
%  sul funzionamento generale dell'engine.
%
%  Pacchetti richiesti nel preambolo del libro (già presenti nel
%  main.tex): \usepackage{booktabs}, \usepackage{amsmath},
%  \usepackage{amssymb}, \usepackage{tikz} con
%    \usetikzlibrary{shapes.geometric,shapes.symbols,shapes.misc,
%      arrows.meta,positioning,calc,fit}
%  (le librerie vanno caricate nel preambolo, non nel capitolo).
%  Consigliati, come per gli altri capitoli: \usepackage[italian]{babel}
%  e \usepackage[T1]{fontenc} (il capitolo compila anche senza).
%
%  Collocazione nel libro: parte "Fondamenti di Unreal Engine",
%  dodicesimo capitolo (numerazione propria 1.12). Nel file principale
%  del libro, subito dopo la riga dell'undicesima lezione:
%
%    % =====================================================================
%  Fondamenti di Unreal Engine, lezione 12 — Il menu principale e il
%  sistema di interazione.
%  Dodicesima e ultima lezione del corso breve di introduzione a Unreal
%  Engine (focus sui Blueprint) con questo docente, online, nel
%  pomeriggio (pausa alle 17:26, ripresa "alle diciassette cinquantacinque,
%  diciotto"). Docente titolare (come nelle lezioni precedenti, nella
%  trascrizione non viene mai chiamato per nome). Da qui in poi il corso
%  prosegue con Davide e il level design.
%  Contenuti, nell'ordine della lezione: le domande d'apertura (recuperare
%  le lezioni o seguire in diretta; l'esame, i Big One dei tre anni e il
%  pezzo da portfolio; il designer che prototipa con i Blueprint; "less is
%  more" e il testo iniziale di Star Wars); gli annunci (ultima lezione,
%  le lezioni con Davide, il server Discord e i suoi canali, le
%  registrazioni fino a settembre); la scelta dell'ordine (prima il menu,
%  poi le interazioni); la dimostrazione del sistema di interazione; il
%  menu: il Widget Blueprint WB_Menu, la mappa vuota M_Menu (Empty Level,
%  il piano con il materiale a griglia), il Canvas Panel, i due Button
%  (Size 300 x 80, ancora al centro, Alignment), i nomi Button_StartGame e
%  Button_QuitGame, i Text dentro i bottoni, lo Style (Normal, Hovered,
%  Pressed), il codice esadecimale sRGB del colore, l'Image dietro ai
%  bottoni e l'ordine della Hierarchy, l'alfa a 0,2; il Level Blueprint
%  (Begin Play, Create Widget, Add to Viewport), il GameMode Override con
%  la Game Mode base e la Default Pawn, il Camera Actor, la Directional
%  Light, Set View Target with Blend (prima cercato sul Camera Manager,
%  poi trovato sul Player Controller); il Graph del menu (Event Construct,
%  Set Show Mouse Cursor, Set Input Mode UI Only, On Clicked, Quit Game,
%  Set Input Mode Game Only, Open Level by Name con ThirdPersonMap); i
%  colori di hover e pressione; il diorama del menu (il bot in Idle, la
%  porta modulare, la luce); i problemi degli studenti (il personaggio che
%  compare ancora, World Settings, il Player Start fuori campo, il target
%  dell'Add to Viewport, il cursore, la spunta del Set Show Mouse Cursor);
%  le interfacce (l'intermediario, il volo di linea) e i Blueprint
%  Component (Actor Component e Scene Component); il componente
%  BPC_RayInteract: la funzione LookForInteractions, Sphere Trace by
%  Channel con Radius 5, Get World Location, il Select con il trucco della
%  variabile Scene Component, Get Owner e Get Component by Class (Camera
%  Component), Is Valid, Get Forward Vector, DistanceTrace a 300, il
%  componente aggiunto al bot, Draw Debug Type, l'Event Tick; il Trace
%  Channel InteractionTrace (Ignore); Break Hit Result, Hit Actor, Does
%  Object Implement Interface; la Blueprint Interface BPI_InteractionRay
%  con le tre funzioni; la pausa (la stamina di Tommaso, il Call che
%  mancava; il rumore della trascrizione); dopo la pausa (la BPI aperta
%  "male"); la variabile Interactable, il diverso (!=), il Validated Get,
%  Interaction End e Interaction Start come messaggi, lo svuotamento
%  della variabile; il problema di tipo di uno studente; l'Event Tick del
%  componente; l'actor BP_InteractableTest (Cube e Widget), Class Settings
%  e Implemented Interfaces, gli eventi dell'interfaccia con tre Print
%  String colorati, la collisione Block su InteractionTrace; IA_Interact
%  (copia di IA_Run, Pressed e Released) e il tasto E nell'IMC; nel bot:
%  Do Once, Interaction (Message), Set Actor Rotation con il solo Yaw da
%  Find Look at Rotation, Play Montage, la variabile Interacting, Disable
%  Input, Stop Movement Immediately, Enable Input, il Reset del Do Once; il
%  Montage dall'animazione "Interact press button"; il widget
%  WB_InteractUI (la cartella UI degli asset, i due Border, Panel e Panel
%  Border, la E centrata); il Widget Component in Screen Space, Draw Size
%  100 x 100, Set Hidden in Game; il problema dei segni di spunta; la luce
%  accesa e spenta con Flip Flop e BP_ActivableLight; la chiusura del
%  corso.
%  Fonte: trascrizione automatica della dodicesima lezione (UE12).
%  Integrazioni dal libro "Blueprints Visual Scripting for Unreal Engine 5"
%  (Romero, Sewell; capitoli 4, 8, 11, 13, 15, 16 e 18) e nozioni tecniche
%  generali su Unreal Engine 5, sempre segnalate come approfondimenti. Il
%  libro non tratta la Default Pawn, il Widget Component in Screen Space,
%  il trucco del Select, i Trace Channel per le interazioni né Stop
%  Movement Immediately: su questi argomenti gli approfondimenti si basano
%  sul funzionamento generale dell'engine.
%
%  Pacchetti richiesti nel preambolo del libro (già presenti nel
%  main.tex): \usepackage{booktabs}, \usepackage{amsmath},
%  \usepackage{amssymb}, \usepackage{tikz} con
%    \usetikzlibrary{shapes.geometric,shapes.symbols,shapes.misc,
%      arrows.meta,positioning,calc,fit}
%  (le librerie vanno caricate nel preambolo, non nel capitolo).
%  Consigliati, come per gli altri capitoli: \usepackage[italian]{babel}
%  e \usepackage[T1]{fontenc} (il capitolo compila anche senza).
%
%  Collocazione nel libro: parte "Fondamenti di Unreal Engine",
%  dodicesimo capitolo (numerazione propria 1.12). Nel file principale
%  del libro, subito dopo la riga dell'undicesima lezione:
%
%    % =====================================================================
%  Fondamenti di Unreal Engine, lezione 12 — Il menu principale e il
%  sistema di interazione.
%  Dodicesima e ultima lezione del corso breve di introduzione a Unreal
%  Engine (focus sui Blueprint) con questo docente, online, nel
%  pomeriggio (pausa alle 17:26, ripresa "alle diciassette cinquantacinque,
%  diciotto"). Docente titolare (come nelle lezioni precedenti, nella
%  trascrizione non viene mai chiamato per nome). Da qui in poi il corso
%  prosegue con Davide e il level design.
%  Contenuti, nell'ordine della lezione: le domande d'apertura (recuperare
%  le lezioni o seguire in diretta; l'esame, i Big One dei tre anni e il
%  pezzo da portfolio; il designer che prototipa con i Blueprint; "less is
%  more" e il testo iniziale di Star Wars); gli annunci (ultima lezione,
%  le lezioni con Davide, il server Discord e i suoi canali, le
%  registrazioni fino a settembre); la scelta dell'ordine (prima il menu,
%  poi le interazioni); la dimostrazione del sistema di interazione; il
%  menu: il Widget Blueprint WB_Menu, la mappa vuota M_Menu (Empty Level,
%  il piano con il materiale a griglia), il Canvas Panel, i due Button
%  (Size 300 x 80, ancora al centro, Alignment), i nomi Button_StartGame e
%  Button_QuitGame, i Text dentro i bottoni, lo Style (Normal, Hovered,
%  Pressed), il codice esadecimale sRGB del colore, l'Image dietro ai
%  bottoni e l'ordine della Hierarchy, l'alfa a 0,2; il Level Blueprint
%  (Begin Play, Create Widget, Add to Viewport), il GameMode Override con
%  la Game Mode base e la Default Pawn, il Camera Actor, la Directional
%  Light, Set View Target with Blend (prima cercato sul Camera Manager,
%  poi trovato sul Player Controller); il Graph del menu (Event Construct,
%  Set Show Mouse Cursor, Set Input Mode UI Only, On Clicked, Quit Game,
%  Set Input Mode Game Only, Open Level by Name con ThirdPersonMap); i
%  colori di hover e pressione; il diorama del menu (il bot in Idle, la
%  porta modulare, la luce); i problemi degli studenti (il personaggio che
%  compare ancora, World Settings, il Player Start fuori campo, il target
%  dell'Add to Viewport, il cursore, la spunta del Set Show Mouse Cursor);
%  le interfacce (l'intermediario, il volo di linea) e i Blueprint
%  Component (Actor Component e Scene Component); il componente
%  BPC_RayInteract: la funzione LookForInteractions, Sphere Trace by
%  Channel con Radius 5, Get World Location, il Select con il trucco della
%  variabile Scene Component, Get Owner e Get Component by Class (Camera
%  Component), Is Valid, Get Forward Vector, DistanceTrace a 300, il
%  componente aggiunto al bot, Draw Debug Type, l'Event Tick; il Trace
%  Channel InteractionTrace (Ignore); Break Hit Result, Hit Actor, Does
%  Object Implement Interface; la Blueprint Interface BPI_InteractionRay
%  con le tre funzioni; la pausa (la stamina di Tommaso, il Call che
%  mancava; il rumore della trascrizione); dopo la pausa (la BPI aperta
%  "male"); la variabile Interactable, il diverso (!=), il Validated Get,
%  Interaction End e Interaction Start come messaggi, lo svuotamento
%  della variabile; il problema di tipo di uno studente; l'Event Tick del
%  componente; l'actor BP_InteractableTest (Cube e Widget), Class Settings
%  e Implemented Interfaces, gli eventi dell'interfaccia con tre Print
%  String colorati, la collisione Block su InteractionTrace; IA_Interact
%  (copia di IA_Run, Pressed e Released) e il tasto E nell'IMC; nel bot:
%  Do Once, Interaction (Message), Set Actor Rotation con il solo Yaw da
%  Find Look at Rotation, Play Montage, la variabile Interacting, Disable
%  Input, Stop Movement Immediately, Enable Input, il Reset del Do Once; il
%  Montage dall'animazione "Interact press button"; il widget
%  WB_InteractUI (la cartella UI degli asset, i due Border, Panel e Panel
%  Border, la E centrata); il Widget Component in Screen Space, Draw Size
%  100 x 100, Set Hidden in Game; il problema dei segni di spunta; la luce
%  accesa e spenta con Flip Flop e BP_ActivableLight; la chiusura del
%  corso.
%  Fonte: trascrizione automatica della dodicesima lezione (UE12).
%  Integrazioni dal libro "Blueprints Visual Scripting for Unreal Engine 5"
%  (Romero, Sewell; capitoli 4, 8, 11, 13, 15, 16 e 18) e nozioni tecniche
%  generali su Unreal Engine 5, sempre segnalate come approfondimenti. Il
%  libro non tratta la Default Pawn, il Widget Component in Screen Space,
%  il trucco del Select, i Trace Channel per le interazioni né Stop
%  Movement Immediately: su questi argomenti gli approfondimenti si basano
%  sul funzionamento generale dell'engine.
%
%  Pacchetti richiesti nel preambolo del libro (già presenti nel
%  main.tex): \usepackage{booktabs}, \usepackage{amsmath},
%  \usepackage{amssymb}, \usepackage{tikz} con
%    \usetikzlibrary{shapes.geometric,shapes.symbols,shapes.misc,
%      arrows.meta,positioning,calc,fit}
%  (le librerie vanno caricate nel preambolo, non nel capitolo).
%  Consigliati, come per gli altri capitoli: \usepackage[italian]{babel}
%  e \usepackage[T1]{fontenc} (il capitolo compila anche senza).
%
%  Collocazione nel libro: parte "Fondamenti di Unreal Engine",
%  dodicesimo capitolo (numerazione propria 1.12). Nel file principale
%  del libro, subito dopo la riga dell'undicesima lezione:
%
%    % =====================================================================
%  Fondamenti di Unreal Engine, lezione 12 — Il menu principale e il
%  sistema di interazione.
%  Dodicesima e ultima lezione del corso breve di introduzione a Unreal
%  Engine (focus sui Blueprint) con questo docente, online, nel
%  pomeriggio (pausa alle 17:26, ripresa "alle diciassette cinquantacinque,
%  diciotto"). Docente titolare (come nelle lezioni precedenti, nella
%  trascrizione non viene mai chiamato per nome). Da qui in poi il corso
%  prosegue con Davide e il level design.
%  Contenuti, nell'ordine della lezione: le domande d'apertura (recuperare
%  le lezioni o seguire in diretta; l'esame, i Big One dei tre anni e il
%  pezzo da portfolio; il designer che prototipa con i Blueprint; "less is
%  more" e il testo iniziale di Star Wars); gli annunci (ultima lezione,
%  le lezioni con Davide, il server Discord e i suoi canali, le
%  registrazioni fino a settembre); la scelta dell'ordine (prima il menu,
%  poi le interazioni); la dimostrazione del sistema di interazione; il
%  menu: il Widget Blueprint WB_Menu, la mappa vuota M_Menu (Empty Level,
%  il piano con il materiale a griglia), il Canvas Panel, i due Button
%  (Size 300 x 80, ancora al centro, Alignment), i nomi Button_StartGame e
%  Button_QuitGame, i Text dentro i bottoni, lo Style (Normal, Hovered,
%  Pressed), il codice esadecimale sRGB del colore, l'Image dietro ai
%  bottoni e l'ordine della Hierarchy, l'alfa a 0,2; il Level Blueprint
%  (Begin Play, Create Widget, Add to Viewport), il GameMode Override con
%  la Game Mode base e la Default Pawn, il Camera Actor, la Directional
%  Light, Set View Target with Blend (prima cercato sul Camera Manager,
%  poi trovato sul Player Controller); il Graph del menu (Event Construct,
%  Set Show Mouse Cursor, Set Input Mode UI Only, On Clicked, Quit Game,
%  Set Input Mode Game Only, Open Level by Name con ThirdPersonMap); i
%  colori di hover e pressione; il diorama del menu (il bot in Idle, la
%  porta modulare, la luce); i problemi degli studenti (il personaggio che
%  compare ancora, World Settings, il Player Start fuori campo, il target
%  dell'Add to Viewport, il cursore, la spunta del Set Show Mouse Cursor);
%  le interfacce (l'intermediario, il volo di linea) e i Blueprint
%  Component (Actor Component e Scene Component); il componente
%  BPC_RayInteract: la funzione LookForInteractions, Sphere Trace by
%  Channel con Radius 5, Get World Location, il Select con il trucco della
%  variabile Scene Component, Get Owner e Get Component by Class (Camera
%  Component), Is Valid, Get Forward Vector, DistanceTrace a 300, il
%  componente aggiunto al bot, Draw Debug Type, l'Event Tick; il Trace
%  Channel InteractionTrace (Ignore); Break Hit Result, Hit Actor, Does
%  Object Implement Interface; la Blueprint Interface BPI_InteractionRay
%  con le tre funzioni; la pausa (la stamina di Tommaso, il Call che
%  mancava; il rumore della trascrizione); dopo la pausa (la BPI aperta
%  "male"); la variabile Interactable, il diverso (!=), il Validated Get,
%  Interaction End e Interaction Start come messaggi, lo svuotamento
%  della variabile; il problema di tipo di uno studente; l'Event Tick del
%  componente; l'actor BP_InteractableTest (Cube e Widget), Class Settings
%  e Implemented Interfaces, gli eventi dell'interfaccia con tre Print
%  String colorati, la collisione Block su InteractionTrace; IA_Interact
%  (copia di IA_Run, Pressed e Released) e il tasto E nell'IMC; nel bot:
%  Do Once, Interaction (Message), Set Actor Rotation con il solo Yaw da
%  Find Look at Rotation, Play Montage, la variabile Interacting, Disable
%  Input, Stop Movement Immediately, Enable Input, il Reset del Do Once; il
%  Montage dall'animazione "Interact press button"; il widget
%  WB_InteractUI (la cartella UI degli asset, i due Border, Panel e Panel
%  Border, la E centrata); il Widget Component in Screen Space, Draw Size
%  100 x 100, Set Hidden in Game; il problema dei segni di spunta; la luce
%  accesa e spenta con Flip Flop e BP_ActivableLight; la chiusura del
%  corso.
%  Fonte: trascrizione automatica della dodicesima lezione (UE12).
%  Integrazioni dal libro "Blueprints Visual Scripting for Unreal Engine 5"
%  (Romero, Sewell; capitoli 4, 8, 11, 13, 15, 16 e 18) e nozioni tecniche
%  generali su Unreal Engine 5, sempre segnalate come approfondimenti. Il
%  libro non tratta la Default Pawn, il Widget Component in Screen Space,
%  il trucco del Select, i Trace Channel per le interazioni né Stop
%  Movement Immediately: su questi argomenti gli approfondimenti si basano
%  sul funzionamento generale dell'engine.
%
%  Pacchetti richiesti nel preambolo del libro (già presenti nel
%  main.tex): \usepackage{booktabs}, \usepackage{amsmath},
%  \usepackage{amssymb}, \usepackage{tikz} con
%    \usetikzlibrary{shapes.geometric,shapes.symbols,shapes.misc,
%      arrows.meta,positioning,calc,fit}
%  (le librerie vanno caricate nel preambolo, non nel capitolo).
%  Consigliati, come per gli altri capitoli: \usepackage[italian]{babel}
%  e \usepackage[T1]{fontenc} (il capitolo compila anche senza).
%
%  Collocazione nel libro: parte "Fondamenti di Unreal Engine",
%  dodicesimo capitolo (numerazione propria 1.12). Nel file principale
%  del libro, subito dopo la riga dell'undicesima lezione:
%
%    \input{UE/UE12_menu_e_sistema_di_interazione}
%
%  Tutte le etichette di questo capitolo hanno il prefisso "ue12:"
%  (ue12:cap:..., ue12:sec:..., ue12:fig:..., ue12:tab:...), e tutti gli
%  stili TikZ hanno il prefisso "ue12" (ue12node, ue12exec, ue12data,
%  ...) per non collidere con chiavi native di TikZ né con quelli delle
%  lezioni precedenti. Il capitolo è autonomo: non contiene \ref verso
%  altri capitoli (i rimandi alle lezioni precedenti sono a parole).
%  Le figure sono ricostruzioni schematiche di quanto mostrato e detto a
%  lezione, non copie delle schermate del docente. Dove la trascrizione
%  è troppo rovinata per ricostruire un dettaglio, il capitolo lo dice
%  invece di inventarlo.
% =====================================================================

\chapter[Menu e sistema di interazione]{Il menu principale e il sistema di
interazione}
\label{ue12:cap:interazione}

Questa è la \textbf{dodicesima lezione} del corso breve di \textbf{introduzione a Unreal
Engine}, ed è l'\textbf{ultima} con questo docente: ``oggi è l'ultima lezione, almeno
mia''. Nella lezione precedente il personaggio, il \textbf{bot}, ha ricevuto la vita, il
danno della torretta, la stamina e le due barre dell'HUD; restavano in sospeso due
argomenti annunciati, il \textbf{menu} e le \textbf{interazioni}. Oggi si fanno entrambi.
``Sono argomenti che abbiamo già fatto in precedenza e adesso li uniamo'': Widget
Blueprint, Level Blueprint, trace, collisioni personalizzate, Montage, timer e attivabili
tornano tutti, messi insieme.

Il primo argomento, il più corto, è un \textbf{menu principale}: una mappa a parte, con un
piccolo \emph{diorama} 3D sullo sfondo, una camera fissa, due bottoni (\emph{gioca} ed
\emph{esci}) e il passaggio alla mappa di gioco. Il secondo, ``piuttosto lungo'', è un
\textbf{sistema di interazione} del tipo ``premi E per interagire'', costruito con due
strumenti nuovi: un \textbf{Blueprint Component}, che si aggiunge a qualunque actor e
``traccia'' dalla camera gli oggetti interagibili, e una \textbf{Blueprint Interface}, il
``porto'' attraverso cui il personaggio dice a un oggetto \emph{sto iniziando a guardarti},
\emph{non ti guardo più}, \emph{interagisco con te}. L'oggetto di prova mostra un piccolo
widget con la lettera E quando è inquadrato; premendo E il personaggio si gira verso di lui
e suona un'animazione; alla fine, lo stesso oggetto accende e spegne una luce.

La lezione è \textbf{online}, come le precedenti, e si apre con una conversazione più
generale: come recuperare le lezioni, che cosa aspetta gli studenti negli anni successivi,
e un consiglio di design, ``\emph{less is more}''. Durante la pausa il docente risolve il
problema della barra della stamina rimasto aperto la volta precedente; si chiude con i
ringraziamenti e con la raccomandazione che riassume il corso: ``ricordatevi di validare
sempre molte variabili''.

Tre avvertenze di lettura, le stesse delle lezioni precedenti. La prima: la trascrizione
è automatica e storpia moltissimi nomi. Unreal è ``Real'' e ``Arrial''; il
\textbf{diorama} è ``diuramino''; il Widget Blueprint \textbf{WB\_Menu} è ``VB menu'' e
``VBIN menu'', e \textbf{WB\_InteractUI} è ``VB Interact UI'' e ``Interacco uai''; il
\textbf{Begin Play} è ``big in play'' e ``back in play''; la \textbf{ThirdPersonMap} è
``Terperson map'' e ``there person''; il \textbf{Scene Component} è ``sin component'' e
``sim component''; il componente \textbf{BPC\_RayInteract} è ``BPC Reinteract'', ``Ray
interact one'' e ``IPC reinteract''; la Blueprint Interface \textbf{BPI\_InteractionRay} è
``VPI interaction'', ``VPA Interaction Array'' e ``bibliotecazione:ray''; l'\textbf{Hit
Actor} è ``it actor'' e ``I Tactor''; il \textbf{Validated Get} è ``variegate'' e ``vari da
get''; \textbf{Does Object Implement Interface} è ``Hacker does object implement
interface''; il \textbf{Draw Debug Type} è ``draw the bug type''; il bot è ``Boto'' e
``Boto2'', e il suo Blueprint ``BiPi Bot''; \textbf{Find Look at Rotation} è
``findLocatrotation''; l'\textbf{Anim Montage} è ``un montaggio un po' brasilero''; i
\emph{Big One} del corso di laurea sono ``biggest one''. ``Mandragata'', invece, non è un
errore: è il romanesco \emph{mandrakata}, ``trucchetto'', e il docente lo usa apposta.
Dove la ricostruzione è sicura il capitolo usa la forma corretta; dove non lo è riporta la
forma della trascrizione e lo dice. La seconda: durante la pausa il riconoscimento vocale
produce un lungo blocco di testo senza senso e ripetuto decine di volte (``grazie a tutti
per l'attenzione che ci ha spiegato'', ``qual è la risposta che ci può essere?''), seguito
da un frammento che sembra l'audio di un video in sottofondo (una ``blue map'' che mostra
``le secret room'' senza sprecare bombe, e ``un'altra monetina''): non appartengono alla
lezione e sono stati ignorati. La terza: buona parte della lezione è fatta di indicazioni
a voce mentre si lavora (``tira fuori'', ``attacca qua'', ``vai giù''), spesso rivolte
allo schermo di uno studente; il capitolo ricostruisce i grafi da ciò che il docente dice
mentre li costruisce, distingue i tentativi dalle soluzioni, completa i nomi esatti dei
nodi e dei pin di Unreal Engine 5, e dichiara i punti in cui la trascrizione lascia un
buco.

\paragraph{Come leggere questo capitolo.} Il capitolo segue l'ordine della lezione. Le
figure ricostruiscono la disposizione del menu, i due grafi che lo fanno funzionare (Level
Blueprint e Widget Blueprint), il ruolo di componente e interfaccia, la funzione che
cerca gli interagibili, il ciclo di vita della variabile Interactable e il grafo
dell'interazione nel personaggio. Le tabelle riassumono gli asset e le variabili della
giornata, gli errori incontrati durante la lezione e i nodi. Qua e là, dei paragrafi di
\textbf{approfondimento} collegano la lezione al libro di Marcos Romero e Brenden Sewell,
\emph{Blueprints Visual Scripting for Unreal Engine 5}, il testo di riferimento per questa
parte del libro, e a qualche nozione tecnica generale su Unreal Engine 5; sono segnalati
come tali e non fanno parte della lezione.

Il percorso del capitolo:
\begin{enumerate}
  \item le \textbf{domande d'apertura}: recuperare le lezioni, l'esame e i Big One, il
        lavoro del designer, \emph{less is more} (sezione~\ref{ue12:sec:inizio});
  \item gli \textbf{annunci}, la scelta dell'ordine e la \textbf{dimostrazione} del
        sistema di interazione (sezione~\ref{ue12:sec:annunci});
  \item il \textbf{menu}: la mappa, il widget, i bottoni e i loro colori
        (sezione~\ref{ue12:sec:menu});
  \item il \textbf{Level Blueprint} del menu: il widget sullo schermo, la Game Mode e la
        camera (sezione~\ref{ue12:sec:levelbp});
  \item il \textbf{Graph del menu}: cursore, input e bottoni; il diorama e i problemi
        degli studenti (sezioni \ref{ue12:sec:menugraph} e~\ref{ue12:sec:diorama});
  \item \textbf{interfacce} e \textbf{Blueprint Component} (sezione~\ref{ue12:sec:concetti});
  \item il componente \textbf{BPC\_RayInteract} e la sua Sphere Trace
        (sezione~\ref{ue12:sec:componente});
  \item il \textbf{Trace Channel} e la \textbf{Blueprint Interface}
        (sezione~\ref{ue12:sec:bpi});
  \item la \textbf{pausa}: la stamina di Tommaso (sezione~\ref{ue12:sec:pausa});
  \item la logica della variabile \textbf{Interactable} (sezione~\ref{ue12:sec:logica});
  \item l'oggetto di prova \textbf{BP\_InteractableTest} (sezione~\ref{ue12:sec:test});
  \item l'\textbf{input} e l'interazione nel personaggio (sezioni \ref{ue12:sec:input}
        e~\ref{ue12:sec:player});
  \item il widget \textbf{``premi E''} (sezione~\ref{ue12:sec:widget});
  \item la \textbf{luce} interattiva e la \textbf{chiusura} del corso (sezioni
        \ref{ue12:sec:luce} e~\ref{ue12:sec:chiusura}).
\end{enumerate}

% ---------------------------------------------------------------------
\section{Prima della lezione: recuperare, l'esame, il designer}
\label{ue12:sec:inizio}

\paragraph{Recuperare o seguire?} In apertura, prima che arrivino gli altri, uno
studente rimasto indietro (``le sto andando pian piano'') chiede se gli convenga
recuperare le lezioni vecchie o seguire questa. (Il nome, nella trascrizione, è incerto:
il docente lo chiama una volta ``Dave'', ma poco dopo ricorda questa conversazione come
quella fatta ``con Lorenzo''.) La risposta: seguire la lezione in diretta ha un vantaggio,
``se c'è un qualche particolare dell'argomento, puoi farmi direttamente una domanda''; poi,
siccome oggi si uniscono argomenti già visti, ``magari ti rifai anche quelle precedenti in
modo tale da stare un po' più al passo''. L'unico prerequisito per l'esercizio di oggi è
``avere il personaggio che si muove'', che ``può guardarsi in giro''; ``puoi anche usare
il personaggio standard, volendo''. Le domande si possono fare anche dopo, sul
\textbf{Discord}: ``mandarmi le domande anche quando non ci vediamo più, perché oggi è
l'ultima lezione''. Lo studente sceglie di seguire e di riguardare la registrazione.

\paragraph{L'esame e i Big One.} Seconda domanda: ci sarà un esame, con un progetto da
consegnare? ``Nel corso no'': questo corso ``è strutturato che sia un'introduzione per
voi''. Il docente chiede che corso farà lo studente (game designer: ``tanto farete sia game
sia level designer'') e descrive il percorso, avvertendo che è quello che ha fatto lui
(``almeno io feci così, adesso non so come sono cambiate le cose''):
\begin{itemize}
  \item al \textbf{primo anno} il progetto principale è un \textbf{gioco da tavolo}:
        ``farete analog game design'';
  \item al \textbf{secondo anno} il \textbf{Big One} è un progetto su Unreal, ``perché
        Unreal è molto più improntato per i designer'';
  \item al \textbf{terzo anno} il progetto più grande (nella trascrizione ``biggest one''),
        che può anche essere la continuazione del secondo, ``se il vostro progetto nel
        secondo anno è andato bene'', da portare fino al rilascio su \textbf{Steam},
        ``con l'Accademia assieme''.
\end{itemize}
In questo senso l'esame c'è: ``sarà farvi un \textbf{pezzo da portfolio}'', ``un gioco che
potete mettere sul vostro portfolio e promuovervi''. ``Quindi, più sapete, meglio è.''

\paragraph{Che cosa fa il designer con Unreal.} ``Io sono designer'': per un designer
Unreal ``servirà principalmente per prototipare le meccaniche'', cioè ``concretizzare
l'idea che tu hai studiato, che tu hai tirato fuori, concretizzarla in editor e vedere se
effettivamente funziona''. Il suo strumento sono i \textbf{Blueprint}, ``quindi non andare
troppo in depth'' nella programmazione: in ambito lavorativo il designer ha ``un
programmatore che lo assiste''; ``poi, quando si sta da soli, si fa di tutto''.

\paragraph{Less is more.} Un consiglio ``da esperienza'': \textbf{less is more}. ``Più è
leggibile un qualcosa anche nella vostra testa a livello di design, più sarà leggibile per
un giocatore.'' L'esempio è ``la regola delle parole a schermo'': il testo che scorre
all'inizio di \emph{Star Wars}. Chi l'ha scritto lo rilegge in venti secondi, perché lo
conosce già; lo spettatore, che deve ``cominciare a creare le connessioni mentali'', ci
mette ``non 20 secondi, ma 50''. Nel gioco vale lo stesso: chi lo ha progettato e testato
pensa ``vabbè, tanto questo lo so ormai'', ma il giocatore ``ci metterà il triplo del
tempo''. ``Più roba fate, più tempo il giocatore dovrà metterci per imparare tutto
quanto.'' ``Credo che \emph{less is more} ve lo faranno fare'', aggiunge, ``ve lo
diranno''; e per qualunque domanda di design, ``possiamo tranquillamente pure parlare''.

\paragraph{Approfondimento: la maledizione della conoscenza.} Il fenomeno descritto dal
docente ha un nome in psicologia, \emph{curse of knowledge}: chi conosce un'informazione
fatica a immaginare come la percepisce chi non la conosce. È il motivo per cui il
\emph{playtest} con persone esterne al team è insostituibile, e per cui le regole di un
gioco andrebbero introdotte una alla volta, ciascuna quando serve (il cosiddetto
\emph{onboarding} progressivo). I numeri del docente (20 secondi contro 50, ``il
triplo'') sono un'immagine, non una misura. Questa è una nozione generale di design, non
della lezione.

% ---------------------------------------------------------------------
\section{Gli annunci, l'ordine della giornata e la dimostrazione}
\label{ue12:sec:annunci}

\paragraph{Che cosa si fa oggi.} Prima che arrivino tutti, il docente chiede agli
studenti di scegliere: oggi ci sono ``sia il sistema di interazioni, che è piuttosto
lungo'', sia il \textbf{menu}, ``che l'altra volta non abbiamo finito, praticamente non
abbiamo neanche iniziato''. Prima il lungo o prima il corto? Gli studenti scelgono
``corto e lungo, così facciamo la pausa prima'' (``andiamo a un caffè''). Uno studente
ricorda un'altra cosa in sospeso: ``durante la pausa sistemare il problema della stamina,
che se tengo premuto non si vede, ma solo quando mollo il tasto si vede'' (è il problema
rimasto aperto alla fine della lezione precedente); ``ricordamelo appena andiamo in
pausa''.

\paragraph{Gli annunci.} Arrivati gli altri, il docente ripete: ``oggi praticamente è la
mia ultima lezione''. Seguiranno ``circa quattro lezioni'' (``adesso non mi ricordo quante
lezioni avete'') con \textbf{Davide}, dedicate alla parte di \textbf{level design}: ``i
tool che vi serviranno per fare level design'' e alcuni approfondimenti, ``credo di
materiali''. ``Ciò non vuol dire che io vi abbandoni'':
\begin{itemize}
  \item il \textbf{canale Discord} resta aperto, anche quando inizierà il corso vero e
        proprio: ``fate domande e dovrebbe arrivare qualcuno in vostro aiuto'';
  \item anche i lavori si possono postare per avere un feedback (il docente cita quello
        di \textbf{Federico La Porta}, ``che è bravissimo'', a cui ha già risposto);
  \item i due canali vanno tenuti separati: quello del corso (``Introduzione a Unreal'')
        per i compiti, ``così almeno non mandate mischiati, perché sennò perdo il filo'',
        e la chat generale ``per qualsiasi comunicazione generica'';
  \item le registrazioni restano disponibili ``fino a settembre, se non sbaglio'', ``quindi
        avete comunque abbastanza tempo per poterle recuperare''.
\end{itemize}

\paragraph{La dimostrazione.} ``Adesso vi faccio vedere solamente un'introduzione, come
nei classici video tutorial.'' Il personaggio è quello della volta scorsa, ``che si muoveva
con stamina e adesso perde vita con la torretta''. Dalla sua testa parte un \textbf{trace}
sottile, ``tipo un cilindrone'': ``questo sarà il nostro sistema che identificherà un
possibile \emph{interactable} tramite la telecamera''. A sinistra, sullo schermo, i
messaggi: se la camera inquadra l'oggetto compare \emph{inizio interazione}, se si
distoglie lo sguardo \emph{fine interazione}; premendo il tasto, il personaggio ``si gira in
automatico'' verso l'oggetto e interagisce. ``Visto che comunque abbiamo creato questo
sistema, è un buon stampo modulare'': in futuro, pulsanti e altri oggetti potranno avere
``l'interazione un po' più precisa'', riusando ciò che si è visto con i pulsanti classici:
``sempre premi pulsantino e fai cosetta, e animazione''.

% ---------------------------------------------------------------------
\section{Il menu: la mappa, il widget e i bottoni}
\label{ue12:sec:menu}

\paragraph{Il Widget Blueprint.} Nella cartella \textbf{UI} del progetto, tasto destro,
\emph{User Interface}, \textbf{Widget Blueprint}, ``come abbiamo creato per l'HUD''. Alla
richiesta della classe base, ``non ci serve una particolare, ce ne serve una generica'':
\textbf{User Widget}. Il nome: \textbf{WB\_Menu}.

\paragraph{Una mappa per il menu.} ``Una cosa importante che io almeno consiglio'': creare
una \textbf{nuova mappa} dedicata al menu, ``quindi differenziare il gioco dal menu''. ``Si
possono fare anche volendo insieme, ma a livello di gestione'' è più comodo separarli, ``ed
è molto più rapido''. Da \textbf{File}, \textbf{New Level}, si sceglie \textbf{Empty
Level}: ``se vi diventa tutto nero, è normale, ragazzi: è un empty level''. Un livello
\emph{Basic}, invece, contiene già ``la directional light, la skybox, tante cose che sono
già prefatte''; qui si parte da zero. Poi:
\begin{itemize}
  \item il menu si può fare ``sia in 2D, con immagini 2D'', sia con elementi 3D, ``una
        specie di diorama''; per cominciare, un \textbf{piano} abbastanza grande (dalle
        forme base: il docente prende prima un cubo, poi si corregge), con un materiale a
        griglia, ``il material base di Unreal'';
  \item si salva la mappa con \textbf{Ctrl+S} nella cartella delle mappe, con il nome
        \textbf{M\_Menu}.
\end{itemize}
(Un dettaglio marginale: il docente nota che qualcosa nel livello risulta \emph{Hidden in
Game}, poi si accorge che non lo è; la trascrizione non dice di che cosa si tratti.)

\paragraph{Il Canvas Panel.} Aperto WB\_Menu, ``la prima cosa, che in questo caso servirà
per adattarlo allo schermo'': il \textbf{Canvas Panel}, ``la classica canvas''. La
\textbf{Hierarchy} si vede ``da qua sotto''. A uno studente che non trova la
\textbf{Palette} il docente mostra il pannello a sinistra dell'editor.

\paragraph{Due bottoni.} ``Adesso andremo a fare due pulsanti molto semplici: uno di
\emph{exit} e uno di \emph{open the game}.'' Nella Palette, la sezione
\textbf{Common} (accanto alle altre: Audio, Input, Panel, Primitive e così via, che la
trascrizione elenca a fatica) contiene il \textbf{Button}; se ne trascinano due al centro
dello schermo. La disposizione passa per qualche tentativo ad alta voce (posizioni e
allineamenti come $-2{,}7$ e $-5$, presto abbandonati: ``ok, mi sono perso''); la versione
pulita, quella che il docente detta alla fine, è questa:
\begin{enumerate}
  \item \textbf{Size}: X \textbf{300}, Y \textbf{80} (``60? No, troppo piccolo'');
  \item \textbf{Anchors} al \textbf{centro}, l'icona centrale del menu delle ancore, ``in
        modo tale che questo bottone sia sempre ancorato al centro''; poi si
        \textbf{azzerano} i valori della \textbf{Position};
  \item \textbf{Alignment} \textbf{0,5} su X e \textbf{1,5} su Y, ``perché così è più
        corretto'': il bottone resta centrato orizzontalmente e si alza di una volta e
        mezza la sua altezza rispetto al centro dello schermo;
  \item per il secondo bottone, ``copiate e incollate letteralmente'' il primo, azzerate
        la posizione e cambiate solo l'allineamento verticale: dopo un ``$-1{,}5$? No,
        troppo'' e un ``$-0{,}5$'', il docente si ferma a \textbf{$-0{,}3$} (``va bene''),
        cioè il secondo bottone sta un po' sotto il centro.
\end{enumerate}
(I valori del secondo bottone, nella trascrizione, sono incerti: il senso è che i due
bottoni stanno uno sopra e uno sotto il centro, con un po' di spazio fra loro.) A uno
studente il copia e incolla non funziona: ``guarda nella Hierarchy: se c'hai due bottoni,
vuol dire che ce ne hai due'', magari sovrapposti; altrimenti, ``metti semplicemente un
altro bottone tramite il Common''.

\paragraph{I nomi.} ``Una cosa che è importante da fare'': nella Hierarchy si rinominano i
due bottoni con \textbf{F2}, \textbf{Button\_StartGame} e \textbf{Button\_QuitGame} (la
trascrizione dice ``Button With Game''). ``Così abbiamo semplicemente perfettamente
nomenclati.''

\paragraph{Il testo dentro il bottone.} ``Una cosa carina che già esiste'': un
\textbf{Text} (dalla sezione Common) trascinato \emph{sopra} un bottone ``si va a incastrare
per benino'', diventa suo figlio e si centra da solo. Nei \textbf{Details} del Text, sezione
\textbf{Content}, il campo \textbf{Text}: ``cancellate e scrivete \emph{Play Game}, o
\emph{Play}, come volete''; nell'altro bottone, \emph{Quit}. Il testo bianco ``non si vede
tantissimo'', ma i colori si cambiano subito dopo.

\paragraph{Lo Style del bottone.} Selezionando il bottone (``mi raccomando''), nei Details
la sezione \textbf{Style} ha tre stati, ciascuno con il suo colore:
\begin{itemize}
  \item \textbf{Normal}: il bottone a riposo;
  \item \textbf{Hovered}: ``vuol dire che ci passi sopra il cursore'';
  \item \textbf{Pressed}: ``se premuto''.
\end{itemize}
(La trascrizione nomina anche i \emph{foreground} e lo stato \emph{Disabled}: nelle
versioni recenti dell'engine lo stile del bottone ha anche i colori del primo piano per
ciascuno stato.) Cliccando sulla barra del colore si apre la ruota dei colori; il docente
mette il Normal \textbf{blu}, ``e avremo questo cambiamento in diretta''.

\paragraph{Copiare un colore esatto.} ``Se vogliamo prendere perfettamente il colore che ci
interessa per spostarlo all'altro bottone'': nella finestra del colore, in basso, ci sono i
campi \textbf{Hex sRGB} e \textbf{Hex Linear}. ``Potete prendere ambi i due'': si copia il
codice con Ctrl+C, si apre lo stesso stato nell'altro bottone e lo si incolla \emph{nello
stesso campo} (``se avete preso sRGB, sempre su sRGB''). È lo stesso principio dei siti
web: ``in HTML, per mettere un colore, i colori sono tutti identificati tramite un codice
colore''.

\paragraph{Un'immagine dietro i bottoni.} Si aggiunge un'\textbf{Image}, allargata dietro i
due bottoni; ma ``come potete vedere, mi copre tutto quanto''. ``Qui ci viene in aiuto la
Hierarchy'': come i livelli di Photoshop, gli elementi del Canvas sono disegnati
nell'ordine dell'elenco, e ``quello che è più sotto'' nell'elenco, stranamente, è quello
``a livello più alto'', disegnato sopra gli altri. L'immagine, aggiunta per ultima, è in
fondo all'elenco e quindi davanti ai bottoni. La si trascina nella Hierarchy fino a far
comparire una \textbf{linea blu} fra il Canvas Panel e il primo bottone, e la si lascia lì:
ora è disegnata per prima, cioè dietro. Poi la si rende trasparente: nel suo
\textbf{Color and Opacity} (che il docente cerca nel Brush), l'\textbf{alfa} a 0,5, ``o
forse un po' più basso'', \textbf{0,2}. ``Questa è la alfa di un colore: l'intensità,
l'opacità.'' Ne viene ``una specie di schermatina piccoletta'', regolata a mano attorno ai
bottoni (figura~\ref{ue12:fig:layout}).

\begin{figure}[htbp]
\centering
\begin{tikzpicture}[font=\footnotesize, x=1cm, y=1cm]
  % schermo
  \draw[very thick, fill=gray!5] (0,0) rectangle (8.0,5.0);
  \node[anchor=north west, font=\scriptsize] at (0.1,4.95) {Canvas Panel (schermo)};
  % immagine trasparente
  \draw[thick, fill=blue!8, opacity=0.9] (2.4,1.15) rectangle (5.6,4.15);
  \node[font=\scriptsize, anchor=north] at (4.0,4.12) {Image, alfa 0,2};
  % bottoni
  \draw[thick, fill=blue!35, rounded corners=2pt] (2.9,2.70) rectangle (5.1,3.30);
  \node at (4.0,3.0) {\textbf{Play Game}};
  \draw[thick, fill=blue!35, rounded corners=2pt] (2.9,1.62) rectangle (5.1,2.22);
  \node at (4.0,1.92) {\textbf{Quit}};
  % ancora al centro
  \draw[red!70!black, thick] (4.0,2.5) circle (0.09);
  \draw[red!70!black, dashed] (0,2.5) -- (8.0,2.5);
  \draw[red!70!black, dashed] (4.0,0) -- (4.0,5.0);
  \node[red!70!black, font=\scriptsize, anchor=west] at (5.7,2.62) {ancora al centro,};
  \node[red!70!black, font=\scriptsize, anchor=west] at (5.7,2.30) {Position 0, 0};
  \node[font=\scriptsize, anchor=west, align=left] at (5.7,3.35) {Alignment\\0,5 / 1,5};
  \node[font=\scriptsize, anchor=west, align=left] at (5.7,1.65) {Alignment\\0,5 / $-0{,}3$\\(circa)};
  % hierarchy
  \begin{scope}[xshift=9.0cm]
    \draw[thick, fill=gray!8] (0,0.6) rectangle (4.2,4.4);
    \node[anchor=north west, font=\scriptsize\bfseries] at (0.1,4.35) {Hierarchy};
    \node[anchor=west] at (0.2,3.6) {[WB\_Menu]};
    \node[anchor=west] at (0.45,3.1) {Canvas Panel};
    \node[anchor=west, fill=blue!10] at (0.7,2.6) {Image};
    \node[anchor=west] at (0.7,2.1) {Button\_StartGame};
    \node[anchor=west] at (0.95,1.65) {Text ``Play Game''};
    \node[anchor=west] at (0.7,1.2) {Button\_QuitGame};
    \node[anchor=west] at (0.95,0.8) {Text ``Quit''};
    \draw[-{Stealth[length=2mm]}, thick] (4.45,3.0) -- node[right, font=\scriptsize,
      align=left]{disegnati\\in ordine:\\l'ultimo\\sta sopra} (4.45,0.9);
  \end{scope}
\end{tikzpicture}
\caption{Il menu nel Designer. I due bottoni sono ancorati al centro dello schermo con la
posizione azzerata; è l'Alignment (il pivot del bottone) a spostarli sopra e sotto il
centro. L'immagine trasparente è stata trascinata in cima alla Hierarchy, così viene
disegnata per prima, dietro ai bottoni. Ricostruzione dalla lezione: i valori del secondo
bottone sono approssimativi.}
\label{ue12:fig:layout}
\end{figure}

\paragraph{Approfondimento: Alignment, ancore e Size To Content.} Con l'ancora al centro e
la posizione a zero, l'\textbf{Alignment} decide quale punto del bottone coincide con il
centro dello schermo: 0 è il bordo sinistro (o superiore), 1 il bordo destro (o inferiore),
0,5 il centro. Valori fuori dall'intervallo da 0 a 1 sono leciti: con Y = 1,5 il punto
d'aggancio sta mezzo bottone \emph{sotto} il bordo inferiore, quindi il bottone sale; con Y
negativo, scende. È il trucco usato a lezione per impilare due bottoni senza toccare la
posizione. La via più ordinata, per menu con molti bottoni, è un \textbf{Vertical Box}
(sezione Panel della Palette) ancorato al centro, con i bottoni dentro: lo spazio fra gli
elementi si regola con il \emph{Padding} di ciascuno e non servono calcoli. Nel libro di
Romero e Sewell (capitolo 8, \emph{WinMenu}) i bottoni sono ancorati al centro e spostati
con la \emph{Position} (per esempio X $-180$, Y 150, Size 360 $\times$ 100), con l'alignment
di default: è la terza via, equivalente. Funzionamento generale dell'UMG, salvo il
riferimento al libro.

% ---------------------------------------------------------------------
\section{Il Level Blueprint del menu: il widget, la Game Mode e la camera}
\label{ue12:sec:levelbp}

\paragraph{Perché il Level Blueprint.} ``Adesso che ce famo con questa roba?'' Il widget
va mostrato appena si apre la mappa del menu, e solo lì. Per questo si usa il
\textbf{Level Blueprint}, ``che abbiamo usato nella prima lezione'': ogni livello ha il
suo, ed è il posto giusto per la logica che appartiene a \emph{quella} mappa, ``perché
vogliamo che questa mappa sia dedicata apposta al nostro menu''. Si apre dalla toolbar
del Level Editor, dal pulsante dei Blueprint (``c'è il cubetto col più''),
\textbf{Open Level Blueprint}.

\paragraph{Il widget sullo schermo.} ``Ci interessa semplicemente usare il Begin Play,
perché vogliamo che si crei già da subito'':
\begin{enumerate}
  \item \textbf{Event BeginPlay} $\rightarrow$ \textbf{Create Widget}, con
        \textbf{Class} = WB\_Menu, ``come abbiamo visto l'altra volta'';
  \item \textbf{Add to Viewport}, con il \textbf{Return Value} del Create Widget nel
        Target. Il docente pensa per un momento di promuovere il widget a variabile, come
        per l'HUD, poi rinuncia: ``no, beh, non lo facciamo, tanto è il level blueprint''.
\end{enumerate}

\paragraph{Il personaggio di troppo.} In Play il menu compare, ma ``è molto strano:
abbiamo il nostro player dietro''. La mappa usa ancora la Game Mode del progetto, che fa
nascere il bot con i suoi controlli. Nel menu, invece, serve ``una pawn senza controlli,
senza niente'', e Unreal ce l'ha già: nelle \textbf{World Settings} della mappa, alla voce
\textbf{GameMode Override}, si sceglie la \textbf{Game Mode base} di Unreal (nella
trascrizione ``la game mode basica che esiste, la default''), che usa come Pawn la
\textbf{Default Pawn}. In Play, ``già è una cosa diversa: non c'è più il nostro character,
sparito, perché in questa mappa ho overridato la modalità Game Mode classica'', quella
``proprio base di Unreal, in cui non c'è character in terza persona''. ``Ma, come potete
vedere, non ho controllo'' (o meglio, come si vedrà, la Default Pawn si muove ancora:
il controllo vero si toglie nel Graph del widget, sezione~\ref{ue12:sec:menugraph}).

\paragraph{La camera fissa.} ``Andiamo a vedere anche un'altra cosa'': nel livello si
aggiunge un \textbf{Camera Actor} (``scrivete \emph{camera actor}, eccola qua''), che sarà
la vista fissa del menu. Insieme, una \textbf{Directional Light} (``e abbiamo messo il
sole, ragazzi''), regolabile a piacere, e un muretto, ``un po' più carino, così capiamo
meglio''. Per dire al gioco di guardare da quella camera, di nuovo nel Level Blueprint:
\begin{itemize}
  \item la \textbf{referenza} alla camera: selezionato il Camera Actor nel livello, nel
        Level Blueprint tasto destro, \textbf{Create a Reference to CameraActor};
  \item il nodo che cambia la vista. Qui il docente sbaglia strada in diretta: prova da
        \emph{Get Player Camera Manager}, non trova il nodo giusto e rimanda (``adesso non
        mi ricordo benissimo, lo vediamo un attimo dopo''). Dopo aver finito il Graph del
        menu, ritrova la soluzione: ``non è che bisogna prendere il camera manager, ma
        bisogna andare sul controller''. Quindi \textbf{Get Player Controller}
        $\rightarrow$ \textbf{Set View Target with Blend}, con la camera nel pin
        \textbf{New View Target}; il \emph{Blend Time} resta al default, ``non ci
        interessa''.
\end{itemize}
``Questo blocco ci permette di blendare una nuova camera'': ``gli diciamo: guarda, questa
è la tua nuova camera, da d'ora in poi'' (figura~\ref{ue12:fig:menugraph}, in alto).

\paragraph{Il Player Start fuori campo.} Ultimo ritocco, emerso aiutando gli studenti: la
Default Pawn ``è una palla'', e se nasce davanti alla camera si vede, o si vede il suo
riflesso. ``Prendete un \textbf{Player Start} e mettetelo fuori campo'', ``così almeno non vi
si vede la pallina''.

\paragraph{Approfondimento: Game Mode, Default Pawn e Player Start.} La \textbf{Game Mode}
è la classe che stabilisce le regole di una mappa e, fra le altre cose, quale
\emph{Pawn} far nascere per il giocatore (\emph{Default Pawn Class}) e con quale
\emph{Player Controller}. Il progetto ne ha una predefinita (in \emph{Project Settings},
\emph{Maps \& Modes}); il \textbf{GameMode Override} delle World Settings la sostituisce
\emph{solo per quella mappa}, ed è il modo standard di avere un menu senza personaggio.
La classe scelta a lezione, \emph{GameModeBase}, usa come Pawn la \textbf{DefaultPawn}:
una sfera con un movimento ``da spettatore'' (WASD e mouse, senza gravità), pensata per
fare prove. Il Pawn nasce in un \textbf{Player Start}; se la mappa non ne ha, nasce
all'origine o nella posizione della camera dell'editor. Un'alternativa ancora più pulita
per un menu è una Game Mode con la \emph{Default Pawn Class} impostata su \emph{None}: il
Player Controller esiste (e quindi esistono il cursore e gli input dell'interfaccia), ma
nessun Pawn nasce. Funzionamento generale dell'engine.

\paragraph{Approfondimento: Set View Target with Blend nel libro.} Romero e Sewell
descrivono il nodo nel capitolo 15 (\emph{Blueprints Tips}): ``è una funzione della classe
Player Controller'', ``usata per passare la vista di gioco fra camere diverse''; ``il
parametro di input New View Target è l'actor da usare come view target, di solito una
camera''. Il loro esempio è proprio in un Level Blueprint: quando il giocatore entra nella
stanza del tesoro, la vista passa alla camera della stanza. Questo spiega l'errore del
docente: il \emph{Player Camera Manager} gestisce la vista, ma la funzione esposta ai
Blueprint per cambiarla appartiene al \emph{Player Controller}. Il \textbf{Blend Time} a 0
fa uno stacco netto; un valore positivo fa una transizione morbida dalla vista precedente,
con una curva scelta in \emph{Blend Func}.

% ---------------------------------------------------------------------
\section{Il Graph del menu: cursore, input e bottoni}
\label{ue12:sec:menugraph}

\paragraph{Event Construct.} Si torna in WB\_Menu, scheda \textbf{Graph}: ``vediamo le
funzionalità del nostro menu quando clicchiamo le cose''. Dei tre eventi già presenti
``non ci serve il Pre Construct, non ci serve il Tick''. Serve l'\textbf{Event Construct}:
``vuol dire quando questo widget è costruito'', ``è un po' il Begin Play'', ``viene
differenziato dal nome, giusto per non confonderti''.

\paragraph{Il cursore e l'input dell'interfaccia.} In Play, per ora, ``mi posso muovere''
(la Default Pawn) e il mouse non si vede. Nell'Event Construct:
\begin{enumerate}
  \item \textbf{Get Player Controller} $\rightarrow$ \textbf{Set Show Mouse Cursor},
        con la spunta attiva: il cursore resta visibile anche cliccando nella finestra di
        gioco. ``Però, ovviamente, posso ancora muovermi'';
  \item \textbf{Set Input Mode UI Only} (``enable input UI'', dice prima il docente, poi
        trova il nodo), con lo stesso Player Controller: ``questo ci permette di
        utilizzare tranquillamente il nostro widget''. Il pin \emph{Widget to Focus} resta
        vuoto: ``no, tanto non ci serve''.
\end{enumerate}
In Play, ora, ``se noi proviamo a muovere lateralmente, non si muove più''. ``Mi
raccomando: questa cosa, quando noi cliccheremo un pulsante, dovrà essere disattivata e
rimessa in un'altra modalità.''

\paragraph{I colori di hover e pressione.} Con il cursore visibile si possono provare gli
stati dello Style. Nel Designer il docente cambia il colore di \textbf{Hovered}: ``voglio
che quando io ci passo sopra con il mouse diventi tipo arancione, sì, un bell'arancione''
(``lo vediamo in preview, è normale''), e lo copia, con il codice esadecimale, sull'altro
bottone. Il risultato, ammette, è ``molto Ikea''. Il \textbf{Pressed} di default diventa
grigio, ``si vede pochissimo'': lo si mette \textbf{rosso}, su entrambi i bottoni.

\paragraph{I bottoni.} Nel Graph, a sinistra, fra le variabili compaiono
\textbf{Button\_QuitGame} e \textbf{Button\_StartGame} (sono variabili perché i bottoni,
nel Designer, hanno \emph{Is Variable} attivo di default). Ogni bottone ha i suoi eventi,
``l'abbiamo visto anche l'altra volta'': \emph{On Pressed}, \emph{On Released},
\emph{On Hovered}, \emph{On Unhovered} e \textbf{On Clicked}. Si prende l'On Clicked di
entrambi.
\begin{itemize}
  \item \textbf{On Clicked (Button\_QuitGame)} $\rightarrow$ \textbf{Quit Game}: ``è
        piuttosto semplice, letteralmente Quit Game è già presente, di standard
        proprio''. (``Adesso non mi ricordo se ti esce proprio dall'editor'': si veda
        l'approfondimento più sotto.)
  \item \textbf{On Clicked (Button\_StartGame)}: prima \textbf{Get Player Controller}
        $\rightarrow$ \textbf{Set Input Mode Game Only}, ``molto importante'', ``in modo
        tale che ritornino gli input al gioco''; poi \textbf{Open Level (by Name)}, con il
        nome della mappa di gioco. Il docente lo cerca nel Content Browser (``era
        MLevel1? Ah, questa no, era M level 2?'') e trova la mappa del template,
        \textbf{ThirdPersonMap}; prova a copiarne il nome con F2 (``non lo fa, va bene''),
        e lo scrive a mano.
\end{itemize}
Tornati nella mappa del menu e salvato tutto, in Play: Play Game ``ci carica la nostra
bellissima mappina'' (figura~\ref{ue12:fig:menugraph}).

\begin{figure}[htbp]
\centering
\begin{tikzpicture}[font=\footnotesize,
  ue12node/.style={draw, thick, rounded corners=4pt, align=center, fill=gray!8,
                  minimum height=0.9cm, text width=2.3cm},
  ue12ev/.style={ue12node, fill=red!12},
  ue12pure/.style={ue12node, fill=green!10, rounded corners=8pt},
  ue12exec/.style={-{Stealth[length=2.5mm]}, very thick, black!80},
  ue12data/.style={-{Stealth[length=2mm]}, thick, blue!60!black}]
  \node[font=\scriptsize\bfseries, anchor=west] at (-1.3,1.0) {Level Blueprint di M\_Menu};
  \node[ue12ev, text width=1.7cm] (bp) at (0,0) {\textbf{Event BeginPlay}};
  \node[ue12node] (cw) at (2.8,0) {\textbf{Create Widget}\\{\scriptsize Class: WB\_Menu}};
  \node[ue12node, text width=2.5cm] (av) at (5.8,0) {\textbf{Add to Viewport}\\{\scriptsize Target: Return
    Value}};
  \node[ue12node, text width=2.6cm] (vt) at (8.9,0) {\textbf{Set View Target with
    Blend}\\{\scriptsize Blend Time 0}};
  \node[ue12pure] (gpc1) at (7.4,-1.8) {Get Player Controller};
  \node[ue12pure, fill=blue!10] (cam) at (10.2,-1.8) {CameraActor\\{\scriptsize
    (referenza)}};
  \draw[ue12exec] (bp) -- (cw); \draw[ue12exec] (cw) -- (av); \draw[ue12exec] (av) -- (vt);
  \draw[ue12data] ([yshift=-3mm]cw.east) -- ([yshift=-3mm]av.west);
  \draw[ue12data] (gpc1.north) -- node[left, font=\scriptsize]{Target} ([xshift=-6mm]vt.south);
  \draw[ue12data] (cam.north) -- node[right, font=\scriptsize]{New View Target}
    ([xshift=6mm]vt.south);
  % widget graph
  \node[font=\scriptsize\bfseries, anchor=west] at (-1.3,-3.2) {Graph di WB\_Menu};
  \node[ue12ev, text width=1.7cm] (con) at (0,-4.2) {\textbf{Event Construct}};
  \node[ue12node] (smc) at (2.8,-4.2) {\textbf{Set Show Mouse Cursor}\\{\scriptsize
    spunta attiva}};
  \node[ue12node] (ui) at (5.7,-4.2) {\textbf{Set Input Mode UI Only}};
  \node[ue12pure] (gpc2) at (4.2,-5.9) {Get Player Controller};
  \draw[ue12exec] (con) -- (smc); \draw[ue12exec] (smc) -- (ui);
  \draw[ue12data] (gpc2.north) -- (smc.south); \draw[ue12data] (gpc2.north) -- (ui.south);
  \node[ue12ev, text width=2.6cm] (oq) at (0.3,-7.6) {\textbf{On Clicked}\\{\scriptsize
    (Button\_QuitGame)}};
  \node[ue12node] (qg) at (3.6,-7.6) {\textbf{Quit Game}};
  \draw[ue12exec] (oq) -- (qg);
  \node[ue12ev, text width=2.6cm] (os) at (0.3,-9.3) {\textbf{On Clicked}\\{\scriptsize
    (Button\_StartGame)}};
  \node[ue12node, text width=2.7cm] (gm) at (3.8,-9.3) {\textbf{Set Input Mode Game Only}};
  \node[ue12node, text width=2.8cm] (ol) at (7.2,-9.3) {\textbf{Open Level (by
    Name)}\\{\scriptsize ThirdPersonMap}};
  \node[ue12pure] (gpc3) at (3.8,-11.0) {Get Player Controller};
  \draw[ue12exec] (os) -- (gm); \draw[ue12exec] (gm) -- (ol);
  \draw[ue12data] (gpc3) -- (gm);
\end{tikzpicture}
\caption{I grafi del menu. In alto, il Level Blueprint della mappa M\_Menu crea il widget,
lo mostra e passa la vista alla camera fissa. In basso, il widget prende il controllo del
cursore quando nasce, esce dal gioco con Quit e, con Play Game, restituisce l'input al
gioco e apre la mappa di gioco. Ricostruzione dalla lezione.}
\label{ue12:fig:menugraph}
\end{figure}

\paragraph{Approfondimento: il menu nel libro.} Romero e Sewell costruiscono un menu molto
simile nel capitolo 8 (\emph{Creating Constraints and Gameplay Objectives}), il
\emph{WinMenu} che compare a fine partita: due bottoni \emph{Btn restart} e \emph{Btn
quit} ancorati al centro, ciascuno con un Text dentro (``font size 48''), e gli eventi
aggiunti dal pulsante \emph{+} accanto a \emph{On Clicked}, in fondo ai Details del bottone.
Il bottone Restart usa \emph{Open Level (by Object Reference)}, con la mappa scelta da un
menu a tendina, seguito da \emph{Remove from Parent}, ``che rimuove il widget dalla
vista''; il bottone Quit usa \emph{Quit Game}, ``così che il giocatore possa chiudere il
gioco''. Il menu viene mostrato dal personaggio, in un Custom Event \emph{End Game}: Set
Game Paused, \emph{SET Show Mouse Cursor} dal Get Player Controller (``per permettere al
giocatore di riprendere il controllo del mouse''), Create Widget, Add to Viewport. Nel
capitolo 15 il libro elenca anche le tre modalità di input:
\begin{itemize}
  \item \textbf{Set Input Mode Game Only}: ``solo il Player Controller riceve gli eventi
        di input'';
  \item \textbf{Set Input Mode UI Only}: ``solo la UI riceve gli eventi di input'';
  \item \textbf{Set Input Mode Game and UI}: ``la UI ha la priorità nel gestire un evento
        di input, ma se la UI non lo gestisce, lo riceve il Player Controller''; l'esempio
        è un negozio, in cui si scelgono le opzioni con il mouse ma ci si può allontanare
        con le frecce.
\end{itemize}
Una differenza utile da notare: \emph{Open Level (by Object Reference)} non si rompe se la
mappa viene rinominata, mentre \emph{Open Level (by Name)}, usato a lezione, richiede il
nome esatto (si può scrivere il nome semplice, \emph{ThirdPersonMap}, o il percorso
completo, se esistono due mappe con lo stesso nome).

\paragraph{Approfondimento: Construct, Quit Game e il cursore nella mappa di gioco.}
Alcune nozioni generali, non della lezione. \textbf{Event Pre Construct} viene eseguito
anche nell'editor, ogni volta che il widget viene ridisegnato nel Designer: serve per
anteprime e impostazioni visive; \textbf{Event Construct} solo in gioco, quando il widget
viene aggiunto allo schermo (e di nuovo se lo si rimuove e lo si aggiunge). \textbf{Quit
Game}, durante il Play in Editor, \emph{termina la sessione di gioco} e riporta
all'editor, senza chiudere l'editor; in un gioco pacchettizzato chiude l'applicazione. Il
dubbio del docente si scioglie così. Infine, \emph{Open Level} distrugge il mondo corrente
con tutti i suoi actor, widget compresi, e crea un Player Controller nuovo: per questo,
nella mappa di gioco, il menu sparisce da solo e il cursore torna nascosto, anche senza un
\emph{Remove from Parent} o un \emph{Set Show Mouse Cursor} a falso. Il \emph{Set Input
Mode Game Only} prima dell'Open Level, raccomandato dal docente, è comunque una buona
abitudine: il caso in cui non è superfluo è quello di un menu che si chiude \emph{senza}
cambiare livello, come un menu di pausa (nel libro, capitolo 11, il \emph{Resume} fa
Remove from Parent, Set Show Mouse Cursor a falso e Set Game Paused a falso).

% ---------------------------------------------------------------------
\section{Il diorama, e i problemi degli studenti}
\label{ue12:sec:diorama}

\paragraph{Il diorama.} Con il menu funzionante, ``passiamo sopra'' all'aspetto. Sullo
sfondo, ``un piccolo diorama, una piccola box, in cui ci andiamo a mettere tutto quello che
vogliamo''. ``Volendo, potete fare tutto un menu interattivo'': gli ostacoli, i Blueprint
creati nel corso, i bottoni, la torretta. Il docente si limita ai modelli:
\begin{itemize}
  \item il \textbf{bot}: dal Content Browser, o la sua Skeletal Mesh, ``o li metto
        direttamente l'animazione in Idle'': trascinando nel livello un asset di animazione,
        Unreal crea un actor con la mesh che la riproduce;
  \item la \textbf{porta} modulare del compito della volta scorsa, \textbf{SM\_Modular\_\allowbreak Wall\_Door}
        (``wall door, eccola qua''), con le ante \textbf{SM\_Door};
  \item qualche ritocco alla \textbf{camera} e alla \textbf{luce}: l'intensità a 10 è
        ``enorme'', poi 5, poi valori molto più bassi (la trascrizione, qui, è quasi
        incomprensibile: si sentono ``0,01'' e ``1'').
\end{itemize}
``Questo caso è più una cosa artistica, ragazzi, quindi potete fare un po' come volete.''
Le funzioni di base, entrare nel gioco e uscire, ci sono. ``Questo è un menu che io di
solito faccio quando faccio i miei prototipi'', cioè ``quando devo semplicemente mostrare
qualche cosa a livello di gameplay''. Un menu delle \emph{opzioni} ``è già un discorso
molto più'' complesso, ``perché bisogna andare a vedere gli slider'': non si tratta, ``perché
tanto poi andiamo a fare qualcosa di più gameplayoso''.

\paragraph{I problemi dei presenti.} Diversi studenti hanno ancora il personaggio nel
menu.
\begin{itemize}
  \item \textbf{La camera.} Ripasso passo per passo: selezionare la camera nel livello
        (``proprio diventa arancione''), Open Level Blueprint, tasto destro, Create a
        Reference to CameraActor (``devi stare dentro il Level Blueprint''); Get Player
        Controller, Set View Target with Blend, la camera nel New View Target.
  \item \textbf{Il personaggio che nasce ancora.} ``Devi passare alla Game Mode
        default'': \textbf{World Settings}, \textbf{GameMode Override}. Uno studente non
        trova la finestra delle World Settings (si apre dal menu \emph{Window} dell'editor,
        se è chiusa). Poi il Player Start fuori dalla camera.
  \item \textbf{L'errore su Add to Viewport.} ``Mi dà un problema proprio su Add to
        Viewport'': ``perché il target di base è Self, quindi non so se devo cambiare il
        target''. Sì: il Target deve essere il \textbf{Return Value} del Create Widget.
        ``Probabilmente devi dare return value'' al target.
  \item \textbf{Il mouse che non si vede.} Va impostato nell'Event Construct del widget:
        Get Player Controller, Set Show Mouse Cursor, Set Input Mode UI Only. Lo studente
        aveva cancellato l'Event Construct (``basta che fai tasto destro'' e lo si ricrea)
        e non aveva messo la \textbf{spunta} nel Set Show Mouse Cursor (``ah, però non
        avevo messo il tick''): il nodo, senza spunta, imposta il cursore a \emph{falso}.
        ``Funziona? Sì.''
\end{itemize}
Chiuso il menu, si torna alla ThirdPersonMap.

% ---------------------------------------------------------------------
\section{Interfacce e Blueprint Component}
\label{ue12:sec:concetti}

``Andiamo sul nostro amico fidato'': il docente apre il suo bot (``ah no, questo è Boto:
dobbiamo andare su Boto2''), cioè BP\_Bot2. ``Oggi parleremo delle interfacce, ragazzi.''

\paragraph{L'interfaccia, un intermediario.} ``Immaginatevi che un'interfaccia è un
intermediario fra due o più Blueprint.'' Il docente mostra quella già pronta nel suo
progetto, nella cartella dei componenti dei Blueprint: un tipo particolare di Blueprint
che contiene solo tre \textbf{funzioni}, \emph{Interaction Start}, \emph{Interaction End}
e \emph{Interaction}. L'immagine: un \textbf{volo di linea}, ``un'interconnessione tra due
o più poli'', che ``trasporta molto spesso le stesse cose, passeggeri di un tipo,
passeggeri di un altro''. ``Quindi è un qualcosa che va a trasportare varie funzioni da una
parte all'altra.'' Nel caso di oggi l'interfaccia comunica tre messaggi: \emph{sto
iniziando un'interazione} (ti sto guardando), \emph{sto finendo un'interazione} (non ti
guardo più), \emph{sto interagendo} (ho premuto il tasto).

\paragraph{Il componente, una funzionalità isolata.} Il lavoro di \emph{vedere se c'è
qualcosa di interagibile} non lo fa il personaggio, ma un \textbf{Blueprint Component}:
``dei componenti che noi possiamo aggiungere all'interno dei nostri actor per dare
funzionalità in modo isolato''. Per esempio: ``questo actor può tracciare degli
interactable''; ``questo actor ha la vita''. La vita, fatta la volta scorsa \emph{dentro}
il personaggio, si potrebbe spostare in un componente, ``perché magari abbiamo anche i
nemici che hanno la vita'', o ``degli oggetti che si possono rompere''. Lo scopo è
``isolare sempre più possibile le funzionalità, e quindi essere anche modulari''. Infatti
il BP\_Bot2 del docente, a parte un input nuovo, ``non ha niente di nuovo'': si gira verso
l'oggetto e suona un Montage, ma ``non ha niente per tracciare'': il trace è tutto nel
componente.

\paragraph{Actor Component o Scene Component.} Si crea una \textbf{Blueprint Class} e, fra
le classi comuni, non un Actor ma un componente. Ce ne sono due:
\begin{itemize}
  \item l'\textbf{Actor Component} ``mette semplicemente la funzionalità'';
  \item lo \textbf{Scene Component} ``è un Actor Component che ha anche una scena'',
        cioè un \textbf{Transform}: ``non so se vi ricordate la Scene, quando l'abbiamo
        messa: il punto vuoto''. ``Almeno si può spostare nello spazio.''
\end{itemize}
Qui serve il secondo, perché il trace è ``un raggio che da quel punto parte e in un punto
finisce'', e il punto di partenza si deve poter spostare. Il docente lo mostra nel suo bot:
spostando il componente nel Blueprint del personaggio, in Play il raggio parte dal nuovo
punto; lo rimette ``al centro della testa, più o meno''. Il nuovo Blueprint:
\textbf{BPC\_RayInteract} (BPC sta per \emph{Blueprint Component}; il docente, che ne ha già
uno, chiama il suo BPC\_RayInteract1).

\paragraph{Il componente aperto.} ``Come potete vedere, presenta delle similitudini con
tutti i Blueprint.'' ``L'unica cosa: non ha una viewport, perché non ha una forma fisica'',
``tra virgolette, perché è una scena, ma è vuota''. ``E ricordiamoci che è un componente:
quindi sarà qualcosa che deve essere per forza attaccato a qualcos'altro.''

\paragraph{Che cos'è una funzione.} Prima di scrivere la logica, una domanda: chi ha già
sentito parlare di \textbf{funzioni} in informatica? Risponde \textbf{Federico}, che ha
studiato da programmatore: ``la funzione è un insieme di istruzioni: chiami la funzione e
fai quelle istruzioni''. ``Esattamente.'' Il docente indica una catena di nodi nel bot: ``gli
eventi non possono essere messi come funzioni, ma questo qua può essere una funzione''. È
``un modo per ordinare anche il lavoro in varie tipologie di funzionalità, che possono
essere magari trasportabili da un punto all'altro'': ``sempre questo senso di modularità''.
(``Grande Federico, sempre miglior assistente.'')

\begin{figure}[htbp]
\centering
\begin{tikzpicture}[font=\footnotesize,
  ue12box/.style={draw, thick, rounded corners=5pt, align=center, inner sep=5pt},
  ue12msg/.style={-{Stealth[length=2.5mm]}, very thick, violet!70!black},
  ue12own/.style={-{Stealth[length=2.2mm]}, thick, dashed, black!70}]
  \node[ue12box, fill=gray!8, minimum width=4.4cm, minimum height=4.6cm] (bot) at (0,-0.3) {};
  \node[anchor=north, font=\bfseries] at (bot.north) {BP\_Bot2};
  \node[ue12box, fill=blue!10, text width=3.6cm] (cam) at (0,0.8) {Camera Component};
  \node[ue12box, fill=orange!18, text width=3.6cm] (comp) at (0,-1.4)
    {\textbf{BPC\_RayInteract}\\{\scriptsize Scene Component: ogni frame una Sphere Trace
    lungo la camera; variabile Interactable}};
  \draw[ue12own] (comp.north) -- node[right, font=\scriptsize, text width=1.9cm,
    align=left]{Get Owner,\\Get Component by Class} (cam.south);
  \node[ue12box, fill=violet!10, text width=3.2cm] (bpi) at (6.6,-0.3)
    {\textbf{BPI\_InteractionRay}\\{\scriptsize Blueprint Interface}\\[2pt]{\scriptsize
    Interaction Start\\Interaction End\\Interaction}};
  \node[ue12box, fill=green!10, text width=3.3cm] (obj1) at (11.6,1.0)
    {\textbf{BP\_InteractableTest}\\{\scriptsize implementa la BPI: Print String, widget
    ``E'', luce}};
  \node[ue12box, fill=green!5, text width=3.3cm] (obj2) at (11.6,-1.6) {{\scriptsize un
    altro oggetto}\\{\scriptsize (bottone, ascensore\dots) con la stessa BPI e un'altra
    risposta}};
  \draw[ue12msg] (bot.east |- cam) -- node[above, font=\scriptsize, text width=2.4cm,
    align=center]{Interaction\\(tasto E)} (bpi.west |- cam);
  \draw[ue12msg] (comp.east) -- node[below, font=\scriptsize, text width=2.4cm,
    align=center]{Start / End\\(guardare, smettere)} (bpi.west |- comp);
  \draw[ue12msg] (bpi.east) -- (obj1.west);
  \draw[ue12msg] (bpi.east) -- (obj2.west);
\end{tikzpicture}
\caption{I ruoli del sistema di interazione. Il componente, attaccato al personaggio, trova
la camera del suo proprietario, traccia ciò che il giocatore guarda e invia i messaggi di
inizio e fine; il personaggio invia il messaggio Interaction quando si preme E. I messaggi
viaggiano attraverso l'interfaccia: chi li riceve decide che cosa fare, e un actor che non
la implementa li ignora. Schema ricostruito dalla lezione.}
\label{ue12:fig:ruoli}
\end{figure}

\paragraph{Approfondimento: i componenti nel libro.} Romero e Sewell dedicano ai componenti
la seconda metà del capitolo 18 (\emph{Creating Blueprint Libraries and Components}).
``Quando creiamo un Actor Blueprint aggiungiamo spesso componenti con una funzionalità
incapsulata, pronta all'uso, come Projectile Movement, Static Mesh e Collision''; ``possiamo
anche creare i nostri Actor Component con i Blueprint''. Nel pannello \emph{Pick Parent
Class} le due classi sono fra quelle \emph{comuni}: ``lo Scene Component è una classe figlia
dell'Actor Component che ha la struttura Transform (location, rotation e scale). Grazie al
Transform, uno Scene Component può essere attaccato a un altro Scene Component.'' E, come a
lezione: ``quando si scrive la logica di un componente nell'Event Graph, si può ottenere la
referenza all'actor che usa il componente con il nodo \textbf{Get Owner}''. I due esempi del
libro mostrano bene la differenza: \emph{BP\_ExpLevelComp} è un Actor Component che gestisce
punti esperienza e livelli (non ha bisogno di una posizione), mentre \emph{BP\_CircularMovComp}
è uno Scene Component che ruota attorno all'actor e a cui si attacca una Static Mesh, per
simulare uno scudo rotante: ``lo Scene Component è usato come posizione di riferimento per
altri componenti''. Il componente della lezione è dello stesso tipo: la sua posizione è il
punto di partenza del raggio.

% ---------------------------------------------------------------------
\section{Il componente BPC\_RayInteract: la Sphere Trace}
\label{ue12:sec:componente}

\paragraph{La funzione LookForInteractions.} Nel componente si crea una nuova funzione
(``eccolo che vi si crea questo pannellino violaceo''), \textbf{LookForInteractions}.
Dentro, una \textbf{Sphere Trace By Channel}. ``Seguitemi per bene.''
\begin{itemize}
  \item \textbf{Radius}: \textbf{5}, ``perché se no non vi apparirà mai''. Con raggio zero
        la Sphere Trace degenera in un raggio sottilissimo (e il debug quasi non si vede);
        5 unità bastano a renderla visibile e un po' più tollerante.
  \item \textbf{Start}: \textbf{Get World Location} del componente stesso, ``la posizione
        nel mondo del nostro Scene Component: dovrà partire da qui''.
  \item \textbf{End}: da costruire, ed è la parte con il trucco.
\end{itemize}

\paragraph{Prima la camera: Get Owner e Get Component by Class.} Il raggio deve seguire lo
sguardo del giocatore, cioè la \emph{rotazione della camera}. Il componente, però, non sa
nulla della camera del personaggio: deve trovarla. Nell'\textbf{Event Graph} del
componente, all'\textbf{Event BeginPlay}:
\begin{enumerate}
  \item \textbf{Get Owner}: ``la referenza del proprietario di questo Scene Component'';
  \item dal proprietario, \textbf{Get Component by Class} (``non \emph{components},
        \emph{component}, perché se no vi dà un array''), con \textbf{Component Class} =
        \textbf{Camera Component}: ``sarà il nostro punto principale da cui andremo a fare il
        nostro trace'';
  \item il risultato si promuove a variabile, \textbf{CameraComponent}, e si collega tutto
        al Begin Play.
\end{enumerate}
``Perdonate, sono 30.000 componenti: un componente su componente su componente.'' In
sintesi: ``in Begin Play io prendo il mio proprietario, prendo la sua camera, mi salvo la
referenza alla sua camera, e poi la useremo all'interno del nostro componente''.

\paragraph{La camera, se c'è.} Il componente deve funzionare anche su un actor
\emph{senza} camera: ``se il nostro proprietario non ha una camera, allora prendi te stesso
come pivot''. Quindi un \textbf{Select} fra due oggetti:
\begin{itemize}
  \item l'indice è un booleano, \textbf{Is Valid} (``la F verde'', la versione \emph{pura}
        della funzione) sulla variabile CameraComponent: ``se il componente camera esiste,
        quindi il nostro actor ha effettivamente una camera, e quindi questa variabile
        CameraComponent ha qualcosa dentro'';
  \item su \textbf{True}, CameraComponent: ``usami la camera come punto di riferimento'';
  \item su \textbf{False}, \textbf{Self}: ``se no, usi te stesso''.
\end{itemize}
Il Select era già comparso con la torretta, per scegliere fra le due bocche di fuoco: lì
con un \emph{Random Bool}, qui con l'Is Valid. ``So che è un po' complesso.''

\paragraph{La ``mandragata'': il tipo del Select.} Il problema è il \emph{tipo} delle
opzioni. ``Il Select praticamente c'è un piccolo problema'': il nodo è jolly (\emph{wildcard})
e assume il tipo della \emph{prima} cosa che gli si collega. Collegando prima la camera, il
Select diventa un Select di \emph{Camera Component Object Reference}, e Self (il componente
BPC\_RayInteract) non entra più; collegando prima Self, il Select diventa un Select di
riferimenti a BPC\_RayInteract, e la camera non entra. Entrambi però \emph{sono} Scene
Component: la camera è uno Scene Component, e lo è anche il nostro componente. Il trucco
(``qui bisogna fare una mandragata, ragazzi''):
\begin{enumerate}
  \item si crea una variabile qualunque di tipo \textbf{Scene Component Object Reference}
        (il docente la chiama \emph{Scene}: ``e poi non ci servirà a niente'');
  \item se ne mette un Get nel Select: ora le opzioni sono di tipo \emph{Scene Component};
  \item si \textbf{cancella} il Get e si collegano CameraComponent e Self, che sono entrambi
        accettati.
\end{enumerate}
``Ecco fatto che abbiamo praticamente fregato Unreal.''

\paragraph{La domanda di Filippo.} ``Non ho capito perché abbiamo creato la variabile Scene
prima.'' Il docente rifà la prova davanti a tutti: un Select nuovo, il cursore sul pin, la
scritta ``Scene Component Object Reference''; poi Self in un Select vuoto, e il nodo diventa
``Self Object Reference'', ``mi farà fare solamente un select fra due self, non due Scene
Component''. ``Per fregare Unreal'' serve ``una variabile che appartiene generalmente allo
Scene Component''. Uno studente prova a trasformare la variabile in Self: ``non è che te lo
trasforma: lo devi cancellare, e poi tasto destro, scrivi \emph{self}''. Filippo insiste: e
se avessimo collegato subito la camera? Prova del docente: ``Select camera component, camera
component object reference, e io, se provo ad attaccare Self, non me lo attacca''. ``Ok,
capito.''

\paragraph{Il punto finale.} Dal risultato del Select, ``quello che abbiamo sempre fatto'':
\begin{enumerate}
  \item \textbf{Get Forward Vector} (della camera, o del componente stesso);
  \item moltiplicato per una distanza: il secondo pin del Multiply si converte in float
        (``tasto destro, mi raccomando: \emph{To Float}''), si promuove a variabile,
        \textbf{DistanceTrace}, con valore \textbf{300};
  \item il prodotto si \textbf{somma} (Add) allo Start, e il risultato va nell'\textbf{End}.
\end{enumerate}

\paragraph{Perché lo Start è il componente.} Una studentessa chiede perché lo Start sia la
World Location: ``in relazione a dove attacchiamo il componente?'' ``Esatto.'' ``Il punto
iniziale deve essere sempre il punto, il Transform, del nostro componente''; ``il punto che
deve diversificare è la camera''. Il docente lo mostra nel bot: guardando in basso il
raggio va in basso, guardando di lato va di lato, ma parte sempre dal componente, ``così
magari posso posizionarlo avanti, posizionarlo indietro, posizionarlo qui, posizionarlo
là'' (figura~\ref{ue12:fig:look}).

\paragraph{Il componente nel personaggio.} Nel BP\_Bot2 (``a bot2, in questo caso, è il
vostro character''), \textbf{Add} e il nome del componente: BPC\_RayInteract compare fra i
componenti del personaggio. Per vederlo si imposta il \textbf{Draw Debug Type} della Sphere
Trace (``lo metto a \emph{For One Frame}''). In Play, però, non succede niente: ``ah,
certo, ecco perché: non facciamo LookForInteractions''. La funzione va chiamata
nell'\textbf{Event Tick} del componente (``ah, nel tick, infatti''), così il controllo
avviene a ogni frame. Con il debug persistente (\emph{Persistent}) il raggio lascia una scia
``gigantesca''; con \emph{For Duration} resta qualche istante; con \emph{For One Frame}
``rimane un po' più carino'', e si vede solo il raggio attuale. (La trascrizione, qui, è
piena di prove e di parole storpiate: ``la duration vado a activate, troppo corto'', ``Lugano'',
``Adobe''; anche una confusione momentanea fra il bot e il bot2 del docente.)

\paragraph{Approfondimento: Sphere Trace, Tick e costo.} Una \textbf{Sphere Trace} è un
\emph{line trace} ``spesso'': fa scorrere una sfera del raggio indicato dallo Start
all'End e restituisce il primo oggetto che la sfera tocca nel canale scelto. Rispetto a un
raggio sottile perdona di più la mira (con 5 unità di raggio, circa 5 cm, conta anche
sfiorare il bordo di un oggetto), ed è il motivo per cui molti sistemi di interazione la
preferiscono. Farla nel \textbf{Tick}, una volta per frame, costa poco: una singola trace di
300 unità è un'operazione leggera, e il sistema la esegue solo per il personaggio del
giocatore. Se gli actor con il componente fossero molti (per esempio dei nemici), conviene
ridurre la frequenza con un timer (come quelli della stamina, a 0,1 secondi) o disattivare
il Tick del componente quando non serve (\emph{Set Component Tick Enabled}). Una Sphere
Trace può stare senza problemi dentro una funzione, perché non ha uno stato interno da
ricordare fra una chiamata e l'altra; il Do Once e il Flip Flop, invece, sì (si veda più
avanti). Funzionamento generale dell'engine.

% ---------------------------------------------------------------------
\section{Il Trace Channel e la Blueprint Interface}
\label{ue12:sec:bpi}

\paragraph{Un canale per le interazioni.} ``Ha colpito qualcosa? Sì, no.'' Prima di
leggere il risultato serve un canale dedicato, ``che ci permette di identificare gli
oggetti che vogliamo che siano interagibili'': ``se no, praticamente, interagirà con
tutto''. ``Come abbiamo visto anche nella nostra torretta'': \textbf{Project Settings},
\emph{Collision}, \textbf{New Trace Channel}, nome \textbf{InteractionTrace}, risposta di
default \textbf{Ignore}. Il docente, che ce l'ha già, non lo rifà. Nella Sphere Trace, il
\textbf{Trace Channel} diventa InteractionTrace: ``mi raccomando, se non vi appare, fate
compila, perché ha bisogno un attimo di compilare il nuovo Trace Channel''.

\paragraph{Che cosa abbiamo colpito.} Dopo la trace, un \textbf{Branch} sul suo Return
Value (``ha colpito qualcosa?'') e un \textbf{Break Hit Result} sull'Out Hit, da cui, ``vi
ricordate, no?'', escono l'actor colpito, il Physical Material ``e tutto''. Il piano:
``useremo l'interfaccia, la Blueprint Interface, per andare a comunicare con questo actor'':
``ha questo porto in cui può accogliere questa tipologia di interazione? Sì? Benissimo,
allora fai l'interazione''. Dall'\textbf{Hit Actor} si chiede \textbf{Does Object
Implement Interface}: ``ehi, questo actor ce l'ha l'interface?'' Ma l'interfaccia non
esiste ancora: ``ci fermiamo qua per adesso, creiamo l'interface, e ci fermiamo''. (Il
docente parla qui anche di promuovere l'Hit Actor a variabile; la variabile di appoggio
vera e propria, \emph{Interactable}, si crea dopo la pausa, sezione~\ref{ue12:sec:logica}.)

\paragraph{La Blueprint Interface.} Nella cartella \emph{Blueprints}, \emph{Components},
tasto destro, \textbf{Blueprint}, ``non Blueprint Class questo, ma Blueprint questo'': il
sottomenu offre \emph{Blueprint Class}, \emph{Blueprint Function Library} e
\textbf{Blueprint Interface}. Il nome: \textbf{BPI\_InteractionRay}. Aperta:
\begin{itemize}
  \item il grafo è \textbf{in sola lettura}: ``non possiamo far niente'', ``cannot create
        new nodes in read only graph''. Un'interfaccia contiene solo \emph{dichiarazioni};
  \item si possono aggiungere \textbf{funzioni}: la prima c'è già e si rinomina (con F2),
        \textbf{Interaction Start}; con \emph{Add Function} le altre due,
        \textbf{Interaction End} e \textbf{Interaction}.
\end{itemize}
``Abbiamo messo che in questo cargo, questo trasporto, ci stanno Interaction Start,
Interaction End e Interaction: quindi queste funzioni possono essere passate tra di loro.''

\paragraph{Il collegamento.} Nel componente, il pin \textbf{Interface} di Does Object
Implement Interface diventa \mbox{BPI\_}\allowbreak InteractionRay. ``Se quel qualcosa che ho colpito ha questa
interfaccia, allora andremo a fare tutto il procedimento di interazione.'' ``Ci sono domande?
Mi sembra di capire no.'' Sono le 17:26: pausa fino alle 17:55, ``diciotto''.

\paragraph{Approfondimento: la Blueprint Interface nel libro.} Romero e Sewell presentano le
interfacce nel capitolo 16 (\emph{Introduction to VR Development}), nel paragrafo
\emph{Blueprint Communication using interfaces}: ``una Blueprint Interface è un tipo speciale
di Blueprint che contiene solo nomi di funzioni e parametri. Si usa per permettere la
comunicazione fra tipi diversi di Blueprint.'' Si crea dal pulsante \emph{Add} del Content
Browser, sottomenu \emph{Blueprints}, \emph{Blueprint Interface}. L'esempio è il template VR:
l'interfaccia \emph{VRInteraction BPI} è aggiunta alla pistola dalle \emph{Class Settings},
categoria \emph{Interfaces}, e la pistola implementa la funzione \emph{Trigger Pressed}:
``siccome la funzione non ha parametri di output, è implementata come un \emph{evento}''. Il
Pawn VR chiama Trigger Pressed sull'actor che tiene in mano, ``e se l'actor ha implementato
l'interfaccia VRInteraction BPI, Trigger Pressed viene eseguito. Non succede niente se l'actor
non l'ha implementata.'' È esattamente lo schema della lezione: il personaggio non sa
\emph{che cosa} sta guardando (un cubo, un bottone, un ascensore) e non deve fare un
\emph{cast} a nessuna classe; manda il messaggio, e ciascun oggetto risponde a modo suo.
Questo è il vantaggio delle interfacce rispetto al cast e all'ereditarietà viste con
BP\_Activable: classi che non hanno niente in comune (un actor, un personaggio, un oggetto
fisico) possono condividere la stessa ``porta''.

% ---------------------------------------------------------------------
\section{La pausa: la stamina di Tommaso}
\label{ue12:sec:pausa}

\paragraph{Il sintomo.} In pausa, come promesso, il problema rimasto aperto la volta
precedente: \textbf{Tommaso} tiene premuto Shift, la barra della stamina non scende, e
``appena mollo, si vede''. Il docente ripercorre il codice condiviso: l'evento StopRunning,
il Clear and Invalidate, il Set Timer by Event e la variabile TimeStamina (``quanto sta?
0,92''; nella trascrizione il valore è questo, ma va preso con cautela: nella lezione
precedente il valore indicato era 0,05), RegainStamina, il Clamp con CurrentStamina e
MaxStamina, l'InitializeStaminaBar nel Begin Play. ``Sembra tutto a posto.''

\paragraph{Il debug.} Il metodo è quello della volta scorsa: un \textbf{Print String} con la
CurrentStamina (``attaccalo al viola'') dopo l'evento di aggiornamento. I valori scendono
correttamente (``500, 500, 300'' dice la trascrizione, forse leggendo la velocità), poi
risalgono. Ma a schermo ci sono ``troppi Print String'': si cercano e si tolgono quelli
vecchi, uno per volta (quello della torretta, uno della piattaforma, ``forse i mulini?'',
``forse il collectable''; ``nessuno è perfetto''). Il comportamento: la barra
\emph{ricarica} visibilmente, ma \emph{non scende}.

\paragraph{La causa.} ``Ah, ecco perché: non chiami l'evento Call Update Stamina.'' Nella
catena di \textbf{StaminaDeduction}, dopo il Set di CurrentStamina, mancava la
\textbf{Call E\_UpdateStamina}: il valore calava, ma nessuno lo comunicava alla barra. La
Call c'era solo in RegainStamina, che parte al rilascio del tasto; per questo la barra si
aggiornava ``solo quando mollo''. Il rimedio: la Call anche in StaminaDeduction, con
CurrentStamina e MaxStamina collegati. ``Perfetto. Ottimo.''

\paragraph{Il rumore.} Il resto della pausa, nella trascrizione, è il blocco di testo
ripetuto e il frammento di video già segnalati nell'introduzione. Alla ripresa, uno studente
chiede aiuto per BPI\_InteractionRay, che dopo un riavvio di Unreal si apre ``male'': il
grafo in sola lettura non si vede. Il docente prova i menu \emph{Window} (Bookmarks, Find
Results, Compiler Results\dots), carica per sbaglio il layout di default dell'editor
(``mi sono caricato il bellissimo layout di Unreal''), e alla fine trova la via più
semplice: ``fai doppio clic sulla funzione'' nel pannello \emph{My Blueprint}, e il suo
grafo si apre. ``Tanto non ti serve'': nell'interfaccia non si scrive nulla.

% ---------------------------------------------------------------------
\section{La logica: la variabile Interactable}
\label{ue12:sec:logica}

``Vai, allora finiamo questo bellissimo sistema di interazioni.'' Si torna in
BPC\_RayInteract, nella funzione LookForInteractions. Il ramo di partenza è quello in cui
la trace ha colpito qualcosa e l'actor colpito implementa BPI\_InteractionRay.

\paragraph{La variabile d'appoggio.} Serve una variabile di tipo \textbf{Actor}
(\emph{Object Reference}), \textbf{Interactable}: ``sarà la nostra variabile d'appoggio'',
cioè l'oggetto interagibile che il giocatore sta guardando \emph{in questo momento}, o
niente. Il suo scopo è evitare di ripetere il lavoro a ogni frame: ``se l'actor che io ho
intercettato con il mio trace è lo stesso di prima, allora non voglio che mi rifai alcune
sezioni del codice''; ``se è lo stesso, non lo svuoto''.

\paragraph{È un oggetto nuovo?} Un secondo \textbf{Branch}, con la condizione \textbf{Hit
Actor} \textbf{diverso da} Interactable: il nodo \emph{Not Equal} (``diverso si fa con il
punto esclamativo e l'uguale''). ``Il punto esclamativo in informatica vuol dire che questo è
un condizionale che ci dice: guarda, questo è diverso da questo, e ci restituisce un
booleano. Se è diverso ci dà True, se sono uguali ci dà False.'' Sul False non si fa nulla:
è lo stesso oggetto di prima.

\paragraph{I messaggi dell'interfaccia.} ``Adesso, importante'': scrivendo
\emph{interaction} nel menu dei nodi, compaiono già le funzioni di BPI\_InteractionRay,
``senza fare niente'': ``ho già la possibilità di poter comunicare e messaggiare con
l'interfaccia''. Sono i nodi \textbf{(Message)}, con l'icona della busta, che si possono
chiamare su qualunque actor. Servono \textbf{Interaction End} e \textbf{Interaction Start}.

\paragraph{Il ramo True: chiudere il vecchio, aprire il nuovo.} Sul True del diverso:
\begin{enumerate}
  \item un \textbf{Validated Get} di Interactable (si mette il Get, tasto destro,
        \emph{Convert to Validated Get}): ``questa variabile Interactable è valida? C'è
        qualcosa dentro?'';
  \item \textbf{Is Valid}: c'era già un oggetto, quindi prima si chiude la vecchia
        interazione, \textbf{Interaction End} con target il \emph{vecchio} Interactable;
        poi \textbf{Set Interactable} = Hit Actor; poi \textbf{Interaction Start} con
        target il \emph{nuovo} Interactable;
  \item \textbf{Is Not Valid}: ``la variabile è vuota'', ``non è mai iniziata
        un'interazione precedente''; si salta l'End e si fa direttamente il Set
        Interactable = Hit Actor e l'\textbf{Interaction Start}.
\end{enumerate}

\paragraph{Il ramo False: non guardo niente.} Se la trace non colpisce nulla, \emph{oppure}
colpisce un actor senza l'interfaccia (``il muro, ad esempio''), il giocatore ha distolto lo
sguardo: ``allora ti finisco l'interazione''. Entrambi i False (quello del Branch sul Return
Value e quello di Does Object Implement Interface) portano a un altro Validated Get di
Interactable:
\begin{itemize}
  \item \textbf{Is Valid}: \textbf{Interaction End} sul vecchio Interactable, poi
        \textbf{Set Interactable} senza collegare nulla al pin, cioè lo si
        \textbf{svuota};
  \item \textbf{Is Not Valid}: non c'era nessuna interazione da chiudere; per sicurezza
        lo si svuota comunque (``semplicemente dirgli: guarda, svuotati'').
\end{itemize}
``Quindi mettere il Set così è come mettere uguale a null?'' ``È nulla, proprio. La
annulliamo: la mettiamo vuota, e quindi diventerà non valida.'' Il caso del primo ramo,
il passaggio diretto da un oggetto all'altro, serve quando due interagibili sono vicini:
``lui finisce con quello precedente e inizia con il nuovo''; di solito, invece, gli oggetti
sono distanti, la variabile si svuota e poi si riempie di nuovo
(figure~\ref{ue12:fig:look} e~\ref{ue12:fig:stati}).

\paragraph{Il problema di tipo.} Uno studente non riesce a collegare l'Hit Actor al Set
Interactable: ``mi dice che non sono compatibili, perché uno è un riferimento oggetto e
l'altro è un componente di scena''. La variabile Interactable era stata creata con il tipo
sbagliato (ed era cambiata per un clic involontario in \emph{Geometry}, ``che non ho
cliccato''). Rimedio: nei Details della variabile, il tipo \textbf{Actor}, \textbf{Object
Reference} (``non quell'Actor: tipo oggetto, riferimento oggetto''), confermare la
modifica del tipo, poi staccare e riattaccare i fili rossi dei nodi che usavano la
variabile, ``era semplicemente sparito''.

\paragraph{Il Tick.} ``Importante, ragazzi: far sì che questa funzione venga fatta su ogni
frame.'' Nell'Event Graph del componente, \textbf{Event Tick} $\rightarrow$
LookForInteractions (se non lo si è già fatto): ``a ogni frame farà il controllo: c'ho
qualcosa di interagibile davanti a me, dove sto puntando la camera? Sì o no?''.

\begin{figure}[htbp]
\centering
\begin{tikzpicture}[font=\footnotesize,
  ue12node/.style={draw, thick, rounded corners=4pt, align=center, fill=gray!8,
                  minimum height=0.85cm, text width=2.4cm},
  ue12dec/.style={draw, thick, diamond, aspect=2.2, align=center, fill=yellow!15,
                  inner sep=1pt, text width=2.0cm, font=\scriptsize},
  ue12msgn/.style={ue12node, fill=violet!12},
  ue12exec/.style={-{Stealth[length=2.5mm]}, very thick, black!80}]
  \node[ue12node, fill=red!12, text width=3.4cm] (st) at (0,0) {\textbf{Sphere Trace By
    Channel}\\{\scriptsize Start: posizione del componente; End: Start + Forward (camera o
    self) $\times$ 300; Radius 5; InteractionTrace}};
  \node[ue12dec] (hit) at (0,-2.2) {colpito qualcosa?};
  \node[ue12dec] (impl) at (0,-4.4) {Hit Actor implementa la BPI?};
  \node[ue12dec] (diff) at (0,-6.6) {Hit Actor $\neq$ Interactable?};
  \node[ue12dec] (val1) at (0,-8.9) {Interactable valido?};
  \node[ue12msgn] (end1) at (-3.6,-10.8) {\textbf{Interaction End}\\{\scriptsize sul vecchio}};
  \node[ue12node] (set1) at (0,-12.4) {\textbf{Set Interactable}\\{\scriptsize = Hit Actor}};
  \node[ue12msgn] (start) at (0,-14.0) {\textbf{Interaction Start}\\{\scriptsize sul nuovo}};
  \node[ue12node, fill=gray!3, text width=2.0cm] (none) at (3.7,-6.6) {niente\\{\scriptsize
    (stesso oggetto)}};
  \node[ue12dec] (val2) at (6.4,-3.3) {Interactable valido?};
  \node[ue12msgn] (end2) at (6.4,-5.6) {\textbf{Interaction End}\\{\scriptsize sul vecchio}};
  \node[ue12node] (clr) at (6.4,-7.6) {\textbf{Set Interactable}\\{\scriptsize vuoto (null)}};
  \draw[ue12exec] (st) -- (hit);
  \draw[ue12exec] (hit) -- node[left, font=\scriptsize]{True} (impl);
  \draw[ue12exec] (impl) -- node[left, font=\scriptsize]{True} (diff);
  \draw[ue12exec] (diff) -- node[left, font=\scriptsize]{True} (val1);
  \draw[ue12exec] (diff) -- node[above, font=\scriptsize]{False} (none);
  \draw[ue12exec] (val1.west) -| node[pos=0.25, above, font=\scriptsize]{Is Valid} (end1.north);
  \draw[ue12exec] (end1.south) |- (set1.west);
  \draw[ue12exec] (val1) -- node[right, font=\scriptsize]{Is Not Valid} (set1);
  \draw[ue12exec] (set1) -- (start);
  \draw[ue12exec] (hit.east) -| node[pos=0.2, above, font=\scriptsize]{False} (val2.north);
  \draw[ue12exec] (impl.east) -- node[above, font=\scriptsize]{False} ++(2.5,0) |-
    (val2.west);
  \draw[ue12exec] (val2) -- node[right, font=\scriptsize]{Is Valid} (end2);
  \draw[ue12exec] (end2) -- (clr);
  \draw[ue12exec] (val2.east) -- ++(0.6,0) |- node[pos=0.3, right, font=\scriptsize, text
    width=1.0cm, align=left]{Is Not Valid} (clr.east);
\end{tikzpicture}
\caption{La funzione LookForInteractions del componente, chiamata nell'Event Tick. Se il
giocatore guarda un nuovo interagibile, si chiude l'interazione con il vecchio (se c'era)
e si apre con il nuovo; se guarda un oggetto non interagibile, o niente, si chiude
l'interazione in corso e si svuota la variabile. I nodi viola sono messaggi
dell'interfaccia. Ricostruzione dalla lezione.}
\label{ue12:fig:look}
\end{figure}

\begin{figure}[htbp]
\centering
\begin{tikzpicture}[font=\footnotesize,
  ue12state/.style={draw, very thick, rounded corners=10pt, align=center,
                    minimum height=1.3cm, text width=3.4cm},
  ue12tr/.style={-{Stealth[length=2.5mm]}, very thick}]
  \node[ue12state, fill=gray!10] (vuoto) at (0,0) {\textbf{Interactable vuoto}\\{\scriptsize
    non guardo niente di interagibile}};
  \node[ue12state, fill=green!12] (a) at (7.0,0) {\textbf{Interactable = A}\\{\scriptsize
    widget ``E'' di A visibile}};
  \node[ue12state, fill=green!6] (b) at (7.0,-3.4) {\textbf{Interactable = B}\\{\scriptsize
    widget ``E'' di B visibile}};
  \draw[ue12tr] ([yshift=3mm]vuoto.east) -- node[above, font=\scriptsize]{guardo A: Start(A)}
    ([yshift=3mm]a.west);
  \draw[ue12tr] ([yshift=-3mm]a.west) -- node[below, font=\scriptsize, text width=3.2cm,
    align=center]{guardo il muro o niente: End(A), svuota} ([yshift=-3mm]vuoto.east);
  \draw[ue12tr] (a.south) -- node[right, font=\scriptsize, text width=3.4cm]{passo
    direttamente a B: End(A), Set, Start(B)} (b.north);
  \draw[ue12tr] (b.west) -| node[pos=0.3, below, font=\scriptsize]{guardo il muro: End(B),
    svuota} (vuoto.south);
  \draw[ue12tr] (a.east) .. controls +(1.6,0.9) and +(1.6,-0.9) .. node[right,
    font=\scriptsize, text width=1.9cm]{guardo ancora A: nulla} (a.east);
\end{tikzpicture}
\caption{La variabile Interactable come stato. Ogni Interaction Start è sempre seguito, prima
o poi, dal suo Interaction End: l'oggetto può contare sulla coppia per mostrare e nascondere
il suggerimento. Schema ricostruito.}
\label{ue12:fig:stati}
\end{figure}

\paragraph{Approfondimento: messaggi, Does Implement Interface e Validated Get.} Tre
precisazioni generali sui nodi usati. Primo: un nodo \emph{(Message)} di un'interfaccia si
può chiamare su qualunque riferimento a oggetto; se l'oggetto non implementa l'interfaccia,
la chiamata semplicemente non fa nulla, senza errori (la frase del libro citata sopra: ``non
succede niente''). Il controllo \textbf{Does Object Implement Interface}, quindi, non serve
a evitare un errore, ma a \emph{distinguere} gli interagibili dagli altri oggetti, così da
chiudere l'interazione quando si guarda un muro. Secondo: il confronto \emph{Not Equal} fra
due riferimenti confronta l'\emph{identità} degli oggetti (è lo stesso actor?), non la loro
classe; due cubi identici sono comunque diversi. Terzo: il \textbf{Validated Get} è un Get
con l'\emph{Is Valid} incorporato, con le due uscite \emph{Is Valid} e \emph{Is Not Valid};
``validare'' un riferimento prima di usarlo è la raccomandazione con cui il docente chiude
il corso, perché un riferimento vuoto o a un actor distrutto (per esempio un interagibile
raccolto o eliminato) produce, se usato, gli errori \emph{Accessed None} a fine Play.
Funzionamento generale dell'engine.

% ---------------------------------------------------------------------
\section{L'oggetto di prova: BP\_InteractableTest}
\label{ue12:sec:test}

\paragraph{L'actor.} ``Adesso creiamo il nostro actor interagibile.'' Nella cartella dei
Blueprint, una Blueprint Class di tipo \textbf{Actor}, ``veramente semplicissimo'':
\textbf{BP\_InteractableTest}. Due componenti: un \textbf{Cube} e un \textbf{Widget}
(``fate widget, che poi faremo''), messo sopra il cubo.

\paragraph{Implementare l'interfaccia.} ``Non si trova'', se si cerca l'interfaccia fra i
componenti: va aggiunta in \textbf{Class Settings}, nella toolbar. Nei Details compaiono le
categorie \emph{Imports} (\emph{Default Namespace}, \emph{Imported Namespaces}) e
\textbf{Interfaces}, con \emph{Inherited Interfaces} e \textbf{Implemented Interfaces}; in
quest'ultima, \textbf{Add}, BPI\_InteractionRay. ``Per vedere che l'hai aggiunto'': nel
pannello My Blueprint compare la sezione \textbf{Interfaces}, con Interaction, Interaction End
e Interaction Start.

\paragraph{Implementare le funzioni.} Trascinare Interaction Start nel grafo non funziona:
l'editor avverte che la funzione ``was not marked as Blueprint Callable'' e non può essere
piazzata così. ``Dobbiamo overrideare la loro funzione, che di base è nulla'': le funzioni
dell'interfaccia ``sono solo degli appoggi'', vuote, e il loro contenuto ``sarà da
overrideare proprio dal livello di design''. Con un doppio clic (o tasto destro,
\emph{Implement Event}) si ottengono tre eventi rossi: \textbf{Event Interaction Start},
\textbf{Event Interaction End} ed \textbf{Event Interaction}. Ciascuno con un
\textbf{Print String}:
\begin{itemize}
  \item Interaction Start: ``inizio interazione'', \textbf{Text Color} verde;
  \item Interaction End: ``fine interazione'', rosso;
  \item Interaction: ``interagendo'', giallo-arancione (``ancora non l'abbiamo usato,
        quindi non lo vedremo, è normale'').
\end{itemize}
Il colore si trova aprendo la tendina del Print String (la freccia in basso del nodo).

\paragraph{La collisione.} ``È importante però una cosa'': il canale InteractionTrace è in
Ignore di default, quindi il cubo, così com'è, non viene visto dal raggio. Nella
\textbf{Collision} del Cube, preset \emph{Custom} e \textbf{InteractionTrace} su
\textbf{Block}. ``Su Cube, tanto la mesh è quella.''

\paragraph{La prova.} Si mette il cubo in scena. In Play: ``inizio interazione, fine
interazione, inizio interazione, fine interazione''. Con la camera in terza persona colpire
il cubo è un po' difficile, ``ma ci sono delle volte in cui diventa verde''. ``Vi funziona?''
``Sì, funziona.'' ``Quindi adesso abbiamo creato il nostro trace per quanto riguarda le
interazioni.'' Il componente ``fa le veci di \emph{vado a controllarti se ci sono delle
interazioni, e gestisco tutta la cosa}''; l'interazione vera, il tasto, la fa il personaggio,
``in modo classico''.

\paragraph{Approfondimento: eventi o funzioni dell'interfaccia.} Le tre funzioni di
BPI\_InteractionRay non hanno \emph{output}: per questo, nell'actor che le implementa,
diventano \emph{eventi} nell'Event Graph (è la stessa regola citata dal libro per la Trigger
Pressed della pistola VR). Se una funzione dell'interfaccia avesse un valore di ritorno (per
esempio un \emph{GetInteractionText} che restituisce ``Apri'' o ``Accendi''), andrebbe
implementata come \emph{funzione}, con un doppio clic nella sezione Interfaces, e potrebbe
essere chiamata per leggere un dato. Un'interfaccia non ha stato né variabili: ogni classe
che la implementa decide da sé che cosa fare, e può anche lasciare vuote le funzioni che non
le servono. Funzionamento generale dell'engine.

% ---------------------------------------------------------------------
\section{L'input: IA\_Interact e il tasto E}
\label{ue12:sec:input}

Nel BP\_Bot2 serve un nuovo input. ``Come si crea un nuovo input? Con l'IMC e con l'Input
Action.''
\begin{enumerate}
  \item Nel Content Browser si va alle Input Action del progetto (cliccando su IA\_Jump,
        \emph{Show in Folder View});
  \item ``basta che fate una copia di IA\_Run'': \textbf{Duplicate}, nome
        \textbf{IA\_Interact}; ``basta che qua dentro abbia i \textbf{Triggers}
        \emph{Pressed} e \emph{Released}'' (come quelli di IA\_Run), da lasciare così;
  \item nell'\textbf{IMC} di default (\emph{IMC\_\allowbreak Default}), il più accanto alle
        \emph{Mappings} aggiunge una riga per IA\_Interact; poi, cliccando sull'icona
        della tastiera e premendo il tasto, si assegna \textbf{E}.
\end{enumerate}
Il bot2 ``eredita la stessa cosa'': usa lo stesso IMC\_Default del template. Le domande dei
presenti sono tutte sulla tendina dell'IMC (``non mi fa mettere l'azione di input'': ``è la
tendina'') e sui trigger (``il primo è premuto e l'altro qual è?'': ``premuto e rilascio va
benissimo così, lasciato così''). Nel personaggio, cercando \emph{IA\_Interact}, compare il
nuovo evento d'input.

\paragraph{Approfondimento: Pressed e Released nello stesso trigger.} Un'Input Action
digitale senza trigger scatta in \emph{Triggered} a ogni frame finché il tasto è premuto; con
il trigger \emph{Pressed} scatta una sola volta alla pressione, con \emph{Released} una
volta al rilascio. Con \emph{entrambi}, come in IA\_Run e quindi in IA\_Interact, l'uscita
\emph{Triggered} scatta due volte per ogni pressione: quando si preme e quando si rilascia.
Nella lezione questo non dà problemi, perché il grafo del personaggio passa da un \emph{Do
Once} che si resetta solo alla fine del Montage: il secondo impulso, al rilascio, arriva
mentre l'animazione è in corso e viene ignorato. Ma se il Montage fosse molto breve, o il
giocatore tenesse premuto a lungo, il rilascio potrebbe far partire una seconda
interazione: per un'interazione ``a pressione'' basta il solo trigger Pressed, o l'uscita
\emph{Started}. La trascrizione non dice quale uscita dell'evento IA\_Interact usi il docente.
Funzionamento generale dell'Enhanced Input.

% ---------------------------------------------------------------------
\section{L'interazione nel personaggio}
\label{ue12:sec:player}

``Allora adesso partiamo a fare quello che ci servirà per indirizzare il nostro player,
cioè farlo ruotare verso il nostro oggetto, e allo stesso modo far sì che l'interazione
parta.''

\paragraph{Un Do Once.} ``Vi dico che è un'azione che deve essere fatta su un \textbf{Do
Once}'': l'interazione è una sequenza (girarsi, far partire il Montage, aspettarne la fine)
e solo ``a fine del montaggio può reiniziare una prossima interazione''. Quindi IA\_Interact
$\rightarrow$ Do Once.

\paragraph{Chi stiamo guardando.} Dal componente \textbf{BPC\_RayInteract} (trascinato dal
pannello dei componenti nel grafo: ``prendilo e trascinalo'') si prende la sua variabile
\textbf{Interactable} e la si converte in \textbf{Validated Get}: ``questo Interactable è
presente? Ci sta o non ci sta?''. Il personaggio non rifà la trace: legge il risultato del
componente.
\begin{itemize}
  \item \textbf{Is Valid}: il messaggio \textbf{Interaction} di BPI\_InteractionRay
        (``tira fuori e fai \emph{interaction}, cioè Interaction Message''), con target
        Interactable: ``e questa è proprio la base: così interagiamo''.
  \item \textbf{Is Not Valid}: \textbf{Reset} del Do Once (aggiunto dopo, quando il docente
        si accorge che senza il reset, premendo E senza guardare nulla, il Do Once resta
        chiuso per sempre: ``se non è valido, resetta il Do Once comunque, in ogni caso'').
\end{itemize}

\paragraph{Girarsi verso l'oggetto.} Dopo il messaggio, \textbf{Set Actor Rotation} con
target \emph{Self} (``siamo già nel player''). Si splitta il pin \emph{New Rotation}, e si
splitta anche l'uscita di un \textbf{Get Actor Rotation}: ``l'unica asse che ci interessa è
la Z'', perché il personaggio deve girarsi solo attorno all'asse verticale (``se
l'interactable sta qua, deve girare così''). Quindi:
\begin{itemize}
  \item X (Roll) e Y (Pitch) restano quelle attuali, dal Get Actor Rotation;
  \item Z (Yaw) viene da un \textbf{Find Look at Rotation}, ``quello che abbiamo fatto
        l'altra volta per la torretta'', con \textbf{Start} = \textbf{Get Actor Location}
        del personaggio (il docente dice prima ``rotation'', poi si corregge: ``location,
        ragazzi, scusatemi'') e \textbf{Target} = Get Actor Location dell'Interactable. Il
        pin di uscita non si lascia splittare (``non me lo splitta''), quindi un
        \textbf{Break Rotator}, e il suo Z va nella Z della nuova rotazione.
\end{itemize}

\paragraph{La prova provvisoria.} Prima del Montage, una prova: la scritta ``interagendo''
compare, ma il personaggio non si gira; e dopo la prima volta non succede più nulla. ``Ah,
giustamente, ragazzi'': il Do Once non si resetta mai. Il docente collega temporaneamente
il Reset alla fine della catena (``lo faccio solamente per farvelo vedere''): ``mi giro di
qua, e si gira verso di lui''. Poi lo stacca, ``perché ci servirà più avanti''.

\paragraph{Il Montage.} Nel Content Browser si cerca \emph{interact}: fra le animazioni del
bot c'è \textbf{Interact Press Button}. ``Vi ricordate come si fa il montaggio?'' Tasto
destro sull'animazione, \textbf{Create}, \textbf{Create AnimMontage}: si crea un asset che
``è sempre la stessa animazione, ma in questo caso è un montaggio''. Nel personaggio,
\textbf{Play Montage}: nel pin \textbf{In Skeletal Mesh Component} la \textbf{Mesh} del
personaggio (``basta che prendi la Mesh da qui e la metti così''), nel \textbf{Montage to
Play} il nuovo Montage (\emph{Interact Press Button Montage}).

\paragraph{La variabile Interacting.} ``Mettiamo, giusto perché in futuro ci possa
servire, una variabile'' booleana, \textbf{Interacting}: ``dobbiamo verificare se stiamo
interagendo, non stiamo interagendo, eccetera. Sempre buono mettere una variabile che dice:
guarda, io sto interagendo''. Prima del Play Montage, Set Interacting a \textbf{vero}; sull'
\textbf{On Completed}, ``quindi è finita l'animazione'', a \textbf{falso}.

\paragraph{Fermare il giocatore.} Durante l'interazione il giocatore non deve muoversi:
\begin{itemize}
  \item prima del Montage, \textbf{Disable Input}, con il \textbf{Get Player Controller}
        nel pin \emph{Player Controller}: ``chi dobbiamo levare? Gli input al controller del
        nostro player'';
  \item lo \textbf{Stop Movement Immediately}, ``andiamo a fermare qualsiasi tipologia
        di movimento''. Il target non viene detto: il nodo appartiene al componente di
        movimento del personaggio, il \emph{Character Movement};
  \item sull'On Completed, \textbf{Enable Input}, con lo stesso Player Controller;
  \item e infine il \textbf{Reset} del Do Once, ``in modo tale che io possa reinteragire di
        nuovo''.
\end{itemize}
In Play: ``io giro di qua, si gira, fa l'animazione, non mi posso muovere''; ``e adesso mi
riposso muovere appena finito il Montage'' (figura~\ref{ue12:fig:player}). ``Ci siete? Vi
funziona? Perfetto.''

\paragraph{I problemi dei presenti.} Uno studente non trova come prendere il componente
(``l'hai aggiunto come components? Prendilo e trascinalo''); altri si perdono i due nodi
successivi, ripetuti dal docente (Get Interactable, Convert to Validated Get, il messaggio
Interaction). Un altro ha il Montage che non parte: manca la Mesh nel pin del Play Montage.
Uno studente, a cui prima non funzionava, conferma che ora va, ma ha lo
schermo ``pieno di \emph{Hello}'' (forse i Print String rimasti con il testo di default),
``quindi li vedo con la coda dell'occhio''.

\begin{figure}[htbp]
\centering
\begin{tikzpicture}[font=\footnotesize,
  ue12node/.style={draw, thick, rounded corners=4pt, align=center, fill=gray!8,
                  minimum height=0.9cm, text width=2.2cm},
  ue12ev/.style={ue12node, fill=red!12},
  ue12msgn/.style={ue12node, fill=violet!12},
  ue12pure/.style={ue12node, fill=green!10, rounded corners=8pt},
  ue12exec/.style={-{Stealth[length=2.5mm]}, very thick, black!80},
  ue12data/.style={-{Stealth[length=2mm]}, thick, blue!60!black}]
  \node[ue12ev, text width=1.9cm] (ia) at (0,-0.0) {\textbf{IA\_Interact}\\{\scriptsize tasto
    E}};
  \node[ue12node, text width=1.5cm] (do) at (2.3,0) {\textbf{Do Once}};
  \node[ue12node, text width=2.4cm] (vg) at (4.9,0) {\textbf{Validated Get}\\{\scriptsize
    BPC\_RayInteract.\\Interactable}};
  \node[ue12msgn] (msg) at (8.6,0) {\textbf{Interaction}\\{\scriptsize (Message) su
    Interactable}};
  \node[ue12node, text width=2.5cm] (rot) at (11.9,0) {\textbf{Set Actor Rotation}\\{\scriptsize
    X, Y attuali; Z dal Look at}};
  \node[ue12pure, text width=3.0cm] (look) at (11.9,-1.85) {\textbf{Find Look at
    Rotation}\\{\scriptsize mia posizione $\rightarrow$ posizione di Interactable}\\
    $\rightarrow$ \textbf{Break Rotator}: Z};
  \node[ue12node, text width=3.5cm] (pre) at (8.0,-4.4) {\textbf{Set Interacting} = vero\\
    \textbf{Disable Input}\\\textbf{Stop Movement Immediately}};
  \node[ue12node, text width=2.8cm] (pm) at (12.0,-4.4) {\textbf{Play Montage}\\{\scriptsize
    Mesh; Interact Press Button Montage}};
  \node[ue12node, text width=3.5cm] (post) at (12.0,-6.8) {\textbf{Set Interacting} =
    falso\\\textbf{Enable Input}\\\textbf{Reset} del Do Once};
  \node[ue12pure, text width=1.9cm] (gpc) at (8.0,-6.8) {Get Player Controller};
  \draw[ue12exec] (ia) -- (do); \draw[ue12exec] (do) -- (vg);
  \draw[ue12exec] (vg) -- node[below, font=\scriptsize]{Is Valid} (msg);
  \draw[ue12exec] (msg) -- (rot);
  \draw[ue12exec] (rot.east) -- ++(0.35,0) -- ++(0,-3.15) -| (pre.north);
  \draw[ue12exec] (pre) -- (pm);
  \draw[ue12exec] (pm) -- node[right, font=\scriptsize]{On Completed} (post);
  \draw[ue12exec] (vg.south) -- ++(0,-1.0) -| node[pos=0.25, above, font=\scriptsize]{Is
    Not Valid: Reset} ([xshift=-3mm]do.south);
  \draw[ue12data] (look) -- (rot);
  \draw[ue12data] (gpc) -- (pre);
  \draw[ue12data] (gpc) -- (post);
\end{tikzpicture}
\caption{L'interazione nel personaggio. Il personaggio non rifà la trace: legge la variabile
Interactable del suo componente, invia il messaggio, si gira solo sull'asse Z verso
l'oggetto, blocca l'input e suona il Montage; a Montage finito restituisce l'input e riapre
il Do Once. Ricostruzione dalla lezione; la posizione esatta del Set Interacting e del
Disable Input nella catena è dedotta dalle parole del docente.}
\label{ue12:fig:player}
\end{figure}

\paragraph{Approfondimento: Do Once, Disable Input e la rotazione.} Romero e Sewell
descrivono il \textbf{Do Once} nel capitolo 13 (\emph{Data Structures and Flow Control}):
``il nodo Do Once esegue il pin di uscita una sola volta. Se viene attivato di nuovo, il pin
di uscita non viene eseguito. Il pin di input Reset deve essere attivato per permettere al
Do Once di eseguire di nuovo il pin di uscita.'' Il loro esempio è un'arma caricata, che si
ricarica dopo 30 secondi con un timer che resetta il Do Once: lo stesso schema
della lezione, con il Montage al posto del timer. Nel capitolo 15 il libro descrive
\textbf{Enable Input} e \textbf{Disable Input}: ``definiscono se un actor deve rispondere agli
eventi di input'', e ``hanno bisogno di un riferimento al Player Controller in uso''. Due
note generali sull'engine. La prima: \emph{Disable Input} chiamato sul personaggio blocca
gli eventi d'input di quel Blueprint (movimento, camera, salto), ma non ferma un movimento
già in corso, cioè la velocità residua del Character Movement; per questo serve anche
\emph{Stop Movement Immediately}, che azzera la velocità sul posto. La seconda: il
\emph{Set Actor Rotation} è istantaneo, quindi il personaggio ``scatta'' verso l'oggetto. Per
una rotazione morbida si può usare un \emph{RInterp To} nel Tick o una Timeline, come per le
porte; e il template Third Person, che ha l'\emph{Orient Rotation to Movement} attivo, non
contrasta la rotazione perché durante l'interazione il personaggio è fermo.

% ---------------------------------------------------------------------
\section{Il widget ``premi E''}
\label{ue12:sec:widget}

``Andiamo semplicemente a rendere tutto un po' più carino'': un piccolo widget che mostri il
tasto da premere sopra l'oggetto.

\paragraph{Gli asset della UI.} Il docente aveva messo nel materiale didattico una cartella
di immagini per la UI, ``la UI asset''. Uno studente non l'ha trovata (``pensavo
fosse sul Drive''): il docente controlla il Drive condiviso e i canali del Discord, scopre che
non era stata caricata e la carica al momento, ``tanto pesa pochi kilobyte''. ``Estraetela e
schiaffatela dentro Unreal'', nella cartella \textbf{UI}. Il trascinamento nel Content
Browser a qualcuno dà ``fallimento''; il ``metodo abbastanza brusco, però va benissimo'' è
copiare la cartella \emph{estratta} con il gestore dei file direttamente nella cartella
\emph{Content/UI} del progetto. ``Adesso funziona.''

\paragraph{WB\_InteractUI.} Nella cartella UI, un nuovo Widget Blueprint, classe User
Widget: \textbf{WB\_InteractUI}. Nel Designer:
\begin{enumerate}
  \item un \textbf{Border}, con l'alfa del colore a \textbf{0}, ridimensionato a un
        quadrato di circa \textbf{96 $\times$ 96} (``94 va bene anche''), e il
        \textbf{Padding} tutto a \textbf{zero};
  \item un secondo \textbf{Border} \emph{dentro} il primo;
  \item le immagini: nel \textbf{Brush} di ciascun Border (``vedi che c'è il brush, apri il
        brush'') si sceglie un'immagine del pacchetto, ``il panel che volete voi'' (ce ne
        sono diversi, Panel 0, 1, 2, 3, 5\dots). Dopo una correzione (``questo è il
        contrario, scusate''), il Border esterno prende \textbf{Panel Border 5}, con l'alfa a
        \textbf{1} e la tinta \textbf{nera}, e quello interno \textbf{Panel 5};
  \item un \textbf{Text} con la lettera \textbf{E}, con \emph{Horizontal Alignment} e
        \emph{Vertical Alignment} al \textbf{centro}, e la dimensione del \textbf{font} a
        35, poi 40--45 (``una bella E'').
\end{enumerate}
``Piccola cosa, ragazzi, prima non l'ho fatto fare a mano'': la dimensione dei Border si può
impostare anche con precisione, dall'\textbf{Image Size} del Brush (96 $\times$ 96).

\paragraph{Il widget che contiene la barra della vita.} Una studentessa crea un nuovo
widget, ma aprendolo trova le barre della vita e della stamina, e non riesce a modificarle.
Rifacendo la creazione insieme, la causa emerge: nella finestra di scelta della classe aveva
preso un widget esistente come genitore (``anteprima widget praticamente ti dà
un'anteprima dei widget che già ci sono''), cioè il suo HUD; il nuovo widget ne
\emph{ereditava} il contenuto, in sola lettura. ``È buona prassi fare quello tutto in
alto'', \textbf{User Widget}.

\paragraph{L'immagine che non si vede.} Un'altra studentessa ha il bordo esterno che non
compare. Il docente controlla tinta, alfa, Image Size, Padding, \emph{Content Opacity}: ``è
a 1, va bene così''. ``Perdonami, purtroppo non saprei dire'': si cancellano i Border e si
rifanno, ``basta funzionare''. Rifatti, funzionano.

\paragraph{Il widget nell'oggetto.} In BP\_InteractableTest, il componente \textbf{Widget}:
\begin{itemize}
  \item \textbf{Space}: \textbf{Screen}, non World, ``in modo tale che ci sia una specie di
        2D, quindi relativo allo schermo, non al mondo'';
  \item \textbf{Draw Size}: \textbf{100 $\times$ 100}, ``un po' più piccolo'';
  \item \textbf{Widget Class}: WB\_InteractUI.
\end{itemize}
In Play, il quadratino con la E ``scala con lo schermo'', ma, se non lo si nasconde,
``rimane per sempre''.

\paragraph{Mostrare e nascondere.} Nel grafo di BP\_InteractableTest:
\begin{itemize}
  \item \textbf{Event BeginPlay} $\rightarrow$ \textbf{Set Hidden in Game} sul Widget, con
        \emph{New Hidden} \textbf{attivo}: nascosto da subito;
  \item \textbf{Event Interaction Start} $\rightarrow$ Set Hidden in Game,
        \textbf{disattivo}: ``allora la faccio apparire'';
  \item \textbf{Event Interaction End} $\rightarrow$ Set Hidden in Game, \textbf{attivo}.
\end{itemize}
Il docente prova anche un'idea per l'Event Interaction (nascondere il widget, un
\textbf{Delay} di un secondo, mostrarlo di nuovo), ma la toglie: ``questo potrebbe dare
errori, ragazzi'', perché il Delay ``va in contrasto'' con Start ed End (se lo sguardo si
sposta durante il secondo d'attesa, il Delay rimostrerebbe un widget che l'End ha appena
nascosto). ``Lascialo così per adesso.''

\paragraph{La spunta al contrario.} Una studentessa vede la E anche quando non guarda il
cubo: ``allora vuol dire che hai sbagliato l'Hidden in Game''. ``Il tick è positivo sull'End
e negativo sullo Start'', ``e anche positivo sul Begin Play''. Dopo la correzione resta un
dubbio, poi chiarito: all'inizio della prova la camera stava già inquadrando il cubo da una
posizione strana, e il widget era giustamente visibile. ``Non lo vedo più. Perfetto,
ragazzi.''

\paragraph{Approfondimento: il Widget Component.} Il \textbf{Widget Component} mostra un
Widget Blueprint attaccato a un actor. In \textbf{World} space il widget è un piano nel
mondo 3D: ha prospettiva, può essere coperto da altri oggetti, si rimpicciolisce con la
distanza, ed è adatto a schermi diegetici (un terminale, un cartello). In \textbf{Screen}
space il widget è disegnato sopra l'immagine, sempre rivolto alla camera, con una dimensione
costante in pixel, e segue sullo schermo la posizione proiettata del componente: è la scelta
tipica per i suggerimenti ``premi E'' e per le barre della vita sopra i nemici. La
\textbf{Draw Size} è la risoluzione in pixel a cui il widget viene disegnato. Per nascondere
il suggerimento, l'alternativa a \emph{Set Hidden in Game} è \emph{Set Visibility} sul
componente; il risultato è lo stesso. Funzionamento generale dell'engine.

% ---------------------------------------------------------------------
\section{Un esempio: la luce che si accende}
\label{ue12:sec:luce}

\paragraph{L'idea.} ``Sappiate che tutti gli oggetti che create'', se implementano
l'interfaccia, ``avete già un qualcosa che vi identifica se state iniziando
un'interazione'', se la state finendo, ``e la tipologia di interazione''. Le funzioni
costruite per i bottoni nelle lezioni passate ``possiamo traslarle sul nostro
interactable'': ``un diverso modo di interagire e di attivare le cose''. Con Davide si
vedranno ``anche delle piccole cose di modeling'': si può modellare un pulsante, dargli
BPI\_InteractionRay, ``e a quel punto magari potete chiamare l'ascensore''. ``Cosa succede se
interagisco? Ci metto quello che voglio: ci metto un timer, ci metto una cosa, accendo una
luce.''

\paragraph{Velocissimo.} ``Facciamo velocissimo'', in BP\_InteractableTest. Dall'Event
Interaction si toglie il Print String, e:
\begin{enumerate}
  \item una variabile del tipo dell'attivabile luce, la classe figlia di BP\_Activable delle
        lezioni sul pulsante (``la Point Light'', \textbf{BP\_ActivableLight}; la
        trascrizione dice ``BP Activables'' e ``inter light''), con i suoi eventi
        \textbf{Activation} e \textbf{Deactivation};
  \item un \textbf{Flip Flop}: ``questo fa A, poi fa B, poi fa di nuovo A: fondamentalmente
        gioca su due'';
  \item su \textbf{A}, \textbf{Activation} della luce; su \textbf{B},
        \textbf{Deactivation};
  \item la variabile va resa modificabile nell'istanza e collegata alla luce del livello
        (``però gli dobbiamo dare quale light'', ``dammi la light, dammi questa'').
\end{enumerate}
``Abbiamo acceso la luce e la rispegniamo.'' ``Questo è un esempio, ragazzi'': ``qualsiasi
actor può essere separato dagli altri con la sua specifica interazione'', e ``quello che la
identifica già ce l'abbiamo''.

\paragraph{Approfondimento: Flip Flop nel libro.} Romero e Sewell (capitolo 13): ``ogni volta
che un nodo Flip Flop viene eseguito, alterna fra due pin di uscita, A e B''; ha anche un
output booleano \emph{Is A}, vero quando sta eseguendo A. Con una nota importante: ``il Flip
Flop non funziona correttamente dentro una funzione, perché ha una variabile interna che
viene cancellata quando la funzione finisce'', e quindi dentro una funzione esegue sempre A.
Lo stesso vale per il Do Once. È il motivo per cui, nella lezione, il Flip Flop sta
nell'Event Graph dell'oggetto e il Do Once nell'Event Graph del personaggio, mentre la
funzione LookForInteractions, che non deve ricordare nulla fra un frame e l'altro (lo stato
è nella variabile Interactable, che è del componente), può essere una funzione.

% ---------------------------------------------------------------------
\section{La chiusura del corso}
\label{ue12:sec:chiusura}

``Ci sono domande, ragazzi? Domande su tutto il corso, in generale.'' Il docente chiede
anche un \textbf{feedback} sul corso, ``adesso o dopo, tranquillamente sul server''. La
risposta di uno studente, scherzosa: ``tutto chiaro, tutto bellissimo, incredibile''. ``Sono
contento.''

Per il futuro, anche ``per chi guarda in differita'': ``se avete un problema, state facendo il
Big One, eccetera, buttate il vostro problema sul server: sia alunni sia professori
provvederanno a rispondervi e darvi una mano''. ``Se ci sono ultime cose che volete dire,
ragazzi, ditele ora, o effettivamente questa volta tacete per sempre.'' E il congedo: ``vi
ringrazio per la vostra attenzione, vi ringrazio per aver partecipato al corso. Spero che i
vostri progettini possano prendere vita per tutte le cose che poi andrete a fare in Unreal.
Per benino, ricordatevi di \textbf{validare sempre} molte variabili.'' ``Da questo punto in
poi'' le lezioni continuano con Davide. ``È stato un piacere.''

\begin{table}[htbp]
\centering
\small
\begin{tabular}{@{}p{3.7cm}p{3.0cm}p{2.9cm}p{4.4cm}@{}}
\toprule
\textbf{Asset o variabile} & \textbf{Dove} & \textbf{Tipo, valore} & \textbf{A che serve} \\
\midrule
M\_Menu & mappe & Empty Level & la mappa del menu \\
WB\_Menu & UI & Widget Blueprint & il menu con i due bottoni \\
Button\_StartGame, Button\_QuitGame & WB\_Menu & Button, 300 $\times$ 80 & gioca; esci \\
CameraActor & M\_Menu & Camera Actor & la vista fissa del menu \\
GameMode Override & World Settings di M\_Menu & Game Mode base & niente personaggio nel menu \\
ThirdPersonMap & mappe & livello & la mappa aperta da Play Game \\
BPC\_RayInteract & Blueprints, Components & Scene Component & traccia gli interagibili \\
CameraComponent & BPC\_RayInteract & Camera Component (ref.) & la camera del proprietario \\
DistanceTrace & BPC\_RayInteract & Float, 300 & la lunghezza del raggio \\
Interactable & BPC\_RayInteract & Actor (ref.) & l'interagibile guardato ora, o vuoto \\
Scene & BPC\_RayInteract & Scene Component (ref.) & solo per il trucco del Select \\
InteractionTrace & Project Settings & Trace Channel, Ignore & il canale delle interazioni \\
BPI\_InteractionRay & Blueprints, Components & Blueprint Interface & Start, End, Interaction \\
BP\_InteractableTest & Blueprints & Actor (Cube, Widget) & l'oggetto di prova \\
IA\_Interact & Input & Input Action, Pressed e Released & l'interazione, tasto E \\
Interacting & BP\_Bot2 & Boolean & vero durante l'interazione \\
Interact Press Button Montage & animazioni del bot & Anim Montage & l'animazione dell'interazione \\
WB\_InteractUI & UI & Widget Blueprint, 96 $\times$ 96 & il suggerimento ``E'' \\
Widget (componente) & BP\_InteractableTest & Screen, 100 $\times$ 100 & mostra WB\_InteractUI \\
Light & BP\_InteractableTest & BP\_ActivableLight (ref.) & la luce accesa e spenta \\
\bottomrule
\end{tabular}
\caption{Gli asset e le variabili della giornata, con i nomi usati dal docente (BPC\_RayInteract
è, nel progetto del docente, BPC\_RayInteract1; il nome della variabile della luce è
incerto).}
\label{ue12:tab:asset}
\end{table}

\begin{table}[htbp]
\centering
\small
\begin{tabular}{@{}p{5.4cm}p{8.6cm}@{}}
\toprule
\textbf{Sintomo} & \textbf{Causa e rimedio} \\
\midrule
La barra della stamina scende solo al rilascio & manca la Call E\_UpdateStamina in
StaminaDeduction: c'era solo in RegainStamina \\
Il secondo bottone incollato non compare & è sovrapposto al primo (controllare la
Hierarchy); oppure aggiungerne uno nuovo dalla Palette \\
L'immagine copre i bottoni & è in fondo alla Hierarchy, cioè disegnata sopra: trascinarla
prima dei bottoni \\
Il personaggio compare dietro il menu & GameMode Override nelle World Settings della
mappa del menu: Game Mode base \\
Si vede la pallina (Default Pawn) & spostare il Player Start fuori dalla camera \\
La camera del menu non si usa & Set View Target with Blend dal Get Player Controller (non
dal Camera Manager), con la referenza alla camera \\
Errore su Add to Viewport & il Target è Self: collegare il Return Value del Create Widget \\
Il cursore non si vede & Set Show Mouse Cursor nell'Event Construct, con la spunta attiva \\
Nel menu ci si muove ancora & Set Input Mode UI Only nell'Event Construct \\
Il Select non accetta Self (o la camera) & il Select prende il tipo del primo collegamento:
mettere prima una variabile Scene Component, poi cancellarla \\
In Play non si vede il raggio & LookForInteractions non chiamata nell'Event Tick; Draw
Debug Type a None \\
InteractionTrace non compare nella trace & compilare dopo aver creato il canale \\
Il raggio ignora il cubo & la mesh non blocca InteractionTrace (default Ignore): Block \\
L'evento dell'interfaccia non si trascina & non è Blueprint Callable: implementarlo come
evento (doppio clic o Implement Event) \\
Hit Actor non si collega a Interactable & Interactable di tipo sbagliato: Actor, Object
Reference \\
Si interagisce una sola volta & il Do Once non viene resettato (On Completed e Is Not
Valid) \\
Il Montage non parte & manca la Mesh nell'In Skeletal Mesh Component \\
Il nuovo widget contiene le barre dell'HUD & creato con l'HUD come genitore: scegliere User
Widget \\
Il bordo del widget non si vede & causa non trovata; rifare i Border \\
La E si vede quando non si guarda & spunte di Set Hidden in Game invertite: attivo nel
Begin Play e nell'End, disattivo nello Start \\
Trascinare la cartella UI dà ``fallimento'' & copiare la cartella estratta in Content/UI dal
gestore dei file \\
La BPI si apre senza grafo & doppio clic sulla funzione in My Blueprint \\
\bottomrule
\end{tabular}
\caption{I problemi incontrati durante la lezione, con la causa e il rimedio.}
\label{ue12:tab:errori}
\end{table}

\begin{table}[htbp]
\centering
\footnotesize
\begin{tabular}{@{}p{5.8cm}p{8.2cm}@{}}
\toprule
\textbf{Nodo, asset o impostazione} & \textbf{Che cosa fa} \\
\midrule
New Level: Empty / Basic & livello vuoto; con luce e cielo già pronti \\
Widget Blueprint (User Widget) & l'asset dell'interfaccia: Designer e Graph \\
Button: Style (Normal, Hovered, Pressed) & i colori del bottone nei tre stati \\
Hex sRGB / Hex Linear & il codice del colore, da copiare \\
Text dentro un Button & l'etichetta del bottone, figlia del bottone \\
Image: Color and Opacity (alfa) & la trasparenza dell'immagine \\
Hierarchy (ordine) & l'ultimo elemento è disegnato sopra \\
Anchors / Position / Alignment / Size & ancora; scostamento; pivot; dimensioni \\
Open Level Blueprint & il Blueprint del livello corrente \\
Create Widget / Add to Viewport & crea il widget; lo mostra \\
World Settings: GameMode Override & la Game Mode per questa mappa \\
Default Pawn / Player Start & la sfera della Game Mode base; dove nasce il Pawn \\
Camera Actor / Directional Light & una camera nel livello; il sole \\
Create a Reference to (actor) & la referenza a un actor del livello nel Level Blueprint \\
Set View Target with Blend & cambia la camera del Player Controller \\
Event Pre Construct / Construct & alla costruzione nel Designer; alla costruzione in gioco \\
Set Show Mouse Cursor & mostra o nasconde il cursore \\
Set Input Mode UI Only / Game Only & l'input va solo all'interfaccia; solo al gioco \\
On Clicked (Button) & l'evento del clic su un bottone \\
Quit Game & esce dal gioco (in editor, chiude il Play) \\
Open Level (by Name) & apre una mappa dal nome \\
Blueprint Interface & solo dichiarazioni di funzioni, in sola lettura \\
Class Settings: Implemented Interfaces & aggiunge un'interfaccia a un Blueprint \\
Interaction Start / End / Interaction (Message) & i messaggi dell'interfaccia \\
Event Interaction Start / End / Interaction & la loro implementazione nell'oggetto \\
Actor Component / Scene Component & componente senza / con Transform \\
Get Owner & l'actor che possiede il componente \\
Get Component by Class & il primo componente di una classe (non l'array) \\
Sphere Trace By Channel & una sfera lungo un segmento: Radius, Draw Debug Type \\
Get World Location / Get Forward Vector & posizione nel mondo; direzione davanti \\
Select (wildcard) & sceglie fra due valori; il tipo è quello del primo collegato \\
Is Valid (pura) / Validated Get & il riferimento esiste?; Get con le due uscite \\
Not Equal (!=) & vero se i due valori sono diversi \\
New Trace Channel (Ignore / Block) & un canale di collisione personalizzato \\
Break Hit Result: Hit Actor & l'actor colpito \\
Does Object Implement Interface & l'oggetto implementa l'interfaccia? \\
Event Tick & a ogni frame \\
Print String: Text Color & il colore del messaggio \\
Input Action (Pressed, Released) / IMC & l'azione e il tasto (E) \\
Do Once (Reset) & esegue una volta, finché non si resetta \\
Set Actor Rotation / Get Actor Rotation & imposta; legge la rotazione dell'actor \\
Find Look at Rotation / Break Rotator & rotazione per guardare un punto; i suoi tre assi \\
Create AnimMontage / Play Montage & un Montage da un'animazione; lo suona (On Completed) \\
Disable Input / Enable Input & toglie; restituisce l'input all'actor \\
Stop Movement Immediately & azzera il movimento del personaggio \\
Border: Brush, Padding, Image Size & il riquadro con immagine; margine interno; dimensione \\
Widget Component: Space, Draw Size, Widget Class & un widget attaccato a un actor \\
Set Hidden in Game (New Hidden) & nasconde o mostra un componente \\
Delay & aspetta prima di proseguire \\
Flip Flop & alterna fra le uscite A e B \\
\bottomrule
\end{tabular}
\caption{I nodi, gli asset e le impostazioni della dodicesima lezione.}
\label{ue12:tab:nodi}
\end{table}

% ---------------------------------------------------------------------
\section*{Punti chiave del capitolo}
\addcontentsline{toc}{section}{Punti chiave del capitolo}

\begin{itemize}
  \item Il designer usa Unreal per \textbf{prototipare} le meccaniche e verificare che
        funzionino; il progetto d'esame degli anni successivi è un \textbf{pezzo da
        portfolio}. \textbf{Less is more}: il giocatore impiega molto più tempo di chi
        ha progettato il gioco per capirne ogni elemento.
  \item Il menu vive in una \textbf{mappa a parte} (Empty Level), con un \textbf{Widget
        Blueprint} (User Widget) creato e mostrato nel \textbf{Level Blueprint} (Create
        Widget, Add to Viewport con il Return Value).
  \item Nel Designer: Canvas Panel, bottoni con l'\textbf{ancora} al centro e
        l'\textbf{Alignment} per impilarli, un \textbf{Text} dentro ogni bottone, i colori
        di \textbf{Normal}, \textbf{Hovered} e \textbf{Pressed} nello Style (copiabili con
        il codice Hex), e l'ordine della \textbf{Hierarchy} per mettere un'immagine
        dietro (l'ultimo elemento sta sopra).
  \item Niente personaggio nel menu: \textbf{GameMode Override} con la Game Mode base
        (Default Pawn) e Player Start fuori campo. La vista fissa: un Camera Actor e
        \textbf{Set View Target with Blend}, dal \textbf{Player Controller}.
  \item Nell'\textbf{Event Construct} del widget: \textbf{Set Show Mouse Cursor} (con la
        spunta) e \textbf{Set Input Mode UI Only}. Sui bottoni, \textbf{On Clicked}:
        \textbf{Quit Game}; \textbf{Set Input Mode Game Only} e \textbf{Open Level (by
        Name)}.
  \item Una \textbf{Blueprint Interface} è un intermediario (``un volo di linea''): solo
        nomi di funzioni, in sola lettura. Chi la implementa (Class Settings,
        Implemented Interfaces) decide che cosa fare; chi non la implementa ignora i
        messaggi.
  \item Un \textbf{Blueprint Component} isola una funzionalità da aggiungere a qualunque
        actor. L'\textbf{Actor Component} non ha posizione; lo \textbf{Scene Component}
        ha un Transform. Il componente trova il suo actor con \textbf{Get Owner} e la
        camera con \textbf{Get Component by Class}.
  \item \textbf{BPC\_RayInteract}: a ogni frame (\textbf{Event Tick}) una \textbf{Sphere
        Trace} (Radius 5) dalla posizione del componente, lungo il Forward Vector della
        camera (o del componente, se la camera non c'è) per 300 unità, sul canale
        \textbf{InteractionTrace} (Ignore di default, Block sugli interagibili).
  \item Il \textbf{Select} prende il tipo del primo collegamento: per scegliere fra la
        camera e Self, prima una variabile \emph{Scene Component}, poi la si cancella (la
        ``mandragata'').
  \item La variabile \textbf{Interactable} ricorda l'oggetto guardato: se cambia,
        \textbf{Interaction End} sul vecchio, Set, \textbf{Interaction Start} sul nuovo;
        se non si guarda niente di interagibile, End e svuota. Sempre con il
        \textbf{Validated Get}.
  \item Nel personaggio: \textbf{IA\_Interact} (tasto E), \textbf{Do Once}, il messaggio
        \textbf{Interaction}, \textbf{Set Actor Rotation} con il solo \textbf{Yaw} da
        \textbf{Find Look at Rotation}, \textbf{Disable Input} e \textbf{Stop Movement
        Immediately}, \textbf{Play Montage}, poi Enable Input e \textbf{Reset} del Do Once
        (anche quando non c'è nulla da interagire).
  \item Il suggerimento ``E'': un \textbf{Widget Component} in \textbf{Screen} space,
        nascosto nel Begin Play e mostrato e nascosto da Interaction Start ed End con
        \textbf{Set Hidden in Game}.
  \item Ogni oggetto può dare la sua risposta all'interazione: per esempio un
        \textbf{Flip Flop} che chiama Activation e Deactivation di una
        \textbf{BP\_ActivableLight}. Flip Flop e Do Once non vanno dentro le funzioni.
  \item La raccomandazione finale del corso: \textbf{validare sempre} le variabili.
\end{itemize}

%
%  Tutte le etichette di questo capitolo hanno il prefisso "ue12:"
%  (ue12:cap:..., ue12:sec:..., ue12:fig:..., ue12:tab:...), e tutti gli
%  stili TikZ hanno il prefisso "ue12" (ue12node, ue12exec, ue12data,
%  ...) per non collidere con chiavi native di TikZ né con quelli delle
%  lezioni precedenti. Il capitolo è autonomo: non contiene \ref verso
%  altri capitoli (i rimandi alle lezioni precedenti sono a parole).
%  Le figure sono ricostruzioni schematiche di quanto mostrato e detto a
%  lezione, non copie delle schermate del docente. Dove la trascrizione
%  è troppo rovinata per ricostruire un dettaglio, il capitolo lo dice
%  invece di inventarlo.
% =====================================================================

\chapter[Menu e sistema di interazione]{Il menu principale e il sistema di
interazione}
\label{ue12:cap:interazione}

Questa è la \textbf{dodicesima lezione} del corso breve di \textbf{introduzione a Unreal
Engine}, ed è l'\textbf{ultima} con questo docente: ``oggi è l'ultima lezione, almeno
mia''. Nella lezione precedente il personaggio, il \textbf{bot}, ha ricevuto la vita, il
danno della torretta, la stamina e le due barre dell'HUD; restavano in sospeso due
argomenti annunciati, il \textbf{menu} e le \textbf{interazioni}. Oggi si fanno entrambi.
``Sono argomenti che abbiamo già fatto in precedenza e adesso li uniamo'': Widget
Blueprint, Level Blueprint, trace, collisioni personalizzate, Montage, timer e attivabili
tornano tutti, messi insieme.

Il primo argomento, il più corto, è un \textbf{menu principale}: una mappa a parte, con un
piccolo \emph{diorama} 3D sullo sfondo, una camera fissa, due bottoni (\emph{gioca} ed
\emph{esci}) e il passaggio alla mappa di gioco. Il secondo, ``piuttosto lungo'', è un
\textbf{sistema di interazione} del tipo ``premi E per interagire'', costruito con due
strumenti nuovi: un \textbf{Blueprint Component}, che si aggiunge a qualunque actor e
``traccia'' dalla camera gli oggetti interagibili, e una \textbf{Blueprint Interface}, il
``porto'' attraverso cui il personaggio dice a un oggetto \emph{sto iniziando a guardarti},
\emph{non ti guardo più}, \emph{interagisco con te}. L'oggetto di prova mostra un piccolo
widget con la lettera E quando è inquadrato; premendo E il personaggio si gira verso di lui
e suona un'animazione; alla fine, lo stesso oggetto accende e spegne una luce.

La lezione è \textbf{online}, come le precedenti, e si apre con una conversazione più
generale: come recuperare le lezioni, che cosa aspetta gli studenti negli anni successivi,
e un consiglio di design, ``\emph{less is more}''. Durante la pausa il docente risolve il
problema della barra della stamina rimasto aperto la volta precedente; si chiude con i
ringraziamenti e con la raccomandazione che riassume il corso: ``ricordatevi di validare
sempre molte variabili''.

Tre avvertenze di lettura, le stesse delle lezioni precedenti. La prima: la trascrizione
è automatica e storpia moltissimi nomi. Unreal è ``Real'' e ``Arrial''; il
\textbf{diorama} è ``diuramino''; il Widget Blueprint \textbf{WB\_Menu} è ``VB menu'' e
``VBIN menu'', e \textbf{WB\_InteractUI} è ``VB Interact UI'' e ``Interacco uai''; il
\textbf{Begin Play} è ``big in play'' e ``back in play''; la \textbf{ThirdPersonMap} è
``Terperson map'' e ``there person''; il \textbf{Scene Component} è ``sin component'' e
``sim component''; il componente \textbf{BPC\_RayInteract} è ``BPC Reinteract'', ``Ray
interact one'' e ``IPC reinteract''; la Blueprint Interface \textbf{BPI\_InteractionRay} è
``VPI interaction'', ``VPA Interaction Array'' e ``bibliotecazione:ray''; l'\textbf{Hit
Actor} è ``it actor'' e ``I Tactor''; il \textbf{Validated Get} è ``variegate'' e ``vari da
get''; \textbf{Does Object Implement Interface} è ``Hacker does object implement
interface''; il \textbf{Draw Debug Type} è ``draw the bug type''; il bot è ``Boto'' e
``Boto2'', e il suo Blueprint ``BiPi Bot''; \textbf{Find Look at Rotation} è
``findLocatrotation''; l'\textbf{Anim Montage} è ``un montaggio un po' brasilero''; i
\emph{Big One} del corso di laurea sono ``biggest one''. ``Mandragata'', invece, non è un
errore: è il romanesco \emph{mandrakata}, ``trucchetto'', e il docente lo usa apposta.
Dove la ricostruzione è sicura il capitolo usa la forma corretta; dove non lo è riporta la
forma della trascrizione e lo dice. La seconda: durante la pausa il riconoscimento vocale
produce un lungo blocco di testo senza senso e ripetuto decine di volte (``grazie a tutti
per l'attenzione che ci ha spiegato'', ``qual è la risposta che ci può essere?''), seguito
da un frammento che sembra l'audio di un video in sottofondo (una ``blue map'' che mostra
``le secret room'' senza sprecare bombe, e ``un'altra monetina''): non appartengono alla
lezione e sono stati ignorati. La terza: buona parte della lezione è fatta di indicazioni
a voce mentre si lavora (``tira fuori'', ``attacca qua'', ``vai giù''), spesso rivolte
allo schermo di uno studente; il capitolo ricostruisce i grafi da ciò che il docente dice
mentre li costruisce, distingue i tentativi dalle soluzioni, completa i nomi esatti dei
nodi e dei pin di Unreal Engine 5, e dichiara i punti in cui la trascrizione lascia un
buco.

\paragraph{Come leggere questo capitolo.} Il capitolo segue l'ordine della lezione. Le
figure ricostruiscono la disposizione del menu, i due grafi che lo fanno funzionare (Level
Blueprint e Widget Blueprint), il ruolo di componente e interfaccia, la funzione che
cerca gli interagibili, il ciclo di vita della variabile Interactable e il grafo
dell'interazione nel personaggio. Le tabelle riassumono gli asset e le variabili della
giornata, gli errori incontrati durante la lezione e i nodi. Qua e là, dei paragrafi di
\textbf{approfondimento} collegano la lezione al libro di Marcos Romero e Brenden Sewell,
\emph{Blueprints Visual Scripting for Unreal Engine 5}, il testo di riferimento per questa
parte del libro, e a qualche nozione tecnica generale su Unreal Engine 5; sono segnalati
come tali e non fanno parte della lezione.

Il percorso del capitolo:
\begin{enumerate}
  \item le \textbf{domande d'apertura}: recuperare le lezioni, l'esame e i Big One, il
        lavoro del designer, \emph{less is more} (sezione~\ref{ue12:sec:inizio});
  \item gli \textbf{annunci}, la scelta dell'ordine e la \textbf{dimostrazione} del
        sistema di interazione (sezione~\ref{ue12:sec:annunci});
  \item il \textbf{menu}: la mappa, il widget, i bottoni e i loro colori
        (sezione~\ref{ue12:sec:menu});
  \item il \textbf{Level Blueprint} del menu: il widget sullo schermo, la Game Mode e la
        camera (sezione~\ref{ue12:sec:levelbp});
  \item il \textbf{Graph del menu}: cursore, input e bottoni; il diorama e i problemi
        degli studenti (sezioni \ref{ue12:sec:menugraph} e~\ref{ue12:sec:diorama});
  \item \textbf{interfacce} e \textbf{Blueprint Component} (sezione~\ref{ue12:sec:concetti});
  \item il componente \textbf{BPC\_RayInteract} e la sua Sphere Trace
        (sezione~\ref{ue12:sec:componente});
  \item il \textbf{Trace Channel} e la \textbf{Blueprint Interface}
        (sezione~\ref{ue12:sec:bpi});
  \item la \textbf{pausa}: la stamina di Tommaso (sezione~\ref{ue12:sec:pausa});
  \item la logica della variabile \textbf{Interactable} (sezione~\ref{ue12:sec:logica});
  \item l'oggetto di prova \textbf{BP\_InteractableTest} (sezione~\ref{ue12:sec:test});
  \item l'\textbf{input} e l'interazione nel personaggio (sezioni \ref{ue12:sec:input}
        e~\ref{ue12:sec:player});
  \item il widget \textbf{``premi E''} (sezione~\ref{ue12:sec:widget});
  \item la \textbf{luce} interattiva e la \textbf{chiusura} del corso (sezioni
        \ref{ue12:sec:luce} e~\ref{ue12:sec:chiusura}).
\end{enumerate}

% ---------------------------------------------------------------------
\section{Prima della lezione: recuperare, l'esame, il designer}
\label{ue12:sec:inizio}

\paragraph{Recuperare o seguire?} In apertura, prima che arrivino gli altri, uno
studente rimasto indietro (``le sto andando pian piano'') chiede se gli convenga
recuperare le lezioni vecchie o seguire questa. (Il nome, nella trascrizione, è incerto:
il docente lo chiama una volta ``Dave'', ma poco dopo ricorda questa conversazione come
quella fatta ``con Lorenzo''.) La risposta: seguire la lezione in diretta ha un vantaggio,
``se c'è un qualche particolare dell'argomento, puoi farmi direttamente una domanda''; poi,
siccome oggi si uniscono argomenti già visti, ``magari ti rifai anche quelle precedenti in
modo tale da stare un po' più al passo''. L'unico prerequisito per l'esercizio di oggi è
``avere il personaggio che si muove'', che ``può guardarsi in giro''; ``puoi anche usare
il personaggio standard, volendo''. Le domande si possono fare anche dopo, sul
\textbf{Discord}: ``mandarmi le domande anche quando non ci vediamo più, perché oggi è
l'ultima lezione''. Lo studente sceglie di seguire e di riguardare la registrazione.

\paragraph{L'esame e i Big One.} Seconda domanda: ci sarà un esame, con un progetto da
consegnare? ``Nel corso no'': questo corso ``è strutturato che sia un'introduzione per
voi''. Il docente chiede che corso farà lo studente (game designer: ``tanto farete sia game
sia level designer'') e descrive il percorso, avvertendo che è quello che ha fatto lui
(``almeno io feci così, adesso non so come sono cambiate le cose''):
\begin{itemize}
  \item al \textbf{primo anno} il progetto principale è un \textbf{gioco da tavolo}:
        ``farete analog game design'';
  \item al \textbf{secondo anno} il \textbf{Big One} è un progetto su Unreal, ``perché
        Unreal è molto più improntato per i designer'';
  \item al \textbf{terzo anno} il progetto più grande (nella trascrizione ``biggest one''),
        che può anche essere la continuazione del secondo, ``se il vostro progetto nel
        secondo anno è andato bene'', da portare fino al rilascio su \textbf{Steam},
        ``con l'Accademia assieme''.
\end{itemize}
In questo senso l'esame c'è: ``sarà farvi un \textbf{pezzo da portfolio}'', ``un gioco che
potete mettere sul vostro portfolio e promuovervi''. ``Quindi, più sapete, meglio è.''

\paragraph{Che cosa fa il designer con Unreal.} ``Io sono designer'': per un designer
Unreal ``servirà principalmente per prototipare le meccaniche'', cioè ``concretizzare
l'idea che tu hai studiato, che tu hai tirato fuori, concretizzarla in editor e vedere se
effettivamente funziona''. Il suo strumento sono i \textbf{Blueprint}, ``quindi non andare
troppo in depth'' nella programmazione: in ambito lavorativo il designer ha ``un
programmatore che lo assiste''; ``poi, quando si sta da soli, si fa di tutto''.

\paragraph{Less is more.} Un consiglio ``da esperienza'': \textbf{less is more}. ``Più è
leggibile un qualcosa anche nella vostra testa a livello di design, più sarà leggibile per
un giocatore.'' L'esempio è ``la regola delle parole a schermo'': il testo che scorre
all'inizio di \emph{Star Wars}. Chi l'ha scritto lo rilegge in venti secondi, perché lo
conosce già; lo spettatore, che deve ``cominciare a creare le connessioni mentali'', ci
mette ``non 20 secondi, ma 50''. Nel gioco vale lo stesso: chi lo ha progettato e testato
pensa ``vabbè, tanto questo lo so ormai'', ma il giocatore ``ci metterà il triplo del
tempo''. ``Più roba fate, più tempo il giocatore dovrà metterci per imparare tutto
quanto.'' ``Credo che \emph{less is more} ve lo faranno fare'', aggiunge, ``ve lo
diranno''; e per qualunque domanda di design, ``possiamo tranquillamente pure parlare''.

\paragraph{Approfondimento: la maledizione della conoscenza.} Il fenomeno descritto dal
docente ha un nome in psicologia, \emph{curse of knowledge}: chi conosce un'informazione
fatica a immaginare come la percepisce chi non la conosce. È il motivo per cui il
\emph{playtest} con persone esterne al team è insostituibile, e per cui le regole di un
gioco andrebbero introdotte una alla volta, ciascuna quando serve (il cosiddetto
\emph{onboarding} progressivo). I numeri del docente (20 secondi contro 50, ``il
triplo'') sono un'immagine, non una misura. Questa è una nozione generale di design, non
della lezione.

% ---------------------------------------------------------------------
\section{Gli annunci, l'ordine della giornata e la dimostrazione}
\label{ue12:sec:annunci}

\paragraph{Che cosa si fa oggi.} Prima che arrivino tutti, il docente chiede agli
studenti di scegliere: oggi ci sono ``sia il sistema di interazioni, che è piuttosto
lungo'', sia il \textbf{menu}, ``che l'altra volta non abbiamo finito, praticamente non
abbiamo neanche iniziato''. Prima il lungo o prima il corto? Gli studenti scelgono
``corto e lungo, così facciamo la pausa prima'' (``andiamo a un caffè''). Uno studente
ricorda un'altra cosa in sospeso: ``durante la pausa sistemare il problema della stamina,
che se tengo premuto non si vede, ma solo quando mollo il tasto si vede'' (è il problema
rimasto aperto alla fine della lezione precedente); ``ricordamelo appena andiamo in
pausa''.

\paragraph{Gli annunci.} Arrivati gli altri, il docente ripete: ``oggi praticamente è la
mia ultima lezione''. Seguiranno ``circa quattro lezioni'' (``adesso non mi ricordo quante
lezioni avete'') con \textbf{Davide}, dedicate alla parte di \textbf{level design}: ``i
tool che vi serviranno per fare level design'' e alcuni approfondimenti, ``credo di
materiali''. ``Ciò non vuol dire che io vi abbandoni'':
\begin{itemize}
  \item il \textbf{canale Discord} resta aperto, anche quando inizierà il corso vero e
        proprio: ``fate domande e dovrebbe arrivare qualcuno in vostro aiuto'';
  \item anche i lavori si possono postare per avere un feedback (il docente cita quello
        di \textbf{Federico La Porta}, ``che è bravissimo'', a cui ha già risposto);
  \item i due canali vanno tenuti separati: quello del corso (``Introduzione a Unreal'')
        per i compiti, ``così almeno non mandate mischiati, perché sennò perdo il filo'',
        e la chat generale ``per qualsiasi comunicazione generica'';
  \item le registrazioni restano disponibili ``fino a settembre, se non sbaglio'', ``quindi
        avete comunque abbastanza tempo per poterle recuperare''.
\end{itemize}

\paragraph{La dimostrazione.} ``Adesso vi faccio vedere solamente un'introduzione, come
nei classici video tutorial.'' Il personaggio è quello della volta scorsa, ``che si muoveva
con stamina e adesso perde vita con la torretta''. Dalla sua testa parte un \textbf{trace}
sottile, ``tipo un cilindrone'': ``questo sarà il nostro sistema che identificherà un
possibile \emph{interactable} tramite la telecamera''. A sinistra, sullo schermo, i
messaggi: se la camera inquadra l'oggetto compare \emph{inizio interazione}, se si
distoglie lo sguardo \emph{fine interazione}; premendo il tasto, il personaggio ``si gira in
automatico'' verso l'oggetto e interagisce. ``Visto che comunque abbiamo creato questo
sistema, è un buon stampo modulare'': in futuro, pulsanti e altri oggetti potranno avere
``l'interazione un po' più precisa'', riusando ciò che si è visto con i pulsanti classici:
``sempre premi pulsantino e fai cosetta, e animazione''.

% ---------------------------------------------------------------------
\section{Il menu: la mappa, il widget e i bottoni}
\label{ue12:sec:menu}

\paragraph{Il Widget Blueprint.} Nella cartella \textbf{UI} del progetto, tasto destro,
\emph{User Interface}, \textbf{Widget Blueprint}, ``come abbiamo creato per l'HUD''. Alla
richiesta della classe base, ``non ci serve una particolare, ce ne serve una generica'':
\textbf{User Widget}. Il nome: \textbf{WB\_Menu}.

\paragraph{Una mappa per il menu.} ``Una cosa importante che io almeno consiglio'': creare
una \textbf{nuova mappa} dedicata al menu, ``quindi differenziare il gioco dal menu''. ``Si
possono fare anche volendo insieme, ma a livello di gestione'' è più comodo separarli, ``ed
è molto più rapido''. Da \textbf{File}, \textbf{New Level}, si sceglie \textbf{Empty
Level}: ``se vi diventa tutto nero, è normale, ragazzi: è un empty level''. Un livello
\emph{Basic}, invece, contiene già ``la directional light, la skybox, tante cose che sono
già prefatte''; qui si parte da zero. Poi:
\begin{itemize}
  \item il menu si può fare ``sia in 2D, con immagini 2D'', sia con elementi 3D, ``una
        specie di diorama''; per cominciare, un \textbf{piano} abbastanza grande (dalle
        forme base: il docente prende prima un cubo, poi si corregge), con un materiale a
        griglia, ``il material base di Unreal'';
  \item si salva la mappa con \textbf{Ctrl+S} nella cartella delle mappe, con il nome
        \textbf{M\_Menu}.
\end{itemize}
(Un dettaglio marginale: il docente nota che qualcosa nel livello risulta \emph{Hidden in
Game}, poi si accorge che non lo è; la trascrizione non dice di che cosa si tratti.)

\paragraph{Il Canvas Panel.} Aperto WB\_Menu, ``la prima cosa, che in questo caso servirà
per adattarlo allo schermo'': il \textbf{Canvas Panel}, ``la classica canvas''. La
\textbf{Hierarchy} si vede ``da qua sotto''. A uno studente che non trova la
\textbf{Palette} il docente mostra il pannello a sinistra dell'editor.

\paragraph{Due bottoni.} ``Adesso andremo a fare due pulsanti molto semplici: uno di
\emph{exit} e uno di \emph{open the game}.'' Nella Palette, la sezione
\textbf{Common} (accanto alle altre: Audio, Input, Panel, Primitive e così via, che la
trascrizione elenca a fatica) contiene il \textbf{Button}; se ne trascinano due al centro
dello schermo. La disposizione passa per qualche tentativo ad alta voce (posizioni e
allineamenti come $-2{,}7$ e $-5$, presto abbandonati: ``ok, mi sono perso''); la versione
pulita, quella che il docente detta alla fine, è questa:
\begin{enumerate}
  \item \textbf{Size}: X \textbf{300}, Y \textbf{80} (``60? No, troppo piccolo'');
  \item \textbf{Anchors} al \textbf{centro}, l'icona centrale del menu delle ancore, ``in
        modo tale che questo bottone sia sempre ancorato al centro''; poi si
        \textbf{azzerano} i valori della \textbf{Position};
  \item \textbf{Alignment} \textbf{0,5} su X e \textbf{1,5} su Y, ``perché così è più
        corretto'': il bottone resta centrato orizzontalmente e si alza di una volta e
        mezza la sua altezza rispetto al centro dello schermo;
  \item per il secondo bottone, ``copiate e incollate letteralmente'' il primo, azzerate
        la posizione e cambiate solo l'allineamento verticale: dopo un ``$-1{,}5$? No,
        troppo'' e un ``$-0{,}5$'', il docente si ferma a \textbf{$-0{,}3$} (``va bene''),
        cioè il secondo bottone sta un po' sotto il centro.
\end{enumerate}
(I valori del secondo bottone, nella trascrizione, sono incerti: il senso è che i due
bottoni stanno uno sopra e uno sotto il centro, con un po' di spazio fra loro.) A uno
studente il copia e incolla non funziona: ``guarda nella Hierarchy: se c'hai due bottoni,
vuol dire che ce ne hai due'', magari sovrapposti; altrimenti, ``metti semplicemente un
altro bottone tramite il Common''.

\paragraph{I nomi.} ``Una cosa che è importante da fare'': nella Hierarchy si rinominano i
due bottoni con \textbf{F2}, \textbf{Button\_StartGame} e \textbf{Button\_QuitGame} (la
trascrizione dice ``Button With Game''). ``Così abbiamo semplicemente perfettamente
nomenclati.''

\paragraph{Il testo dentro il bottone.} ``Una cosa carina che già esiste'': un
\textbf{Text} (dalla sezione Common) trascinato \emph{sopra} un bottone ``si va a incastrare
per benino'', diventa suo figlio e si centra da solo. Nei \textbf{Details} del Text, sezione
\textbf{Content}, il campo \textbf{Text}: ``cancellate e scrivete \emph{Play Game}, o
\emph{Play}, come volete''; nell'altro bottone, \emph{Quit}. Il testo bianco ``non si vede
tantissimo'', ma i colori si cambiano subito dopo.

\paragraph{Lo Style del bottone.} Selezionando il bottone (``mi raccomando''), nei Details
la sezione \textbf{Style} ha tre stati, ciascuno con il suo colore:
\begin{itemize}
  \item \textbf{Normal}: il bottone a riposo;
  \item \textbf{Hovered}: ``vuol dire che ci passi sopra il cursore'';
  \item \textbf{Pressed}: ``se premuto''.
\end{itemize}
(La trascrizione nomina anche i \emph{foreground} e lo stato \emph{Disabled}: nelle
versioni recenti dell'engine lo stile del bottone ha anche i colori del primo piano per
ciascuno stato.) Cliccando sulla barra del colore si apre la ruota dei colori; il docente
mette il Normal \textbf{blu}, ``e avremo questo cambiamento in diretta''.

\paragraph{Copiare un colore esatto.} ``Se vogliamo prendere perfettamente il colore che ci
interessa per spostarlo all'altro bottone'': nella finestra del colore, in basso, ci sono i
campi \textbf{Hex sRGB} e \textbf{Hex Linear}. ``Potete prendere ambi i due'': si copia il
codice con Ctrl+C, si apre lo stesso stato nell'altro bottone e lo si incolla \emph{nello
stesso campo} (``se avete preso sRGB, sempre su sRGB''). È lo stesso principio dei siti
web: ``in HTML, per mettere un colore, i colori sono tutti identificati tramite un codice
colore''.

\paragraph{Un'immagine dietro i bottoni.} Si aggiunge un'\textbf{Image}, allargata dietro i
due bottoni; ma ``come potete vedere, mi copre tutto quanto''. ``Qui ci viene in aiuto la
Hierarchy'': come i livelli di Photoshop, gli elementi del Canvas sono disegnati
nell'ordine dell'elenco, e ``quello che è più sotto'' nell'elenco, stranamente, è quello
``a livello più alto'', disegnato sopra gli altri. L'immagine, aggiunta per ultima, è in
fondo all'elenco e quindi davanti ai bottoni. La si trascina nella Hierarchy fino a far
comparire una \textbf{linea blu} fra il Canvas Panel e il primo bottone, e la si lascia lì:
ora è disegnata per prima, cioè dietro. Poi la si rende trasparente: nel suo
\textbf{Color and Opacity} (che il docente cerca nel Brush), l'\textbf{alfa} a 0,5, ``o
forse un po' più basso'', \textbf{0,2}. ``Questa è la alfa di un colore: l'intensità,
l'opacità.'' Ne viene ``una specie di schermatina piccoletta'', regolata a mano attorno ai
bottoni (figura~\ref{ue12:fig:layout}).

\begin{figure}[htbp]
\centering
\begin{tikzpicture}[font=\footnotesize, x=1cm, y=1cm]
  % schermo
  \draw[very thick, fill=gray!5] (0,0) rectangle (8.0,5.0);
  \node[anchor=north west, font=\scriptsize] at (0.1,4.95) {Canvas Panel (schermo)};
  % immagine trasparente
  \draw[thick, fill=blue!8, opacity=0.9] (2.4,1.15) rectangle (5.6,4.15);
  \node[font=\scriptsize, anchor=north] at (4.0,4.12) {Image, alfa 0,2};
  % bottoni
  \draw[thick, fill=blue!35, rounded corners=2pt] (2.9,2.70) rectangle (5.1,3.30);
  \node at (4.0,3.0) {\textbf{Play Game}};
  \draw[thick, fill=blue!35, rounded corners=2pt] (2.9,1.62) rectangle (5.1,2.22);
  \node at (4.0,1.92) {\textbf{Quit}};
  % ancora al centro
  \draw[red!70!black, thick] (4.0,2.5) circle (0.09);
  \draw[red!70!black, dashed] (0,2.5) -- (8.0,2.5);
  \draw[red!70!black, dashed] (4.0,0) -- (4.0,5.0);
  \node[red!70!black, font=\scriptsize, anchor=west] at (5.7,2.62) {ancora al centro,};
  \node[red!70!black, font=\scriptsize, anchor=west] at (5.7,2.30) {Position 0, 0};
  \node[font=\scriptsize, anchor=west, align=left] at (5.7,3.35) {Alignment\\0,5 / 1,5};
  \node[font=\scriptsize, anchor=west, align=left] at (5.7,1.65) {Alignment\\0,5 / $-0{,}3$\\(circa)};
  % hierarchy
  \begin{scope}[xshift=9.0cm]
    \draw[thick, fill=gray!8] (0,0.6) rectangle (4.2,4.4);
    \node[anchor=north west, font=\scriptsize\bfseries] at (0.1,4.35) {Hierarchy};
    \node[anchor=west] at (0.2,3.6) {[WB\_Menu]};
    \node[anchor=west] at (0.45,3.1) {Canvas Panel};
    \node[anchor=west, fill=blue!10] at (0.7,2.6) {Image};
    \node[anchor=west] at (0.7,2.1) {Button\_StartGame};
    \node[anchor=west] at (0.95,1.65) {Text ``Play Game''};
    \node[anchor=west] at (0.7,1.2) {Button\_QuitGame};
    \node[anchor=west] at (0.95,0.8) {Text ``Quit''};
    \draw[-{Stealth[length=2mm]}, thick] (4.45,3.0) -- node[right, font=\scriptsize,
      align=left]{disegnati\\in ordine:\\l'ultimo\\sta sopra} (4.45,0.9);
  \end{scope}
\end{tikzpicture}
\caption{Il menu nel Designer. I due bottoni sono ancorati al centro dello schermo con la
posizione azzerata; è l'Alignment (il pivot del bottone) a spostarli sopra e sotto il
centro. L'immagine trasparente è stata trascinata in cima alla Hierarchy, così viene
disegnata per prima, dietro ai bottoni. Ricostruzione dalla lezione: i valori del secondo
bottone sono approssimativi.}
\label{ue12:fig:layout}
\end{figure}

\paragraph{Approfondimento: Alignment, ancore e Size To Content.} Con l'ancora al centro e
la posizione a zero, l'\textbf{Alignment} decide quale punto del bottone coincide con il
centro dello schermo: 0 è il bordo sinistro (o superiore), 1 il bordo destro (o inferiore),
0,5 il centro. Valori fuori dall'intervallo da 0 a 1 sono leciti: con Y = 1,5 il punto
d'aggancio sta mezzo bottone \emph{sotto} il bordo inferiore, quindi il bottone sale; con Y
negativo, scende. È il trucco usato a lezione per impilare due bottoni senza toccare la
posizione. La via più ordinata, per menu con molti bottoni, è un \textbf{Vertical Box}
(sezione Panel della Palette) ancorato al centro, con i bottoni dentro: lo spazio fra gli
elementi si regola con il \emph{Padding} di ciascuno e non servono calcoli. Nel libro di
Romero e Sewell (capitolo 8, \emph{WinMenu}) i bottoni sono ancorati al centro e spostati
con la \emph{Position} (per esempio X $-180$, Y 150, Size 360 $\times$ 100), con l'alignment
di default: è la terza via, equivalente. Funzionamento generale dell'UMG, salvo il
riferimento al libro.

% ---------------------------------------------------------------------
\section{Il Level Blueprint del menu: il widget, la Game Mode e la camera}
\label{ue12:sec:levelbp}

\paragraph{Perché il Level Blueprint.} ``Adesso che ce famo con questa roba?'' Il widget
va mostrato appena si apre la mappa del menu, e solo lì. Per questo si usa il
\textbf{Level Blueprint}, ``che abbiamo usato nella prima lezione'': ogni livello ha il
suo, ed è il posto giusto per la logica che appartiene a \emph{quella} mappa, ``perché
vogliamo che questa mappa sia dedicata apposta al nostro menu''. Si apre dalla toolbar
del Level Editor, dal pulsante dei Blueprint (``c'è il cubetto col più''),
\textbf{Open Level Blueprint}.

\paragraph{Il widget sullo schermo.} ``Ci interessa semplicemente usare il Begin Play,
perché vogliamo che si crei già da subito'':
\begin{enumerate}
  \item \textbf{Event BeginPlay} $\rightarrow$ \textbf{Create Widget}, con
        \textbf{Class} = WB\_Menu, ``come abbiamo visto l'altra volta'';
  \item \textbf{Add to Viewport}, con il \textbf{Return Value} del Create Widget nel
        Target. Il docente pensa per un momento di promuovere il widget a variabile, come
        per l'HUD, poi rinuncia: ``no, beh, non lo facciamo, tanto è il level blueprint''.
\end{enumerate}

\paragraph{Il personaggio di troppo.} In Play il menu compare, ma ``è molto strano:
abbiamo il nostro player dietro''. La mappa usa ancora la Game Mode del progetto, che fa
nascere il bot con i suoi controlli. Nel menu, invece, serve ``una pawn senza controlli,
senza niente'', e Unreal ce l'ha già: nelle \textbf{World Settings} della mappa, alla voce
\textbf{GameMode Override}, si sceglie la \textbf{Game Mode base} di Unreal (nella
trascrizione ``la game mode basica che esiste, la default''), che usa come Pawn la
\textbf{Default Pawn}. In Play, ``già è una cosa diversa: non c'è più il nostro character,
sparito, perché in questa mappa ho overridato la modalità Game Mode classica'', quella
``proprio base di Unreal, in cui non c'è character in terza persona''. ``Ma, come potete
vedere, non ho controllo'' (o meglio, come si vedrà, la Default Pawn si muove ancora:
il controllo vero si toglie nel Graph del widget, sezione~\ref{ue12:sec:menugraph}).

\paragraph{La camera fissa.} ``Andiamo a vedere anche un'altra cosa'': nel livello si
aggiunge un \textbf{Camera Actor} (``scrivete \emph{camera actor}, eccola qua''), che sarà
la vista fissa del menu. Insieme, una \textbf{Directional Light} (``e abbiamo messo il
sole, ragazzi''), regolabile a piacere, e un muretto, ``un po' più carino, così capiamo
meglio''. Per dire al gioco di guardare da quella camera, di nuovo nel Level Blueprint:
\begin{itemize}
  \item la \textbf{referenza} alla camera: selezionato il Camera Actor nel livello, nel
        Level Blueprint tasto destro, \textbf{Create a Reference to CameraActor};
  \item il nodo che cambia la vista. Qui il docente sbaglia strada in diretta: prova da
        \emph{Get Player Camera Manager}, non trova il nodo giusto e rimanda (``adesso non
        mi ricordo benissimo, lo vediamo un attimo dopo''). Dopo aver finito il Graph del
        menu, ritrova la soluzione: ``non è che bisogna prendere il camera manager, ma
        bisogna andare sul controller''. Quindi \textbf{Get Player Controller}
        $\rightarrow$ \textbf{Set View Target with Blend}, con la camera nel pin
        \textbf{New View Target}; il \emph{Blend Time} resta al default, ``non ci
        interessa''.
\end{itemize}
``Questo blocco ci permette di blendare una nuova camera'': ``gli diciamo: guarda, questa
è la tua nuova camera, da d'ora in poi'' (figura~\ref{ue12:fig:menugraph}, in alto).

\paragraph{Il Player Start fuori campo.} Ultimo ritocco, emerso aiutando gli studenti: la
Default Pawn ``è una palla'', e se nasce davanti alla camera si vede, o si vede il suo
riflesso. ``Prendete un \textbf{Player Start} e mettetelo fuori campo'', ``così almeno non vi
si vede la pallina''.

\paragraph{Approfondimento: Game Mode, Default Pawn e Player Start.} La \textbf{Game Mode}
è la classe che stabilisce le regole di una mappa e, fra le altre cose, quale
\emph{Pawn} far nascere per il giocatore (\emph{Default Pawn Class}) e con quale
\emph{Player Controller}. Il progetto ne ha una predefinita (in \emph{Project Settings},
\emph{Maps \& Modes}); il \textbf{GameMode Override} delle World Settings la sostituisce
\emph{solo per quella mappa}, ed è il modo standard di avere un menu senza personaggio.
La classe scelta a lezione, \emph{GameModeBase}, usa come Pawn la \textbf{DefaultPawn}:
una sfera con un movimento ``da spettatore'' (WASD e mouse, senza gravità), pensata per
fare prove. Il Pawn nasce in un \textbf{Player Start}; se la mappa non ne ha, nasce
all'origine o nella posizione della camera dell'editor. Un'alternativa ancora più pulita
per un menu è una Game Mode con la \emph{Default Pawn Class} impostata su \emph{None}: il
Player Controller esiste (e quindi esistono il cursore e gli input dell'interfaccia), ma
nessun Pawn nasce. Funzionamento generale dell'engine.

\paragraph{Approfondimento: Set View Target with Blend nel libro.} Romero e Sewell
descrivono il nodo nel capitolo 15 (\emph{Blueprints Tips}): ``è una funzione della classe
Player Controller'', ``usata per passare la vista di gioco fra camere diverse''; ``il
parametro di input New View Target è l'actor da usare come view target, di solito una
camera''. Il loro esempio è proprio in un Level Blueprint: quando il giocatore entra nella
stanza del tesoro, la vista passa alla camera della stanza. Questo spiega l'errore del
docente: il \emph{Player Camera Manager} gestisce la vista, ma la funzione esposta ai
Blueprint per cambiarla appartiene al \emph{Player Controller}. Il \textbf{Blend Time} a 0
fa uno stacco netto; un valore positivo fa una transizione morbida dalla vista precedente,
con una curva scelta in \emph{Blend Func}.

% ---------------------------------------------------------------------
\section{Il Graph del menu: cursore, input e bottoni}
\label{ue12:sec:menugraph}

\paragraph{Event Construct.} Si torna in WB\_Menu, scheda \textbf{Graph}: ``vediamo le
funzionalità del nostro menu quando clicchiamo le cose''. Dei tre eventi già presenti
``non ci serve il Pre Construct, non ci serve il Tick''. Serve l'\textbf{Event Construct}:
``vuol dire quando questo widget è costruito'', ``è un po' il Begin Play'', ``viene
differenziato dal nome, giusto per non confonderti''.

\paragraph{Il cursore e l'input dell'interfaccia.} In Play, per ora, ``mi posso muovere''
(la Default Pawn) e il mouse non si vede. Nell'Event Construct:
\begin{enumerate}
  \item \textbf{Get Player Controller} $\rightarrow$ \textbf{Set Show Mouse Cursor},
        con la spunta attiva: il cursore resta visibile anche cliccando nella finestra di
        gioco. ``Però, ovviamente, posso ancora muovermi'';
  \item \textbf{Set Input Mode UI Only} (``enable input UI'', dice prima il docente, poi
        trova il nodo), con lo stesso Player Controller: ``questo ci permette di
        utilizzare tranquillamente il nostro widget''. Il pin \emph{Widget to Focus} resta
        vuoto: ``no, tanto non ci serve''.
\end{enumerate}
In Play, ora, ``se noi proviamo a muovere lateralmente, non si muove più''. ``Mi
raccomando: questa cosa, quando noi cliccheremo un pulsante, dovrà essere disattivata e
rimessa in un'altra modalità.''

\paragraph{I colori di hover e pressione.} Con il cursore visibile si possono provare gli
stati dello Style. Nel Designer il docente cambia il colore di \textbf{Hovered}: ``voglio
che quando io ci passo sopra con il mouse diventi tipo arancione, sì, un bell'arancione''
(``lo vediamo in preview, è normale''), e lo copia, con il codice esadecimale, sull'altro
bottone. Il risultato, ammette, è ``molto Ikea''. Il \textbf{Pressed} di default diventa
grigio, ``si vede pochissimo'': lo si mette \textbf{rosso}, su entrambi i bottoni.

\paragraph{I bottoni.} Nel Graph, a sinistra, fra le variabili compaiono
\textbf{Button\_QuitGame} e \textbf{Button\_StartGame} (sono variabili perché i bottoni,
nel Designer, hanno \emph{Is Variable} attivo di default). Ogni bottone ha i suoi eventi,
``l'abbiamo visto anche l'altra volta'': \emph{On Pressed}, \emph{On Released},
\emph{On Hovered}, \emph{On Unhovered} e \textbf{On Clicked}. Si prende l'On Clicked di
entrambi.
\begin{itemize}
  \item \textbf{On Clicked (Button\_QuitGame)} $\rightarrow$ \textbf{Quit Game}: ``è
        piuttosto semplice, letteralmente Quit Game è già presente, di standard
        proprio''. (``Adesso non mi ricordo se ti esce proprio dall'editor'': si veda
        l'approfondimento più sotto.)
  \item \textbf{On Clicked (Button\_StartGame)}: prima \textbf{Get Player Controller}
        $\rightarrow$ \textbf{Set Input Mode Game Only}, ``molto importante'', ``in modo
        tale che ritornino gli input al gioco''; poi \textbf{Open Level (by Name)}, con il
        nome della mappa di gioco. Il docente lo cerca nel Content Browser (``era
        MLevel1? Ah, questa no, era M level 2?'') e trova la mappa del template,
        \textbf{ThirdPersonMap}; prova a copiarne il nome con F2 (``non lo fa, va bene''),
        e lo scrive a mano.
\end{itemize}
Tornati nella mappa del menu e salvato tutto, in Play: Play Game ``ci carica la nostra
bellissima mappina'' (figura~\ref{ue12:fig:menugraph}).

\begin{figure}[htbp]
\centering
\begin{tikzpicture}[font=\footnotesize,
  ue12node/.style={draw, thick, rounded corners=4pt, align=center, fill=gray!8,
                  minimum height=0.9cm, text width=2.3cm},
  ue12ev/.style={ue12node, fill=red!12},
  ue12pure/.style={ue12node, fill=green!10, rounded corners=8pt},
  ue12exec/.style={-{Stealth[length=2.5mm]}, very thick, black!80},
  ue12data/.style={-{Stealth[length=2mm]}, thick, blue!60!black}]
  \node[font=\scriptsize\bfseries, anchor=west] at (-1.3,1.0) {Level Blueprint di M\_Menu};
  \node[ue12ev, text width=1.7cm] (bp) at (0,0) {\textbf{Event BeginPlay}};
  \node[ue12node] (cw) at (2.8,0) {\textbf{Create Widget}\\{\scriptsize Class: WB\_Menu}};
  \node[ue12node, text width=2.5cm] (av) at (5.8,0) {\textbf{Add to Viewport}\\{\scriptsize Target: Return
    Value}};
  \node[ue12node, text width=2.6cm] (vt) at (8.9,0) {\textbf{Set View Target with
    Blend}\\{\scriptsize Blend Time 0}};
  \node[ue12pure] (gpc1) at (7.4,-1.8) {Get Player Controller};
  \node[ue12pure, fill=blue!10] (cam) at (10.2,-1.8) {CameraActor\\{\scriptsize
    (referenza)}};
  \draw[ue12exec] (bp) -- (cw); \draw[ue12exec] (cw) -- (av); \draw[ue12exec] (av) -- (vt);
  \draw[ue12data] ([yshift=-3mm]cw.east) -- ([yshift=-3mm]av.west);
  \draw[ue12data] (gpc1.north) -- node[left, font=\scriptsize]{Target} ([xshift=-6mm]vt.south);
  \draw[ue12data] (cam.north) -- node[right, font=\scriptsize]{New View Target}
    ([xshift=6mm]vt.south);
  % widget graph
  \node[font=\scriptsize\bfseries, anchor=west] at (-1.3,-3.2) {Graph di WB\_Menu};
  \node[ue12ev, text width=1.7cm] (con) at (0,-4.2) {\textbf{Event Construct}};
  \node[ue12node] (smc) at (2.8,-4.2) {\textbf{Set Show Mouse Cursor}\\{\scriptsize
    spunta attiva}};
  \node[ue12node] (ui) at (5.7,-4.2) {\textbf{Set Input Mode UI Only}};
  \node[ue12pure] (gpc2) at (4.2,-5.9) {Get Player Controller};
  \draw[ue12exec] (con) -- (smc); \draw[ue12exec] (smc) -- (ui);
  \draw[ue12data] (gpc2.north) -- (smc.south); \draw[ue12data] (gpc2.north) -- (ui.south);
  \node[ue12ev, text width=2.6cm] (oq) at (0.3,-7.6) {\textbf{On Clicked}\\{\scriptsize
    (Button\_QuitGame)}};
  \node[ue12node] (qg) at (3.6,-7.6) {\textbf{Quit Game}};
  \draw[ue12exec] (oq) -- (qg);
  \node[ue12ev, text width=2.6cm] (os) at (0.3,-9.3) {\textbf{On Clicked}\\{\scriptsize
    (Button\_StartGame)}};
  \node[ue12node, text width=2.7cm] (gm) at (3.8,-9.3) {\textbf{Set Input Mode Game Only}};
  \node[ue12node, text width=2.8cm] (ol) at (7.2,-9.3) {\textbf{Open Level (by
    Name)}\\{\scriptsize ThirdPersonMap}};
  \node[ue12pure] (gpc3) at (3.8,-11.0) {Get Player Controller};
  \draw[ue12exec] (os) -- (gm); \draw[ue12exec] (gm) -- (ol);
  \draw[ue12data] (gpc3) -- (gm);
\end{tikzpicture}
\caption{I grafi del menu. In alto, il Level Blueprint della mappa M\_Menu crea il widget,
lo mostra e passa la vista alla camera fissa. In basso, il widget prende il controllo del
cursore quando nasce, esce dal gioco con Quit e, con Play Game, restituisce l'input al
gioco e apre la mappa di gioco. Ricostruzione dalla lezione.}
\label{ue12:fig:menugraph}
\end{figure}

\paragraph{Approfondimento: il menu nel libro.} Romero e Sewell costruiscono un menu molto
simile nel capitolo 8 (\emph{Creating Constraints and Gameplay Objectives}), il
\emph{WinMenu} che compare a fine partita: due bottoni \emph{Btn restart} e \emph{Btn
quit} ancorati al centro, ciascuno con un Text dentro (``font size 48''), e gli eventi
aggiunti dal pulsante \emph{+} accanto a \emph{On Clicked}, in fondo ai Details del bottone.
Il bottone Restart usa \emph{Open Level (by Object Reference)}, con la mappa scelta da un
menu a tendina, seguito da \emph{Remove from Parent}, ``che rimuove il widget dalla
vista''; il bottone Quit usa \emph{Quit Game}, ``così che il giocatore possa chiudere il
gioco''. Il menu viene mostrato dal personaggio, in un Custom Event \emph{End Game}: Set
Game Paused, \emph{SET Show Mouse Cursor} dal Get Player Controller (``per permettere al
giocatore di riprendere il controllo del mouse''), Create Widget, Add to Viewport. Nel
capitolo 15 il libro elenca anche le tre modalità di input:
\begin{itemize}
  \item \textbf{Set Input Mode Game Only}: ``solo il Player Controller riceve gli eventi
        di input'';
  \item \textbf{Set Input Mode UI Only}: ``solo la UI riceve gli eventi di input'';
  \item \textbf{Set Input Mode Game and UI}: ``la UI ha la priorità nel gestire un evento
        di input, ma se la UI non lo gestisce, lo riceve il Player Controller''; l'esempio
        è un negozio, in cui si scelgono le opzioni con il mouse ma ci si può allontanare
        con le frecce.
\end{itemize}
Una differenza utile da notare: \emph{Open Level (by Object Reference)} non si rompe se la
mappa viene rinominata, mentre \emph{Open Level (by Name)}, usato a lezione, richiede il
nome esatto (si può scrivere il nome semplice, \emph{ThirdPersonMap}, o il percorso
completo, se esistono due mappe con lo stesso nome).

\paragraph{Approfondimento: Construct, Quit Game e il cursore nella mappa di gioco.}
Alcune nozioni generali, non della lezione. \textbf{Event Pre Construct} viene eseguito
anche nell'editor, ogni volta che il widget viene ridisegnato nel Designer: serve per
anteprime e impostazioni visive; \textbf{Event Construct} solo in gioco, quando il widget
viene aggiunto allo schermo (e di nuovo se lo si rimuove e lo si aggiunge). \textbf{Quit
Game}, durante il Play in Editor, \emph{termina la sessione di gioco} e riporta
all'editor, senza chiudere l'editor; in un gioco pacchettizzato chiude l'applicazione. Il
dubbio del docente si scioglie così. Infine, \emph{Open Level} distrugge il mondo corrente
con tutti i suoi actor, widget compresi, e crea un Player Controller nuovo: per questo,
nella mappa di gioco, il menu sparisce da solo e il cursore torna nascosto, anche senza un
\emph{Remove from Parent} o un \emph{Set Show Mouse Cursor} a falso. Il \emph{Set Input
Mode Game Only} prima dell'Open Level, raccomandato dal docente, è comunque una buona
abitudine: il caso in cui non è superfluo è quello di un menu che si chiude \emph{senza}
cambiare livello, come un menu di pausa (nel libro, capitolo 11, il \emph{Resume} fa
Remove from Parent, Set Show Mouse Cursor a falso e Set Game Paused a falso).

% ---------------------------------------------------------------------
\section{Il diorama, e i problemi degli studenti}
\label{ue12:sec:diorama}

\paragraph{Il diorama.} Con il menu funzionante, ``passiamo sopra'' all'aspetto. Sullo
sfondo, ``un piccolo diorama, una piccola box, in cui ci andiamo a mettere tutto quello che
vogliamo''. ``Volendo, potete fare tutto un menu interattivo'': gli ostacoli, i Blueprint
creati nel corso, i bottoni, la torretta. Il docente si limita ai modelli:
\begin{itemize}
  \item il \textbf{bot}: dal Content Browser, o la sua Skeletal Mesh, ``o li metto
        direttamente l'animazione in Idle'': trascinando nel livello un asset di animazione,
        Unreal crea un actor con la mesh che la riproduce;
  \item la \textbf{porta} modulare del compito della volta scorsa, \textbf{SM\_Modular\_\allowbreak Wall\_Door}
        (``wall door, eccola qua''), con le ante \textbf{SM\_Door};
  \item qualche ritocco alla \textbf{camera} e alla \textbf{luce}: l'intensità a 10 è
        ``enorme'', poi 5, poi valori molto più bassi (la trascrizione, qui, è quasi
        incomprensibile: si sentono ``0,01'' e ``1'').
\end{itemize}
``Questo caso è più una cosa artistica, ragazzi, quindi potete fare un po' come volete.''
Le funzioni di base, entrare nel gioco e uscire, ci sono. ``Questo è un menu che io di
solito faccio quando faccio i miei prototipi'', cioè ``quando devo semplicemente mostrare
qualche cosa a livello di gameplay''. Un menu delle \emph{opzioni} ``è già un discorso
molto più'' complesso, ``perché bisogna andare a vedere gli slider'': non si tratta, ``perché
tanto poi andiamo a fare qualcosa di più gameplayoso''.

\paragraph{I problemi dei presenti.} Diversi studenti hanno ancora il personaggio nel
menu.
\begin{itemize}
  \item \textbf{La camera.} Ripasso passo per passo: selezionare la camera nel livello
        (``proprio diventa arancione''), Open Level Blueprint, tasto destro, Create a
        Reference to CameraActor (``devi stare dentro il Level Blueprint''); Get Player
        Controller, Set View Target with Blend, la camera nel New View Target.
  \item \textbf{Il personaggio che nasce ancora.} ``Devi passare alla Game Mode
        default'': \textbf{World Settings}, \textbf{GameMode Override}. Uno studente non
        trova la finestra delle World Settings (si apre dal menu \emph{Window} dell'editor,
        se è chiusa). Poi il Player Start fuori dalla camera.
  \item \textbf{L'errore su Add to Viewport.} ``Mi dà un problema proprio su Add to
        Viewport'': ``perché il target di base è Self, quindi non so se devo cambiare il
        target''. Sì: il Target deve essere il \textbf{Return Value} del Create Widget.
        ``Probabilmente devi dare return value'' al target.
  \item \textbf{Il mouse che non si vede.} Va impostato nell'Event Construct del widget:
        Get Player Controller, Set Show Mouse Cursor, Set Input Mode UI Only. Lo studente
        aveva cancellato l'Event Construct (``basta che fai tasto destro'' e lo si ricrea)
        e non aveva messo la \textbf{spunta} nel Set Show Mouse Cursor (``ah, però non
        avevo messo il tick''): il nodo, senza spunta, imposta il cursore a \emph{falso}.
        ``Funziona? Sì.''
\end{itemize}
Chiuso il menu, si torna alla ThirdPersonMap.

% ---------------------------------------------------------------------
\section{Interfacce e Blueprint Component}
\label{ue12:sec:concetti}

``Andiamo sul nostro amico fidato'': il docente apre il suo bot (``ah no, questo è Boto:
dobbiamo andare su Boto2''), cioè BP\_Bot2. ``Oggi parleremo delle interfacce, ragazzi.''

\paragraph{L'interfaccia, un intermediario.} ``Immaginatevi che un'interfaccia è un
intermediario fra due o più Blueprint.'' Il docente mostra quella già pronta nel suo
progetto, nella cartella dei componenti dei Blueprint: un tipo particolare di Blueprint
che contiene solo tre \textbf{funzioni}, \emph{Interaction Start}, \emph{Interaction End}
e \emph{Interaction}. L'immagine: un \textbf{volo di linea}, ``un'interconnessione tra due
o più poli'', che ``trasporta molto spesso le stesse cose, passeggeri di un tipo,
passeggeri di un altro''. ``Quindi è un qualcosa che va a trasportare varie funzioni da una
parte all'altra.'' Nel caso di oggi l'interfaccia comunica tre messaggi: \emph{sto
iniziando un'interazione} (ti sto guardando), \emph{sto finendo un'interazione} (non ti
guardo più), \emph{sto interagendo} (ho premuto il tasto).

\paragraph{Il componente, una funzionalità isolata.} Il lavoro di \emph{vedere se c'è
qualcosa di interagibile} non lo fa il personaggio, ma un \textbf{Blueprint Component}:
``dei componenti che noi possiamo aggiungere all'interno dei nostri actor per dare
funzionalità in modo isolato''. Per esempio: ``questo actor può tracciare degli
interactable''; ``questo actor ha la vita''. La vita, fatta la volta scorsa \emph{dentro}
il personaggio, si potrebbe spostare in un componente, ``perché magari abbiamo anche i
nemici che hanno la vita'', o ``degli oggetti che si possono rompere''. Lo scopo è
``isolare sempre più possibile le funzionalità, e quindi essere anche modulari''. Infatti
il BP\_Bot2 del docente, a parte un input nuovo, ``non ha niente di nuovo'': si gira verso
l'oggetto e suona un Montage, ma ``non ha niente per tracciare'': il trace è tutto nel
componente.

\paragraph{Actor Component o Scene Component.} Si crea una \textbf{Blueprint Class} e, fra
le classi comuni, non un Actor ma un componente. Ce ne sono due:
\begin{itemize}
  \item l'\textbf{Actor Component} ``mette semplicemente la funzionalità'';
  \item lo \textbf{Scene Component} ``è un Actor Component che ha anche una scena'',
        cioè un \textbf{Transform}: ``non so se vi ricordate la Scene, quando l'abbiamo
        messa: il punto vuoto''. ``Almeno si può spostare nello spazio.''
\end{itemize}
Qui serve il secondo, perché il trace è ``un raggio che da quel punto parte e in un punto
finisce'', e il punto di partenza si deve poter spostare. Il docente lo mostra nel suo bot:
spostando il componente nel Blueprint del personaggio, in Play il raggio parte dal nuovo
punto; lo rimette ``al centro della testa, più o meno''. Il nuovo Blueprint:
\textbf{BPC\_RayInteract} (BPC sta per \emph{Blueprint Component}; il docente, che ne ha già
uno, chiama il suo BPC\_RayInteract1).

\paragraph{Il componente aperto.} ``Come potete vedere, presenta delle similitudini con
tutti i Blueprint.'' ``L'unica cosa: non ha una viewport, perché non ha una forma fisica'',
``tra virgolette, perché è una scena, ma è vuota''. ``E ricordiamoci che è un componente:
quindi sarà qualcosa che deve essere per forza attaccato a qualcos'altro.''

\paragraph{Che cos'è una funzione.} Prima di scrivere la logica, una domanda: chi ha già
sentito parlare di \textbf{funzioni} in informatica? Risponde \textbf{Federico}, che ha
studiato da programmatore: ``la funzione è un insieme di istruzioni: chiami la funzione e
fai quelle istruzioni''. ``Esattamente.'' Il docente indica una catena di nodi nel bot: ``gli
eventi non possono essere messi come funzioni, ma questo qua può essere una funzione''. È
``un modo per ordinare anche il lavoro in varie tipologie di funzionalità, che possono
essere magari trasportabili da un punto all'altro'': ``sempre questo senso di modularità''.
(``Grande Federico, sempre miglior assistente.'')

\begin{figure}[htbp]
\centering
\begin{tikzpicture}[font=\footnotesize,
  ue12box/.style={draw, thick, rounded corners=5pt, align=center, inner sep=5pt},
  ue12msg/.style={-{Stealth[length=2.5mm]}, very thick, violet!70!black},
  ue12own/.style={-{Stealth[length=2.2mm]}, thick, dashed, black!70}]
  \node[ue12box, fill=gray!8, minimum width=4.4cm, minimum height=4.6cm] (bot) at (0,-0.3) {};
  \node[anchor=north, font=\bfseries] at (bot.north) {BP\_Bot2};
  \node[ue12box, fill=blue!10, text width=3.6cm] (cam) at (0,0.8) {Camera Component};
  \node[ue12box, fill=orange!18, text width=3.6cm] (comp) at (0,-1.4)
    {\textbf{BPC\_RayInteract}\\{\scriptsize Scene Component: ogni frame una Sphere Trace
    lungo la camera; variabile Interactable}};
  \draw[ue12own] (comp.north) -- node[right, font=\scriptsize, text width=1.9cm,
    align=left]{Get Owner,\\Get Component by Class} (cam.south);
  \node[ue12box, fill=violet!10, text width=3.2cm] (bpi) at (6.6,-0.3)
    {\textbf{BPI\_InteractionRay}\\{\scriptsize Blueprint Interface}\\[2pt]{\scriptsize
    Interaction Start\\Interaction End\\Interaction}};
  \node[ue12box, fill=green!10, text width=3.3cm] (obj1) at (11.6,1.0)
    {\textbf{BP\_InteractableTest}\\{\scriptsize implementa la BPI: Print String, widget
    ``E'', luce}};
  \node[ue12box, fill=green!5, text width=3.3cm] (obj2) at (11.6,-1.6) {{\scriptsize un
    altro oggetto}\\{\scriptsize (bottone, ascensore\dots) con la stessa BPI e un'altra
    risposta}};
  \draw[ue12msg] (bot.east |- cam) -- node[above, font=\scriptsize, text width=2.4cm,
    align=center]{Interaction\\(tasto E)} (bpi.west |- cam);
  \draw[ue12msg] (comp.east) -- node[below, font=\scriptsize, text width=2.4cm,
    align=center]{Start / End\\(guardare, smettere)} (bpi.west |- comp);
  \draw[ue12msg] (bpi.east) -- (obj1.west);
  \draw[ue12msg] (bpi.east) -- (obj2.west);
\end{tikzpicture}
\caption{I ruoli del sistema di interazione. Il componente, attaccato al personaggio, trova
la camera del suo proprietario, traccia ciò che il giocatore guarda e invia i messaggi di
inizio e fine; il personaggio invia il messaggio Interaction quando si preme E. I messaggi
viaggiano attraverso l'interfaccia: chi li riceve decide che cosa fare, e un actor che non
la implementa li ignora. Schema ricostruito dalla lezione.}
\label{ue12:fig:ruoli}
\end{figure}

\paragraph{Approfondimento: i componenti nel libro.} Romero e Sewell dedicano ai componenti
la seconda metà del capitolo 18 (\emph{Creating Blueprint Libraries and Components}).
``Quando creiamo un Actor Blueprint aggiungiamo spesso componenti con una funzionalità
incapsulata, pronta all'uso, come Projectile Movement, Static Mesh e Collision''; ``possiamo
anche creare i nostri Actor Component con i Blueprint''. Nel pannello \emph{Pick Parent
Class} le due classi sono fra quelle \emph{comuni}: ``lo Scene Component è una classe figlia
dell'Actor Component che ha la struttura Transform (location, rotation e scale). Grazie al
Transform, uno Scene Component può essere attaccato a un altro Scene Component.'' E, come a
lezione: ``quando si scrive la logica di un componente nell'Event Graph, si può ottenere la
referenza all'actor che usa il componente con il nodo \textbf{Get Owner}''. I due esempi del
libro mostrano bene la differenza: \emph{BP\_ExpLevelComp} è un Actor Component che gestisce
punti esperienza e livelli (non ha bisogno di una posizione), mentre \emph{BP\_CircularMovComp}
è uno Scene Component che ruota attorno all'actor e a cui si attacca una Static Mesh, per
simulare uno scudo rotante: ``lo Scene Component è usato come posizione di riferimento per
altri componenti''. Il componente della lezione è dello stesso tipo: la sua posizione è il
punto di partenza del raggio.

% ---------------------------------------------------------------------
\section{Il componente BPC\_RayInteract: la Sphere Trace}
\label{ue12:sec:componente}

\paragraph{La funzione LookForInteractions.} Nel componente si crea una nuova funzione
(``eccolo che vi si crea questo pannellino violaceo''), \textbf{LookForInteractions}.
Dentro, una \textbf{Sphere Trace By Channel}. ``Seguitemi per bene.''
\begin{itemize}
  \item \textbf{Radius}: \textbf{5}, ``perché se no non vi apparirà mai''. Con raggio zero
        la Sphere Trace degenera in un raggio sottilissimo (e il debug quasi non si vede);
        5 unità bastano a renderla visibile e un po' più tollerante.
  \item \textbf{Start}: \textbf{Get World Location} del componente stesso, ``la posizione
        nel mondo del nostro Scene Component: dovrà partire da qui''.
  \item \textbf{End}: da costruire, ed è la parte con il trucco.
\end{itemize}

\paragraph{Prima la camera: Get Owner e Get Component by Class.} Il raggio deve seguire lo
sguardo del giocatore, cioè la \emph{rotazione della camera}. Il componente, però, non sa
nulla della camera del personaggio: deve trovarla. Nell'\textbf{Event Graph} del
componente, all'\textbf{Event BeginPlay}:
\begin{enumerate}
  \item \textbf{Get Owner}: ``la referenza del proprietario di questo Scene Component'';
  \item dal proprietario, \textbf{Get Component by Class} (``non \emph{components},
        \emph{component}, perché se no vi dà un array''), con \textbf{Component Class} =
        \textbf{Camera Component}: ``sarà il nostro punto principale da cui andremo a fare il
        nostro trace'';
  \item il risultato si promuove a variabile, \textbf{CameraComponent}, e si collega tutto
        al Begin Play.
\end{enumerate}
``Perdonate, sono 30.000 componenti: un componente su componente su componente.'' In
sintesi: ``in Begin Play io prendo il mio proprietario, prendo la sua camera, mi salvo la
referenza alla sua camera, e poi la useremo all'interno del nostro componente''.

\paragraph{La camera, se c'è.} Il componente deve funzionare anche su un actor
\emph{senza} camera: ``se il nostro proprietario non ha una camera, allora prendi te stesso
come pivot''. Quindi un \textbf{Select} fra due oggetti:
\begin{itemize}
  \item l'indice è un booleano, \textbf{Is Valid} (``la F verde'', la versione \emph{pura}
        della funzione) sulla variabile CameraComponent: ``se il componente camera esiste,
        quindi il nostro actor ha effettivamente una camera, e quindi questa variabile
        CameraComponent ha qualcosa dentro'';
  \item su \textbf{True}, CameraComponent: ``usami la camera come punto di riferimento'';
  \item su \textbf{False}, \textbf{Self}: ``se no, usi te stesso''.
\end{itemize}
Il Select era già comparso con la torretta, per scegliere fra le due bocche di fuoco: lì
con un \emph{Random Bool}, qui con l'Is Valid. ``So che è un po' complesso.''

\paragraph{La ``mandragata'': il tipo del Select.} Il problema è il \emph{tipo} delle
opzioni. ``Il Select praticamente c'è un piccolo problema'': il nodo è jolly (\emph{wildcard})
e assume il tipo della \emph{prima} cosa che gli si collega. Collegando prima la camera, il
Select diventa un Select di \emph{Camera Component Object Reference}, e Self (il componente
BPC\_RayInteract) non entra più; collegando prima Self, il Select diventa un Select di
riferimenti a BPC\_RayInteract, e la camera non entra. Entrambi però \emph{sono} Scene
Component: la camera è uno Scene Component, e lo è anche il nostro componente. Il trucco
(``qui bisogna fare una mandragata, ragazzi''):
\begin{enumerate}
  \item si crea una variabile qualunque di tipo \textbf{Scene Component Object Reference}
        (il docente la chiama \emph{Scene}: ``e poi non ci servirà a niente'');
  \item se ne mette un Get nel Select: ora le opzioni sono di tipo \emph{Scene Component};
  \item si \textbf{cancella} il Get e si collegano CameraComponent e Self, che sono entrambi
        accettati.
\end{enumerate}
``Ecco fatto che abbiamo praticamente fregato Unreal.''

\paragraph{La domanda di Filippo.} ``Non ho capito perché abbiamo creato la variabile Scene
prima.'' Il docente rifà la prova davanti a tutti: un Select nuovo, il cursore sul pin, la
scritta ``Scene Component Object Reference''; poi Self in un Select vuoto, e il nodo diventa
``Self Object Reference'', ``mi farà fare solamente un select fra due self, non due Scene
Component''. ``Per fregare Unreal'' serve ``una variabile che appartiene generalmente allo
Scene Component''. Uno studente prova a trasformare la variabile in Self: ``non è che te lo
trasforma: lo devi cancellare, e poi tasto destro, scrivi \emph{self}''. Filippo insiste: e
se avessimo collegato subito la camera? Prova del docente: ``Select camera component, camera
component object reference, e io, se provo ad attaccare Self, non me lo attacca''. ``Ok,
capito.''

\paragraph{Il punto finale.} Dal risultato del Select, ``quello che abbiamo sempre fatto'':
\begin{enumerate}
  \item \textbf{Get Forward Vector} (della camera, o del componente stesso);
  \item moltiplicato per una distanza: il secondo pin del Multiply si converte in float
        (``tasto destro, mi raccomando: \emph{To Float}''), si promuove a variabile,
        \textbf{DistanceTrace}, con valore \textbf{300};
  \item il prodotto si \textbf{somma} (Add) allo Start, e il risultato va nell'\textbf{End}.
\end{enumerate}

\paragraph{Perché lo Start è il componente.} Una studentessa chiede perché lo Start sia la
World Location: ``in relazione a dove attacchiamo il componente?'' ``Esatto.'' ``Il punto
iniziale deve essere sempre il punto, il Transform, del nostro componente''; ``il punto che
deve diversificare è la camera''. Il docente lo mostra nel bot: guardando in basso il
raggio va in basso, guardando di lato va di lato, ma parte sempre dal componente, ``così
magari posso posizionarlo avanti, posizionarlo indietro, posizionarlo qui, posizionarlo
là'' (figura~\ref{ue12:fig:look}).

\paragraph{Il componente nel personaggio.} Nel BP\_Bot2 (``a bot2, in questo caso, è il
vostro character''), \textbf{Add} e il nome del componente: BPC\_RayInteract compare fra i
componenti del personaggio. Per vederlo si imposta il \textbf{Draw Debug Type} della Sphere
Trace (``lo metto a \emph{For One Frame}''). In Play, però, non succede niente: ``ah,
certo, ecco perché: non facciamo LookForInteractions''. La funzione va chiamata
nell'\textbf{Event Tick} del componente (``ah, nel tick, infatti''), così il controllo
avviene a ogni frame. Con il debug persistente (\emph{Persistent}) il raggio lascia una scia
``gigantesca''; con \emph{For Duration} resta qualche istante; con \emph{For One Frame}
``rimane un po' più carino'', e si vede solo il raggio attuale. (La trascrizione, qui, è
piena di prove e di parole storpiate: ``la duration vado a activate, troppo corto'', ``Lugano'',
``Adobe''; anche una confusione momentanea fra il bot e il bot2 del docente.)

\paragraph{Approfondimento: Sphere Trace, Tick e costo.} Una \textbf{Sphere Trace} è un
\emph{line trace} ``spesso'': fa scorrere una sfera del raggio indicato dallo Start
all'End e restituisce il primo oggetto che la sfera tocca nel canale scelto. Rispetto a un
raggio sottile perdona di più la mira (con 5 unità di raggio, circa 5 cm, conta anche
sfiorare il bordo di un oggetto), ed è il motivo per cui molti sistemi di interazione la
preferiscono. Farla nel \textbf{Tick}, una volta per frame, costa poco: una singola trace di
300 unità è un'operazione leggera, e il sistema la esegue solo per il personaggio del
giocatore. Se gli actor con il componente fossero molti (per esempio dei nemici), conviene
ridurre la frequenza con un timer (come quelli della stamina, a 0,1 secondi) o disattivare
il Tick del componente quando non serve (\emph{Set Component Tick Enabled}). Una Sphere
Trace può stare senza problemi dentro una funzione, perché non ha uno stato interno da
ricordare fra una chiamata e l'altra; il Do Once e il Flip Flop, invece, sì (si veda più
avanti). Funzionamento generale dell'engine.

% ---------------------------------------------------------------------
\section{Il Trace Channel e la Blueprint Interface}
\label{ue12:sec:bpi}

\paragraph{Un canale per le interazioni.} ``Ha colpito qualcosa? Sì, no.'' Prima di
leggere il risultato serve un canale dedicato, ``che ci permette di identificare gli
oggetti che vogliamo che siano interagibili'': ``se no, praticamente, interagirà con
tutto''. ``Come abbiamo visto anche nella nostra torretta'': \textbf{Project Settings},
\emph{Collision}, \textbf{New Trace Channel}, nome \textbf{InteractionTrace}, risposta di
default \textbf{Ignore}. Il docente, che ce l'ha già, non lo rifà. Nella Sphere Trace, il
\textbf{Trace Channel} diventa InteractionTrace: ``mi raccomando, se non vi appare, fate
compila, perché ha bisogno un attimo di compilare il nuovo Trace Channel''.

\paragraph{Che cosa abbiamo colpito.} Dopo la trace, un \textbf{Branch} sul suo Return
Value (``ha colpito qualcosa?'') e un \textbf{Break Hit Result} sull'Out Hit, da cui, ``vi
ricordate, no?'', escono l'actor colpito, il Physical Material ``e tutto''. Il piano:
``useremo l'interfaccia, la Blueprint Interface, per andare a comunicare con questo actor'':
``ha questo porto in cui può accogliere questa tipologia di interazione? Sì? Benissimo,
allora fai l'interazione''. Dall'\textbf{Hit Actor} si chiede \textbf{Does Object
Implement Interface}: ``ehi, questo actor ce l'ha l'interface?'' Ma l'interfaccia non
esiste ancora: ``ci fermiamo qua per adesso, creiamo l'interface, e ci fermiamo''. (Il
docente parla qui anche di promuovere l'Hit Actor a variabile; la variabile di appoggio
vera e propria, \emph{Interactable}, si crea dopo la pausa, sezione~\ref{ue12:sec:logica}.)

\paragraph{La Blueprint Interface.} Nella cartella \emph{Blueprints}, \emph{Components},
tasto destro, \textbf{Blueprint}, ``non Blueprint Class questo, ma Blueprint questo'': il
sottomenu offre \emph{Blueprint Class}, \emph{Blueprint Function Library} e
\textbf{Blueprint Interface}. Il nome: \textbf{BPI\_InteractionRay}. Aperta:
\begin{itemize}
  \item il grafo è \textbf{in sola lettura}: ``non possiamo far niente'', ``cannot create
        new nodes in read only graph''. Un'interfaccia contiene solo \emph{dichiarazioni};
  \item si possono aggiungere \textbf{funzioni}: la prima c'è già e si rinomina (con F2),
        \textbf{Interaction Start}; con \emph{Add Function} le altre due,
        \textbf{Interaction End} e \textbf{Interaction}.
\end{itemize}
``Abbiamo messo che in questo cargo, questo trasporto, ci stanno Interaction Start,
Interaction End e Interaction: quindi queste funzioni possono essere passate tra di loro.''

\paragraph{Il collegamento.} Nel componente, il pin \textbf{Interface} di Does Object
Implement Interface diventa \mbox{BPI\_}\allowbreak InteractionRay. ``Se quel qualcosa che ho colpito ha questa
interfaccia, allora andremo a fare tutto il procedimento di interazione.'' ``Ci sono domande?
Mi sembra di capire no.'' Sono le 17:26: pausa fino alle 17:55, ``diciotto''.

\paragraph{Approfondimento: la Blueprint Interface nel libro.} Romero e Sewell presentano le
interfacce nel capitolo 16 (\emph{Introduction to VR Development}), nel paragrafo
\emph{Blueprint Communication using interfaces}: ``una Blueprint Interface è un tipo speciale
di Blueprint che contiene solo nomi di funzioni e parametri. Si usa per permettere la
comunicazione fra tipi diversi di Blueprint.'' Si crea dal pulsante \emph{Add} del Content
Browser, sottomenu \emph{Blueprints}, \emph{Blueprint Interface}. L'esempio è il template VR:
l'interfaccia \emph{VRInteraction BPI} è aggiunta alla pistola dalle \emph{Class Settings},
categoria \emph{Interfaces}, e la pistola implementa la funzione \emph{Trigger Pressed}:
``siccome la funzione non ha parametri di output, è implementata come un \emph{evento}''. Il
Pawn VR chiama Trigger Pressed sull'actor che tiene in mano, ``e se l'actor ha implementato
l'interfaccia VRInteraction BPI, Trigger Pressed viene eseguito. Non succede niente se l'actor
non l'ha implementata.'' È esattamente lo schema della lezione: il personaggio non sa
\emph{che cosa} sta guardando (un cubo, un bottone, un ascensore) e non deve fare un
\emph{cast} a nessuna classe; manda il messaggio, e ciascun oggetto risponde a modo suo.
Questo è il vantaggio delle interfacce rispetto al cast e all'ereditarietà viste con
BP\_Activable: classi che non hanno niente in comune (un actor, un personaggio, un oggetto
fisico) possono condividere la stessa ``porta''.

% ---------------------------------------------------------------------
\section{La pausa: la stamina di Tommaso}
\label{ue12:sec:pausa}

\paragraph{Il sintomo.} In pausa, come promesso, il problema rimasto aperto la volta
precedente: \textbf{Tommaso} tiene premuto Shift, la barra della stamina non scende, e
``appena mollo, si vede''. Il docente ripercorre il codice condiviso: l'evento StopRunning,
il Clear and Invalidate, il Set Timer by Event e la variabile TimeStamina (``quanto sta?
0,92''; nella trascrizione il valore è questo, ma va preso con cautela: nella lezione
precedente il valore indicato era 0,05), RegainStamina, il Clamp con CurrentStamina e
MaxStamina, l'InitializeStaminaBar nel Begin Play. ``Sembra tutto a posto.''

\paragraph{Il debug.} Il metodo è quello della volta scorsa: un \textbf{Print String} con la
CurrentStamina (``attaccalo al viola'') dopo l'evento di aggiornamento. I valori scendono
correttamente (``500, 500, 300'' dice la trascrizione, forse leggendo la velocità), poi
risalgono. Ma a schermo ci sono ``troppi Print String'': si cercano e si tolgono quelli
vecchi, uno per volta (quello della torretta, uno della piattaforma, ``forse i mulini?'',
``forse il collectable''; ``nessuno è perfetto''). Il comportamento: la barra
\emph{ricarica} visibilmente, ma \emph{non scende}.

\paragraph{La causa.} ``Ah, ecco perché: non chiami l'evento Call Update Stamina.'' Nella
catena di \textbf{StaminaDeduction}, dopo il Set di CurrentStamina, mancava la
\textbf{Call E\_UpdateStamina}: il valore calava, ma nessuno lo comunicava alla barra. La
Call c'era solo in RegainStamina, che parte al rilascio del tasto; per questo la barra si
aggiornava ``solo quando mollo''. Il rimedio: la Call anche in StaminaDeduction, con
CurrentStamina e MaxStamina collegati. ``Perfetto. Ottimo.''

\paragraph{Il rumore.} Il resto della pausa, nella trascrizione, è il blocco di testo
ripetuto e il frammento di video già segnalati nell'introduzione. Alla ripresa, uno studente
chiede aiuto per BPI\_InteractionRay, che dopo un riavvio di Unreal si apre ``male'': il
grafo in sola lettura non si vede. Il docente prova i menu \emph{Window} (Bookmarks, Find
Results, Compiler Results\dots), carica per sbaglio il layout di default dell'editor
(``mi sono caricato il bellissimo layout di Unreal''), e alla fine trova la via più
semplice: ``fai doppio clic sulla funzione'' nel pannello \emph{My Blueprint}, e il suo
grafo si apre. ``Tanto non ti serve'': nell'interfaccia non si scrive nulla.

% ---------------------------------------------------------------------
\section{La logica: la variabile Interactable}
\label{ue12:sec:logica}

``Vai, allora finiamo questo bellissimo sistema di interazioni.'' Si torna in
BPC\_RayInteract, nella funzione LookForInteractions. Il ramo di partenza è quello in cui
la trace ha colpito qualcosa e l'actor colpito implementa BPI\_InteractionRay.

\paragraph{La variabile d'appoggio.} Serve una variabile di tipo \textbf{Actor}
(\emph{Object Reference}), \textbf{Interactable}: ``sarà la nostra variabile d'appoggio'',
cioè l'oggetto interagibile che il giocatore sta guardando \emph{in questo momento}, o
niente. Il suo scopo è evitare di ripetere il lavoro a ogni frame: ``se l'actor che io ho
intercettato con il mio trace è lo stesso di prima, allora non voglio che mi rifai alcune
sezioni del codice''; ``se è lo stesso, non lo svuoto''.

\paragraph{È un oggetto nuovo?} Un secondo \textbf{Branch}, con la condizione \textbf{Hit
Actor} \textbf{diverso da} Interactable: il nodo \emph{Not Equal} (``diverso si fa con il
punto esclamativo e l'uguale''). ``Il punto esclamativo in informatica vuol dire che questo è
un condizionale che ci dice: guarda, questo è diverso da questo, e ci restituisce un
booleano. Se è diverso ci dà True, se sono uguali ci dà False.'' Sul False non si fa nulla:
è lo stesso oggetto di prima.

\paragraph{I messaggi dell'interfaccia.} ``Adesso, importante'': scrivendo
\emph{interaction} nel menu dei nodi, compaiono già le funzioni di BPI\_InteractionRay,
``senza fare niente'': ``ho già la possibilità di poter comunicare e messaggiare con
l'interfaccia''. Sono i nodi \textbf{(Message)}, con l'icona della busta, che si possono
chiamare su qualunque actor. Servono \textbf{Interaction End} e \textbf{Interaction Start}.

\paragraph{Il ramo True: chiudere il vecchio, aprire il nuovo.} Sul True del diverso:
\begin{enumerate}
  \item un \textbf{Validated Get} di Interactable (si mette il Get, tasto destro,
        \emph{Convert to Validated Get}): ``questa variabile Interactable è valida? C'è
        qualcosa dentro?'';
  \item \textbf{Is Valid}: c'era già un oggetto, quindi prima si chiude la vecchia
        interazione, \textbf{Interaction End} con target il \emph{vecchio} Interactable;
        poi \textbf{Set Interactable} = Hit Actor; poi \textbf{Interaction Start} con
        target il \emph{nuovo} Interactable;
  \item \textbf{Is Not Valid}: ``la variabile è vuota'', ``non è mai iniziata
        un'interazione precedente''; si salta l'End e si fa direttamente il Set
        Interactable = Hit Actor e l'\textbf{Interaction Start}.
\end{enumerate}

\paragraph{Il ramo False: non guardo niente.} Se la trace non colpisce nulla, \emph{oppure}
colpisce un actor senza l'interfaccia (``il muro, ad esempio''), il giocatore ha distolto lo
sguardo: ``allora ti finisco l'interazione''. Entrambi i False (quello del Branch sul Return
Value e quello di Does Object Implement Interface) portano a un altro Validated Get di
Interactable:
\begin{itemize}
  \item \textbf{Is Valid}: \textbf{Interaction End} sul vecchio Interactable, poi
        \textbf{Set Interactable} senza collegare nulla al pin, cioè lo si
        \textbf{svuota};
  \item \textbf{Is Not Valid}: non c'era nessuna interazione da chiudere; per sicurezza
        lo si svuota comunque (``semplicemente dirgli: guarda, svuotati'').
\end{itemize}
``Quindi mettere il Set così è come mettere uguale a null?'' ``È nulla, proprio. La
annulliamo: la mettiamo vuota, e quindi diventerà non valida.'' Il caso del primo ramo,
il passaggio diretto da un oggetto all'altro, serve quando due interagibili sono vicini:
``lui finisce con quello precedente e inizia con il nuovo''; di solito, invece, gli oggetti
sono distanti, la variabile si svuota e poi si riempie di nuovo
(figure~\ref{ue12:fig:look} e~\ref{ue12:fig:stati}).

\paragraph{Il problema di tipo.} Uno studente non riesce a collegare l'Hit Actor al Set
Interactable: ``mi dice che non sono compatibili, perché uno è un riferimento oggetto e
l'altro è un componente di scena''. La variabile Interactable era stata creata con il tipo
sbagliato (ed era cambiata per un clic involontario in \emph{Geometry}, ``che non ho
cliccato''). Rimedio: nei Details della variabile, il tipo \textbf{Actor}, \textbf{Object
Reference} (``non quell'Actor: tipo oggetto, riferimento oggetto''), confermare la
modifica del tipo, poi staccare e riattaccare i fili rossi dei nodi che usavano la
variabile, ``era semplicemente sparito''.

\paragraph{Il Tick.} ``Importante, ragazzi: far sì che questa funzione venga fatta su ogni
frame.'' Nell'Event Graph del componente, \textbf{Event Tick} $\rightarrow$
LookForInteractions (se non lo si è già fatto): ``a ogni frame farà il controllo: c'ho
qualcosa di interagibile davanti a me, dove sto puntando la camera? Sì o no?''.

\begin{figure}[htbp]
\centering
\begin{tikzpicture}[font=\footnotesize,
  ue12node/.style={draw, thick, rounded corners=4pt, align=center, fill=gray!8,
                  minimum height=0.85cm, text width=2.4cm},
  ue12dec/.style={draw, thick, diamond, aspect=2.2, align=center, fill=yellow!15,
                  inner sep=1pt, text width=2.0cm, font=\scriptsize},
  ue12msgn/.style={ue12node, fill=violet!12},
  ue12exec/.style={-{Stealth[length=2.5mm]}, very thick, black!80}]
  \node[ue12node, fill=red!12, text width=3.4cm] (st) at (0,0) {\textbf{Sphere Trace By
    Channel}\\{\scriptsize Start: posizione del componente; End: Start + Forward (camera o
    self) $\times$ 300; Radius 5; InteractionTrace}};
  \node[ue12dec] (hit) at (0,-2.2) {colpito qualcosa?};
  \node[ue12dec] (impl) at (0,-4.4) {Hit Actor implementa la BPI?};
  \node[ue12dec] (diff) at (0,-6.6) {Hit Actor $\neq$ Interactable?};
  \node[ue12dec] (val1) at (0,-8.9) {Interactable valido?};
  \node[ue12msgn] (end1) at (-3.6,-10.8) {\textbf{Interaction End}\\{\scriptsize sul vecchio}};
  \node[ue12node] (set1) at (0,-12.4) {\textbf{Set Interactable}\\{\scriptsize = Hit Actor}};
  \node[ue12msgn] (start) at (0,-14.0) {\textbf{Interaction Start}\\{\scriptsize sul nuovo}};
  \node[ue12node, fill=gray!3, text width=2.0cm] (none) at (3.7,-6.6) {niente\\{\scriptsize
    (stesso oggetto)}};
  \node[ue12dec] (val2) at (6.4,-3.3) {Interactable valido?};
  \node[ue12msgn] (end2) at (6.4,-5.6) {\textbf{Interaction End}\\{\scriptsize sul vecchio}};
  \node[ue12node] (clr) at (6.4,-7.6) {\textbf{Set Interactable}\\{\scriptsize vuoto (null)}};
  \draw[ue12exec] (st) -- (hit);
  \draw[ue12exec] (hit) -- node[left, font=\scriptsize]{True} (impl);
  \draw[ue12exec] (impl) -- node[left, font=\scriptsize]{True} (diff);
  \draw[ue12exec] (diff) -- node[left, font=\scriptsize]{True} (val1);
  \draw[ue12exec] (diff) -- node[above, font=\scriptsize]{False} (none);
  \draw[ue12exec] (val1.west) -| node[pos=0.25, above, font=\scriptsize]{Is Valid} (end1.north);
  \draw[ue12exec] (end1.south) |- (set1.west);
  \draw[ue12exec] (val1) -- node[right, font=\scriptsize]{Is Not Valid} (set1);
  \draw[ue12exec] (set1) -- (start);
  \draw[ue12exec] (hit.east) -| node[pos=0.2, above, font=\scriptsize]{False} (val2.north);
  \draw[ue12exec] (impl.east) -- node[above, font=\scriptsize]{False} ++(2.5,0) |-
    (val2.west);
  \draw[ue12exec] (val2) -- node[right, font=\scriptsize]{Is Valid} (end2);
  \draw[ue12exec] (end2) -- (clr);
  \draw[ue12exec] (val2.east) -- ++(0.6,0) |- node[pos=0.3, right, font=\scriptsize, text
    width=1.0cm, align=left]{Is Not Valid} (clr.east);
\end{tikzpicture}
\caption{La funzione LookForInteractions del componente, chiamata nell'Event Tick. Se il
giocatore guarda un nuovo interagibile, si chiude l'interazione con il vecchio (se c'era)
e si apre con il nuovo; se guarda un oggetto non interagibile, o niente, si chiude
l'interazione in corso e si svuota la variabile. I nodi viola sono messaggi
dell'interfaccia. Ricostruzione dalla lezione.}
\label{ue12:fig:look}
\end{figure}

\begin{figure}[htbp]
\centering
\begin{tikzpicture}[font=\footnotesize,
  ue12state/.style={draw, very thick, rounded corners=10pt, align=center,
                    minimum height=1.3cm, text width=3.4cm},
  ue12tr/.style={-{Stealth[length=2.5mm]}, very thick}]
  \node[ue12state, fill=gray!10] (vuoto) at (0,0) {\textbf{Interactable vuoto}\\{\scriptsize
    non guardo niente di interagibile}};
  \node[ue12state, fill=green!12] (a) at (7.0,0) {\textbf{Interactable = A}\\{\scriptsize
    widget ``E'' di A visibile}};
  \node[ue12state, fill=green!6] (b) at (7.0,-3.4) {\textbf{Interactable = B}\\{\scriptsize
    widget ``E'' di B visibile}};
  \draw[ue12tr] ([yshift=3mm]vuoto.east) -- node[above, font=\scriptsize]{guardo A: Start(A)}
    ([yshift=3mm]a.west);
  \draw[ue12tr] ([yshift=-3mm]a.west) -- node[below, font=\scriptsize, text width=3.2cm,
    align=center]{guardo il muro o niente: End(A), svuota} ([yshift=-3mm]vuoto.east);
  \draw[ue12tr] (a.south) -- node[right, font=\scriptsize, text width=3.4cm]{passo
    direttamente a B: End(A), Set, Start(B)} (b.north);
  \draw[ue12tr] (b.west) -| node[pos=0.3, below, font=\scriptsize]{guardo il muro: End(B),
    svuota} (vuoto.south);
  \draw[ue12tr] (a.east) .. controls +(1.6,0.9) and +(1.6,-0.9) .. node[right,
    font=\scriptsize, text width=1.9cm]{guardo ancora A: nulla} (a.east);
\end{tikzpicture}
\caption{La variabile Interactable come stato. Ogni Interaction Start è sempre seguito, prima
o poi, dal suo Interaction End: l'oggetto può contare sulla coppia per mostrare e nascondere
il suggerimento. Schema ricostruito.}
\label{ue12:fig:stati}
\end{figure}

\paragraph{Approfondimento: messaggi, Does Implement Interface e Validated Get.} Tre
precisazioni generali sui nodi usati. Primo: un nodo \emph{(Message)} di un'interfaccia si
può chiamare su qualunque riferimento a oggetto; se l'oggetto non implementa l'interfaccia,
la chiamata semplicemente non fa nulla, senza errori (la frase del libro citata sopra: ``non
succede niente''). Il controllo \textbf{Does Object Implement Interface}, quindi, non serve
a evitare un errore, ma a \emph{distinguere} gli interagibili dagli altri oggetti, così da
chiudere l'interazione quando si guarda un muro. Secondo: il confronto \emph{Not Equal} fra
due riferimenti confronta l'\emph{identità} degli oggetti (è lo stesso actor?), non la loro
classe; due cubi identici sono comunque diversi. Terzo: il \textbf{Validated Get} è un Get
con l'\emph{Is Valid} incorporato, con le due uscite \emph{Is Valid} e \emph{Is Not Valid};
``validare'' un riferimento prima di usarlo è la raccomandazione con cui il docente chiude
il corso, perché un riferimento vuoto o a un actor distrutto (per esempio un interagibile
raccolto o eliminato) produce, se usato, gli errori \emph{Accessed None} a fine Play.
Funzionamento generale dell'engine.

% ---------------------------------------------------------------------
\section{L'oggetto di prova: BP\_InteractableTest}
\label{ue12:sec:test}

\paragraph{L'actor.} ``Adesso creiamo il nostro actor interagibile.'' Nella cartella dei
Blueprint, una Blueprint Class di tipo \textbf{Actor}, ``veramente semplicissimo'':
\textbf{BP\_InteractableTest}. Due componenti: un \textbf{Cube} e un \textbf{Widget}
(``fate widget, che poi faremo''), messo sopra il cubo.

\paragraph{Implementare l'interfaccia.} ``Non si trova'', se si cerca l'interfaccia fra i
componenti: va aggiunta in \textbf{Class Settings}, nella toolbar. Nei Details compaiono le
categorie \emph{Imports} (\emph{Default Namespace}, \emph{Imported Namespaces}) e
\textbf{Interfaces}, con \emph{Inherited Interfaces} e \textbf{Implemented Interfaces}; in
quest'ultima, \textbf{Add}, BPI\_InteractionRay. ``Per vedere che l'hai aggiunto'': nel
pannello My Blueprint compare la sezione \textbf{Interfaces}, con Interaction, Interaction End
e Interaction Start.

\paragraph{Implementare le funzioni.} Trascinare Interaction Start nel grafo non funziona:
l'editor avverte che la funzione ``was not marked as Blueprint Callable'' e non può essere
piazzata così. ``Dobbiamo overrideare la loro funzione, che di base è nulla'': le funzioni
dell'interfaccia ``sono solo degli appoggi'', vuote, e il loro contenuto ``sarà da
overrideare proprio dal livello di design''. Con un doppio clic (o tasto destro,
\emph{Implement Event}) si ottengono tre eventi rossi: \textbf{Event Interaction Start},
\textbf{Event Interaction End} ed \textbf{Event Interaction}. Ciascuno con un
\textbf{Print String}:
\begin{itemize}
  \item Interaction Start: ``inizio interazione'', \textbf{Text Color} verde;
  \item Interaction End: ``fine interazione'', rosso;
  \item Interaction: ``interagendo'', giallo-arancione (``ancora non l'abbiamo usato,
        quindi non lo vedremo, è normale'').
\end{itemize}
Il colore si trova aprendo la tendina del Print String (la freccia in basso del nodo).

\paragraph{La collisione.} ``È importante però una cosa'': il canale InteractionTrace è in
Ignore di default, quindi il cubo, così com'è, non viene visto dal raggio. Nella
\textbf{Collision} del Cube, preset \emph{Custom} e \textbf{InteractionTrace} su
\textbf{Block}. ``Su Cube, tanto la mesh è quella.''

\paragraph{La prova.} Si mette il cubo in scena. In Play: ``inizio interazione, fine
interazione, inizio interazione, fine interazione''. Con la camera in terza persona colpire
il cubo è un po' difficile, ``ma ci sono delle volte in cui diventa verde''. ``Vi funziona?''
``Sì, funziona.'' ``Quindi adesso abbiamo creato il nostro trace per quanto riguarda le
interazioni.'' Il componente ``fa le veci di \emph{vado a controllarti se ci sono delle
interazioni, e gestisco tutta la cosa}''; l'interazione vera, il tasto, la fa il personaggio,
``in modo classico''.

\paragraph{Approfondimento: eventi o funzioni dell'interfaccia.} Le tre funzioni di
BPI\_InteractionRay non hanno \emph{output}: per questo, nell'actor che le implementa,
diventano \emph{eventi} nell'Event Graph (è la stessa regola citata dal libro per la Trigger
Pressed della pistola VR). Se una funzione dell'interfaccia avesse un valore di ritorno (per
esempio un \emph{GetInteractionText} che restituisce ``Apri'' o ``Accendi''), andrebbe
implementata come \emph{funzione}, con un doppio clic nella sezione Interfaces, e potrebbe
essere chiamata per leggere un dato. Un'interfaccia non ha stato né variabili: ogni classe
che la implementa decide da sé che cosa fare, e può anche lasciare vuote le funzioni che non
le servono. Funzionamento generale dell'engine.

% ---------------------------------------------------------------------
\section{L'input: IA\_Interact e il tasto E}
\label{ue12:sec:input}

Nel BP\_Bot2 serve un nuovo input. ``Come si crea un nuovo input? Con l'IMC e con l'Input
Action.''
\begin{enumerate}
  \item Nel Content Browser si va alle Input Action del progetto (cliccando su IA\_Jump,
        \emph{Show in Folder View});
  \item ``basta che fate una copia di IA\_Run'': \textbf{Duplicate}, nome
        \textbf{IA\_Interact}; ``basta che qua dentro abbia i \textbf{Triggers}
        \emph{Pressed} e \emph{Released}'' (come quelli di IA\_Run), da lasciare così;
  \item nell'\textbf{IMC} di default (\emph{IMC\_\allowbreak Default}), il più accanto alle
        \emph{Mappings} aggiunge una riga per IA\_Interact; poi, cliccando sull'icona
        della tastiera e premendo il tasto, si assegna \textbf{E}.
\end{enumerate}
Il bot2 ``eredita la stessa cosa'': usa lo stesso IMC\_Default del template. Le domande dei
presenti sono tutte sulla tendina dell'IMC (``non mi fa mettere l'azione di input'': ``è la
tendina'') e sui trigger (``il primo è premuto e l'altro qual è?'': ``premuto e rilascio va
benissimo così, lasciato così''). Nel personaggio, cercando \emph{IA\_Interact}, compare il
nuovo evento d'input.

\paragraph{Approfondimento: Pressed e Released nello stesso trigger.} Un'Input Action
digitale senza trigger scatta in \emph{Triggered} a ogni frame finché il tasto è premuto; con
il trigger \emph{Pressed} scatta una sola volta alla pressione, con \emph{Released} una
volta al rilascio. Con \emph{entrambi}, come in IA\_Run e quindi in IA\_Interact, l'uscita
\emph{Triggered} scatta due volte per ogni pressione: quando si preme e quando si rilascia.
Nella lezione questo non dà problemi, perché il grafo del personaggio passa da un \emph{Do
Once} che si resetta solo alla fine del Montage: il secondo impulso, al rilascio, arriva
mentre l'animazione è in corso e viene ignorato. Ma se il Montage fosse molto breve, o il
giocatore tenesse premuto a lungo, il rilascio potrebbe far partire una seconda
interazione: per un'interazione ``a pressione'' basta il solo trigger Pressed, o l'uscita
\emph{Started}. La trascrizione non dice quale uscita dell'evento IA\_Interact usi il docente.
Funzionamento generale dell'Enhanced Input.

% ---------------------------------------------------------------------
\section{L'interazione nel personaggio}
\label{ue12:sec:player}

``Allora adesso partiamo a fare quello che ci servirà per indirizzare il nostro player,
cioè farlo ruotare verso il nostro oggetto, e allo stesso modo far sì che l'interazione
parta.''

\paragraph{Un Do Once.} ``Vi dico che è un'azione che deve essere fatta su un \textbf{Do
Once}'': l'interazione è una sequenza (girarsi, far partire il Montage, aspettarne la fine)
e solo ``a fine del montaggio può reiniziare una prossima interazione''. Quindi IA\_Interact
$\rightarrow$ Do Once.

\paragraph{Chi stiamo guardando.} Dal componente \textbf{BPC\_RayInteract} (trascinato dal
pannello dei componenti nel grafo: ``prendilo e trascinalo'') si prende la sua variabile
\textbf{Interactable} e la si converte in \textbf{Validated Get}: ``questo Interactable è
presente? Ci sta o non ci sta?''. Il personaggio non rifà la trace: legge il risultato del
componente.
\begin{itemize}
  \item \textbf{Is Valid}: il messaggio \textbf{Interaction} di BPI\_InteractionRay
        (``tira fuori e fai \emph{interaction}, cioè Interaction Message''), con target
        Interactable: ``e questa è proprio la base: così interagiamo''.
  \item \textbf{Is Not Valid}: \textbf{Reset} del Do Once (aggiunto dopo, quando il docente
        si accorge che senza il reset, premendo E senza guardare nulla, il Do Once resta
        chiuso per sempre: ``se non è valido, resetta il Do Once comunque, in ogni caso'').
\end{itemize}

\paragraph{Girarsi verso l'oggetto.} Dopo il messaggio, \textbf{Set Actor Rotation} con
target \emph{Self} (``siamo già nel player''). Si splitta il pin \emph{New Rotation}, e si
splitta anche l'uscita di un \textbf{Get Actor Rotation}: ``l'unica asse che ci interessa è
la Z'', perché il personaggio deve girarsi solo attorno all'asse verticale (``se
l'interactable sta qua, deve girare così''). Quindi:
\begin{itemize}
  \item X (Roll) e Y (Pitch) restano quelle attuali, dal Get Actor Rotation;
  \item Z (Yaw) viene da un \textbf{Find Look at Rotation}, ``quello che abbiamo fatto
        l'altra volta per la torretta'', con \textbf{Start} = \textbf{Get Actor Location}
        del personaggio (il docente dice prima ``rotation'', poi si corregge: ``location,
        ragazzi, scusatemi'') e \textbf{Target} = Get Actor Location dell'Interactable. Il
        pin di uscita non si lascia splittare (``non me lo splitta''), quindi un
        \textbf{Break Rotator}, e il suo Z va nella Z della nuova rotazione.
\end{itemize}

\paragraph{La prova provvisoria.} Prima del Montage, una prova: la scritta ``interagendo''
compare, ma il personaggio non si gira; e dopo la prima volta non succede più nulla. ``Ah,
giustamente, ragazzi'': il Do Once non si resetta mai. Il docente collega temporaneamente
il Reset alla fine della catena (``lo faccio solamente per farvelo vedere''): ``mi giro di
qua, e si gira verso di lui''. Poi lo stacca, ``perché ci servirà più avanti''.

\paragraph{Il Montage.} Nel Content Browser si cerca \emph{interact}: fra le animazioni del
bot c'è \textbf{Interact Press Button}. ``Vi ricordate come si fa il montaggio?'' Tasto
destro sull'animazione, \textbf{Create}, \textbf{Create AnimMontage}: si crea un asset che
``è sempre la stessa animazione, ma in questo caso è un montaggio''. Nel personaggio,
\textbf{Play Montage}: nel pin \textbf{In Skeletal Mesh Component} la \textbf{Mesh} del
personaggio (``basta che prendi la Mesh da qui e la metti così''), nel \textbf{Montage to
Play} il nuovo Montage (\emph{Interact Press Button Montage}).

\paragraph{La variabile Interacting.} ``Mettiamo, giusto perché in futuro ci possa
servire, una variabile'' booleana, \textbf{Interacting}: ``dobbiamo verificare se stiamo
interagendo, non stiamo interagendo, eccetera. Sempre buono mettere una variabile che dice:
guarda, io sto interagendo''. Prima del Play Montage, Set Interacting a \textbf{vero}; sull'
\textbf{On Completed}, ``quindi è finita l'animazione'', a \textbf{falso}.

\paragraph{Fermare il giocatore.} Durante l'interazione il giocatore non deve muoversi:
\begin{itemize}
  \item prima del Montage, \textbf{Disable Input}, con il \textbf{Get Player Controller}
        nel pin \emph{Player Controller}: ``chi dobbiamo levare? Gli input al controller del
        nostro player'';
  \item lo \textbf{Stop Movement Immediately}, ``andiamo a fermare qualsiasi tipologia
        di movimento''. Il target non viene detto: il nodo appartiene al componente di
        movimento del personaggio, il \emph{Character Movement};
  \item sull'On Completed, \textbf{Enable Input}, con lo stesso Player Controller;
  \item e infine il \textbf{Reset} del Do Once, ``in modo tale che io possa reinteragire di
        nuovo''.
\end{itemize}
In Play: ``io giro di qua, si gira, fa l'animazione, non mi posso muovere''; ``e adesso mi
riposso muovere appena finito il Montage'' (figura~\ref{ue12:fig:player}). ``Ci siete? Vi
funziona? Perfetto.''

\paragraph{I problemi dei presenti.} Uno studente non trova come prendere il componente
(``l'hai aggiunto come components? Prendilo e trascinalo''); altri si perdono i due nodi
successivi, ripetuti dal docente (Get Interactable, Convert to Validated Get, il messaggio
Interaction). Un altro ha il Montage che non parte: manca la Mesh nel pin del Play Montage.
Uno studente, a cui prima non funzionava, conferma che ora va, ma ha lo
schermo ``pieno di \emph{Hello}'' (forse i Print String rimasti con il testo di default),
``quindi li vedo con la coda dell'occhio''.

\begin{figure}[htbp]
\centering
\begin{tikzpicture}[font=\footnotesize,
  ue12node/.style={draw, thick, rounded corners=4pt, align=center, fill=gray!8,
                  minimum height=0.9cm, text width=2.2cm},
  ue12ev/.style={ue12node, fill=red!12},
  ue12msgn/.style={ue12node, fill=violet!12},
  ue12pure/.style={ue12node, fill=green!10, rounded corners=8pt},
  ue12exec/.style={-{Stealth[length=2.5mm]}, very thick, black!80},
  ue12data/.style={-{Stealth[length=2mm]}, thick, blue!60!black}]
  \node[ue12ev, text width=1.9cm] (ia) at (0,-0.0) {\textbf{IA\_Interact}\\{\scriptsize tasto
    E}};
  \node[ue12node, text width=1.5cm] (do) at (2.3,0) {\textbf{Do Once}};
  \node[ue12node, text width=2.4cm] (vg) at (4.9,0) {\textbf{Validated Get}\\{\scriptsize
    BPC\_RayInteract.\\Interactable}};
  \node[ue12msgn] (msg) at (8.6,0) {\textbf{Interaction}\\{\scriptsize (Message) su
    Interactable}};
  \node[ue12node, text width=2.5cm] (rot) at (11.9,0) {\textbf{Set Actor Rotation}\\{\scriptsize
    X, Y attuali; Z dal Look at}};
  \node[ue12pure, text width=3.0cm] (look) at (11.9,-1.85) {\textbf{Find Look at
    Rotation}\\{\scriptsize mia posizione $\rightarrow$ posizione di Interactable}\\
    $\rightarrow$ \textbf{Break Rotator}: Z};
  \node[ue12node, text width=3.5cm] (pre) at (8.0,-4.4) {\textbf{Set Interacting} = vero\\
    \textbf{Disable Input}\\\textbf{Stop Movement Immediately}};
  \node[ue12node, text width=2.8cm] (pm) at (12.0,-4.4) {\textbf{Play Montage}\\{\scriptsize
    Mesh; Interact Press Button Montage}};
  \node[ue12node, text width=3.5cm] (post) at (12.0,-6.8) {\textbf{Set Interacting} =
    falso\\\textbf{Enable Input}\\\textbf{Reset} del Do Once};
  \node[ue12pure, text width=1.9cm] (gpc) at (8.0,-6.8) {Get Player Controller};
  \draw[ue12exec] (ia) -- (do); \draw[ue12exec] (do) -- (vg);
  \draw[ue12exec] (vg) -- node[below, font=\scriptsize]{Is Valid} (msg);
  \draw[ue12exec] (msg) -- (rot);
  \draw[ue12exec] (rot.east) -- ++(0.35,0) -- ++(0,-3.15) -| (pre.north);
  \draw[ue12exec] (pre) -- (pm);
  \draw[ue12exec] (pm) -- node[right, font=\scriptsize]{On Completed} (post);
  \draw[ue12exec] (vg.south) -- ++(0,-1.0) -| node[pos=0.25, above, font=\scriptsize]{Is
    Not Valid: Reset} ([xshift=-3mm]do.south);
  \draw[ue12data] (look) -- (rot);
  \draw[ue12data] (gpc) -- (pre);
  \draw[ue12data] (gpc) -- (post);
\end{tikzpicture}
\caption{L'interazione nel personaggio. Il personaggio non rifà la trace: legge la variabile
Interactable del suo componente, invia il messaggio, si gira solo sull'asse Z verso
l'oggetto, blocca l'input e suona il Montage; a Montage finito restituisce l'input e riapre
il Do Once. Ricostruzione dalla lezione; la posizione esatta del Set Interacting e del
Disable Input nella catena è dedotta dalle parole del docente.}
\label{ue12:fig:player}
\end{figure}

\paragraph{Approfondimento: Do Once, Disable Input e la rotazione.} Romero e Sewell
descrivono il \textbf{Do Once} nel capitolo 13 (\emph{Data Structures and Flow Control}):
``il nodo Do Once esegue il pin di uscita una sola volta. Se viene attivato di nuovo, il pin
di uscita non viene eseguito. Il pin di input Reset deve essere attivato per permettere al
Do Once di eseguire di nuovo il pin di uscita.'' Il loro esempio è un'arma caricata, che si
ricarica dopo 30 secondi con un timer che resetta il Do Once: lo stesso schema
della lezione, con il Montage al posto del timer. Nel capitolo 15 il libro descrive
\textbf{Enable Input} e \textbf{Disable Input}: ``definiscono se un actor deve rispondere agli
eventi di input'', e ``hanno bisogno di un riferimento al Player Controller in uso''. Due
note generali sull'engine. La prima: \emph{Disable Input} chiamato sul personaggio blocca
gli eventi d'input di quel Blueprint (movimento, camera, salto), ma non ferma un movimento
già in corso, cioè la velocità residua del Character Movement; per questo serve anche
\emph{Stop Movement Immediately}, che azzera la velocità sul posto. La seconda: il
\emph{Set Actor Rotation} è istantaneo, quindi il personaggio ``scatta'' verso l'oggetto. Per
una rotazione morbida si può usare un \emph{RInterp To} nel Tick o una Timeline, come per le
porte; e il template Third Person, che ha l'\emph{Orient Rotation to Movement} attivo, non
contrasta la rotazione perché durante l'interazione il personaggio è fermo.

% ---------------------------------------------------------------------
\section{Il widget ``premi E''}
\label{ue12:sec:widget}

``Andiamo semplicemente a rendere tutto un po' più carino'': un piccolo widget che mostri il
tasto da premere sopra l'oggetto.

\paragraph{Gli asset della UI.} Il docente aveva messo nel materiale didattico una cartella
di immagini per la UI, ``la UI asset''. Uno studente non l'ha trovata (``pensavo
fosse sul Drive''): il docente controlla il Drive condiviso e i canali del Discord, scopre che
non era stata caricata e la carica al momento, ``tanto pesa pochi kilobyte''. ``Estraetela e
schiaffatela dentro Unreal'', nella cartella \textbf{UI}. Il trascinamento nel Content
Browser a qualcuno dà ``fallimento''; il ``metodo abbastanza brusco, però va benissimo'' è
copiare la cartella \emph{estratta} con il gestore dei file direttamente nella cartella
\emph{Content/UI} del progetto. ``Adesso funziona.''

\paragraph{WB\_InteractUI.} Nella cartella UI, un nuovo Widget Blueprint, classe User
Widget: \textbf{WB\_InteractUI}. Nel Designer:
\begin{enumerate}
  \item un \textbf{Border}, con l'alfa del colore a \textbf{0}, ridimensionato a un
        quadrato di circa \textbf{96 $\times$ 96} (``94 va bene anche''), e il
        \textbf{Padding} tutto a \textbf{zero};
  \item un secondo \textbf{Border} \emph{dentro} il primo;
  \item le immagini: nel \textbf{Brush} di ciascun Border (``vedi che c'è il brush, apri il
        brush'') si sceglie un'immagine del pacchetto, ``il panel che volete voi'' (ce ne
        sono diversi, Panel 0, 1, 2, 3, 5\dots). Dopo una correzione (``questo è il
        contrario, scusate''), il Border esterno prende \textbf{Panel Border 5}, con l'alfa a
        \textbf{1} e la tinta \textbf{nera}, e quello interno \textbf{Panel 5};
  \item un \textbf{Text} con la lettera \textbf{E}, con \emph{Horizontal Alignment} e
        \emph{Vertical Alignment} al \textbf{centro}, e la dimensione del \textbf{font} a
        35, poi 40--45 (``una bella E'').
\end{enumerate}
``Piccola cosa, ragazzi, prima non l'ho fatto fare a mano'': la dimensione dei Border si può
impostare anche con precisione, dall'\textbf{Image Size} del Brush (96 $\times$ 96).

\paragraph{Il widget che contiene la barra della vita.} Una studentessa crea un nuovo
widget, ma aprendolo trova le barre della vita e della stamina, e non riesce a modificarle.
Rifacendo la creazione insieme, la causa emerge: nella finestra di scelta della classe aveva
preso un widget esistente come genitore (``anteprima widget praticamente ti dà
un'anteprima dei widget che già ci sono''), cioè il suo HUD; il nuovo widget ne
\emph{ereditava} il contenuto, in sola lettura. ``È buona prassi fare quello tutto in
alto'', \textbf{User Widget}.

\paragraph{L'immagine che non si vede.} Un'altra studentessa ha il bordo esterno che non
compare. Il docente controlla tinta, alfa, Image Size, Padding, \emph{Content Opacity}: ``è
a 1, va bene così''. ``Perdonami, purtroppo non saprei dire'': si cancellano i Border e si
rifanno, ``basta funzionare''. Rifatti, funzionano.

\paragraph{Il widget nell'oggetto.} In BP\_InteractableTest, il componente \textbf{Widget}:
\begin{itemize}
  \item \textbf{Space}: \textbf{Screen}, non World, ``in modo tale che ci sia una specie di
        2D, quindi relativo allo schermo, non al mondo'';
  \item \textbf{Draw Size}: \textbf{100 $\times$ 100}, ``un po' più piccolo'';
  \item \textbf{Widget Class}: WB\_InteractUI.
\end{itemize}
In Play, il quadratino con la E ``scala con lo schermo'', ma, se non lo si nasconde,
``rimane per sempre''.

\paragraph{Mostrare e nascondere.} Nel grafo di BP\_InteractableTest:
\begin{itemize}
  \item \textbf{Event BeginPlay} $\rightarrow$ \textbf{Set Hidden in Game} sul Widget, con
        \emph{New Hidden} \textbf{attivo}: nascosto da subito;
  \item \textbf{Event Interaction Start} $\rightarrow$ Set Hidden in Game,
        \textbf{disattivo}: ``allora la faccio apparire'';
  \item \textbf{Event Interaction End} $\rightarrow$ Set Hidden in Game, \textbf{attivo}.
\end{itemize}
Il docente prova anche un'idea per l'Event Interaction (nascondere il widget, un
\textbf{Delay} di un secondo, mostrarlo di nuovo), ma la toglie: ``questo potrebbe dare
errori, ragazzi'', perché il Delay ``va in contrasto'' con Start ed End (se lo sguardo si
sposta durante il secondo d'attesa, il Delay rimostrerebbe un widget che l'End ha appena
nascosto). ``Lascialo così per adesso.''

\paragraph{La spunta al contrario.} Una studentessa vede la E anche quando non guarda il
cubo: ``allora vuol dire che hai sbagliato l'Hidden in Game''. ``Il tick è positivo sull'End
e negativo sullo Start'', ``e anche positivo sul Begin Play''. Dopo la correzione resta un
dubbio, poi chiarito: all'inizio della prova la camera stava già inquadrando il cubo da una
posizione strana, e il widget era giustamente visibile. ``Non lo vedo più. Perfetto,
ragazzi.''

\paragraph{Approfondimento: il Widget Component.} Il \textbf{Widget Component} mostra un
Widget Blueprint attaccato a un actor. In \textbf{World} space il widget è un piano nel
mondo 3D: ha prospettiva, può essere coperto da altri oggetti, si rimpicciolisce con la
distanza, ed è adatto a schermi diegetici (un terminale, un cartello). In \textbf{Screen}
space il widget è disegnato sopra l'immagine, sempre rivolto alla camera, con una dimensione
costante in pixel, e segue sullo schermo la posizione proiettata del componente: è la scelta
tipica per i suggerimenti ``premi E'' e per le barre della vita sopra i nemici. La
\textbf{Draw Size} è la risoluzione in pixel a cui il widget viene disegnato. Per nascondere
il suggerimento, l'alternativa a \emph{Set Hidden in Game} è \emph{Set Visibility} sul
componente; il risultato è lo stesso. Funzionamento generale dell'engine.

% ---------------------------------------------------------------------
\section{Un esempio: la luce che si accende}
\label{ue12:sec:luce}

\paragraph{L'idea.} ``Sappiate che tutti gli oggetti che create'', se implementano
l'interfaccia, ``avete già un qualcosa che vi identifica se state iniziando
un'interazione'', se la state finendo, ``e la tipologia di interazione''. Le funzioni
costruite per i bottoni nelle lezioni passate ``possiamo traslarle sul nostro
interactable'': ``un diverso modo di interagire e di attivare le cose''. Con Davide si
vedranno ``anche delle piccole cose di modeling'': si può modellare un pulsante, dargli
BPI\_InteractionRay, ``e a quel punto magari potete chiamare l'ascensore''. ``Cosa succede se
interagisco? Ci metto quello che voglio: ci metto un timer, ci metto una cosa, accendo una
luce.''

\paragraph{Velocissimo.} ``Facciamo velocissimo'', in BP\_InteractableTest. Dall'Event
Interaction si toglie il Print String, e:
\begin{enumerate}
  \item una variabile del tipo dell'attivabile luce, la classe figlia di BP\_Activable delle
        lezioni sul pulsante (``la Point Light'', \textbf{BP\_ActivableLight}; la
        trascrizione dice ``BP Activables'' e ``inter light''), con i suoi eventi
        \textbf{Activation} e \textbf{Deactivation};
  \item un \textbf{Flip Flop}: ``questo fa A, poi fa B, poi fa di nuovo A: fondamentalmente
        gioca su due'';
  \item su \textbf{A}, \textbf{Activation} della luce; su \textbf{B},
        \textbf{Deactivation};
  \item la variabile va resa modificabile nell'istanza e collegata alla luce del livello
        (``però gli dobbiamo dare quale light'', ``dammi la light, dammi questa'').
\end{enumerate}
``Abbiamo acceso la luce e la rispegniamo.'' ``Questo è un esempio, ragazzi'': ``qualsiasi
actor può essere separato dagli altri con la sua specifica interazione'', e ``quello che la
identifica già ce l'abbiamo''.

\paragraph{Approfondimento: Flip Flop nel libro.} Romero e Sewell (capitolo 13): ``ogni volta
che un nodo Flip Flop viene eseguito, alterna fra due pin di uscita, A e B''; ha anche un
output booleano \emph{Is A}, vero quando sta eseguendo A. Con una nota importante: ``il Flip
Flop non funziona correttamente dentro una funzione, perché ha una variabile interna che
viene cancellata quando la funzione finisce'', e quindi dentro una funzione esegue sempre A.
Lo stesso vale per il Do Once. È il motivo per cui, nella lezione, il Flip Flop sta
nell'Event Graph dell'oggetto e il Do Once nell'Event Graph del personaggio, mentre la
funzione LookForInteractions, che non deve ricordare nulla fra un frame e l'altro (lo stato
è nella variabile Interactable, che è del componente), può essere una funzione.

% ---------------------------------------------------------------------
\section{La chiusura del corso}
\label{ue12:sec:chiusura}

``Ci sono domande, ragazzi? Domande su tutto il corso, in generale.'' Il docente chiede
anche un \textbf{feedback} sul corso, ``adesso o dopo, tranquillamente sul server''. La
risposta di uno studente, scherzosa: ``tutto chiaro, tutto bellissimo, incredibile''. ``Sono
contento.''

Per il futuro, anche ``per chi guarda in differita'': ``se avete un problema, state facendo il
Big One, eccetera, buttate il vostro problema sul server: sia alunni sia professori
provvederanno a rispondervi e darvi una mano''. ``Se ci sono ultime cose che volete dire,
ragazzi, ditele ora, o effettivamente questa volta tacete per sempre.'' E il congedo: ``vi
ringrazio per la vostra attenzione, vi ringrazio per aver partecipato al corso. Spero che i
vostri progettini possano prendere vita per tutte le cose che poi andrete a fare in Unreal.
Per benino, ricordatevi di \textbf{validare sempre} molte variabili.'' ``Da questo punto in
poi'' le lezioni continuano con Davide. ``È stato un piacere.''

\begin{table}[htbp]
\centering
\small
\begin{tabular}{@{}p{3.7cm}p{3.0cm}p{2.9cm}p{4.4cm}@{}}
\toprule
\textbf{Asset o variabile} & \textbf{Dove} & \textbf{Tipo, valore} & \textbf{A che serve} \\
\midrule
M\_Menu & mappe & Empty Level & la mappa del menu \\
WB\_Menu & UI & Widget Blueprint & il menu con i due bottoni \\
Button\_StartGame, Button\_QuitGame & WB\_Menu & Button, 300 $\times$ 80 & gioca; esci \\
CameraActor & M\_Menu & Camera Actor & la vista fissa del menu \\
GameMode Override & World Settings di M\_Menu & Game Mode base & niente personaggio nel menu \\
ThirdPersonMap & mappe & livello & la mappa aperta da Play Game \\
BPC\_RayInteract & Blueprints, Components & Scene Component & traccia gli interagibili \\
CameraComponent & BPC\_RayInteract & Camera Component (ref.) & la camera del proprietario \\
DistanceTrace & BPC\_RayInteract & Float, 300 & la lunghezza del raggio \\
Interactable & BPC\_RayInteract & Actor (ref.) & l'interagibile guardato ora, o vuoto \\
Scene & BPC\_RayInteract & Scene Component (ref.) & solo per il trucco del Select \\
InteractionTrace & Project Settings & Trace Channel, Ignore & il canale delle interazioni \\
BPI\_InteractionRay & Blueprints, Components & Blueprint Interface & Start, End, Interaction \\
BP\_InteractableTest & Blueprints & Actor (Cube, Widget) & l'oggetto di prova \\
IA\_Interact & Input & Input Action, Pressed e Released & l'interazione, tasto E \\
Interacting & BP\_Bot2 & Boolean & vero durante l'interazione \\
Interact Press Button Montage & animazioni del bot & Anim Montage & l'animazione dell'interazione \\
WB\_InteractUI & UI & Widget Blueprint, 96 $\times$ 96 & il suggerimento ``E'' \\
Widget (componente) & BP\_InteractableTest & Screen, 100 $\times$ 100 & mostra WB\_InteractUI \\
Light & BP\_InteractableTest & BP\_ActivableLight (ref.) & la luce accesa e spenta \\
\bottomrule
\end{tabular}
\caption{Gli asset e le variabili della giornata, con i nomi usati dal docente (BPC\_RayInteract
è, nel progetto del docente, BPC\_RayInteract1; il nome della variabile della luce è
incerto).}
\label{ue12:tab:asset}
\end{table}

\begin{table}[htbp]
\centering
\small
\begin{tabular}{@{}p{5.4cm}p{8.6cm}@{}}
\toprule
\textbf{Sintomo} & \textbf{Causa e rimedio} \\
\midrule
La barra della stamina scende solo al rilascio & manca la Call E\_UpdateStamina in
StaminaDeduction: c'era solo in RegainStamina \\
Il secondo bottone incollato non compare & è sovrapposto al primo (controllare la
Hierarchy); oppure aggiungerne uno nuovo dalla Palette \\
L'immagine copre i bottoni & è in fondo alla Hierarchy, cioè disegnata sopra: trascinarla
prima dei bottoni \\
Il personaggio compare dietro il menu & GameMode Override nelle World Settings della
mappa del menu: Game Mode base \\
Si vede la pallina (Default Pawn) & spostare il Player Start fuori dalla camera \\
La camera del menu non si usa & Set View Target with Blend dal Get Player Controller (non
dal Camera Manager), con la referenza alla camera \\
Errore su Add to Viewport & il Target è Self: collegare il Return Value del Create Widget \\
Il cursore non si vede & Set Show Mouse Cursor nell'Event Construct, con la spunta attiva \\
Nel menu ci si muove ancora & Set Input Mode UI Only nell'Event Construct \\
Il Select non accetta Self (o la camera) & il Select prende il tipo del primo collegamento:
mettere prima una variabile Scene Component, poi cancellarla \\
In Play non si vede il raggio & LookForInteractions non chiamata nell'Event Tick; Draw
Debug Type a None \\
InteractionTrace non compare nella trace & compilare dopo aver creato il canale \\
Il raggio ignora il cubo & la mesh non blocca InteractionTrace (default Ignore): Block \\
L'evento dell'interfaccia non si trascina & non è Blueprint Callable: implementarlo come
evento (doppio clic o Implement Event) \\
Hit Actor non si collega a Interactable & Interactable di tipo sbagliato: Actor, Object
Reference \\
Si interagisce una sola volta & il Do Once non viene resettato (On Completed e Is Not
Valid) \\
Il Montage non parte & manca la Mesh nell'In Skeletal Mesh Component \\
Il nuovo widget contiene le barre dell'HUD & creato con l'HUD come genitore: scegliere User
Widget \\
Il bordo del widget non si vede & causa non trovata; rifare i Border \\
La E si vede quando non si guarda & spunte di Set Hidden in Game invertite: attivo nel
Begin Play e nell'End, disattivo nello Start \\
Trascinare la cartella UI dà ``fallimento'' & copiare la cartella estratta in Content/UI dal
gestore dei file \\
La BPI si apre senza grafo & doppio clic sulla funzione in My Blueprint \\
\bottomrule
\end{tabular}
\caption{I problemi incontrati durante la lezione, con la causa e il rimedio.}
\label{ue12:tab:errori}
\end{table}

\begin{table}[htbp]
\centering
\footnotesize
\begin{tabular}{@{}p{5.8cm}p{8.2cm}@{}}
\toprule
\textbf{Nodo, asset o impostazione} & \textbf{Che cosa fa} \\
\midrule
New Level: Empty / Basic & livello vuoto; con luce e cielo già pronti \\
Widget Blueprint (User Widget) & l'asset dell'interfaccia: Designer e Graph \\
Button: Style (Normal, Hovered, Pressed) & i colori del bottone nei tre stati \\
Hex sRGB / Hex Linear & il codice del colore, da copiare \\
Text dentro un Button & l'etichetta del bottone, figlia del bottone \\
Image: Color and Opacity (alfa) & la trasparenza dell'immagine \\
Hierarchy (ordine) & l'ultimo elemento è disegnato sopra \\
Anchors / Position / Alignment / Size & ancora; scostamento; pivot; dimensioni \\
Open Level Blueprint & il Blueprint del livello corrente \\
Create Widget / Add to Viewport & crea il widget; lo mostra \\
World Settings: GameMode Override & la Game Mode per questa mappa \\
Default Pawn / Player Start & la sfera della Game Mode base; dove nasce il Pawn \\
Camera Actor / Directional Light & una camera nel livello; il sole \\
Create a Reference to (actor) & la referenza a un actor del livello nel Level Blueprint \\
Set View Target with Blend & cambia la camera del Player Controller \\
Event Pre Construct / Construct & alla costruzione nel Designer; alla costruzione in gioco \\
Set Show Mouse Cursor & mostra o nasconde il cursore \\
Set Input Mode UI Only / Game Only & l'input va solo all'interfaccia; solo al gioco \\
On Clicked (Button) & l'evento del clic su un bottone \\
Quit Game & esce dal gioco (in editor, chiude il Play) \\
Open Level (by Name) & apre una mappa dal nome \\
Blueprint Interface & solo dichiarazioni di funzioni, in sola lettura \\
Class Settings: Implemented Interfaces & aggiunge un'interfaccia a un Blueprint \\
Interaction Start / End / Interaction (Message) & i messaggi dell'interfaccia \\
Event Interaction Start / End / Interaction & la loro implementazione nell'oggetto \\
Actor Component / Scene Component & componente senza / con Transform \\
Get Owner & l'actor che possiede il componente \\
Get Component by Class & il primo componente di una classe (non l'array) \\
Sphere Trace By Channel & una sfera lungo un segmento: Radius, Draw Debug Type \\
Get World Location / Get Forward Vector & posizione nel mondo; direzione davanti \\
Select (wildcard) & sceglie fra due valori; il tipo è quello del primo collegato \\
Is Valid (pura) / Validated Get & il riferimento esiste?; Get con le due uscite \\
Not Equal (!=) & vero se i due valori sono diversi \\
New Trace Channel (Ignore / Block) & un canale di collisione personalizzato \\
Break Hit Result: Hit Actor & l'actor colpito \\
Does Object Implement Interface & l'oggetto implementa l'interfaccia? \\
Event Tick & a ogni frame \\
Print String: Text Color & il colore del messaggio \\
Input Action (Pressed, Released) / IMC & l'azione e il tasto (E) \\
Do Once (Reset) & esegue una volta, finché non si resetta \\
Set Actor Rotation / Get Actor Rotation & imposta; legge la rotazione dell'actor \\
Find Look at Rotation / Break Rotator & rotazione per guardare un punto; i suoi tre assi \\
Create AnimMontage / Play Montage & un Montage da un'animazione; lo suona (On Completed) \\
Disable Input / Enable Input & toglie; restituisce l'input all'actor \\
Stop Movement Immediately & azzera il movimento del personaggio \\
Border: Brush, Padding, Image Size & il riquadro con immagine; margine interno; dimensione \\
Widget Component: Space, Draw Size, Widget Class & un widget attaccato a un actor \\
Set Hidden in Game (New Hidden) & nasconde o mostra un componente \\
Delay & aspetta prima di proseguire \\
Flip Flop & alterna fra le uscite A e B \\
\bottomrule
\end{tabular}
\caption{I nodi, gli asset e le impostazioni della dodicesima lezione.}
\label{ue12:tab:nodi}
\end{table}

% ---------------------------------------------------------------------
\section*{Punti chiave del capitolo}
\addcontentsline{toc}{section}{Punti chiave del capitolo}

\begin{itemize}
  \item Il designer usa Unreal per \textbf{prototipare} le meccaniche e verificare che
        funzionino; il progetto d'esame degli anni successivi è un \textbf{pezzo da
        portfolio}. \textbf{Less is more}: il giocatore impiega molto più tempo di chi
        ha progettato il gioco per capirne ogni elemento.
  \item Il menu vive in una \textbf{mappa a parte} (Empty Level), con un \textbf{Widget
        Blueprint} (User Widget) creato e mostrato nel \textbf{Level Blueprint} (Create
        Widget, Add to Viewport con il Return Value).
  \item Nel Designer: Canvas Panel, bottoni con l'\textbf{ancora} al centro e
        l'\textbf{Alignment} per impilarli, un \textbf{Text} dentro ogni bottone, i colori
        di \textbf{Normal}, \textbf{Hovered} e \textbf{Pressed} nello Style (copiabili con
        il codice Hex), e l'ordine della \textbf{Hierarchy} per mettere un'immagine
        dietro (l'ultimo elemento sta sopra).
  \item Niente personaggio nel menu: \textbf{GameMode Override} con la Game Mode base
        (Default Pawn) e Player Start fuori campo. La vista fissa: un Camera Actor e
        \textbf{Set View Target with Blend}, dal \textbf{Player Controller}.
  \item Nell'\textbf{Event Construct} del widget: \textbf{Set Show Mouse Cursor} (con la
        spunta) e \textbf{Set Input Mode UI Only}. Sui bottoni, \textbf{On Clicked}:
        \textbf{Quit Game}; \textbf{Set Input Mode Game Only} e \textbf{Open Level (by
        Name)}.
  \item Una \textbf{Blueprint Interface} è un intermediario (``un volo di linea''): solo
        nomi di funzioni, in sola lettura. Chi la implementa (Class Settings,
        Implemented Interfaces) decide che cosa fare; chi non la implementa ignora i
        messaggi.
  \item Un \textbf{Blueprint Component} isola una funzionalità da aggiungere a qualunque
        actor. L'\textbf{Actor Component} non ha posizione; lo \textbf{Scene Component}
        ha un Transform. Il componente trova il suo actor con \textbf{Get Owner} e la
        camera con \textbf{Get Component by Class}.
  \item \textbf{BPC\_RayInteract}: a ogni frame (\textbf{Event Tick}) una \textbf{Sphere
        Trace} (Radius 5) dalla posizione del componente, lungo il Forward Vector della
        camera (o del componente, se la camera non c'è) per 300 unità, sul canale
        \textbf{InteractionTrace} (Ignore di default, Block sugli interagibili).
  \item Il \textbf{Select} prende il tipo del primo collegamento: per scegliere fra la
        camera e Self, prima una variabile \emph{Scene Component}, poi la si cancella (la
        ``mandragata'').
  \item La variabile \textbf{Interactable} ricorda l'oggetto guardato: se cambia,
        \textbf{Interaction End} sul vecchio, Set, \textbf{Interaction Start} sul nuovo;
        se non si guarda niente di interagibile, End e svuota. Sempre con il
        \textbf{Validated Get}.
  \item Nel personaggio: \textbf{IA\_Interact} (tasto E), \textbf{Do Once}, il messaggio
        \textbf{Interaction}, \textbf{Set Actor Rotation} con il solo \textbf{Yaw} da
        \textbf{Find Look at Rotation}, \textbf{Disable Input} e \textbf{Stop Movement
        Immediately}, \textbf{Play Montage}, poi Enable Input e \textbf{Reset} del Do Once
        (anche quando non c'è nulla da interagire).
  \item Il suggerimento ``E'': un \textbf{Widget Component} in \textbf{Screen} space,
        nascosto nel Begin Play e mostrato e nascosto da Interaction Start ed End con
        \textbf{Set Hidden in Game}.
  \item Ogni oggetto può dare la sua risposta all'interazione: per esempio un
        \textbf{Flip Flop} che chiama Activation e Deactivation di una
        \textbf{BP\_ActivableLight}. Flip Flop e Do Once non vanno dentro le funzioni.
  \item La raccomandazione finale del corso: \textbf{validare sempre} le variabili.
\end{itemize}

%
%  Tutte le etichette di questo capitolo hanno il prefisso "ue12:"
%  (ue12:cap:..., ue12:sec:..., ue12:fig:..., ue12:tab:...), e tutti gli
%  stili TikZ hanno il prefisso "ue12" (ue12node, ue12exec, ue12data,
%  ...) per non collidere con chiavi native di TikZ né con quelli delle
%  lezioni precedenti. Il capitolo è autonomo: non contiene \ref verso
%  altri capitoli (i rimandi alle lezioni precedenti sono a parole).
%  Le figure sono ricostruzioni schematiche di quanto mostrato e detto a
%  lezione, non copie delle schermate del docente. Dove la trascrizione
%  è troppo rovinata per ricostruire un dettaglio, il capitolo lo dice
%  invece di inventarlo.
% =====================================================================

\chapter[Menu e sistema di interazione]{Il menu principale e il sistema di
interazione}
\label{ue12:cap:interazione}

Questa è la \textbf{dodicesima lezione} del corso breve di \textbf{introduzione a Unreal
Engine}, ed è l'\textbf{ultima} con questo docente: ``oggi è l'ultima lezione, almeno
mia''. Nella lezione precedente il personaggio, il \textbf{bot}, ha ricevuto la vita, il
danno della torretta, la stamina e le due barre dell'HUD; restavano in sospeso due
argomenti annunciati, il \textbf{menu} e le \textbf{interazioni}. Oggi si fanno entrambi.
``Sono argomenti che abbiamo già fatto in precedenza e adesso li uniamo'': Widget
Blueprint, Level Blueprint, trace, collisioni personalizzate, Montage, timer e attivabili
tornano tutti, messi insieme.

Il primo argomento, il più corto, è un \textbf{menu principale}: una mappa a parte, con un
piccolo \emph{diorama} 3D sullo sfondo, una camera fissa, due bottoni (\emph{gioca} ed
\emph{esci}) e il passaggio alla mappa di gioco. Il secondo, ``piuttosto lungo'', è un
\textbf{sistema di interazione} del tipo ``premi E per interagire'', costruito con due
strumenti nuovi: un \textbf{Blueprint Component}, che si aggiunge a qualunque actor e
``traccia'' dalla camera gli oggetti interagibili, e una \textbf{Blueprint Interface}, il
``porto'' attraverso cui il personaggio dice a un oggetto \emph{sto iniziando a guardarti},
\emph{non ti guardo più}, \emph{interagisco con te}. L'oggetto di prova mostra un piccolo
widget con la lettera E quando è inquadrato; premendo E il personaggio si gira verso di lui
e suona un'animazione; alla fine, lo stesso oggetto accende e spegne una luce.

La lezione è \textbf{online}, come le precedenti, e si apre con una conversazione più
generale: come recuperare le lezioni, che cosa aspetta gli studenti negli anni successivi,
e un consiglio di design, ``\emph{less is more}''. Durante la pausa il docente risolve il
problema della barra della stamina rimasto aperto la volta precedente; si chiude con i
ringraziamenti e con la raccomandazione che riassume il corso: ``ricordatevi di validare
sempre molte variabili''.

Tre avvertenze di lettura, le stesse delle lezioni precedenti. La prima: la trascrizione
è automatica e storpia moltissimi nomi. Unreal è ``Real'' e ``Arrial''; il
\textbf{diorama} è ``diuramino''; il Widget Blueprint \textbf{WB\_Menu} è ``VB menu'' e
``VBIN menu'', e \textbf{WB\_InteractUI} è ``VB Interact UI'' e ``Interacco uai''; il
\textbf{Begin Play} è ``big in play'' e ``back in play''; la \textbf{ThirdPersonMap} è
``Terperson map'' e ``there person''; il \textbf{Scene Component} è ``sin component'' e
``sim component''; il componente \textbf{BPC\_RayInteract} è ``BPC Reinteract'', ``Ray
interact one'' e ``IPC reinteract''; la Blueprint Interface \textbf{BPI\_InteractionRay} è
``VPI interaction'', ``VPA Interaction Array'' e ``bibliotecazione:ray''; l'\textbf{Hit
Actor} è ``it actor'' e ``I Tactor''; il \textbf{Validated Get} è ``variegate'' e ``vari da
get''; \textbf{Does Object Implement Interface} è ``Hacker does object implement
interface''; il \textbf{Draw Debug Type} è ``draw the bug type''; il bot è ``Boto'' e
``Boto2'', e il suo Blueprint ``BiPi Bot''; \textbf{Find Look at Rotation} è
``findLocatrotation''; l'\textbf{Anim Montage} è ``un montaggio un po' brasilero''; i
\emph{Big One} del corso di laurea sono ``biggest one''. ``Mandragata'', invece, non è un
errore: è il romanesco \emph{mandrakata}, ``trucchetto'', e il docente lo usa apposta.
Dove la ricostruzione è sicura il capitolo usa la forma corretta; dove non lo è riporta la
forma della trascrizione e lo dice. La seconda: durante la pausa il riconoscimento vocale
produce un lungo blocco di testo senza senso e ripetuto decine di volte (``grazie a tutti
per l'attenzione che ci ha spiegato'', ``qual è la risposta che ci può essere?''), seguito
da un frammento che sembra l'audio di un video in sottofondo (una ``blue map'' che mostra
``le secret room'' senza sprecare bombe, e ``un'altra monetina''): non appartengono alla
lezione e sono stati ignorati. La terza: buona parte della lezione è fatta di indicazioni
a voce mentre si lavora (``tira fuori'', ``attacca qua'', ``vai giù''), spesso rivolte
allo schermo di uno studente; il capitolo ricostruisce i grafi da ciò che il docente dice
mentre li costruisce, distingue i tentativi dalle soluzioni, completa i nomi esatti dei
nodi e dei pin di Unreal Engine 5, e dichiara i punti in cui la trascrizione lascia un
buco.

\paragraph{Come leggere questo capitolo.} Il capitolo segue l'ordine della lezione. Le
figure ricostruiscono la disposizione del menu, i due grafi che lo fanno funzionare (Level
Blueprint e Widget Blueprint), il ruolo di componente e interfaccia, la funzione che
cerca gli interagibili, il ciclo di vita della variabile Interactable e il grafo
dell'interazione nel personaggio. Le tabelle riassumono gli asset e le variabili della
giornata, gli errori incontrati durante la lezione e i nodi. Qua e là, dei paragrafi di
\textbf{approfondimento} collegano la lezione al libro di Marcos Romero e Brenden Sewell,
\emph{Blueprints Visual Scripting for Unreal Engine 5}, il testo di riferimento per questa
parte del libro, e a qualche nozione tecnica generale su Unreal Engine 5; sono segnalati
come tali e non fanno parte della lezione.

Il percorso del capitolo:
\begin{enumerate}
  \item le \textbf{domande d'apertura}: recuperare le lezioni, l'esame e i Big One, il
        lavoro del designer, \emph{less is more} (sezione~\ref{ue12:sec:inizio});
  \item gli \textbf{annunci}, la scelta dell'ordine e la \textbf{dimostrazione} del
        sistema di interazione (sezione~\ref{ue12:sec:annunci});
  \item il \textbf{menu}: la mappa, il widget, i bottoni e i loro colori
        (sezione~\ref{ue12:sec:menu});
  \item il \textbf{Level Blueprint} del menu: il widget sullo schermo, la Game Mode e la
        camera (sezione~\ref{ue12:sec:levelbp});
  \item il \textbf{Graph del menu}: cursore, input e bottoni; il diorama e i problemi
        degli studenti (sezioni \ref{ue12:sec:menugraph} e~\ref{ue12:sec:diorama});
  \item \textbf{interfacce} e \textbf{Blueprint Component} (sezione~\ref{ue12:sec:concetti});
  \item il componente \textbf{BPC\_RayInteract} e la sua Sphere Trace
        (sezione~\ref{ue12:sec:componente});
  \item il \textbf{Trace Channel} e la \textbf{Blueprint Interface}
        (sezione~\ref{ue12:sec:bpi});
  \item la \textbf{pausa}: la stamina di Tommaso (sezione~\ref{ue12:sec:pausa});
  \item la logica della variabile \textbf{Interactable} (sezione~\ref{ue12:sec:logica});
  \item l'oggetto di prova \textbf{BP\_InteractableTest} (sezione~\ref{ue12:sec:test});
  \item l'\textbf{input} e l'interazione nel personaggio (sezioni \ref{ue12:sec:input}
        e~\ref{ue12:sec:player});
  \item il widget \textbf{``premi E''} (sezione~\ref{ue12:sec:widget});
  \item la \textbf{luce} interattiva e la \textbf{chiusura} del corso (sezioni
        \ref{ue12:sec:luce} e~\ref{ue12:sec:chiusura}).
\end{enumerate}

% ---------------------------------------------------------------------
\section{Prima della lezione: recuperare, l'esame, il designer}
\label{ue12:sec:inizio}

\paragraph{Recuperare o seguire?} In apertura, prima che arrivino gli altri, uno
studente rimasto indietro (``le sto andando pian piano'') chiede se gli convenga
recuperare le lezioni vecchie o seguire questa. (Il nome, nella trascrizione, è incerto:
il docente lo chiama una volta ``Dave'', ma poco dopo ricorda questa conversazione come
quella fatta ``con Lorenzo''.) La risposta: seguire la lezione in diretta ha un vantaggio,
``se c'è un qualche particolare dell'argomento, puoi farmi direttamente una domanda''; poi,
siccome oggi si uniscono argomenti già visti, ``magari ti rifai anche quelle precedenti in
modo tale da stare un po' più al passo''. L'unico prerequisito per l'esercizio di oggi è
``avere il personaggio che si muove'', che ``può guardarsi in giro''; ``puoi anche usare
il personaggio standard, volendo''. Le domande si possono fare anche dopo, sul
\textbf{Discord}: ``mandarmi le domande anche quando non ci vediamo più, perché oggi è
l'ultima lezione''. Lo studente sceglie di seguire e di riguardare la registrazione.

\paragraph{L'esame e i Big One.} Seconda domanda: ci sarà un esame, con un progetto da
consegnare? ``Nel corso no'': questo corso ``è strutturato che sia un'introduzione per
voi''. Il docente chiede che corso farà lo studente (game designer: ``tanto farete sia game
sia level designer'') e descrive il percorso, avvertendo che è quello che ha fatto lui
(``almeno io feci così, adesso non so come sono cambiate le cose''):
\begin{itemize}
  \item al \textbf{primo anno} il progetto principale è un \textbf{gioco da tavolo}:
        ``farete analog game design'';
  \item al \textbf{secondo anno} il \textbf{Big One} è un progetto su Unreal, ``perché
        Unreal è molto più improntato per i designer'';
  \item al \textbf{terzo anno} il progetto più grande (nella trascrizione ``biggest one''),
        che può anche essere la continuazione del secondo, ``se il vostro progetto nel
        secondo anno è andato bene'', da portare fino al rilascio su \textbf{Steam},
        ``con l'Accademia assieme''.
\end{itemize}
In questo senso l'esame c'è: ``sarà farvi un \textbf{pezzo da portfolio}'', ``un gioco che
potete mettere sul vostro portfolio e promuovervi''. ``Quindi, più sapete, meglio è.''

\paragraph{Che cosa fa il designer con Unreal.} ``Io sono designer'': per un designer
Unreal ``servirà principalmente per prototipare le meccaniche'', cioè ``concretizzare
l'idea che tu hai studiato, che tu hai tirato fuori, concretizzarla in editor e vedere se
effettivamente funziona''. Il suo strumento sono i \textbf{Blueprint}, ``quindi non andare
troppo in depth'' nella programmazione: in ambito lavorativo il designer ha ``un
programmatore che lo assiste''; ``poi, quando si sta da soli, si fa di tutto''.

\paragraph{Less is more.} Un consiglio ``da esperienza'': \textbf{less is more}. ``Più è
leggibile un qualcosa anche nella vostra testa a livello di design, più sarà leggibile per
un giocatore.'' L'esempio è ``la regola delle parole a schermo'': il testo che scorre
all'inizio di \emph{Star Wars}. Chi l'ha scritto lo rilegge in venti secondi, perché lo
conosce già; lo spettatore, che deve ``cominciare a creare le connessioni mentali'', ci
mette ``non 20 secondi, ma 50''. Nel gioco vale lo stesso: chi lo ha progettato e testato
pensa ``vabbè, tanto questo lo so ormai'', ma il giocatore ``ci metterà il triplo del
tempo''. ``Più roba fate, più tempo il giocatore dovrà metterci per imparare tutto
quanto.'' ``Credo che \emph{less is more} ve lo faranno fare'', aggiunge, ``ve lo
diranno''; e per qualunque domanda di design, ``possiamo tranquillamente pure parlare''.

\paragraph{Approfondimento: la maledizione della conoscenza.} Il fenomeno descritto dal
docente ha un nome in psicologia, \emph{curse of knowledge}: chi conosce un'informazione
fatica a immaginare come la percepisce chi non la conosce. È il motivo per cui il
\emph{playtest} con persone esterne al team è insostituibile, e per cui le regole di un
gioco andrebbero introdotte una alla volta, ciascuna quando serve (il cosiddetto
\emph{onboarding} progressivo). I numeri del docente (20 secondi contro 50, ``il
triplo'') sono un'immagine, non una misura. Questa è una nozione generale di design, non
della lezione.

% ---------------------------------------------------------------------
\section{Gli annunci, l'ordine della giornata e la dimostrazione}
\label{ue12:sec:annunci}

\paragraph{Che cosa si fa oggi.} Prima che arrivino tutti, il docente chiede agli
studenti di scegliere: oggi ci sono ``sia il sistema di interazioni, che è piuttosto
lungo'', sia il \textbf{menu}, ``che l'altra volta non abbiamo finito, praticamente non
abbiamo neanche iniziato''. Prima il lungo o prima il corto? Gli studenti scelgono
``corto e lungo, così facciamo la pausa prima'' (``andiamo a un caffè''). Uno studente
ricorda un'altra cosa in sospeso: ``durante la pausa sistemare il problema della stamina,
che se tengo premuto non si vede, ma solo quando mollo il tasto si vede'' (è il problema
rimasto aperto alla fine della lezione precedente); ``ricordamelo appena andiamo in
pausa''.

\paragraph{Gli annunci.} Arrivati gli altri, il docente ripete: ``oggi praticamente è la
mia ultima lezione''. Seguiranno ``circa quattro lezioni'' (``adesso non mi ricordo quante
lezioni avete'') con \textbf{Davide}, dedicate alla parte di \textbf{level design}: ``i
tool che vi serviranno per fare level design'' e alcuni approfondimenti, ``credo di
materiali''. ``Ciò non vuol dire che io vi abbandoni'':
\begin{itemize}
  \item il \textbf{canale Discord} resta aperto, anche quando inizierà il corso vero e
        proprio: ``fate domande e dovrebbe arrivare qualcuno in vostro aiuto'';
  \item anche i lavori si possono postare per avere un feedback (il docente cita quello
        di \textbf{Federico La Porta}, ``che è bravissimo'', a cui ha già risposto);
  \item i due canali vanno tenuti separati: quello del corso (``Introduzione a Unreal'')
        per i compiti, ``così almeno non mandate mischiati, perché sennò perdo il filo'',
        e la chat generale ``per qualsiasi comunicazione generica'';
  \item le registrazioni restano disponibili ``fino a settembre, se non sbaglio'', ``quindi
        avete comunque abbastanza tempo per poterle recuperare''.
\end{itemize}

\paragraph{La dimostrazione.} ``Adesso vi faccio vedere solamente un'introduzione, come
nei classici video tutorial.'' Il personaggio è quello della volta scorsa, ``che si muoveva
con stamina e adesso perde vita con la torretta''. Dalla sua testa parte un \textbf{trace}
sottile, ``tipo un cilindrone'': ``questo sarà il nostro sistema che identificherà un
possibile \emph{interactable} tramite la telecamera''. A sinistra, sullo schermo, i
messaggi: se la camera inquadra l'oggetto compare \emph{inizio interazione}, se si
distoglie lo sguardo \emph{fine interazione}; premendo il tasto, il personaggio ``si gira in
automatico'' verso l'oggetto e interagisce. ``Visto che comunque abbiamo creato questo
sistema, è un buon stampo modulare'': in futuro, pulsanti e altri oggetti potranno avere
``l'interazione un po' più precisa'', riusando ciò che si è visto con i pulsanti classici:
``sempre premi pulsantino e fai cosetta, e animazione''.

% ---------------------------------------------------------------------
\section{Il menu: la mappa, il widget e i bottoni}
\label{ue12:sec:menu}

\paragraph{Il Widget Blueprint.} Nella cartella \textbf{UI} del progetto, tasto destro,
\emph{User Interface}, \textbf{Widget Blueprint}, ``come abbiamo creato per l'HUD''. Alla
richiesta della classe base, ``non ci serve una particolare, ce ne serve una generica'':
\textbf{User Widget}. Il nome: \textbf{WB\_Menu}.

\paragraph{Una mappa per il menu.} ``Una cosa importante che io almeno consiglio'': creare
una \textbf{nuova mappa} dedicata al menu, ``quindi differenziare il gioco dal menu''. ``Si
possono fare anche volendo insieme, ma a livello di gestione'' è più comodo separarli, ``ed
è molto più rapido''. Da \textbf{File}, \textbf{New Level}, si sceglie \textbf{Empty
Level}: ``se vi diventa tutto nero, è normale, ragazzi: è un empty level''. Un livello
\emph{Basic}, invece, contiene già ``la directional light, la skybox, tante cose che sono
già prefatte''; qui si parte da zero. Poi:
\begin{itemize}
  \item il menu si può fare ``sia in 2D, con immagini 2D'', sia con elementi 3D, ``una
        specie di diorama''; per cominciare, un \textbf{piano} abbastanza grande (dalle
        forme base: il docente prende prima un cubo, poi si corregge), con un materiale a
        griglia, ``il material base di Unreal'';
  \item si salva la mappa con \textbf{Ctrl+S} nella cartella delle mappe, con il nome
        \textbf{M\_Menu}.
\end{itemize}
(Un dettaglio marginale: il docente nota che qualcosa nel livello risulta \emph{Hidden in
Game}, poi si accorge che non lo è; la trascrizione non dice di che cosa si tratti.)

\paragraph{Il Canvas Panel.} Aperto WB\_Menu, ``la prima cosa, che in questo caso servirà
per adattarlo allo schermo'': il \textbf{Canvas Panel}, ``la classica canvas''. La
\textbf{Hierarchy} si vede ``da qua sotto''. A uno studente che non trova la
\textbf{Palette} il docente mostra il pannello a sinistra dell'editor.

\paragraph{Due bottoni.} ``Adesso andremo a fare due pulsanti molto semplici: uno di
\emph{exit} e uno di \emph{open the game}.'' Nella Palette, la sezione
\textbf{Common} (accanto alle altre: Audio, Input, Panel, Primitive e così via, che la
trascrizione elenca a fatica) contiene il \textbf{Button}; se ne trascinano due al centro
dello schermo. La disposizione passa per qualche tentativo ad alta voce (posizioni e
allineamenti come $-2{,}7$ e $-5$, presto abbandonati: ``ok, mi sono perso''); la versione
pulita, quella che il docente detta alla fine, è questa:
\begin{enumerate}
  \item \textbf{Size}: X \textbf{300}, Y \textbf{80} (``60? No, troppo piccolo'');
  \item \textbf{Anchors} al \textbf{centro}, l'icona centrale del menu delle ancore, ``in
        modo tale che questo bottone sia sempre ancorato al centro''; poi si
        \textbf{azzerano} i valori della \textbf{Position};
  \item \textbf{Alignment} \textbf{0,5} su X e \textbf{1,5} su Y, ``perché così è più
        corretto'': il bottone resta centrato orizzontalmente e si alza di una volta e
        mezza la sua altezza rispetto al centro dello schermo;
  \item per il secondo bottone, ``copiate e incollate letteralmente'' il primo, azzerate
        la posizione e cambiate solo l'allineamento verticale: dopo un ``$-1{,}5$? No,
        troppo'' e un ``$-0{,}5$'', il docente si ferma a \textbf{$-0{,}3$} (``va bene''),
        cioè il secondo bottone sta un po' sotto il centro.
\end{enumerate}
(I valori del secondo bottone, nella trascrizione, sono incerti: il senso è che i due
bottoni stanno uno sopra e uno sotto il centro, con un po' di spazio fra loro.) A uno
studente il copia e incolla non funziona: ``guarda nella Hierarchy: se c'hai due bottoni,
vuol dire che ce ne hai due'', magari sovrapposti; altrimenti, ``metti semplicemente un
altro bottone tramite il Common''.

\paragraph{I nomi.} ``Una cosa che è importante da fare'': nella Hierarchy si rinominano i
due bottoni con \textbf{F2}, \textbf{Button\_StartGame} e \textbf{Button\_QuitGame} (la
trascrizione dice ``Button With Game''). ``Così abbiamo semplicemente perfettamente
nomenclati.''

\paragraph{Il testo dentro il bottone.} ``Una cosa carina che già esiste'': un
\textbf{Text} (dalla sezione Common) trascinato \emph{sopra} un bottone ``si va a incastrare
per benino'', diventa suo figlio e si centra da solo. Nei \textbf{Details} del Text, sezione
\textbf{Content}, il campo \textbf{Text}: ``cancellate e scrivete \emph{Play Game}, o
\emph{Play}, come volete''; nell'altro bottone, \emph{Quit}. Il testo bianco ``non si vede
tantissimo'', ma i colori si cambiano subito dopo.

\paragraph{Lo Style del bottone.} Selezionando il bottone (``mi raccomando''), nei Details
la sezione \textbf{Style} ha tre stati, ciascuno con il suo colore:
\begin{itemize}
  \item \textbf{Normal}: il bottone a riposo;
  \item \textbf{Hovered}: ``vuol dire che ci passi sopra il cursore'';
  \item \textbf{Pressed}: ``se premuto''.
\end{itemize}
(La trascrizione nomina anche i \emph{foreground} e lo stato \emph{Disabled}: nelle
versioni recenti dell'engine lo stile del bottone ha anche i colori del primo piano per
ciascuno stato.) Cliccando sulla barra del colore si apre la ruota dei colori; il docente
mette il Normal \textbf{blu}, ``e avremo questo cambiamento in diretta''.

\paragraph{Copiare un colore esatto.} ``Se vogliamo prendere perfettamente il colore che ci
interessa per spostarlo all'altro bottone'': nella finestra del colore, in basso, ci sono i
campi \textbf{Hex sRGB} e \textbf{Hex Linear}. ``Potete prendere ambi i due'': si copia il
codice con Ctrl+C, si apre lo stesso stato nell'altro bottone e lo si incolla \emph{nello
stesso campo} (``se avete preso sRGB, sempre su sRGB''). È lo stesso principio dei siti
web: ``in HTML, per mettere un colore, i colori sono tutti identificati tramite un codice
colore''.

\paragraph{Un'immagine dietro i bottoni.} Si aggiunge un'\textbf{Image}, allargata dietro i
due bottoni; ma ``come potete vedere, mi copre tutto quanto''. ``Qui ci viene in aiuto la
Hierarchy'': come i livelli di Photoshop, gli elementi del Canvas sono disegnati
nell'ordine dell'elenco, e ``quello che è più sotto'' nell'elenco, stranamente, è quello
``a livello più alto'', disegnato sopra gli altri. L'immagine, aggiunta per ultima, è in
fondo all'elenco e quindi davanti ai bottoni. La si trascina nella Hierarchy fino a far
comparire una \textbf{linea blu} fra il Canvas Panel e il primo bottone, e la si lascia lì:
ora è disegnata per prima, cioè dietro. Poi la si rende trasparente: nel suo
\textbf{Color and Opacity} (che il docente cerca nel Brush), l'\textbf{alfa} a 0,5, ``o
forse un po' più basso'', \textbf{0,2}. ``Questa è la alfa di un colore: l'intensità,
l'opacità.'' Ne viene ``una specie di schermatina piccoletta'', regolata a mano attorno ai
bottoni (figura~\ref{ue12:fig:layout}).

\begin{figure}[htbp]
\centering
\begin{tikzpicture}[font=\footnotesize, x=1cm, y=1cm]
  % schermo
  \draw[very thick, fill=gray!5] (0,0) rectangle (8.0,5.0);
  \node[anchor=north west, font=\scriptsize] at (0.1,4.95) {Canvas Panel (schermo)};
  % immagine trasparente
  \draw[thick, fill=blue!8, opacity=0.9] (2.4,1.15) rectangle (5.6,4.15);
  \node[font=\scriptsize, anchor=north] at (4.0,4.12) {Image, alfa 0,2};
  % bottoni
  \draw[thick, fill=blue!35, rounded corners=2pt] (2.9,2.70) rectangle (5.1,3.30);
  \node at (4.0,3.0) {\textbf{Play Game}};
  \draw[thick, fill=blue!35, rounded corners=2pt] (2.9,1.62) rectangle (5.1,2.22);
  \node at (4.0,1.92) {\textbf{Quit}};
  % ancora al centro
  \draw[red!70!black, thick] (4.0,2.5) circle (0.09);
  \draw[red!70!black, dashed] (0,2.5) -- (8.0,2.5);
  \draw[red!70!black, dashed] (4.0,0) -- (4.0,5.0);
  \node[red!70!black, font=\scriptsize, anchor=west] at (5.7,2.62) {ancora al centro,};
  \node[red!70!black, font=\scriptsize, anchor=west] at (5.7,2.30) {Position 0, 0};
  \node[font=\scriptsize, anchor=west, align=left] at (5.7,3.35) {Alignment\\0,5 / 1,5};
  \node[font=\scriptsize, anchor=west, align=left] at (5.7,1.65) {Alignment\\0,5 / $-0{,}3$\\(circa)};
  % hierarchy
  \begin{scope}[xshift=9.0cm]
    \draw[thick, fill=gray!8] (0,0.6) rectangle (4.2,4.4);
    \node[anchor=north west, font=\scriptsize\bfseries] at (0.1,4.35) {Hierarchy};
    \node[anchor=west] at (0.2,3.6) {[WB\_Menu]};
    \node[anchor=west] at (0.45,3.1) {Canvas Panel};
    \node[anchor=west, fill=blue!10] at (0.7,2.6) {Image};
    \node[anchor=west] at (0.7,2.1) {Button\_StartGame};
    \node[anchor=west] at (0.95,1.65) {Text ``Play Game''};
    \node[anchor=west] at (0.7,1.2) {Button\_QuitGame};
    \node[anchor=west] at (0.95,0.8) {Text ``Quit''};
    \draw[-{Stealth[length=2mm]}, thick] (4.45,3.0) -- node[right, font=\scriptsize,
      align=left]{disegnati\\in ordine:\\l'ultimo\\sta sopra} (4.45,0.9);
  \end{scope}
\end{tikzpicture}
\caption{Il menu nel Designer. I due bottoni sono ancorati al centro dello schermo con la
posizione azzerata; è l'Alignment (il pivot del bottone) a spostarli sopra e sotto il
centro. L'immagine trasparente è stata trascinata in cima alla Hierarchy, così viene
disegnata per prima, dietro ai bottoni. Ricostruzione dalla lezione: i valori del secondo
bottone sono approssimativi.}
\label{ue12:fig:layout}
\end{figure}

\paragraph{Approfondimento: Alignment, ancore e Size To Content.} Con l'ancora al centro e
la posizione a zero, l'\textbf{Alignment} decide quale punto del bottone coincide con il
centro dello schermo: 0 è il bordo sinistro (o superiore), 1 il bordo destro (o inferiore),
0,5 il centro. Valori fuori dall'intervallo da 0 a 1 sono leciti: con Y = 1,5 il punto
d'aggancio sta mezzo bottone \emph{sotto} il bordo inferiore, quindi il bottone sale; con Y
negativo, scende. È il trucco usato a lezione per impilare due bottoni senza toccare la
posizione. La via più ordinata, per menu con molti bottoni, è un \textbf{Vertical Box}
(sezione Panel della Palette) ancorato al centro, con i bottoni dentro: lo spazio fra gli
elementi si regola con il \emph{Padding} di ciascuno e non servono calcoli. Nel libro di
Romero e Sewell (capitolo 8, \emph{WinMenu}) i bottoni sono ancorati al centro e spostati
con la \emph{Position} (per esempio X $-180$, Y 150, Size 360 $\times$ 100), con l'alignment
di default: è la terza via, equivalente. Funzionamento generale dell'UMG, salvo il
riferimento al libro.

% ---------------------------------------------------------------------
\section{Il Level Blueprint del menu: il widget, la Game Mode e la camera}
\label{ue12:sec:levelbp}

\paragraph{Perché il Level Blueprint.} ``Adesso che ce famo con questa roba?'' Il widget
va mostrato appena si apre la mappa del menu, e solo lì. Per questo si usa il
\textbf{Level Blueprint}, ``che abbiamo usato nella prima lezione'': ogni livello ha il
suo, ed è il posto giusto per la logica che appartiene a \emph{quella} mappa, ``perché
vogliamo che questa mappa sia dedicata apposta al nostro menu''. Si apre dalla toolbar
del Level Editor, dal pulsante dei Blueprint (``c'è il cubetto col più''),
\textbf{Open Level Blueprint}.

\paragraph{Il widget sullo schermo.} ``Ci interessa semplicemente usare il Begin Play,
perché vogliamo che si crei già da subito'':
\begin{enumerate}
  \item \textbf{Event BeginPlay} $\rightarrow$ \textbf{Create Widget}, con
        \textbf{Class} = WB\_Menu, ``come abbiamo visto l'altra volta'';
  \item \textbf{Add to Viewport}, con il \textbf{Return Value} del Create Widget nel
        Target. Il docente pensa per un momento di promuovere il widget a variabile, come
        per l'HUD, poi rinuncia: ``no, beh, non lo facciamo, tanto è il level blueprint''.
\end{enumerate}

\paragraph{Il personaggio di troppo.} In Play il menu compare, ma ``è molto strano:
abbiamo il nostro player dietro''. La mappa usa ancora la Game Mode del progetto, che fa
nascere il bot con i suoi controlli. Nel menu, invece, serve ``una pawn senza controlli,
senza niente'', e Unreal ce l'ha già: nelle \textbf{World Settings} della mappa, alla voce
\textbf{GameMode Override}, si sceglie la \textbf{Game Mode base} di Unreal (nella
trascrizione ``la game mode basica che esiste, la default''), che usa come Pawn la
\textbf{Default Pawn}. In Play, ``già è una cosa diversa: non c'è più il nostro character,
sparito, perché in questa mappa ho overridato la modalità Game Mode classica'', quella
``proprio base di Unreal, in cui non c'è character in terza persona''. ``Ma, come potete
vedere, non ho controllo'' (o meglio, come si vedrà, la Default Pawn si muove ancora:
il controllo vero si toglie nel Graph del widget, sezione~\ref{ue12:sec:menugraph}).

\paragraph{La camera fissa.} ``Andiamo a vedere anche un'altra cosa'': nel livello si
aggiunge un \textbf{Camera Actor} (``scrivete \emph{camera actor}, eccola qua''), che sarà
la vista fissa del menu. Insieme, una \textbf{Directional Light} (``e abbiamo messo il
sole, ragazzi''), regolabile a piacere, e un muretto, ``un po' più carino, così capiamo
meglio''. Per dire al gioco di guardare da quella camera, di nuovo nel Level Blueprint:
\begin{itemize}
  \item la \textbf{referenza} alla camera: selezionato il Camera Actor nel livello, nel
        Level Blueprint tasto destro, \textbf{Create a Reference to CameraActor};
  \item il nodo che cambia la vista. Qui il docente sbaglia strada in diretta: prova da
        \emph{Get Player Camera Manager}, non trova il nodo giusto e rimanda (``adesso non
        mi ricordo benissimo, lo vediamo un attimo dopo''). Dopo aver finito il Graph del
        menu, ritrova la soluzione: ``non è che bisogna prendere il camera manager, ma
        bisogna andare sul controller''. Quindi \textbf{Get Player Controller}
        $\rightarrow$ \textbf{Set View Target with Blend}, con la camera nel pin
        \textbf{New View Target}; il \emph{Blend Time} resta al default, ``non ci
        interessa''.
\end{itemize}
``Questo blocco ci permette di blendare una nuova camera'': ``gli diciamo: guarda, questa
è la tua nuova camera, da d'ora in poi'' (figura~\ref{ue12:fig:menugraph}, in alto).

\paragraph{Il Player Start fuori campo.} Ultimo ritocco, emerso aiutando gli studenti: la
Default Pawn ``è una palla'', e se nasce davanti alla camera si vede, o si vede il suo
riflesso. ``Prendete un \textbf{Player Start} e mettetelo fuori campo'', ``così almeno non vi
si vede la pallina''.

\paragraph{Approfondimento: Game Mode, Default Pawn e Player Start.} La \textbf{Game Mode}
è la classe che stabilisce le regole di una mappa e, fra le altre cose, quale
\emph{Pawn} far nascere per il giocatore (\emph{Default Pawn Class}) e con quale
\emph{Player Controller}. Il progetto ne ha una predefinita (in \emph{Project Settings},
\emph{Maps \& Modes}); il \textbf{GameMode Override} delle World Settings la sostituisce
\emph{solo per quella mappa}, ed è il modo standard di avere un menu senza personaggio.
La classe scelta a lezione, \emph{GameModeBase}, usa come Pawn la \textbf{DefaultPawn}:
una sfera con un movimento ``da spettatore'' (WASD e mouse, senza gravità), pensata per
fare prove. Il Pawn nasce in un \textbf{Player Start}; se la mappa non ne ha, nasce
all'origine o nella posizione della camera dell'editor. Un'alternativa ancora più pulita
per un menu è una Game Mode con la \emph{Default Pawn Class} impostata su \emph{None}: il
Player Controller esiste (e quindi esistono il cursore e gli input dell'interfaccia), ma
nessun Pawn nasce. Funzionamento generale dell'engine.

\paragraph{Approfondimento: Set View Target with Blend nel libro.} Romero e Sewell
descrivono il nodo nel capitolo 15 (\emph{Blueprints Tips}): ``è una funzione della classe
Player Controller'', ``usata per passare la vista di gioco fra camere diverse''; ``il
parametro di input New View Target è l'actor da usare come view target, di solito una
camera''. Il loro esempio è proprio in un Level Blueprint: quando il giocatore entra nella
stanza del tesoro, la vista passa alla camera della stanza. Questo spiega l'errore del
docente: il \emph{Player Camera Manager} gestisce la vista, ma la funzione esposta ai
Blueprint per cambiarla appartiene al \emph{Player Controller}. Il \textbf{Blend Time} a 0
fa uno stacco netto; un valore positivo fa una transizione morbida dalla vista precedente,
con una curva scelta in \emph{Blend Func}.

% ---------------------------------------------------------------------
\section{Il Graph del menu: cursore, input e bottoni}
\label{ue12:sec:menugraph}

\paragraph{Event Construct.} Si torna in WB\_Menu, scheda \textbf{Graph}: ``vediamo le
funzionalità del nostro menu quando clicchiamo le cose''. Dei tre eventi già presenti
``non ci serve il Pre Construct, non ci serve il Tick''. Serve l'\textbf{Event Construct}:
``vuol dire quando questo widget è costruito'', ``è un po' il Begin Play'', ``viene
differenziato dal nome, giusto per non confonderti''.

\paragraph{Il cursore e l'input dell'interfaccia.} In Play, per ora, ``mi posso muovere''
(la Default Pawn) e il mouse non si vede. Nell'Event Construct:
\begin{enumerate}
  \item \textbf{Get Player Controller} $\rightarrow$ \textbf{Set Show Mouse Cursor},
        con la spunta attiva: il cursore resta visibile anche cliccando nella finestra di
        gioco. ``Però, ovviamente, posso ancora muovermi'';
  \item \textbf{Set Input Mode UI Only} (``enable input UI'', dice prima il docente, poi
        trova il nodo), con lo stesso Player Controller: ``questo ci permette di
        utilizzare tranquillamente il nostro widget''. Il pin \emph{Widget to Focus} resta
        vuoto: ``no, tanto non ci serve''.
\end{enumerate}
In Play, ora, ``se noi proviamo a muovere lateralmente, non si muove più''. ``Mi
raccomando: questa cosa, quando noi cliccheremo un pulsante, dovrà essere disattivata e
rimessa in un'altra modalità.''

\paragraph{I colori di hover e pressione.} Con il cursore visibile si possono provare gli
stati dello Style. Nel Designer il docente cambia il colore di \textbf{Hovered}: ``voglio
che quando io ci passo sopra con il mouse diventi tipo arancione, sì, un bell'arancione''
(``lo vediamo in preview, è normale''), e lo copia, con il codice esadecimale, sull'altro
bottone. Il risultato, ammette, è ``molto Ikea''. Il \textbf{Pressed} di default diventa
grigio, ``si vede pochissimo'': lo si mette \textbf{rosso}, su entrambi i bottoni.

\paragraph{I bottoni.} Nel Graph, a sinistra, fra le variabili compaiono
\textbf{Button\_QuitGame} e \textbf{Button\_StartGame} (sono variabili perché i bottoni,
nel Designer, hanno \emph{Is Variable} attivo di default). Ogni bottone ha i suoi eventi,
``l'abbiamo visto anche l'altra volta'': \emph{On Pressed}, \emph{On Released},
\emph{On Hovered}, \emph{On Unhovered} e \textbf{On Clicked}. Si prende l'On Clicked di
entrambi.
\begin{itemize}
  \item \textbf{On Clicked (Button\_QuitGame)} $\rightarrow$ \textbf{Quit Game}: ``è
        piuttosto semplice, letteralmente Quit Game è già presente, di standard
        proprio''. (``Adesso non mi ricordo se ti esce proprio dall'editor'': si veda
        l'approfondimento più sotto.)
  \item \textbf{On Clicked (Button\_StartGame)}: prima \textbf{Get Player Controller}
        $\rightarrow$ \textbf{Set Input Mode Game Only}, ``molto importante'', ``in modo
        tale che ritornino gli input al gioco''; poi \textbf{Open Level (by Name)}, con il
        nome della mappa di gioco. Il docente lo cerca nel Content Browser (``era
        MLevel1? Ah, questa no, era M level 2?'') e trova la mappa del template,
        \textbf{ThirdPersonMap}; prova a copiarne il nome con F2 (``non lo fa, va bene''),
        e lo scrive a mano.
\end{itemize}
Tornati nella mappa del menu e salvato tutto, in Play: Play Game ``ci carica la nostra
bellissima mappina'' (figura~\ref{ue12:fig:menugraph}).

\begin{figure}[htbp]
\centering
\begin{tikzpicture}[font=\footnotesize,
  ue12node/.style={draw, thick, rounded corners=4pt, align=center, fill=gray!8,
                  minimum height=0.9cm, text width=2.3cm},
  ue12ev/.style={ue12node, fill=red!12},
  ue12pure/.style={ue12node, fill=green!10, rounded corners=8pt},
  ue12exec/.style={-{Stealth[length=2.5mm]}, very thick, black!80},
  ue12data/.style={-{Stealth[length=2mm]}, thick, blue!60!black}]
  \node[font=\scriptsize\bfseries, anchor=west] at (-1.3,1.0) {Level Blueprint di M\_Menu};
  \node[ue12ev, text width=1.7cm] (bp) at (0,0) {\textbf{Event BeginPlay}};
  \node[ue12node] (cw) at (2.8,0) {\textbf{Create Widget}\\{\scriptsize Class: WB\_Menu}};
  \node[ue12node, text width=2.5cm] (av) at (5.8,0) {\textbf{Add to Viewport}\\{\scriptsize Target: Return
    Value}};
  \node[ue12node, text width=2.6cm] (vt) at (8.9,0) {\textbf{Set View Target with
    Blend}\\{\scriptsize Blend Time 0}};
  \node[ue12pure] (gpc1) at (7.4,-1.8) {Get Player Controller};
  \node[ue12pure, fill=blue!10] (cam) at (10.2,-1.8) {CameraActor\\{\scriptsize
    (referenza)}};
  \draw[ue12exec] (bp) -- (cw); \draw[ue12exec] (cw) -- (av); \draw[ue12exec] (av) -- (vt);
  \draw[ue12data] ([yshift=-3mm]cw.east) -- ([yshift=-3mm]av.west);
  \draw[ue12data] (gpc1.north) -- node[left, font=\scriptsize]{Target} ([xshift=-6mm]vt.south);
  \draw[ue12data] (cam.north) -- node[right, font=\scriptsize]{New View Target}
    ([xshift=6mm]vt.south);
  % widget graph
  \node[font=\scriptsize\bfseries, anchor=west] at (-1.3,-3.2) {Graph di WB\_Menu};
  \node[ue12ev, text width=1.7cm] (con) at (0,-4.2) {\textbf{Event Construct}};
  \node[ue12node] (smc) at (2.8,-4.2) {\textbf{Set Show Mouse Cursor}\\{\scriptsize
    spunta attiva}};
  \node[ue12node] (ui) at (5.7,-4.2) {\textbf{Set Input Mode UI Only}};
  \node[ue12pure] (gpc2) at (4.2,-5.9) {Get Player Controller};
  \draw[ue12exec] (con) -- (smc); \draw[ue12exec] (smc) -- (ui);
  \draw[ue12data] (gpc2.north) -- (smc.south); \draw[ue12data] (gpc2.north) -- (ui.south);
  \node[ue12ev, text width=2.6cm] (oq) at (0.3,-7.6) {\textbf{On Clicked}\\{\scriptsize
    (Button\_QuitGame)}};
  \node[ue12node] (qg) at (3.6,-7.6) {\textbf{Quit Game}};
  \draw[ue12exec] (oq) -- (qg);
  \node[ue12ev, text width=2.6cm] (os) at (0.3,-9.3) {\textbf{On Clicked}\\{\scriptsize
    (Button\_StartGame)}};
  \node[ue12node, text width=2.7cm] (gm) at (3.8,-9.3) {\textbf{Set Input Mode Game Only}};
  \node[ue12node, text width=2.8cm] (ol) at (7.2,-9.3) {\textbf{Open Level (by
    Name)}\\{\scriptsize ThirdPersonMap}};
  \node[ue12pure] (gpc3) at (3.8,-11.0) {Get Player Controller};
  \draw[ue12exec] (os) -- (gm); \draw[ue12exec] (gm) -- (ol);
  \draw[ue12data] (gpc3) -- (gm);
\end{tikzpicture}
\caption{I grafi del menu. In alto, il Level Blueprint della mappa M\_Menu crea il widget,
lo mostra e passa la vista alla camera fissa. In basso, il widget prende il controllo del
cursore quando nasce, esce dal gioco con Quit e, con Play Game, restituisce l'input al
gioco e apre la mappa di gioco. Ricostruzione dalla lezione.}
\label{ue12:fig:menugraph}
\end{figure}

\paragraph{Approfondimento: il menu nel libro.} Romero e Sewell costruiscono un menu molto
simile nel capitolo 8 (\emph{Creating Constraints and Gameplay Objectives}), il
\emph{WinMenu} che compare a fine partita: due bottoni \emph{Btn restart} e \emph{Btn
quit} ancorati al centro, ciascuno con un Text dentro (``font size 48''), e gli eventi
aggiunti dal pulsante \emph{+} accanto a \emph{On Clicked}, in fondo ai Details del bottone.
Il bottone Restart usa \emph{Open Level (by Object Reference)}, con la mappa scelta da un
menu a tendina, seguito da \emph{Remove from Parent}, ``che rimuove il widget dalla
vista''; il bottone Quit usa \emph{Quit Game}, ``così che il giocatore possa chiudere il
gioco''. Il menu viene mostrato dal personaggio, in un Custom Event \emph{End Game}: Set
Game Paused, \emph{SET Show Mouse Cursor} dal Get Player Controller (``per permettere al
giocatore di riprendere il controllo del mouse''), Create Widget, Add to Viewport. Nel
capitolo 15 il libro elenca anche le tre modalità di input:
\begin{itemize}
  \item \textbf{Set Input Mode Game Only}: ``solo il Player Controller riceve gli eventi
        di input'';
  \item \textbf{Set Input Mode UI Only}: ``solo la UI riceve gli eventi di input'';
  \item \textbf{Set Input Mode Game and UI}: ``la UI ha la priorità nel gestire un evento
        di input, ma se la UI non lo gestisce, lo riceve il Player Controller''; l'esempio
        è un negozio, in cui si scelgono le opzioni con il mouse ma ci si può allontanare
        con le frecce.
\end{itemize}
Una differenza utile da notare: \emph{Open Level (by Object Reference)} non si rompe se la
mappa viene rinominata, mentre \emph{Open Level (by Name)}, usato a lezione, richiede il
nome esatto (si può scrivere il nome semplice, \emph{ThirdPersonMap}, o il percorso
completo, se esistono due mappe con lo stesso nome).

\paragraph{Approfondimento: Construct, Quit Game e il cursore nella mappa di gioco.}
Alcune nozioni generali, non della lezione. \textbf{Event Pre Construct} viene eseguito
anche nell'editor, ogni volta che il widget viene ridisegnato nel Designer: serve per
anteprime e impostazioni visive; \textbf{Event Construct} solo in gioco, quando il widget
viene aggiunto allo schermo (e di nuovo se lo si rimuove e lo si aggiunge). \textbf{Quit
Game}, durante il Play in Editor, \emph{termina la sessione di gioco} e riporta
all'editor, senza chiudere l'editor; in un gioco pacchettizzato chiude l'applicazione. Il
dubbio del docente si scioglie così. Infine, \emph{Open Level} distrugge il mondo corrente
con tutti i suoi actor, widget compresi, e crea un Player Controller nuovo: per questo,
nella mappa di gioco, il menu sparisce da solo e il cursore torna nascosto, anche senza un
\emph{Remove from Parent} o un \emph{Set Show Mouse Cursor} a falso. Il \emph{Set Input
Mode Game Only} prima dell'Open Level, raccomandato dal docente, è comunque una buona
abitudine: il caso in cui non è superfluo è quello di un menu che si chiude \emph{senza}
cambiare livello, come un menu di pausa (nel libro, capitolo 11, il \emph{Resume} fa
Remove from Parent, Set Show Mouse Cursor a falso e Set Game Paused a falso).

% ---------------------------------------------------------------------
\section{Il diorama, e i problemi degli studenti}
\label{ue12:sec:diorama}

\paragraph{Il diorama.} Con il menu funzionante, ``passiamo sopra'' all'aspetto. Sullo
sfondo, ``un piccolo diorama, una piccola box, in cui ci andiamo a mettere tutto quello che
vogliamo''. ``Volendo, potete fare tutto un menu interattivo'': gli ostacoli, i Blueprint
creati nel corso, i bottoni, la torretta. Il docente si limita ai modelli:
\begin{itemize}
  \item il \textbf{bot}: dal Content Browser, o la sua Skeletal Mesh, ``o li metto
        direttamente l'animazione in Idle'': trascinando nel livello un asset di animazione,
        Unreal crea un actor con la mesh che la riproduce;
  \item la \textbf{porta} modulare del compito della volta scorsa, \textbf{SM\_Modular\_\allowbreak Wall\_Door}
        (``wall door, eccola qua''), con le ante \textbf{SM\_Door};
  \item qualche ritocco alla \textbf{camera} e alla \textbf{luce}: l'intensità a 10 è
        ``enorme'', poi 5, poi valori molto più bassi (la trascrizione, qui, è quasi
        incomprensibile: si sentono ``0,01'' e ``1'').
\end{itemize}
``Questo caso è più una cosa artistica, ragazzi, quindi potete fare un po' come volete.''
Le funzioni di base, entrare nel gioco e uscire, ci sono. ``Questo è un menu che io di
solito faccio quando faccio i miei prototipi'', cioè ``quando devo semplicemente mostrare
qualche cosa a livello di gameplay''. Un menu delle \emph{opzioni} ``è già un discorso
molto più'' complesso, ``perché bisogna andare a vedere gli slider'': non si tratta, ``perché
tanto poi andiamo a fare qualcosa di più gameplayoso''.

\paragraph{I problemi dei presenti.} Diversi studenti hanno ancora il personaggio nel
menu.
\begin{itemize}
  \item \textbf{La camera.} Ripasso passo per passo: selezionare la camera nel livello
        (``proprio diventa arancione''), Open Level Blueprint, tasto destro, Create a
        Reference to CameraActor (``devi stare dentro il Level Blueprint''); Get Player
        Controller, Set View Target with Blend, la camera nel New View Target.
  \item \textbf{Il personaggio che nasce ancora.} ``Devi passare alla Game Mode
        default'': \textbf{World Settings}, \textbf{GameMode Override}. Uno studente non
        trova la finestra delle World Settings (si apre dal menu \emph{Window} dell'editor,
        se è chiusa). Poi il Player Start fuori dalla camera.
  \item \textbf{L'errore su Add to Viewport.} ``Mi dà un problema proprio su Add to
        Viewport'': ``perché il target di base è Self, quindi non so se devo cambiare il
        target''. Sì: il Target deve essere il \textbf{Return Value} del Create Widget.
        ``Probabilmente devi dare return value'' al target.
  \item \textbf{Il mouse che non si vede.} Va impostato nell'Event Construct del widget:
        Get Player Controller, Set Show Mouse Cursor, Set Input Mode UI Only. Lo studente
        aveva cancellato l'Event Construct (``basta che fai tasto destro'' e lo si ricrea)
        e non aveva messo la \textbf{spunta} nel Set Show Mouse Cursor (``ah, però non
        avevo messo il tick''): il nodo, senza spunta, imposta il cursore a \emph{falso}.
        ``Funziona? Sì.''
\end{itemize}
Chiuso il menu, si torna alla ThirdPersonMap.

% ---------------------------------------------------------------------
\section{Interfacce e Blueprint Component}
\label{ue12:sec:concetti}

``Andiamo sul nostro amico fidato'': il docente apre il suo bot (``ah no, questo è Boto:
dobbiamo andare su Boto2''), cioè BP\_Bot2. ``Oggi parleremo delle interfacce, ragazzi.''

\paragraph{L'interfaccia, un intermediario.} ``Immaginatevi che un'interfaccia è un
intermediario fra due o più Blueprint.'' Il docente mostra quella già pronta nel suo
progetto, nella cartella dei componenti dei Blueprint: un tipo particolare di Blueprint
che contiene solo tre \textbf{funzioni}, \emph{Interaction Start}, \emph{Interaction End}
e \emph{Interaction}. L'immagine: un \textbf{volo di linea}, ``un'interconnessione tra due
o più poli'', che ``trasporta molto spesso le stesse cose, passeggeri di un tipo,
passeggeri di un altro''. ``Quindi è un qualcosa che va a trasportare varie funzioni da una
parte all'altra.'' Nel caso di oggi l'interfaccia comunica tre messaggi: \emph{sto
iniziando un'interazione} (ti sto guardando), \emph{sto finendo un'interazione} (non ti
guardo più), \emph{sto interagendo} (ho premuto il tasto).

\paragraph{Il componente, una funzionalità isolata.} Il lavoro di \emph{vedere se c'è
qualcosa di interagibile} non lo fa il personaggio, ma un \textbf{Blueprint Component}:
``dei componenti che noi possiamo aggiungere all'interno dei nostri actor per dare
funzionalità in modo isolato''. Per esempio: ``questo actor può tracciare degli
interactable''; ``questo actor ha la vita''. La vita, fatta la volta scorsa \emph{dentro}
il personaggio, si potrebbe spostare in un componente, ``perché magari abbiamo anche i
nemici che hanno la vita'', o ``degli oggetti che si possono rompere''. Lo scopo è
``isolare sempre più possibile le funzionalità, e quindi essere anche modulari''. Infatti
il BP\_Bot2 del docente, a parte un input nuovo, ``non ha niente di nuovo'': si gira verso
l'oggetto e suona un Montage, ma ``non ha niente per tracciare'': il trace è tutto nel
componente.

\paragraph{Actor Component o Scene Component.} Si crea una \textbf{Blueprint Class} e, fra
le classi comuni, non un Actor ma un componente. Ce ne sono due:
\begin{itemize}
  \item l'\textbf{Actor Component} ``mette semplicemente la funzionalità'';
  \item lo \textbf{Scene Component} ``è un Actor Component che ha anche una scena'',
        cioè un \textbf{Transform}: ``non so se vi ricordate la Scene, quando l'abbiamo
        messa: il punto vuoto''. ``Almeno si può spostare nello spazio.''
\end{itemize}
Qui serve il secondo, perché il trace è ``un raggio che da quel punto parte e in un punto
finisce'', e il punto di partenza si deve poter spostare. Il docente lo mostra nel suo bot:
spostando il componente nel Blueprint del personaggio, in Play il raggio parte dal nuovo
punto; lo rimette ``al centro della testa, più o meno''. Il nuovo Blueprint:
\textbf{BPC\_RayInteract} (BPC sta per \emph{Blueprint Component}; il docente, che ne ha già
uno, chiama il suo BPC\_RayInteract1).

\paragraph{Il componente aperto.} ``Come potete vedere, presenta delle similitudini con
tutti i Blueprint.'' ``L'unica cosa: non ha una viewport, perché non ha una forma fisica'',
``tra virgolette, perché è una scena, ma è vuota''. ``E ricordiamoci che è un componente:
quindi sarà qualcosa che deve essere per forza attaccato a qualcos'altro.''

\paragraph{Che cos'è una funzione.} Prima di scrivere la logica, una domanda: chi ha già
sentito parlare di \textbf{funzioni} in informatica? Risponde \textbf{Federico}, che ha
studiato da programmatore: ``la funzione è un insieme di istruzioni: chiami la funzione e
fai quelle istruzioni''. ``Esattamente.'' Il docente indica una catena di nodi nel bot: ``gli
eventi non possono essere messi come funzioni, ma questo qua può essere una funzione''. È
``un modo per ordinare anche il lavoro in varie tipologie di funzionalità, che possono
essere magari trasportabili da un punto all'altro'': ``sempre questo senso di modularità''.
(``Grande Federico, sempre miglior assistente.'')

\begin{figure}[htbp]
\centering
\begin{tikzpicture}[font=\footnotesize,
  ue12box/.style={draw, thick, rounded corners=5pt, align=center, inner sep=5pt},
  ue12msg/.style={-{Stealth[length=2.5mm]}, very thick, violet!70!black},
  ue12own/.style={-{Stealth[length=2.2mm]}, thick, dashed, black!70}]
  \node[ue12box, fill=gray!8, minimum width=4.4cm, minimum height=4.6cm] (bot) at (0,-0.3) {};
  \node[anchor=north, font=\bfseries] at (bot.north) {BP\_Bot2};
  \node[ue12box, fill=blue!10, text width=3.6cm] (cam) at (0,0.8) {Camera Component};
  \node[ue12box, fill=orange!18, text width=3.6cm] (comp) at (0,-1.4)
    {\textbf{BPC\_RayInteract}\\{\scriptsize Scene Component: ogni frame una Sphere Trace
    lungo la camera; variabile Interactable}};
  \draw[ue12own] (comp.north) -- node[right, font=\scriptsize, text width=1.9cm,
    align=left]{Get Owner,\\Get Component by Class} (cam.south);
  \node[ue12box, fill=violet!10, text width=3.2cm] (bpi) at (6.6,-0.3)
    {\textbf{BPI\_InteractionRay}\\{\scriptsize Blueprint Interface}\\[2pt]{\scriptsize
    Interaction Start\\Interaction End\\Interaction}};
  \node[ue12box, fill=green!10, text width=3.3cm] (obj1) at (11.6,1.0)
    {\textbf{BP\_InteractableTest}\\{\scriptsize implementa la BPI: Print String, widget
    ``E'', luce}};
  \node[ue12box, fill=green!5, text width=3.3cm] (obj2) at (11.6,-1.6) {{\scriptsize un
    altro oggetto}\\{\scriptsize (bottone, ascensore\dots) con la stessa BPI e un'altra
    risposta}};
  \draw[ue12msg] (bot.east |- cam) -- node[above, font=\scriptsize, text width=2.4cm,
    align=center]{Interaction\\(tasto E)} (bpi.west |- cam);
  \draw[ue12msg] (comp.east) -- node[below, font=\scriptsize, text width=2.4cm,
    align=center]{Start / End\\(guardare, smettere)} (bpi.west |- comp);
  \draw[ue12msg] (bpi.east) -- (obj1.west);
  \draw[ue12msg] (bpi.east) -- (obj2.west);
\end{tikzpicture}
\caption{I ruoli del sistema di interazione. Il componente, attaccato al personaggio, trova
la camera del suo proprietario, traccia ciò che il giocatore guarda e invia i messaggi di
inizio e fine; il personaggio invia il messaggio Interaction quando si preme E. I messaggi
viaggiano attraverso l'interfaccia: chi li riceve decide che cosa fare, e un actor che non
la implementa li ignora. Schema ricostruito dalla lezione.}
\label{ue12:fig:ruoli}
\end{figure}

\paragraph{Approfondimento: i componenti nel libro.} Romero e Sewell dedicano ai componenti
la seconda metà del capitolo 18 (\emph{Creating Blueprint Libraries and Components}).
``Quando creiamo un Actor Blueprint aggiungiamo spesso componenti con una funzionalità
incapsulata, pronta all'uso, come Projectile Movement, Static Mesh e Collision''; ``possiamo
anche creare i nostri Actor Component con i Blueprint''. Nel pannello \emph{Pick Parent
Class} le due classi sono fra quelle \emph{comuni}: ``lo Scene Component è una classe figlia
dell'Actor Component che ha la struttura Transform (location, rotation e scale). Grazie al
Transform, uno Scene Component può essere attaccato a un altro Scene Component.'' E, come a
lezione: ``quando si scrive la logica di un componente nell'Event Graph, si può ottenere la
referenza all'actor che usa il componente con il nodo \textbf{Get Owner}''. I due esempi del
libro mostrano bene la differenza: \emph{BP\_ExpLevelComp} è un Actor Component che gestisce
punti esperienza e livelli (non ha bisogno di una posizione), mentre \emph{BP\_CircularMovComp}
è uno Scene Component che ruota attorno all'actor e a cui si attacca una Static Mesh, per
simulare uno scudo rotante: ``lo Scene Component è usato come posizione di riferimento per
altri componenti''. Il componente della lezione è dello stesso tipo: la sua posizione è il
punto di partenza del raggio.

% ---------------------------------------------------------------------
\section{Il componente BPC\_RayInteract: la Sphere Trace}
\label{ue12:sec:componente}

\paragraph{La funzione LookForInteractions.} Nel componente si crea una nuova funzione
(``eccolo che vi si crea questo pannellino violaceo''), \textbf{LookForInteractions}.
Dentro, una \textbf{Sphere Trace By Channel}. ``Seguitemi per bene.''
\begin{itemize}
  \item \textbf{Radius}: \textbf{5}, ``perché se no non vi apparirà mai''. Con raggio zero
        la Sphere Trace degenera in un raggio sottilissimo (e il debug quasi non si vede);
        5 unità bastano a renderla visibile e un po' più tollerante.
  \item \textbf{Start}: \textbf{Get World Location} del componente stesso, ``la posizione
        nel mondo del nostro Scene Component: dovrà partire da qui''.
  \item \textbf{End}: da costruire, ed è la parte con il trucco.
\end{itemize}

\paragraph{Prima la camera: Get Owner e Get Component by Class.} Il raggio deve seguire lo
sguardo del giocatore, cioè la \emph{rotazione della camera}. Il componente, però, non sa
nulla della camera del personaggio: deve trovarla. Nell'\textbf{Event Graph} del
componente, all'\textbf{Event BeginPlay}:
\begin{enumerate}
  \item \textbf{Get Owner}: ``la referenza del proprietario di questo Scene Component'';
  \item dal proprietario, \textbf{Get Component by Class} (``non \emph{components},
        \emph{component}, perché se no vi dà un array''), con \textbf{Component Class} =
        \textbf{Camera Component}: ``sarà il nostro punto principale da cui andremo a fare il
        nostro trace'';
  \item il risultato si promuove a variabile, \textbf{CameraComponent}, e si collega tutto
        al Begin Play.
\end{enumerate}
``Perdonate, sono 30.000 componenti: un componente su componente su componente.'' In
sintesi: ``in Begin Play io prendo il mio proprietario, prendo la sua camera, mi salvo la
referenza alla sua camera, e poi la useremo all'interno del nostro componente''.

\paragraph{La camera, se c'è.} Il componente deve funzionare anche su un actor
\emph{senza} camera: ``se il nostro proprietario non ha una camera, allora prendi te stesso
come pivot''. Quindi un \textbf{Select} fra due oggetti:
\begin{itemize}
  \item l'indice è un booleano, \textbf{Is Valid} (``la F verde'', la versione \emph{pura}
        della funzione) sulla variabile CameraComponent: ``se il componente camera esiste,
        quindi il nostro actor ha effettivamente una camera, e quindi questa variabile
        CameraComponent ha qualcosa dentro'';
  \item su \textbf{True}, CameraComponent: ``usami la camera come punto di riferimento'';
  \item su \textbf{False}, \textbf{Self}: ``se no, usi te stesso''.
\end{itemize}
Il Select era già comparso con la torretta, per scegliere fra le due bocche di fuoco: lì
con un \emph{Random Bool}, qui con l'Is Valid. ``So che è un po' complesso.''

\paragraph{La ``mandragata'': il tipo del Select.} Il problema è il \emph{tipo} delle
opzioni. ``Il Select praticamente c'è un piccolo problema'': il nodo è jolly (\emph{wildcard})
e assume il tipo della \emph{prima} cosa che gli si collega. Collegando prima la camera, il
Select diventa un Select di \emph{Camera Component Object Reference}, e Self (il componente
BPC\_RayInteract) non entra più; collegando prima Self, il Select diventa un Select di
riferimenti a BPC\_RayInteract, e la camera non entra. Entrambi però \emph{sono} Scene
Component: la camera è uno Scene Component, e lo è anche il nostro componente. Il trucco
(``qui bisogna fare una mandragata, ragazzi''):
\begin{enumerate}
  \item si crea una variabile qualunque di tipo \textbf{Scene Component Object Reference}
        (il docente la chiama \emph{Scene}: ``e poi non ci servirà a niente'');
  \item se ne mette un Get nel Select: ora le opzioni sono di tipo \emph{Scene Component};
  \item si \textbf{cancella} il Get e si collegano CameraComponent e Self, che sono entrambi
        accettati.
\end{enumerate}
``Ecco fatto che abbiamo praticamente fregato Unreal.''

\paragraph{La domanda di Filippo.} ``Non ho capito perché abbiamo creato la variabile Scene
prima.'' Il docente rifà la prova davanti a tutti: un Select nuovo, il cursore sul pin, la
scritta ``Scene Component Object Reference''; poi Self in un Select vuoto, e il nodo diventa
``Self Object Reference'', ``mi farà fare solamente un select fra due self, non due Scene
Component''. ``Per fregare Unreal'' serve ``una variabile che appartiene generalmente allo
Scene Component''. Uno studente prova a trasformare la variabile in Self: ``non è che te lo
trasforma: lo devi cancellare, e poi tasto destro, scrivi \emph{self}''. Filippo insiste: e
se avessimo collegato subito la camera? Prova del docente: ``Select camera component, camera
component object reference, e io, se provo ad attaccare Self, non me lo attacca''. ``Ok,
capito.''

\paragraph{Il punto finale.} Dal risultato del Select, ``quello che abbiamo sempre fatto'':
\begin{enumerate}
  \item \textbf{Get Forward Vector} (della camera, o del componente stesso);
  \item moltiplicato per una distanza: il secondo pin del Multiply si converte in float
        (``tasto destro, mi raccomando: \emph{To Float}''), si promuove a variabile,
        \textbf{DistanceTrace}, con valore \textbf{300};
  \item il prodotto si \textbf{somma} (Add) allo Start, e il risultato va nell'\textbf{End}.
\end{enumerate}

\paragraph{Perché lo Start è il componente.} Una studentessa chiede perché lo Start sia la
World Location: ``in relazione a dove attacchiamo il componente?'' ``Esatto.'' ``Il punto
iniziale deve essere sempre il punto, il Transform, del nostro componente''; ``il punto che
deve diversificare è la camera''. Il docente lo mostra nel bot: guardando in basso il
raggio va in basso, guardando di lato va di lato, ma parte sempre dal componente, ``così
magari posso posizionarlo avanti, posizionarlo indietro, posizionarlo qui, posizionarlo
là'' (figura~\ref{ue12:fig:look}).

\paragraph{Il componente nel personaggio.} Nel BP\_Bot2 (``a bot2, in questo caso, è il
vostro character''), \textbf{Add} e il nome del componente: BPC\_RayInteract compare fra i
componenti del personaggio. Per vederlo si imposta il \textbf{Draw Debug Type} della Sphere
Trace (``lo metto a \emph{For One Frame}''). In Play, però, non succede niente: ``ah,
certo, ecco perché: non facciamo LookForInteractions''. La funzione va chiamata
nell'\textbf{Event Tick} del componente (``ah, nel tick, infatti''), così il controllo
avviene a ogni frame. Con il debug persistente (\emph{Persistent}) il raggio lascia una scia
``gigantesca''; con \emph{For Duration} resta qualche istante; con \emph{For One Frame}
``rimane un po' più carino'', e si vede solo il raggio attuale. (La trascrizione, qui, è
piena di prove e di parole storpiate: ``la duration vado a activate, troppo corto'', ``Lugano'',
``Adobe''; anche una confusione momentanea fra il bot e il bot2 del docente.)

\paragraph{Approfondimento: Sphere Trace, Tick e costo.} Una \textbf{Sphere Trace} è un
\emph{line trace} ``spesso'': fa scorrere una sfera del raggio indicato dallo Start
all'End e restituisce il primo oggetto che la sfera tocca nel canale scelto. Rispetto a un
raggio sottile perdona di più la mira (con 5 unità di raggio, circa 5 cm, conta anche
sfiorare il bordo di un oggetto), ed è il motivo per cui molti sistemi di interazione la
preferiscono. Farla nel \textbf{Tick}, una volta per frame, costa poco: una singola trace di
300 unità è un'operazione leggera, e il sistema la esegue solo per il personaggio del
giocatore. Se gli actor con il componente fossero molti (per esempio dei nemici), conviene
ridurre la frequenza con un timer (come quelli della stamina, a 0,1 secondi) o disattivare
il Tick del componente quando non serve (\emph{Set Component Tick Enabled}). Una Sphere
Trace può stare senza problemi dentro una funzione, perché non ha uno stato interno da
ricordare fra una chiamata e l'altra; il Do Once e il Flip Flop, invece, sì (si veda più
avanti). Funzionamento generale dell'engine.

% ---------------------------------------------------------------------
\section{Il Trace Channel e la Blueprint Interface}
\label{ue12:sec:bpi}

\paragraph{Un canale per le interazioni.} ``Ha colpito qualcosa? Sì, no.'' Prima di
leggere il risultato serve un canale dedicato, ``che ci permette di identificare gli
oggetti che vogliamo che siano interagibili'': ``se no, praticamente, interagirà con
tutto''. ``Come abbiamo visto anche nella nostra torretta'': \textbf{Project Settings},
\emph{Collision}, \textbf{New Trace Channel}, nome \textbf{InteractionTrace}, risposta di
default \textbf{Ignore}. Il docente, che ce l'ha già, non lo rifà. Nella Sphere Trace, il
\textbf{Trace Channel} diventa InteractionTrace: ``mi raccomando, se non vi appare, fate
compila, perché ha bisogno un attimo di compilare il nuovo Trace Channel''.

\paragraph{Che cosa abbiamo colpito.} Dopo la trace, un \textbf{Branch} sul suo Return
Value (``ha colpito qualcosa?'') e un \textbf{Break Hit Result} sull'Out Hit, da cui, ``vi
ricordate, no?'', escono l'actor colpito, il Physical Material ``e tutto''. Il piano:
``useremo l'interfaccia, la Blueprint Interface, per andare a comunicare con questo actor'':
``ha questo porto in cui può accogliere questa tipologia di interazione? Sì? Benissimo,
allora fai l'interazione''. Dall'\textbf{Hit Actor} si chiede \textbf{Does Object
Implement Interface}: ``ehi, questo actor ce l'ha l'interface?'' Ma l'interfaccia non
esiste ancora: ``ci fermiamo qua per adesso, creiamo l'interface, e ci fermiamo''. (Il
docente parla qui anche di promuovere l'Hit Actor a variabile; la variabile di appoggio
vera e propria, \emph{Interactable}, si crea dopo la pausa, sezione~\ref{ue12:sec:logica}.)

\paragraph{La Blueprint Interface.} Nella cartella \emph{Blueprints}, \emph{Components},
tasto destro, \textbf{Blueprint}, ``non Blueprint Class questo, ma Blueprint questo'': il
sottomenu offre \emph{Blueprint Class}, \emph{Blueprint Function Library} e
\textbf{Blueprint Interface}. Il nome: \textbf{BPI\_InteractionRay}. Aperta:
\begin{itemize}
  \item il grafo è \textbf{in sola lettura}: ``non possiamo far niente'', ``cannot create
        new nodes in read only graph''. Un'interfaccia contiene solo \emph{dichiarazioni};
  \item si possono aggiungere \textbf{funzioni}: la prima c'è già e si rinomina (con F2),
        \textbf{Interaction Start}; con \emph{Add Function} le altre due,
        \textbf{Interaction End} e \textbf{Interaction}.
\end{itemize}
``Abbiamo messo che in questo cargo, questo trasporto, ci stanno Interaction Start,
Interaction End e Interaction: quindi queste funzioni possono essere passate tra di loro.''

\paragraph{Il collegamento.} Nel componente, il pin \textbf{Interface} di Does Object
Implement Interface diventa \mbox{BPI\_}\allowbreak InteractionRay. ``Se quel qualcosa che ho colpito ha questa
interfaccia, allora andremo a fare tutto il procedimento di interazione.'' ``Ci sono domande?
Mi sembra di capire no.'' Sono le 17:26: pausa fino alle 17:55, ``diciotto''.

\paragraph{Approfondimento: la Blueprint Interface nel libro.} Romero e Sewell presentano le
interfacce nel capitolo 16 (\emph{Introduction to VR Development}), nel paragrafo
\emph{Blueprint Communication using interfaces}: ``una Blueprint Interface è un tipo speciale
di Blueprint che contiene solo nomi di funzioni e parametri. Si usa per permettere la
comunicazione fra tipi diversi di Blueprint.'' Si crea dal pulsante \emph{Add} del Content
Browser, sottomenu \emph{Blueprints}, \emph{Blueprint Interface}. L'esempio è il template VR:
l'interfaccia \emph{VRInteraction BPI} è aggiunta alla pistola dalle \emph{Class Settings},
categoria \emph{Interfaces}, e la pistola implementa la funzione \emph{Trigger Pressed}:
``siccome la funzione non ha parametri di output, è implementata come un \emph{evento}''. Il
Pawn VR chiama Trigger Pressed sull'actor che tiene in mano, ``e se l'actor ha implementato
l'interfaccia VRInteraction BPI, Trigger Pressed viene eseguito. Non succede niente se l'actor
non l'ha implementata.'' È esattamente lo schema della lezione: il personaggio non sa
\emph{che cosa} sta guardando (un cubo, un bottone, un ascensore) e non deve fare un
\emph{cast} a nessuna classe; manda il messaggio, e ciascun oggetto risponde a modo suo.
Questo è il vantaggio delle interfacce rispetto al cast e all'ereditarietà viste con
BP\_Activable: classi che non hanno niente in comune (un actor, un personaggio, un oggetto
fisico) possono condividere la stessa ``porta''.

% ---------------------------------------------------------------------
\section{La pausa: la stamina di Tommaso}
\label{ue12:sec:pausa}

\paragraph{Il sintomo.} In pausa, come promesso, il problema rimasto aperto la volta
precedente: \textbf{Tommaso} tiene premuto Shift, la barra della stamina non scende, e
``appena mollo, si vede''. Il docente ripercorre il codice condiviso: l'evento StopRunning,
il Clear and Invalidate, il Set Timer by Event e la variabile TimeStamina (``quanto sta?
0,92''; nella trascrizione il valore è questo, ma va preso con cautela: nella lezione
precedente il valore indicato era 0,05), RegainStamina, il Clamp con CurrentStamina e
MaxStamina, l'InitializeStaminaBar nel Begin Play. ``Sembra tutto a posto.''

\paragraph{Il debug.} Il metodo è quello della volta scorsa: un \textbf{Print String} con la
CurrentStamina (``attaccalo al viola'') dopo l'evento di aggiornamento. I valori scendono
correttamente (``500, 500, 300'' dice la trascrizione, forse leggendo la velocità), poi
risalgono. Ma a schermo ci sono ``troppi Print String'': si cercano e si tolgono quelli
vecchi, uno per volta (quello della torretta, uno della piattaforma, ``forse i mulini?'',
``forse il collectable''; ``nessuno è perfetto''). Il comportamento: la barra
\emph{ricarica} visibilmente, ma \emph{non scende}.

\paragraph{La causa.} ``Ah, ecco perché: non chiami l'evento Call Update Stamina.'' Nella
catena di \textbf{StaminaDeduction}, dopo il Set di CurrentStamina, mancava la
\textbf{Call E\_UpdateStamina}: il valore calava, ma nessuno lo comunicava alla barra. La
Call c'era solo in RegainStamina, che parte al rilascio del tasto; per questo la barra si
aggiornava ``solo quando mollo''. Il rimedio: la Call anche in StaminaDeduction, con
CurrentStamina e MaxStamina collegati. ``Perfetto. Ottimo.''

\paragraph{Il rumore.} Il resto della pausa, nella trascrizione, è il blocco di testo
ripetuto e il frammento di video già segnalati nell'introduzione. Alla ripresa, uno studente
chiede aiuto per BPI\_InteractionRay, che dopo un riavvio di Unreal si apre ``male'': il
grafo in sola lettura non si vede. Il docente prova i menu \emph{Window} (Bookmarks, Find
Results, Compiler Results\dots), carica per sbaglio il layout di default dell'editor
(``mi sono caricato il bellissimo layout di Unreal''), e alla fine trova la via più
semplice: ``fai doppio clic sulla funzione'' nel pannello \emph{My Blueprint}, e il suo
grafo si apre. ``Tanto non ti serve'': nell'interfaccia non si scrive nulla.

% ---------------------------------------------------------------------
\section{La logica: la variabile Interactable}
\label{ue12:sec:logica}

``Vai, allora finiamo questo bellissimo sistema di interazioni.'' Si torna in
BPC\_RayInteract, nella funzione LookForInteractions. Il ramo di partenza è quello in cui
la trace ha colpito qualcosa e l'actor colpito implementa BPI\_InteractionRay.

\paragraph{La variabile d'appoggio.} Serve una variabile di tipo \textbf{Actor}
(\emph{Object Reference}), \textbf{Interactable}: ``sarà la nostra variabile d'appoggio'',
cioè l'oggetto interagibile che il giocatore sta guardando \emph{in questo momento}, o
niente. Il suo scopo è evitare di ripetere il lavoro a ogni frame: ``se l'actor che io ho
intercettato con il mio trace è lo stesso di prima, allora non voglio che mi rifai alcune
sezioni del codice''; ``se è lo stesso, non lo svuoto''.

\paragraph{È un oggetto nuovo?} Un secondo \textbf{Branch}, con la condizione \textbf{Hit
Actor} \textbf{diverso da} Interactable: il nodo \emph{Not Equal} (``diverso si fa con il
punto esclamativo e l'uguale''). ``Il punto esclamativo in informatica vuol dire che questo è
un condizionale che ci dice: guarda, questo è diverso da questo, e ci restituisce un
booleano. Se è diverso ci dà True, se sono uguali ci dà False.'' Sul False non si fa nulla:
è lo stesso oggetto di prima.

\paragraph{I messaggi dell'interfaccia.} ``Adesso, importante'': scrivendo
\emph{interaction} nel menu dei nodi, compaiono già le funzioni di BPI\_InteractionRay,
``senza fare niente'': ``ho già la possibilità di poter comunicare e messaggiare con
l'interfaccia''. Sono i nodi \textbf{(Message)}, con l'icona della busta, che si possono
chiamare su qualunque actor. Servono \textbf{Interaction End} e \textbf{Interaction Start}.

\paragraph{Il ramo True: chiudere il vecchio, aprire il nuovo.} Sul True del diverso:
\begin{enumerate}
  \item un \textbf{Validated Get} di Interactable (si mette il Get, tasto destro,
        \emph{Convert to Validated Get}): ``questa variabile Interactable è valida? C'è
        qualcosa dentro?'';
  \item \textbf{Is Valid}: c'era già un oggetto, quindi prima si chiude la vecchia
        interazione, \textbf{Interaction End} con target il \emph{vecchio} Interactable;
        poi \textbf{Set Interactable} = Hit Actor; poi \textbf{Interaction Start} con
        target il \emph{nuovo} Interactable;
  \item \textbf{Is Not Valid}: ``la variabile è vuota'', ``non è mai iniziata
        un'interazione precedente''; si salta l'End e si fa direttamente il Set
        Interactable = Hit Actor e l'\textbf{Interaction Start}.
\end{enumerate}

\paragraph{Il ramo False: non guardo niente.} Se la trace non colpisce nulla, \emph{oppure}
colpisce un actor senza l'interfaccia (``il muro, ad esempio''), il giocatore ha distolto lo
sguardo: ``allora ti finisco l'interazione''. Entrambi i False (quello del Branch sul Return
Value e quello di Does Object Implement Interface) portano a un altro Validated Get di
Interactable:
\begin{itemize}
  \item \textbf{Is Valid}: \textbf{Interaction End} sul vecchio Interactable, poi
        \textbf{Set Interactable} senza collegare nulla al pin, cioè lo si
        \textbf{svuota};
  \item \textbf{Is Not Valid}: non c'era nessuna interazione da chiudere; per sicurezza
        lo si svuota comunque (``semplicemente dirgli: guarda, svuotati'').
\end{itemize}
``Quindi mettere il Set così è come mettere uguale a null?'' ``È nulla, proprio. La
annulliamo: la mettiamo vuota, e quindi diventerà non valida.'' Il caso del primo ramo,
il passaggio diretto da un oggetto all'altro, serve quando due interagibili sono vicini:
``lui finisce con quello precedente e inizia con il nuovo''; di solito, invece, gli oggetti
sono distanti, la variabile si svuota e poi si riempie di nuovo
(figure~\ref{ue12:fig:look} e~\ref{ue12:fig:stati}).

\paragraph{Il problema di tipo.} Uno studente non riesce a collegare l'Hit Actor al Set
Interactable: ``mi dice che non sono compatibili, perché uno è un riferimento oggetto e
l'altro è un componente di scena''. La variabile Interactable era stata creata con il tipo
sbagliato (ed era cambiata per un clic involontario in \emph{Geometry}, ``che non ho
cliccato''). Rimedio: nei Details della variabile, il tipo \textbf{Actor}, \textbf{Object
Reference} (``non quell'Actor: tipo oggetto, riferimento oggetto''), confermare la
modifica del tipo, poi staccare e riattaccare i fili rossi dei nodi che usavano la
variabile, ``era semplicemente sparito''.

\paragraph{Il Tick.} ``Importante, ragazzi: far sì che questa funzione venga fatta su ogni
frame.'' Nell'Event Graph del componente, \textbf{Event Tick} $\rightarrow$
LookForInteractions (se non lo si è già fatto): ``a ogni frame farà il controllo: c'ho
qualcosa di interagibile davanti a me, dove sto puntando la camera? Sì o no?''.

\begin{figure}[htbp]
\centering
\begin{tikzpicture}[font=\footnotesize,
  ue12node/.style={draw, thick, rounded corners=4pt, align=center, fill=gray!8,
                  minimum height=0.85cm, text width=2.4cm},
  ue12dec/.style={draw, thick, diamond, aspect=2.2, align=center, fill=yellow!15,
                  inner sep=1pt, text width=2.0cm, font=\scriptsize},
  ue12msgn/.style={ue12node, fill=violet!12},
  ue12exec/.style={-{Stealth[length=2.5mm]}, very thick, black!80}]
  \node[ue12node, fill=red!12, text width=3.4cm] (st) at (0,0) {\textbf{Sphere Trace By
    Channel}\\{\scriptsize Start: posizione del componente; End: Start + Forward (camera o
    self) $\times$ 300; Radius 5; InteractionTrace}};
  \node[ue12dec] (hit) at (0,-2.2) {colpito qualcosa?};
  \node[ue12dec] (impl) at (0,-4.4) {Hit Actor implementa la BPI?};
  \node[ue12dec] (diff) at (0,-6.6) {Hit Actor $\neq$ Interactable?};
  \node[ue12dec] (val1) at (0,-8.9) {Interactable valido?};
  \node[ue12msgn] (end1) at (-3.6,-10.8) {\textbf{Interaction End}\\{\scriptsize sul vecchio}};
  \node[ue12node] (set1) at (0,-12.4) {\textbf{Set Interactable}\\{\scriptsize = Hit Actor}};
  \node[ue12msgn] (start) at (0,-14.0) {\textbf{Interaction Start}\\{\scriptsize sul nuovo}};
  \node[ue12node, fill=gray!3, text width=2.0cm] (none) at (3.7,-6.6) {niente\\{\scriptsize
    (stesso oggetto)}};
  \node[ue12dec] (val2) at (6.4,-3.3) {Interactable valido?};
  \node[ue12msgn] (end2) at (6.4,-5.6) {\textbf{Interaction End}\\{\scriptsize sul vecchio}};
  \node[ue12node] (clr) at (6.4,-7.6) {\textbf{Set Interactable}\\{\scriptsize vuoto (null)}};
  \draw[ue12exec] (st) -- (hit);
  \draw[ue12exec] (hit) -- node[left, font=\scriptsize]{True} (impl);
  \draw[ue12exec] (impl) -- node[left, font=\scriptsize]{True} (diff);
  \draw[ue12exec] (diff) -- node[left, font=\scriptsize]{True} (val1);
  \draw[ue12exec] (diff) -- node[above, font=\scriptsize]{False} (none);
  \draw[ue12exec] (val1.west) -| node[pos=0.25, above, font=\scriptsize]{Is Valid} (end1.north);
  \draw[ue12exec] (end1.south) |- (set1.west);
  \draw[ue12exec] (val1) -- node[right, font=\scriptsize]{Is Not Valid} (set1);
  \draw[ue12exec] (set1) -- (start);
  \draw[ue12exec] (hit.east) -| node[pos=0.2, above, font=\scriptsize]{False} (val2.north);
  \draw[ue12exec] (impl.east) -- node[above, font=\scriptsize]{False} ++(2.5,0) |-
    (val2.west);
  \draw[ue12exec] (val2) -- node[right, font=\scriptsize]{Is Valid} (end2);
  \draw[ue12exec] (end2) -- (clr);
  \draw[ue12exec] (val2.east) -- ++(0.6,0) |- node[pos=0.3, right, font=\scriptsize, text
    width=1.0cm, align=left]{Is Not Valid} (clr.east);
\end{tikzpicture}
\caption{La funzione LookForInteractions del componente, chiamata nell'Event Tick. Se il
giocatore guarda un nuovo interagibile, si chiude l'interazione con il vecchio (se c'era)
e si apre con il nuovo; se guarda un oggetto non interagibile, o niente, si chiude
l'interazione in corso e si svuota la variabile. I nodi viola sono messaggi
dell'interfaccia. Ricostruzione dalla lezione.}
\label{ue12:fig:look}
\end{figure}

\begin{figure}[htbp]
\centering
\begin{tikzpicture}[font=\footnotesize,
  ue12state/.style={draw, very thick, rounded corners=10pt, align=center,
                    minimum height=1.3cm, text width=3.4cm},
  ue12tr/.style={-{Stealth[length=2.5mm]}, very thick}]
  \node[ue12state, fill=gray!10] (vuoto) at (0,0) {\textbf{Interactable vuoto}\\{\scriptsize
    non guardo niente di interagibile}};
  \node[ue12state, fill=green!12] (a) at (7.0,0) {\textbf{Interactable = A}\\{\scriptsize
    widget ``E'' di A visibile}};
  \node[ue12state, fill=green!6] (b) at (7.0,-3.4) {\textbf{Interactable = B}\\{\scriptsize
    widget ``E'' di B visibile}};
  \draw[ue12tr] ([yshift=3mm]vuoto.east) -- node[above, font=\scriptsize]{guardo A: Start(A)}
    ([yshift=3mm]a.west);
  \draw[ue12tr] ([yshift=-3mm]a.west) -- node[below, font=\scriptsize, text width=3.2cm,
    align=center]{guardo il muro o niente: End(A), svuota} ([yshift=-3mm]vuoto.east);
  \draw[ue12tr] (a.south) -- node[right, font=\scriptsize, text width=3.4cm]{passo
    direttamente a B: End(A), Set, Start(B)} (b.north);
  \draw[ue12tr] (b.west) -| node[pos=0.3, below, font=\scriptsize]{guardo il muro: End(B),
    svuota} (vuoto.south);
  \draw[ue12tr] (a.east) .. controls +(1.6,0.9) and +(1.6,-0.9) .. node[right,
    font=\scriptsize, text width=1.9cm]{guardo ancora A: nulla} (a.east);
\end{tikzpicture}
\caption{La variabile Interactable come stato. Ogni Interaction Start è sempre seguito, prima
o poi, dal suo Interaction End: l'oggetto può contare sulla coppia per mostrare e nascondere
il suggerimento. Schema ricostruito.}
\label{ue12:fig:stati}
\end{figure}

\paragraph{Approfondimento: messaggi, Does Implement Interface e Validated Get.} Tre
precisazioni generali sui nodi usati. Primo: un nodo \emph{(Message)} di un'interfaccia si
può chiamare su qualunque riferimento a oggetto; se l'oggetto non implementa l'interfaccia,
la chiamata semplicemente non fa nulla, senza errori (la frase del libro citata sopra: ``non
succede niente''). Il controllo \textbf{Does Object Implement Interface}, quindi, non serve
a evitare un errore, ma a \emph{distinguere} gli interagibili dagli altri oggetti, così da
chiudere l'interazione quando si guarda un muro. Secondo: il confronto \emph{Not Equal} fra
due riferimenti confronta l'\emph{identità} degli oggetti (è lo stesso actor?), non la loro
classe; due cubi identici sono comunque diversi. Terzo: il \textbf{Validated Get} è un Get
con l'\emph{Is Valid} incorporato, con le due uscite \emph{Is Valid} e \emph{Is Not Valid};
``validare'' un riferimento prima di usarlo è la raccomandazione con cui il docente chiude
il corso, perché un riferimento vuoto o a un actor distrutto (per esempio un interagibile
raccolto o eliminato) produce, se usato, gli errori \emph{Accessed None} a fine Play.
Funzionamento generale dell'engine.

% ---------------------------------------------------------------------
\section{L'oggetto di prova: BP\_InteractableTest}
\label{ue12:sec:test}

\paragraph{L'actor.} ``Adesso creiamo il nostro actor interagibile.'' Nella cartella dei
Blueprint, una Blueprint Class di tipo \textbf{Actor}, ``veramente semplicissimo'':
\textbf{BP\_InteractableTest}. Due componenti: un \textbf{Cube} e un \textbf{Widget}
(``fate widget, che poi faremo''), messo sopra il cubo.

\paragraph{Implementare l'interfaccia.} ``Non si trova'', se si cerca l'interfaccia fra i
componenti: va aggiunta in \textbf{Class Settings}, nella toolbar. Nei Details compaiono le
categorie \emph{Imports} (\emph{Default Namespace}, \emph{Imported Namespaces}) e
\textbf{Interfaces}, con \emph{Inherited Interfaces} e \textbf{Implemented Interfaces}; in
quest'ultima, \textbf{Add}, BPI\_InteractionRay. ``Per vedere che l'hai aggiunto'': nel
pannello My Blueprint compare la sezione \textbf{Interfaces}, con Interaction, Interaction End
e Interaction Start.

\paragraph{Implementare le funzioni.} Trascinare Interaction Start nel grafo non funziona:
l'editor avverte che la funzione ``was not marked as Blueprint Callable'' e non può essere
piazzata così. ``Dobbiamo overrideare la loro funzione, che di base è nulla'': le funzioni
dell'interfaccia ``sono solo degli appoggi'', vuote, e il loro contenuto ``sarà da
overrideare proprio dal livello di design''. Con un doppio clic (o tasto destro,
\emph{Implement Event}) si ottengono tre eventi rossi: \textbf{Event Interaction Start},
\textbf{Event Interaction End} ed \textbf{Event Interaction}. Ciascuno con un
\textbf{Print String}:
\begin{itemize}
  \item Interaction Start: ``inizio interazione'', \textbf{Text Color} verde;
  \item Interaction End: ``fine interazione'', rosso;
  \item Interaction: ``interagendo'', giallo-arancione (``ancora non l'abbiamo usato,
        quindi non lo vedremo, è normale'').
\end{itemize}
Il colore si trova aprendo la tendina del Print String (la freccia in basso del nodo).

\paragraph{La collisione.} ``È importante però una cosa'': il canale InteractionTrace è in
Ignore di default, quindi il cubo, così com'è, non viene visto dal raggio. Nella
\textbf{Collision} del Cube, preset \emph{Custom} e \textbf{InteractionTrace} su
\textbf{Block}. ``Su Cube, tanto la mesh è quella.''

\paragraph{La prova.} Si mette il cubo in scena. In Play: ``inizio interazione, fine
interazione, inizio interazione, fine interazione''. Con la camera in terza persona colpire
il cubo è un po' difficile, ``ma ci sono delle volte in cui diventa verde''. ``Vi funziona?''
``Sì, funziona.'' ``Quindi adesso abbiamo creato il nostro trace per quanto riguarda le
interazioni.'' Il componente ``fa le veci di \emph{vado a controllarti se ci sono delle
interazioni, e gestisco tutta la cosa}''; l'interazione vera, il tasto, la fa il personaggio,
``in modo classico''.

\paragraph{Approfondimento: eventi o funzioni dell'interfaccia.} Le tre funzioni di
BPI\_InteractionRay non hanno \emph{output}: per questo, nell'actor che le implementa,
diventano \emph{eventi} nell'Event Graph (è la stessa regola citata dal libro per la Trigger
Pressed della pistola VR). Se una funzione dell'interfaccia avesse un valore di ritorno (per
esempio un \emph{GetInteractionText} che restituisce ``Apri'' o ``Accendi''), andrebbe
implementata come \emph{funzione}, con un doppio clic nella sezione Interfaces, e potrebbe
essere chiamata per leggere un dato. Un'interfaccia non ha stato né variabili: ogni classe
che la implementa decide da sé che cosa fare, e può anche lasciare vuote le funzioni che non
le servono. Funzionamento generale dell'engine.

% ---------------------------------------------------------------------
\section{L'input: IA\_Interact e il tasto E}
\label{ue12:sec:input}

Nel BP\_Bot2 serve un nuovo input. ``Come si crea un nuovo input? Con l'IMC e con l'Input
Action.''
\begin{enumerate}
  \item Nel Content Browser si va alle Input Action del progetto (cliccando su IA\_Jump,
        \emph{Show in Folder View});
  \item ``basta che fate una copia di IA\_Run'': \textbf{Duplicate}, nome
        \textbf{IA\_Interact}; ``basta che qua dentro abbia i \textbf{Triggers}
        \emph{Pressed} e \emph{Released}'' (come quelli di IA\_Run), da lasciare così;
  \item nell'\textbf{IMC} di default (\emph{IMC\_\allowbreak Default}), il più accanto alle
        \emph{Mappings} aggiunge una riga per IA\_Interact; poi, cliccando sull'icona
        della tastiera e premendo il tasto, si assegna \textbf{E}.
\end{enumerate}
Il bot2 ``eredita la stessa cosa'': usa lo stesso IMC\_Default del template. Le domande dei
presenti sono tutte sulla tendina dell'IMC (``non mi fa mettere l'azione di input'': ``è la
tendina'') e sui trigger (``il primo è premuto e l'altro qual è?'': ``premuto e rilascio va
benissimo così, lasciato così''). Nel personaggio, cercando \emph{IA\_Interact}, compare il
nuovo evento d'input.

\paragraph{Approfondimento: Pressed e Released nello stesso trigger.} Un'Input Action
digitale senza trigger scatta in \emph{Triggered} a ogni frame finché il tasto è premuto; con
il trigger \emph{Pressed} scatta una sola volta alla pressione, con \emph{Released} una
volta al rilascio. Con \emph{entrambi}, come in IA\_Run e quindi in IA\_Interact, l'uscita
\emph{Triggered} scatta due volte per ogni pressione: quando si preme e quando si rilascia.
Nella lezione questo non dà problemi, perché il grafo del personaggio passa da un \emph{Do
Once} che si resetta solo alla fine del Montage: il secondo impulso, al rilascio, arriva
mentre l'animazione è in corso e viene ignorato. Ma se il Montage fosse molto breve, o il
giocatore tenesse premuto a lungo, il rilascio potrebbe far partire una seconda
interazione: per un'interazione ``a pressione'' basta il solo trigger Pressed, o l'uscita
\emph{Started}. La trascrizione non dice quale uscita dell'evento IA\_Interact usi il docente.
Funzionamento generale dell'Enhanced Input.

% ---------------------------------------------------------------------
\section{L'interazione nel personaggio}
\label{ue12:sec:player}

``Allora adesso partiamo a fare quello che ci servirà per indirizzare il nostro player,
cioè farlo ruotare verso il nostro oggetto, e allo stesso modo far sì che l'interazione
parta.''

\paragraph{Un Do Once.} ``Vi dico che è un'azione che deve essere fatta su un \textbf{Do
Once}'': l'interazione è una sequenza (girarsi, far partire il Montage, aspettarne la fine)
e solo ``a fine del montaggio può reiniziare una prossima interazione''. Quindi IA\_Interact
$\rightarrow$ Do Once.

\paragraph{Chi stiamo guardando.} Dal componente \textbf{BPC\_RayInteract} (trascinato dal
pannello dei componenti nel grafo: ``prendilo e trascinalo'') si prende la sua variabile
\textbf{Interactable} e la si converte in \textbf{Validated Get}: ``questo Interactable è
presente? Ci sta o non ci sta?''. Il personaggio non rifà la trace: legge il risultato del
componente.
\begin{itemize}
  \item \textbf{Is Valid}: il messaggio \textbf{Interaction} di BPI\_InteractionRay
        (``tira fuori e fai \emph{interaction}, cioè Interaction Message''), con target
        Interactable: ``e questa è proprio la base: così interagiamo''.
  \item \textbf{Is Not Valid}: \textbf{Reset} del Do Once (aggiunto dopo, quando il docente
        si accorge che senza il reset, premendo E senza guardare nulla, il Do Once resta
        chiuso per sempre: ``se non è valido, resetta il Do Once comunque, in ogni caso'').
\end{itemize}

\paragraph{Girarsi verso l'oggetto.} Dopo il messaggio, \textbf{Set Actor Rotation} con
target \emph{Self} (``siamo già nel player''). Si splitta il pin \emph{New Rotation}, e si
splitta anche l'uscita di un \textbf{Get Actor Rotation}: ``l'unica asse che ci interessa è
la Z'', perché il personaggio deve girarsi solo attorno all'asse verticale (``se
l'interactable sta qua, deve girare così''). Quindi:
\begin{itemize}
  \item X (Roll) e Y (Pitch) restano quelle attuali, dal Get Actor Rotation;
  \item Z (Yaw) viene da un \textbf{Find Look at Rotation}, ``quello che abbiamo fatto
        l'altra volta per la torretta'', con \textbf{Start} = \textbf{Get Actor Location}
        del personaggio (il docente dice prima ``rotation'', poi si corregge: ``location,
        ragazzi, scusatemi'') e \textbf{Target} = Get Actor Location dell'Interactable. Il
        pin di uscita non si lascia splittare (``non me lo splitta''), quindi un
        \textbf{Break Rotator}, e il suo Z va nella Z della nuova rotazione.
\end{itemize}

\paragraph{La prova provvisoria.} Prima del Montage, una prova: la scritta ``interagendo''
compare, ma il personaggio non si gira; e dopo la prima volta non succede più nulla. ``Ah,
giustamente, ragazzi'': il Do Once non si resetta mai. Il docente collega temporaneamente
il Reset alla fine della catena (``lo faccio solamente per farvelo vedere''): ``mi giro di
qua, e si gira verso di lui''. Poi lo stacca, ``perché ci servirà più avanti''.

\paragraph{Il Montage.} Nel Content Browser si cerca \emph{interact}: fra le animazioni del
bot c'è \textbf{Interact Press Button}. ``Vi ricordate come si fa il montaggio?'' Tasto
destro sull'animazione, \textbf{Create}, \textbf{Create AnimMontage}: si crea un asset che
``è sempre la stessa animazione, ma in questo caso è un montaggio''. Nel personaggio,
\textbf{Play Montage}: nel pin \textbf{In Skeletal Mesh Component} la \textbf{Mesh} del
personaggio (``basta che prendi la Mesh da qui e la metti così''), nel \textbf{Montage to
Play} il nuovo Montage (\emph{Interact Press Button Montage}).

\paragraph{La variabile Interacting.} ``Mettiamo, giusto perché in futuro ci possa
servire, una variabile'' booleana, \textbf{Interacting}: ``dobbiamo verificare se stiamo
interagendo, non stiamo interagendo, eccetera. Sempre buono mettere una variabile che dice:
guarda, io sto interagendo''. Prima del Play Montage, Set Interacting a \textbf{vero}; sull'
\textbf{On Completed}, ``quindi è finita l'animazione'', a \textbf{falso}.

\paragraph{Fermare il giocatore.} Durante l'interazione il giocatore non deve muoversi:
\begin{itemize}
  \item prima del Montage, \textbf{Disable Input}, con il \textbf{Get Player Controller}
        nel pin \emph{Player Controller}: ``chi dobbiamo levare? Gli input al controller del
        nostro player'';
  \item lo \textbf{Stop Movement Immediately}, ``andiamo a fermare qualsiasi tipologia
        di movimento''. Il target non viene detto: il nodo appartiene al componente di
        movimento del personaggio, il \emph{Character Movement};
  \item sull'On Completed, \textbf{Enable Input}, con lo stesso Player Controller;
  \item e infine il \textbf{Reset} del Do Once, ``in modo tale che io possa reinteragire di
        nuovo''.
\end{itemize}
In Play: ``io giro di qua, si gira, fa l'animazione, non mi posso muovere''; ``e adesso mi
riposso muovere appena finito il Montage'' (figura~\ref{ue12:fig:player}). ``Ci siete? Vi
funziona? Perfetto.''

\paragraph{I problemi dei presenti.} Uno studente non trova come prendere il componente
(``l'hai aggiunto come components? Prendilo e trascinalo''); altri si perdono i due nodi
successivi, ripetuti dal docente (Get Interactable, Convert to Validated Get, il messaggio
Interaction). Un altro ha il Montage che non parte: manca la Mesh nel pin del Play Montage.
Uno studente, a cui prima non funzionava, conferma che ora va, ma ha lo
schermo ``pieno di \emph{Hello}'' (forse i Print String rimasti con il testo di default),
``quindi li vedo con la coda dell'occhio''.

\begin{figure}[htbp]
\centering
\begin{tikzpicture}[font=\footnotesize,
  ue12node/.style={draw, thick, rounded corners=4pt, align=center, fill=gray!8,
                  minimum height=0.9cm, text width=2.2cm},
  ue12ev/.style={ue12node, fill=red!12},
  ue12msgn/.style={ue12node, fill=violet!12},
  ue12pure/.style={ue12node, fill=green!10, rounded corners=8pt},
  ue12exec/.style={-{Stealth[length=2.5mm]}, very thick, black!80},
  ue12data/.style={-{Stealth[length=2mm]}, thick, blue!60!black}]
  \node[ue12ev, text width=1.9cm] (ia) at (0,-0.0) {\textbf{IA\_Interact}\\{\scriptsize tasto
    E}};
  \node[ue12node, text width=1.5cm] (do) at (2.3,0) {\textbf{Do Once}};
  \node[ue12node, text width=2.4cm] (vg) at (4.9,0) {\textbf{Validated Get}\\{\scriptsize
    BPC\_RayInteract.\\Interactable}};
  \node[ue12msgn] (msg) at (8.6,0) {\textbf{Interaction}\\{\scriptsize (Message) su
    Interactable}};
  \node[ue12node, text width=2.5cm] (rot) at (11.9,0) {\textbf{Set Actor Rotation}\\{\scriptsize
    X, Y attuali; Z dal Look at}};
  \node[ue12pure, text width=3.0cm] (look) at (11.9,-1.85) {\textbf{Find Look at
    Rotation}\\{\scriptsize mia posizione $\rightarrow$ posizione di Interactable}\\
    $\rightarrow$ \textbf{Break Rotator}: Z};
  \node[ue12node, text width=3.5cm] (pre) at (8.0,-4.4) {\textbf{Set Interacting} = vero\\
    \textbf{Disable Input}\\\textbf{Stop Movement Immediately}};
  \node[ue12node, text width=2.8cm] (pm) at (12.0,-4.4) {\textbf{Play Montage}\\{\scriptsize
    Mesh; Interact Press Button Montage}};
  \node[ue12node, text width=3.5cm] (post) at (12.0,-6.8) {\textbf{Set Interacting} =
    falso\\\textbf{Enable Input}\\\textbf{Reset} del Do Once};
  \node[ue12pure, text width=1.9cm] (gpc) at (8.0,-6.8) {Get Player Controller};
  \draw[ue12exec] (ia) -- (do); \draw[ue12exec] (do) -- (vg);
  \draw[ue12exec] (vg) -- node[below, font=\scriptsize]{Is Valid} (msg);
  \draw[ue12exec] (msg) -- (rot);
  \draw[ue12exec] (rot.east) -- ++(0.35,0) -- ++(0,-3.15) -| (pre.north);
  \draw[ue12exec] (pre) -- (pm);
  \draw[ue12exec] (pm) -- node[right, font=\scriptsize]{On Completed} (post);
  \draw[ue12exec] (vg.south) -- ++(0,-1.0) -| node[pos=0.25, above, font=\scriptsize]{Is
    Not Valid: Reset} ([xshift=-3mm]do.south);
  \draw[ue12data] (look) -- (rot);
  \draw[ue12data] (gpc) -- (pre);
  \draw[ue12data] (gpc) -- (post);
\end{tikzpicture}
\caption{L'interazione nel personaggio. Il personaggio non rifà la trace: legge la variabile
Interactable del suo componente, invia il messaggio, si gira solo sull'asse Z verso
l'oggetto, blocca l'input e suona il Montage; a Montage finito restituisce l'input e riapre
il Do Once. Ricostruzione dalla lezione; la posizione esatta del Set Interacting e del
Disable Input nella catena è dedotta dalle parole del docente.}
\label{ue12:fig:player}
\end{figure}

\paragraph{Approfondimento: Do Once, Disable Input e la rotazione.} Romero e Sewell
descrivono il \textbf{Do Once} nel capitolo 13 (\emph{Data Structures and Flow Control}):
``il nodo Do Once esegue il pin di uscita una sola volta. Se viene attivato di nuovo, il pin
di uscita non viene eseguito. Il pin di input Reset deve essere attivato per permettere al
Do Once di eseguire di nuovo il pin di uscita.'' Il loro esempio è un'arma caricata, che si
ricarica dopo 30 secondi con un timer che resetta il Do Once: lo stesso schema
della lezione, con il Montage al posto del timer. Nel capitolo 15 il libro descrive
\textbf{Enable Input} e \textbf{Disable Input}: ``definiscono se un actor deve rispondere agli
eventi di input'', e ``hanno bisogno di un riferimento al Player Controller in uso''. Due
note generali sull'engine. La prima: \emph{Disable Input} chiamato sul personaggio blocca
gli eventi d'input di quel Blueprint (movimento, camera, salto), ma non ferma un movimento
già in corso, cioè la velocità residua del Character Movement; per questo serve anche
\emph{Stop Movement Immediately}, che azzera la velocità sul posto. La seconda: il
\emph{Set Actor Rotation} è istantaneo, quindi il personaggio ``scatta'' verso l'oggetto. Per
una rotazione morbida si può usare un \emph{RInterp To} nel Tick o una Timeline, come per le
porte; e il template Third Person, che ha l'\emph{Orient Rotation to Movement} attivo, non
contrasta la rotazione perché durante l'interazione il personaggio è fermo.

% ---------------------------------------------------------------------
\section{Il widget ``premi E''}
\label{ue12:sec:widget}

``Andiamo semplicemente a rendere tutto un po' più carino'': un piccolo widget che mostri il
tasto da premere sopra l'oggetto.

\paragraph{Gli asset della UI.} Il docente aveva messo nel materiale didattico una cartella
di immagini per la UI, ``la UI asset''. Uno studente non l'ha trovata (``pensavo
fosse sul Drive''): il docente controlla il Drive condiviso e i canali del Discord, scopre che
non era stata caricata e la carica al momento, ``tanto pesa pochi kilobyte''. ``Estraetela e
schiaffatela dentro Unreal'', nella cartella \textbf{UI}. Il trascinamento nel Content
Browser a qualcuno dà ``fallimento''; il ``metodo abbastanza brusco, però va benissimo'' è
copiare la cartella \emph{estratta} con il gestore dei file direttamente nella cartella
\emph{Content/UI} del progetto. ``Adesso funziona.''

\paragraph{WB\_InteractUI.} Nella cartella UI, un nuovo Widget Blueprint, classe User
Widget: \textbf{WB\_InteractUI}. Nel Designer:
\begin{enumerate}
  \item un \textbf{Border}, con l'alfa del colore a \textbf{0}, ridimensionato a un
        quadrato di circa \textbf{96 $\times$ 96} (``94 va bene anche''), e il
        \textbf{Padding} tutto a \textbf{zero};
  \item un secondo \textbf{Border} \emph{dentro} il primo;
  \item le immagini: nel \textbf{Brush} di ciascun Border (``vedi che c'è il brush, apri il
        brush'') si sceglie un'immagine del pacchetto, ``il panel che volete voi'' (ce ne
        sono diversi, Panel 0, 1, 2, 3, 5\dots). Dopo una correzione (``questo è il
        contrario, scusate''), il Border esterno prende \textbf{Panel Border 5}, con l'alfa a
        \textbf{1} e la tinta \textbf{nera}, e quello interno \textbf{Panel 5};
  \item un \textbf{Text} con la lettera \textbf{E}, con \emph{Horizontal Alignment} e
        \emph{Vertical Alignment} al \textbf{centro}, e la dimensione del \textbf{font} a
        35, poi 40--45 (``una bella E'').
\end{enumerate}
``Piccola cosa, ragazzi, prima non l'ho fatto fare a mano'': la dimensione dei Border si può
impostare anche con precisione, dall'\textbf{Image Size} del Brush (96 $\times$ 96).

\paragraph{Il widget che contiene la barra della vita.} Una studentessa crea un nuovo
widget, ma aprendolo trova le barre della vita e della stamina, e non riesce a modificarle.
Rifacendo la creazione insieme, la causa emerge: nella finestra di scelta della classe aveva
preso un widget esistente come genitore (``anteprima widget praticamente ti dà
un'anteprima dei widget che già ci sono''), cioè il suo HUD; il nuovo widget ne
\emph{ereditava} il contenuto, in sola lettura. ``È buona prassi fare quello tutto in
alto'', \textbf{User Widget}.

\paragraph{L'immagine che non si vede.} Un'altra studentessa ha il bordo esterno che non
compare. Il docente controlla tinta, alfa, Image Size, Padding, \emph{Content Opacity}: ``è
a 1, va bene così''. ``Perdonami, purtroppo non saprei dire'': si cancellano i Border e si
rifanno, ``basta funzionare''. Rifatti, funzionano.

\paragraph{Il widget nell'oggetto.} In BP\_InteractableTest, il componente \textbf{Widget}:
\begin{itemize}
  \item \textbf{Space}: \textbf{Screen}, non World, ``in modo tale che ci sia una specie di
        2D, quindi relativo allo schermo, non al mondo'';
  \item \textbf{Draw Size}: \textbf{100 $\times$ 100}, ``un po' più piccolo'';
  \item \textbf{Widget Class}: WB\_InteractUI.
\end{itemize}
In Play, il quadratino con la E ``scala con lo schermo'', ma, se non lo si nasconde,
``rimane per sempre''.

\paragraph{Mostrare e nascondere.} Nel grafo di BP\_InteractableTest:
\begin{itemize}
  \item \textbf{Event BeginPlay} $\rightarrow$ \textbf{Set Hidden in Game} sul Widget, con
        \emph{New Hidden} \textbf{attivo}: nascosto da subito;
  \item \textbf{Event Interaction Start} $\rightarrow$ Set Hidden in Game,
        \textbf{disattivo}: ``allora la faccio apparire'';
  \item \textbf{Event Interaction End} $\rightarrow$ Set Hidden in Game, \textbf{attivo}.
\end{itemize}
Il docente prova anche un'idea per l'Event Interaction (nascondere il widget, un
\textbf{Delay} di un secondo, mostrarlo di nuovo), ma la toglie: ``questo potrebbe dare
errori, ragazzi'', perché il Delay ``va in contrasto'' con Start ed End (se lo sguardo si
sposta durante il secondo d'attesa, il Delay rimostrerebbe un widget che l'End ha appena
nascosto). ``Lascialo così per adesso.''

\paragraph{La spunta al contrario.} Una studentessa vede la E anche quando non guarda il
cubo: ``allora vuol dire che hai sbagliato l'Hidden in Game''. ``Il tick è positivo sull'End
e negativo sullo Start'', ``e anche positivo sul Begin Play''. Dopo la correzione resta un
dubbio, poi chiarito: all'inizio della prova la camera stava già inquadrando il cubo da una
posizione strana, e il widget era giustamente visibile. ``Non lo vedo più. Perfetto,
ragazzi.''

\paragraph{Approfondimento: il Widget Component.} Il \textbf{Widget Component} mostra un
Widget Blueprint attaccato a un actor. In \textbf{World} space il widget è un piano nel
mondo 3D: ha prospettiva, può essere coperto da altri oggetti, si rimpicciolisce con la
distanza, ed è adatto a schermi diegetici (un terminale, un cartello). In \textbf{Screen}
space il widget è disegnato sopra l'immagine, sempre rivolto alla camera, con una dimensione
costante in pixel, e segue sullo schermo la posizione proiettata del componente: è la scelta
tipica per i suggerimenti ``premi E'' e per le barre della vita sopra i nemici. La
\textbf{Draw Size} è la risoluzione in pixel a cui il widget viene disegnato. Per nascondere
il suggerimento, l'alternativa a \emph{Set Hidden in Game} è \emph{Set Visibility} sul
componente; il risultato è lo stesso. Funzionamento generale dell'engine.

% ---------------------------------------------------------------------
\section{Un esempio: la luce che si accende}
\label{ue12:sec:luce}

\paragraph{L'idea.} ``Sappiate che tutti gli oggetti che create'', se implementano
l'interfaccia, ``avete già un qualcosa che vi identifica se state iniziando
un'interazione'', se la state finendo, ``e la tipologia di interazione''. Le funzioni
costruite per i bottoni nelle lezioni passate ``possiamo traslarle sul nostro
interactable'': ``un diverso modo di interagire e di attivare le cose''. Con Davide si
vedranno ``anche delle piccole cose di modeling'': si può modellare un pulsante, dargli
BPI\_InteractionRay, ``e a quel punto magari potete chiamare l'ascensore''. ``Cosa succede se
interagisco? Ci metto quello che voglio: ci metto un timer, ci metto una cosa, accendo una
luce.''

\paragraph{Velocissimo.} ``Facciamo velocissimo'', in BP\_InteractableTest. Dall'Event
Interaction si toglie il Print String, e:
\begin{enumerate}
  \item una variabile del tipo dell'attivabile luce, la classe figlia di BP\_Activable delle
        lezioni sul pulsante (``la Point Light'', \textbf{BP\_ActivableLight}; la
        trascrizione dice ``BP Activables'' e ``inter light''), con i suoi eventi
        \textbf{Activation} e \textbf{Deactivation};
  \item un \textbf{Flip Flop}: ``questo fa A, poi fa B, poi fa di nuovo A: fondamentalmente
        gioca su due'';
  \item su \textbf{A}, \textbf{Activation} della luce; su \textbf{B},
        \textbf{Deactivation};
  \item la variabile va resa modificabile nell'istanza e collegata alla luce del livello
        (``però gli dobbiamo dare quale light'', ``dammi la light, dammi questa'').
\end{enumerate}
``Abbiamo acceso la luce e la rispegniamo.'' ``Questo è un esempio, ragazzi'': ``qualsiasi
actor può essere separato dagli altri con la sua specifica interazione'', e ``quello che la
identifica già ce l'abbiamo''.

\paragraph{Approfondimento: Flip Flop nel libro.} Romero e Sewell (capitolo 13): ``ogni volta
che un nodo Flip Flop viene eseguito, alterna fra due pin di uscita, A e B''; ha anche un
output booleano \emph{Is A}, vero quando sta eseguendo A. Con una nota importante: ``il Flip
Flop non funziona correttamente dentro una funzione, perché ha una variabile interna che
viene cancellata quando la funzione finisce'', e quindi dentro una funzione esegue sempre A.
Lo stesso vale per il Do Once. È il motivo per cui, nella lezione, il Flip Flop sta
nell'Event Graph dell'oggetto e il Do Once nell'Event Graph del personaggio, mentre la
funzione LookForInteractions, che non deve ricordare nulla fra un frame e l'altro (lo stato
è nella variabile Interactable, che è del componente), può essere una funzione.

% ---------------------------------------------------------------------
\section{La chiusura del corso}
\label{ue12:sec:chiusura}

``Ci sono domande, ragazzi? Domande su tutto il corso, in generale.'' Il docente chiede
anche un \textbf{feedback} sul corso, ``adesso o dopo, tranquillamente sul server''. La
risposta di uno studente, scherzosa: ``tutto chiaro, tutto bellissimo, incredibile''. ``Sono
contento.''

Per il futuro, anche ``per chi guarda in differita'': ``se avete un problema, state facendo il
Big One, eccetera, buttate il vostro problema sul server: sia alunni sia professori
provvederanno a rispondervi e darvi una mano''. ``Se ci sono ultime cose che volete dire,
ragazzi, ditele ora, o effettivamente questa volta tacete per sempre.'' E il congedo: ``vi
ringrazio per la vostra attenzione, vi ringrazio per aver partecipato al corso. Spero che i
vostri progettini possano prendere vita per tutte le cose che poi andrete a fare in Unreal.
Per benino, ricordatevi di \textbf{validare sempre} molte variabili.'' ``Da questo punto in
poi'' le lezioni continuano con Davide. ``È stato un piacere.''

\begin{table}[htbp]
\centering
\small
\begin{tabular}{@{}p{3.7cm}p{3.0cm}p{2.9cm}p{4.4cm}@{}}
\toprule
\textbf{Asset o variabile} & \textbf{Dove} & \textbf{Tipo, valore} & \textbf{A che serve} \\
\midrule
M\_Menu & mappe & Empty Level & la mappa del menu \\
WB\_Menu & UI & Widget Blueprint & il menu con i due bottoni \\
Button\_StartGame, Button\_QuitGame & WB\_Menu & Button, 300 $\times$ 80 & gioca; esci \\
CameraActor & M\_Menu & Camera Actor & la vista fissa del menu \\
GameMode Override & World Settings di M\_Menu & Game Mode base & niente personaggio nel menu \\
ThirdPersonMap & mappe & livello & la mappa aperta da Play Game \\
BPC\_RayInteract & Blueprints, Components & Scene Component & traccia gli interagibili \\
CameraComponent & BPC\_RayInteract & Camera Component (ref.) & la camera del proprietario \\
DistanceTrace & BPC\_RayInteract & Float, 300 & la lunghezza del raggio \\
Interactable & BPC\_RayInteract & Actor (ref.) & l'interagibile guardato ora, o vuoto \\
Scene & BPC\_RayInteract & Scene Component (ref.) & solo per il trucco del Select \\
InteractionTrace & Project Settings & Trace Channel, Ignore & il canale delle interazioni \\
BPI\_InteractionRay & Blueprints, Components & Blueprint Interface & Start, End, Interaction \\
BP\_InteractableTest & Blueprints & Actor (Cube, Widget) & l'oggetto di prova \\
IA\_Interact & Input & Input Action, Pressed e Released & l'interazione, tasto E \\
Interacting & BP\_Bot2 & Boolean & vero durante l'interazione \\
Interact Press Button Montage & animazioni del bot & Anim Montage & l'animazione dell'interazione \\
WB\_InteractUI & UI & Widget Blueprint, 96 $\times$ 96 & il suggerimento ``E'' \\
Widget (componente) & BP\_InteractableTest & Screen, 100 $\times$ 100 & mostra WB\_InteractUI \\
Light & BP\_InteractableTest & BP\_ActivableLight (ref.) & la luce accesa e spenta \\
\bottomrule
\end{tabular}
\caption{Gli asset e le variabili della giornata, con i nomi usati dal docente (BPC\_RayInteract
è, nel progetto del docente, BPC\_RayInteract1; il nome della variabile della luce è
incerto).}
\label{ue12:tab:asset}
\end{table}

\begin{table}[htbp]
\centering
\small
\begin{tabular}{@{}p{5.4cm}p{8.6cm}@{}}
\toprule
\textbf{Sintomo} & \textbf{Causa e rimedio} \\
\midrule
La barra della stamina scende solo al rilascio & manca la Call E\_UpdateStamina in
StaminaDeduction: c'era solo in RegainStamina \\
Il secondo bottone incollato non compare & è sovrapposto al primo (controllare la
Hierarchy); oppure aggiungerne uno nuovo dalla Palette \\
L'immagine copre i bottoni & è in fondo alla Hierarchy, cioè disegnata sopra: trascinarla
prima dei bottoni \\
Il personaggio compare dietro il menu & GameMode Override nelle World Settings della
mappa del menu: Game Mode base \\
Si vede la pallina (Default Pawn) & spostare il Player Start fuori dalla camera \\
La camera del menu non si usa & Set View Target with Blend dal Get Player Controller (non
dal Camera Manager), con la referenza alla camera \\
Errore su Add to Viewport & il Target è Self: collegare il Return Value del Create Widget \\
Il cursore non si vede & Set Show Mouse Cursor nell'Event Construct, con la spunta attiva \\
Nel menu ci si muove ancora & Set Input Mode UI Only nell'Event Construct \\
Il Select non accetta Self (o la camera) & il Select prende il tipo del primo collegamento:
mettere prima una variabile Scene Component, poi cancellarla \\
In Play non si vede il raggio & LookForInteractions non chiamata nell'Event Tick; Draw
Debug Type a None \\
InteractionTrace non compare nella trace & compilare dopo aver creato il canale \\
Il raggio ignora il cubo & la mesh non blocca InteractionTrace (default Ignore): Block \\
L'evento dell'interfaccia non si trascina & non è Blueprint Callable: implementarlo come
evento (doppio clic o Implement Event) \\
Hit Actor non si collega a Interactable & Interactable di tipo sbagliato: Actor, Object
Reference \\
Si interagisce una sola volta & il Do Once non viene resettato (On Completed e Is Not
Valid) \\
Il Montage non parte & manca la Mesh nell'In Skeletal Mesh Component \\
Il nuovo widget contiene le barre dell'HUD & creato con l'HUD come genitore: scegliere User
Widget \\
Il bordo del widget non si vede & causa non trovata; rifare i Border \\
La E si vede quando non si guarda & spunte di Set Hidden in Game invertite: attivo nel
Begin Play e nell'End, disattivo nello Start \\
Trascinare la cartella UI dà ``fallimento'' & copiare la cartella estratta in Content/UI dal
gestore dei file \\
La BPI si apre senza grafo & doppio clic sulla funzione in My Blueprint \\
\bottomrule
\end{tabular}
\caption{I problemi incontrati durante la lezione, con la causa e il rimedio.}
\label{ue12:tab:errori}
\end{table}

\begin{table}[htbp]
\centering
\footnotesize
\begin{tabular}{@{}p{5.8cm}p{8.2cm}@{}}
\toprule
\textbf{Nodo, asset o impostazione} & \textbf{Che cosa fa} \\
\midrule
New Level: Empty / Basic & livello vuoto; con luce e cielo già pronti \\
Widget Blueprint (User Widget) & l'asset dell'interfaccia: Designer e Graph \\
Button: Style (Normal, Hovered, Pressed) & i colori del bottone nei tre stati \\
Hex sRGB / Hex Linear & il codice del colore, da copiare \\
Text dentro un Button & l'etichetta del bottone, figlia del bottone \\
Image: Color and Opacity (alfa) & la trasparenza dell'immagine \\
Hierarchy (ordine) & l'ultimo elemento è disegnato sopra \\
Anchors / Position / Alignment / Size & ancora; scostamento; pivot; dimensioni \\
Open Level Blueprint & il Blueprint del livello corrente \\
Create Widget / Add to Viewport & crea il widget; lo mostra \\
World Settings: GameMode Override & la Game Mode per questa mappa \\
Default Pawn / Player Start & la sfera della Game Mode base; dove nasce il Pawn \\
Camera Actor / Directional Light & una camera nel livello; il sole \\
Create a Reference to (actor) & la referenza a un actor del livello nel Level Blueprint \\
Set View Target with Blend & cambia la camera del Player Controller \\
Event Pre Construct / Construct & alla costruzione nel Designer; alla costruzione in gioco \\
Set Show Mouse Cursor & mostra o nasconde il cursore \\
Set Input Mode UI Only / Game Only & l'input va solo all'interfaccia; solo al gioco \\
On Clicked (Button) & l'evento del clic su un bottone \\
Quit Game & esce dal gioco (in editor, chiude il Play) \\
Open Level (by Name) & apre una mappa dal nome \\
Blueprint Interface & solo dichiarazioni di funzioni, in sola lettura \\
Class Settings: Implemented Interfaces & aggiunge un'interfaccia a un Blueprint \\
Interaction Start / End / Interaction (Message) & i messaggi dell'interfaccia \\
Event Interaction Start / End / Interaction & la loro implementazione nell'oggetto \\
Actor Component / Scene Component & componente senza / con Transform \\
Get Owner & l'actor che possiede il componente \\
Get Component by Class & il primo componente di una classe (non l'array) \\
Sphere Trace By Channel & una sfera lungo un segmento: Radius, Draw Debug Type \\
Get World Location / Get Forward Vector & posizione nel mondo; direzione davanti \\
Select (wildcard) & sceglie fra due valori; il tipo è quello del primo collegato \\
Is Valid (pura) / Validated Get & il riferimento esiste?; Get con le due uscite \\
Not Equal (!=) & vero se i due valori sono diversi \\
New Trace Channel (Ignore / Block) & un canale di collisione personalizzato \\
Break Hit Result: Hit Actor & l'actor colpito \\
Does Object Implement Interface & l'oggetto implementa l'interfaccia? \\
Event Tick & a ogni frame \\
Print String: Text Color & il colore del messaggio \\
Input Action (Pressed, Released) / IMC & l'azione e il tasto (E) \\
Do Once (Reset) & esegue una volta, finché non si resetta \\
Set Actor Rotation / Get Actor Rotation & imposta; legge la rotazione dell'actor \\
Find Look at Rotation / Break Rotator & rotazione per guardare un punto; i suoi tre assi \\
Create AnimMontage / Play Montage & un Montage da un'animazione; lo suona (On Completed) \\
Disable Input / Enable Input & toglie; restituisce l'input all'actor \\
Stop Movement Immediately & azzera il movimento del personaggio \\
Border: Brush, Padding, Image Size & il riquadro con immagine; margine interno; dimensione \\
Widget Component: Space, Draw Size, Widget Class & un widget attaccato a un actor \\
Set Hidden in Game (New Hidden) & nasconde o mostra un componente \\
Delay & aspetta prima di proseguire \\
Flip Flop & alterna fra le uscite A e B \\
\bottomrule
\end{tabular}
\caption{I nodi, gli asset e le impostazioni della dodicesima lezione.}
\label{ue12:tab:nodi}
\end{table}

% ---------------------------------------------------------------------
\section*{Punti chiave del capitolo}
\addcontentsline{toc}{section}{Punti chiave del capitolo}

\begin{itemize}
  \item Il designer usa Unreal per \textbf{prototipare} le meccaniche e verificare che
        funzionino; il progetto d'esame degli anni successivi è un \textbf{pezzo da
        portfolio}. \textbf{Less is more}: il giocatore impiega molto più tempo di chi
        ha progettato il gioco per capirne ogni elemento.
  \item Il menu vive in una \textbf{mappa a parte} (Empty Level), con un \textbf{Widget
        Blueprint} (User Widget) creato e mostrato nel \textbf{Level Blueprint} (Create
        Widget, Add to Viewport con il Return Value).
  \item Nel Designer: Canvas Panel, bottoni con l'\textbf{ancora} al centro e
        l'\textbf{Alignment} per impilarli, un \textbf{Text} dentro ogni bottone, i colori
        di \textbf{Normal}, \textbf{Hovered} e \textbf{Pressed} nello Style (copiabili con
        il codice Hex), e l'ordine della \textbf{Hierarchy} per mettere un'immagine
        dietro (l'ultimo elemento sta sopra).
  \item Niente personaggio nel menu: \textbf{GameMode Override} con la Game Mode base
        (Default Pawn) e Player Start fuori campo. La vista fissa: un Camera Actor e
        \textbf{Set View Target with Blend}, dal \textbf{Player Controller}.
  \item Nell'\textbf{Event Construct} del widget: \textbf{Set Show Mouse Cursor} (con la
        spunta) e \textbf{Set Input Mode UI Only}. Sui bottoni, \textbf{On Clicked}:
        \textbf{Quit Game}; \textbf{Set Input Mode Game Only} e \textbf{Open Level (by
        Name)}.
  \item Una \textbf{Blueprint Interface} è un intermediario (``un volo di linea''): solo
        nomi di funzioni, in sola lettura. Chi la implementa (Class Settings,
        Implemented Interfaces) decide che cosa fare; chi non la implementa ignora i
        messaggi.
  \item Un \textbf{Blueprint Component} isola una funzionalità da aggiungere a qualunque
        actor. L'\textbf{Actor Component} non ha posizione; lo \textbf{Scene Component}
        ha un Transform. Il componente trova il suo actor con \textbf{Get Owner} e la
        camera con \textbf{Get Component by Class}.
  \item \textbf{BPC\_RayInteract}: a ogni frame (\textbf{Event Tick}) una \textbf{Sphere
        Trace} (Radius 5) dalla posizione del componente, lungo il Forward Vector della
        camera (o del componente, se la camera non c'è) per 300 unità, sul canale
        \textbf{InteractionTrace} (Ignore di default, Block sugli interagibili).
  \item Il \textbf{Select} prende il tipo del primo collegamento: per scegliere fra la
        camera e Self, prima una variabile \emph{Scene Component}, poi la si cancella (la
        ``mandragata'').
  \item La variabile \textbf{Interactable} ricorda l'oggetto guardato: se cambia,
        \textbf{Interaction End} sul vecchio, Set, \textbf{Interaction Start} sul nuovo;
        se non si guarda niente di interagibile, End e svuota. Sempre con il
        \textbf{Validated Get}.
  \item Nel personaggio: \textbf{IA\_Interact} (tasto E), \textbf{Do Once}, il messaggio
        \textbf{Interaction}, \textbf{Set Actor Rotation} con il solo \textbf{Yaw} da
        \textbf{Find Look at Rotation}, \textbf{Disable Input} e \textbf{Stop Movement
        Immediately}, \textbf{Play Montage}, poi Enable Input e \textbf{Reset} del Do Once
        (anche quando non c'è nulla da interagire).
  \item Il suggerimento ``E'': un \textbf{Widget Component} in \textbf{Screen} space,
        nascosto nel Begin Play e mostrato e nascosto da Interaction Start ed End con
        \textbf{Set Hidden in Game}.
  \item Ogni oggetto può dare la sua risposta all'interazione: per esempio un
        \textbf{Flip Flop} che chiama Activation e Deactivation di una
        \textbf{BP\_ActivableLight}. Flip Flop e Do Once non vanno dentro le funzioni.
  \item La raccomandazione finale del corso: \textbf{validare sempre} le variabili.
\end{itemize}

%
%  Tutte le etichette di questo capitolo hanno il prefisso "ue12:"
%  (ue12:cap:..., ue12:sec:..., ue12:fig:..., ue12:tab:...), e tutti gli
%  stili TikZ hanno il prefisso "ue12" (ue12node, ue12exec, ue12data,
%  ...) per non collidere con chiavi native di TikZ né con quelli delle
%  lezioni precedenti. Il capitolo è autonomo: non contiene \ref verso
%  altri capitoli (i rimandi alle lezioni precedenti sono a parole).
%  Le figure sono ricostruzioni schematiche di quanto mostrato e detto a
%  lezione, non copie delle schermate del docente. Dove la trascrizione
%  è troppo rovinata per ricostruire un dettaglio, il capitolo lo dice
%  invece di inventarlo.
% =====================================================================

\chapter[Menu e sistema di interazione]{Il menu principale e il sistema di
interazione}
\label{ue12:cap:interazione}

Questa è la \textbf{dodicesima lezione} del corso breve di \textbf{introduzione a Unreal
Engine}, ed è l'\textbf{ultima} con questo docente: ``oggi è l'ultima lezione, almeno
mia''. Nella lezione precedente il personaggio, il \textbf{bot}, ha ricevuto la vita, il
danno della torretta, la stamina e le due barre dell'HUD; restavano in sospeso due
argomenti annunciati, il \textbf{menu} e le \textbf{interazioni}. Oggi si fanno entrambi.
``Sono argomenti che abbiamo già fatto in precedenza e adesso li uniamo'': Widget
Blueprint, Level Blueprint, trace, collisioni personalizzate, Montage, timer e attivabili
tornano tutti, messi insieme.

Il primo argomento, il più corto, è un \textbf{menu principale}: una mappa a parte, con un
piccolo \emph{diorama} 3D sullo sfondo, una camera fissa, due bottoni (\emph{gioca} ed
\emph{esci}) e il passaggio alla mappa di gioco. Il secondo, ``piuttosto lungo'', è un
\textbf{sistema di interazione} del tipo ``premi E per interagire'', costruito con due
strumenti nuovi: un \textbf{Blueprint Component}, che si aggiunge a qualunque actor e
``traccia'' dalla camera gli oggetti interagibili, e una \textbf{Blueprint Interface}, il
``porto'' attraverso cui il personaggio dice a un oggetto \emph{sto iniziando a guardarti},
\emph{non ti guardo più}, \emph{interagisco con te}. L'oggetto di prova mostra un piccolo
widget con la lettera E quando è inquadrato; premendo E il personaggio si gira verso di lui
e suona un'animazione; alla fine, lo stesso oggetto accende e spegne una luce.

La lezione è \textbf{online}, come le precedenti, e si apre con una conversazione più
generale: come recuperare le lezioni, che cosa aspetta gli studenti negli anni successivi,
e un consiglio di design, ``\emph{less is more}''. Durante la pausa il docente risolve il
problema della barra della stamina rimasto aperto la volta precedente; si chiude con i
ringraziamenti e con la raccomandazione che riassume il corso: ``ricordatevi di validare
sempre molte variabili''.

Tre avvertenze di lettura, le stesse delle lezioni precedenti. La prima: la trascrizione
è automatica e storpia moltissimi nomi. Unreal è ``Real'' e ``Arrial''; il
\textbf{diorama} è ``diuramino''; il Widget Blueprint \textbf{WB\_Menu} è ``VB menu'' e
``VBIN menu'', e \textbf{WB\_InteractUI} è ``VB Interact UI'' e ``Interacco uai''; il
\textbf{Begin Play} è ``big in play'' e ``back in play''; la \textbf{ThirdPersonMap} è
``Terperson map'' e ``there person''; il \textbf{Scene Component} è ``sin component'' e
``sim component''; il componente \textbf{BPC\_RayInteract} è ``BPC Reinteract'', ``Ray
interact one'' e ``IPC reinteract''; la Blueprint Interface \textbf{BPI\_InteractionRay} è
``VPI interaction'', ``VPA Interaction Array'' e ``bibliotecazione:ray''; l'\textbf{Hit
Actor} è ``it actor'' e ``I Tactor''; il \textbf{Validated Get} è ``variegate'' e ``vari da
get''; \textbf{Does Object Implement Interface} è ``Hacker does object implement
interface''; il \textbf{Draw Debug Type} è ``draw the bug type''; il bot è ``Boto'' e
``Boto2'', e il suo Blueprint ``BiPi Bot''; \textbf{Find Look at Rotation} è
``findLocatrotation''; l'\textbf{Anim Montage} è ``un montaggio un po' brasilero''; i
\emph{Big One} del corso di laurea sono ``biggest one''. ``Mandragata'', invece, non è un
errore: è il romanesco \emph{mandrakata}, ``trucchetto'', e il docente lo usa apposta.
Dove la ricostruzione è sicura il capitolo usa la forma corretta; dove non lo è riporta la
forma della trascrizione e lo dice. La seconda: durante la pausa il riconoscimento vocale
produce un lungo blocco di testo senza senso e ripetuto decine di volte (``grazie a tutti
per l'attenzione che ci ha spiegato'', ``qual è la risposta che ci può essere?''), seguito
da un frammento che sembra l'audio di un video in sottofondo (una ``blue map'' che mostra
``le secret room'' senza sprecare bombe, e ``un'altra monetina''): non appartengono alla
lezione e sono stati ignorati. La terza: buona parte della lezione è fatta di indicazioni
a voce mentre si lavora (``tira fuori'', ``attacca qua'', ``vai giù''), spesso rivolte
allo schermo di uno studente; il capitolo ricostruisce i grafi da ciò che il docente dice
mentre li costruisce, distingue i tentativi dalle soluzioni, completa i nomi esatti dei
nodi e dei pin di Unreal Engine 5, e dichiara i punti in cui la trascrizione lascia un
buco.

\paragraph{Come leggere questo capitolo.} Il capitolo segue l'ordine della lezione. Le
figure ricostruiscono la disposizione del menu, i due grafi che lo fanno funzionare (Level
Blueprint e Widget Blueprint), il ruolo di componente e interfaccia, la funzione che
cerca gli interagibili, il ciclo di vita della variabile Interactable e il grafo
dell'interazione nel personaggio. Le tabelle riassumono gli asset e le variabili della
giornata, gli errori incontrati durante la lezione e i nodi. Qua e là, dei paragrafi di
\textbf{approfondimento} collegano la lezione al libro di Marcos Romero e Brenden Sewell,
\emph{Blueprints Visual Scripting for Unreal Engine 5}, il testo di riferimento per questa
parte del libro, e a qualche nozione tecnica generale su Unreal Engine 5; sono segnalati
come tali e non fanno parte della lezione.

Il percorso del capitolo:
\begin{enumerate}
  \item le \textbf{domande d'apertura}: recuperare le lezioni, l'esame e i Big One, il
        lavoro del designer, \emph{less is more} (sezione~\ref{ue12:sec:inizio});
  \item gli \textbf{annunci}, la scelta dell'ordine e la \textbf{dimostrazione} del
        sistema di interazione (sezione~\ref{ue12:sec:annunci});
  \item il \textbf{menu}: la mappa, il widget, i bottoni e i loro colori
        (sezione~\ref{ue12:sec:menu});
  \item il \textbf{Level Blueprint} del menu: il widget sullo schermo, la Game Mode e la
        camera (sezione~\ref{ue12:sec:levelbp});
  \item il \textbf{Graph del menu}: cursore, input e bottoni; il diorama e i problemi
        degli studenti (sezioni \ref{ue12:sec:menugraph} e~\ref{ue12:sec:diorama});
  \item \textbf{interfacce} e \textbf{Blueprint Component} (sezione~\ref{ue12:sec:concetti});
  \item il componente \textbf{BPC\_RayInteract} e la sua Sphere Trace
        (sezione~\ref{ue12:sec:componente});
  \item il \textbf{Trace Channel} e la \textbf{Blueprint Interface}
        (sezione~\ref{ue12:sec:bpi});
  \item la \textbf{pausa}: la stamina di Tommaso (sezione~\ref{ue12:sec:pausa});
  \item la logica della variabile \textbf{Interactable} (sezione~\ref{ue12:sec:logica});
  \item l'oggetto di prova \textbf{BP\_InteractableTest} (sezione~\ref{ue12:sec:test});
  \item l'\textbf{input} e l'interazione nel personaggio (sezioni \ref{ue12:sec:input}
        e~\ref{ue12:sec:player});
  \item il widget \textbf{``premi E''} (sezione~\ref{ue12:sec:widget});
  \item la \textbf{luce} interattiva e la \textbf{chiusura} del corso (sezioni
        \ref{ue12:sec:luce} e~\ref{ue12:sec:chiusura}).
\end{enumerate}

% ---------------------------------------------------------------------
\section{Prima della lezione: recuperare, l'esame, il designer}
\label{ue12:sec:inizio}

\paragraph{Recuperare o seguire?} In apertura, prima che arrivino gli altri, uno
studente rimasto indietro (``le sto andando pian piano'') chiede se gli convenga
recuperare le lezioni vecchie o seguire questa. (Il nome, nella trascrizione, è incerto:
il docente lo chiama una volta ``Dave'', ma poco dopo ricorda questa conversazione come
quella fatta ``con Lorenzo''.) La risposta: seguire la lezione in diretta ha un vantaggio,
``se c'è un qualche particolare dell'argomento, puoi farmi direttamente una domanda''; poi,
siccome oggi si uniscono argomenti già visti, ``magari ti rifai anche quelle precedenti in
modo tale da stare un po' più al passo''. L'unico prerequisito per l'esercizio di oggi è
``avere il personaggio che si muove'', che ``può guardarsi in giro''; ``puoi anche usare
il personaggio standard, volendo''. Le domande si possono fare anche dopo, sul
\textbf{Discord}: ``mandarmi le domande anche quando non ci vediamo più, perché oggi è
l'ultima lezione''. Lo studente sceglie di seguire e di riguardare la registrazione.

\paragraph{L'esame e i Big One.} Seconda domanda: ci sarà un esame, con un progetto da
consegnare? ``Nel corso no'': questo corso ``è strutturato che sia un'introduzione per
voi''. Il docente chiede che corso farà lo studente (game designer: ``tanto farete sia game
sia level designer'') e descrive il percorso, avvertendo che è quello che ha fatto lui
(``almeno io feci così, adesso non so come sono cambiate le cose''):
\begin{itemize}
  \item al \textbf{primo anno} il progetto principale è un \textbf{gioco da tavolo}:
        ``farete analog game design'';
  \item al \textbf{secondo anno} il \textbf{Big One} è un progetto su Unreal, ``perché
        Unreal è molto più improntato per i designer'';
  \item al \textbf{terzo anno} il progetto più grande (nella trascrizione ``biggest one''),
        che può anche essere la continuazione del secondo, ``se il vostro progetto nel
        secondo anno è andato bene'', da portare fino al rilascio su \textbf{Steam},
        ``con l'Accademia assieme''.
\end{itemize}
In questo senso l'esame c'è: ``sarà farvi un \textbf{pezzo da portfolio}'', ``un gioco che
potete mettere sul vostro portfolio e promuovervi''. ``Quindi, più sapete, meglio è.''

\paragraph{Che cosa fa il designer con Unreal.} ``Io sono designer'': per un designer
Unreal ``servirà principalmente per prototipare le meccaniche'', cioè ``concretizzare
l'idea che tu hai studiato, che tu hai tirato fuori, concretizzarla in editor e vedere se
effettivamente funziona''. Il suo strumento sono i \textbf{Blueprint}, ``quindi non andare
troppo in depth'' nella programmazione: in ambito lavorativo il designer ha ``un
programmatore che lo assiste''; ``poi, quando si sta da soli, si fa di tutto''.

\paragraph{Less is more.} Un consiglio ``da esperienza'': \textbf{less is more}. ``Più è
leggibile un qualcosa anche nella vostra testa a livello di design, più sarà leggibile per
un giocatore.'' L'esempio è ``la regola delle parole a schermo'': il testo che scorre
all'inizio di \emph{Star Wars}. Chi l'ha scritto lo rilegge in venti secondi, perché lo
conosce già; lo spettatore, che deve ``cominciare a creare le connessioni mentali'', ci
mette ``non 20 secondi, ma 50''. Nel gioco vale lo stesso: chi lo ha progettato e testato
pensa ``vabbè, tanto questo lo so ormai'', ma il giocatore ``ci metterà il triplo del
tempo''. ``Più roba fate, più tempo il giocatore dovrà metterci per imparare tutto
quanto.'' ``Credo che \emph{less is more} ve lo faranno fare'', aggiunge, ``ve lo
diranno''; e per qualunque domanda di design, ``possiamo tranquillamente pure parlare''.

\paragraph{Approfondimento: la maledizione della conoscenza.} Il fenomeno descritto dal
docente ha un nome in psicologia, \emph{curse of knowledge}: chi conosce un'informazione
fatica a immaginare come la percepisce chi non la conosce. È il motivo per cui il
\emph{playtest} con persone esterne al team è insostituibile, e per cui le regole di un
gioco andrebbero introdotte una alla volta, ciascuna quando serve (il cosiddetto
\emph{onboarding} progressivo). I numeri del docente (20 secondi contro 50, ``il
triplo'') sono un'immagine, non una misura. Questa è una nozione generale di design, non
della lezione.

% ---------------------------------------------------------------------
\section{Gli annunci, l'ordine della giornata e la dimostrazione}
\label{ue12:sec:annunci}

\paragraph{Che cosa si fa oggi.} Prima che arrivino tutti, il docente chiede agli
studenti di scegliere: oggi ci sono ``sia il sistema di interazioni, che è piuttosto
lungo'', sia il \textbf{menu}, ``che l'altra volta non abbiamo finito, praticamente non
abbiamo neanche iniziato''. Prima il lungo o prima il corto? Gli studenti scelgono
``corto e lungo, così facciamo la pausa prima'' (``andiamo a un caffè''). Uno studente
ricorda un'altra cosa in sospeso: ``durante la pausa sistemare il problema della stamina,
che se tengo premuto non si vede, ma solo quando mollo il tasto si vede'' (è il problema
rimasto aperto alla fine della lezione precedente); ``ricordamelo appena andiamo in
pausa''.

\paragraph{Gli annunci.} Arrivati gli altri, il docente ripete: ``oggi praticamente è la
mia ultima lezione''. Seguiranno ``circa quattro lezioni'' (``adesso non mi ricordo quante
lezioni avete'') con \textbf{Davide}, dedicate alla parte di \textbf{level design}: ``i
tool che vi serviranno per fare level design'' e alcuni approfondimenti, ``credo di
materiali''. ``Ciò non vuol dire che io vi abbandoni'':
\begin{itemize}
  \item il \textbf{canale Discord} resta aperto, anche quando inizierà il corso vero e
        proprio: ``fate domande e dovrebbe arrivare qualcuno in vostro aiuto'';
  \item anche i lavori si possono postare per avere un feedback (il docente cita quello
        di \textbf{Federico La Porta}, ``che è bravissimo'', a cui ha già risposto);
  \item i due canali vanno tenuti separati: quello del corso (``Introduzione a Unreal'')
        per i compiti, ``così almeno non mandate mischiati, perché sennò perdo il filo'',
        e la chat generale ``per qualsiasi comunicazione generica'';
  \item le registrazioni restano disponibili ``fino a settembre, se non sbaglio'', ``quindi
        avete comunque abbastanza tempo per poterle recuperare''.
\end{itemize}

\paragraph{La dimostrazione.} ``Adesso vi faccio vedere solamente un'introduzione, come
nei classici video tutorial.'' Il personaggio è quello della volta scorsa, ``che si muoveva
con stamina e adesso perde vita con la torretta''. Dalla sua testa parte un \textbf{trace}
sottile, ``tipo un cilindrone'': ``questo sarà il nostro sistema che identificherà un
possibile \emph{interactable} tramite la telecamera''. A sinistra, sullo schermo, i
messaggi: se la camera inquadra l'oggetto compare \emph{inizio interazione}, se si
distoglie lo sguardo \emph{fine interazione}; premendo il tasto, il personaggio ``si gira in
automatico'' verso l'oggetto e interagisce. ``Visto che comunque abbiamo creato questo
sistema, è un buon stampo modulare'': in futuro, pulsanti e altri oggetti potranno avere
``l'interazione un po' più precisa'', riusando ciò che si è visto con i pulsanti classici:
``sempre premi pulsantino e fai cosetta, e animazione''.

% ---------------------------------------------------------------------
\section{Il menu: la mappa, il widget e i bottoni}
\label{ue12:sec:menu}

\paragraph{Il Widget Blueprint.} Nella cartella \textbf{UI} del progetto, tasto destro,
\emph{User Interface}, \textbf{Widget Blueprint}, ``come abbiamo creato per l'HUD''. Alla
richiesta della classe base, ``non ci serve una particolare, ce ne serve una generica'':
\textbf{User Widget}. Il nome: \textbf{WB\_Menu}.

\paragraph{Una mappa per il menu.} ``Una cosa importante che io almeno consiglio'': creare
una \textbf{nuova mappa} dedicata al menu, ``quindi differenziare il gioco dal menu''. ``Si
possono fare anche volendo insieme, ma a livello di gestione'' è più comodo separarli, ``ed
è molto più rapido''. Da \textbf{File}, \textbf{New Level}, si sceglie \textbf{Empty
Level}: ``se vi diventa tutto nero, è normale, ragazzi: è un empty level''. Un livello
\emph{Basic}, invece, contiene già ``la directional light, la skybox, tante cose che sono
già prefatte''; qui si parte da zero. Poi:
\begin{itemize}
  \item il menu si può fare ``sia in 2D, con immagini 2D'', sia con elementi 3D, ``una
        specie di diorama''; per cominciare, un \textbf{piano} abbastanza grande (dalle
        forme base: il docente prende prima un cubo, poi si corregge), con un materiale a
        griglia, ``il material base di Unreal'';
  \item si salva la mappa con \textbf{Ctrl+S} nella cartella delle mappe, con il nome
        \textbf{M\_Menu}.
\end{itemize}
(Un dettaglio marginale: il docente nota che qualcosa nel livello risulta \emph{Hidden in
Game}, poi si accorge che non lo è; la trascrizione non dice di che cosa si tratti.)

\paragraph{Il Canvas Panel.} Aperto WB\_Menu, ``la prima cosa, che in questo caso servirà
per adattarlo allo schermo'': il \textbf{Canvas Panel}, ``la classica canvas''. La
\textbf{Hierarchy} si vede ``da qua sotto''. A uno studente che non trova la
\textbf{Palette} il docente mostra il pannello a sinistra dell'editor.

\paragraph{Due bottoni.} ``Adesso andremo a fare due pulsanti molto semplici: uno di
\emph{exit} e uno di \emph{open the game}.'' Nella Palette, la sezione
\textbf{Common} (accanto alle altre: Audio, Input, Panel, Primitive e così via, che la
trascrizione elenca a fatica) contiene il \textbf{Button}; se ne trascinano due al centro
dello schermo. La disposizione passa per qualche tentativo ad alta voce (posizioni e
allineamenti come $-2{,}7$ e $-5$, presto abbandonati: ``ok, mi sono perso''); la versione
pulita, quella che il docente detta alla fine, è questa:
\begin{enumerate}
  \item \textbf{Size}: X \textbf{300}, Y \textbf{80} (``60? No, troppo piccolo'');
  \item \textbf{Anchors} al \textbf{centro}, l'icona centrale del menu delle ancore, ``in
        modo tale che questo bottone sia sempre ancorato al centro''; poi si
        \textbf{azzerano} i valori della \textbf{Position};
  \item \textbf{Alignment} \textbf{0,5} su X e \textbf{1,5} su Y, ``perché così è più
        corretto'': il bottone resta centrato orizzontalmente e si alza di una volta e
        mezza la sua altezza rispetto al centro dello schermo;
  \item per il secondo bottone, ``copiate e incollate letteralmente'' il primo, azzerate
        la posizione e cambiate solo l'allineamento verticale: dopo un ``$-1{,}5$? No,
        troppo'' e un ``$-0{,}5$'', il docente si ferma a \textbf{$-0{,}3$} (``va bene''),
        cioè il secondo bottone sta un po' sotto il centro.
\end{enumerate}
(I valori del secondo bottone, nella trascrizione, sono incerti: il senso è che i due
bottoni stanno uno sopra e uno sotto il centro, con un po' di spazio fra loro.) A uno
studente il copia e incolla non funziona: ``guarda nella Hierarchy: se c'hai due bottoni,
vuol dire che ce ne hai due'', magari sovrapposti; altrimenti, ``metti semplicemente un
altro bottone tramite il Common''.

\paragraph{I nomi.} ``Una cosa che è importante da fare'': nella Hierarchy si rinominano i
due bottoni con \textbf{F2}, \textbf{Button\_StartGame} e \textbf{Button\_QuitGame} (la
trascrizione dice ``Button With Game''). ``Così abbiamo semplicemente perfettamente
nomenclati.''

\paragraph{Il testo dentro il bottone.} ``Una cosa carina che già esiste'': un
\textbf{Text} (dalla sezione Common) trascinato \emph{sopra} un bottone ``si va a incastrare
per benino'', diventa suo figlio e si centra da solo. Nei \textbf{Details} del Text, sezione
\textbf{Content}, il campo \textbf{Text}: ``cancellate e scrivete \emph{Play Game}, o
\emph{Play}, come volete''; nell'altro bottone, \emph{Quit}. Il testo bianco ``non si vede
tantissimo'', ma i colori si cambiano subito dopo.

\paragraph{Lo Style del bottone.} Selezionando il bottone (``mi raccomando''), nei Details
la sezione \textbf{Style} ha tre stati, ciascuno con il suo colore:
\begin{itemize}
  \item \textbf{Normal}: il bottone a riposo;
  \item \textbf{Hovered}: ``vuol dire che ci passi sopra il cursore'';
  \item \textbf{Pressed}: ``se premuto''.
\end{itemize}
(La trascrizione nomina anche i \emph{foreground} e lo stato \emph{Disabled}: nelle
versioni recenti dell'engine lo stile del bottone ha anche i colori del primo piano per
ciascuno stato.) Cliccando sulla barra del colore si apre la ruota dei colori; il docente
mette il Normal \textbf{blu}, ``e avremo questo cambiamento in diretta''.

\paragraph{Copiare un colore esatto.} ``Se vogliamo prendere perfettamente il colore che ci
interessa per spostarlo all'altro bottone'': nella finestra del colore, in basso, ci sono i
campi \textbf{Hex sRGB} e \textbf{Hex Linear}. ``Potete prendere ambi i due'': si copia il
codice con Ctrl+C, si apre lo stesso stato nell'altro bottone e lo si incolla \emph{nello
stesso campo} (``se avete preso sRGB, sempre su sRGB''). È lo stesso principio dei siti
web: ``in HTML, per mettere un colore, i colori sono tutti identificati tramite un codice
colore''.

\paragraph{Un'immagine dietro i bottoni.} Si aggiunge un'\textbf{Image}, allargata dietro i
due bottoni; ma ``come potete vedere, mi copre tutto quanto''. ``Qui ci viene in aiuto la
Hierarchy'': come i livelli di Photoshop, gli elementi del Canvas sono disegnati
nell'ordine dell'elenco, e ``quello che è più sotto'' nell'elenco, stranamente, è quello
``a livello più alto'', disegnato sopra gli altri. L'immagine, aggiunta per ultima, è in
fondo all'elenco e quindi davanti ai bottoni. La si trascina nella Hierarchy fino a far
comparire una \textbf{linea blu} fra il Canvas Panel e il primo bottone, e la si lascia lì:
ora è disegnata per prima, cioè dietro. Poi la si rende trasparente: nel suo
\textbf{Color and Opacity} (che il docente cerca nel Brush), l'\textbf{alfa} a 0,5, ``o
forse un po' più basso'', \textbf{0,2}. ``Questa è la alfa di un colore: l'intensità,
l'opacità.'' Ne viene ``una specie di schermatina piccoletta'', regolata a mano attorno ai
bottoni (figura~\ref{ue12:fig:layout}).

\begin{figure}[htbp]
\centering
\begin{tikzpicture}[font=\footnotesize, x=1cm, y=1cm]
  % schermo
  \draw[very thick, fill=gray!5] (0,0) rectangle (8.0,5.0);
  \node[anchor=north west, font=\scriptsize] at (0.1,4.95) {Canvas Panel (schermo)};
  % immagine trasparente
  \draw[thick, fill=blue!8, opacity=0.9] (2.4,1.15) rectangle (5.6,4.15);
  \node[font=\scriptsize, anchor=north] at (4.0,4.12) {Image, alfa 0,2};
  % bottoni
  \draw[thick, fill=blue!35, rounded corners=2pt] (2.9,2.70) rectangle (5.1,3.30);
  \node at (4.0,3.0) {\textbf{Play Game}};
  \draw[thick, fill=blue!35, rounded corners=2pt] (2.9,1.62) rectangle (5.1,2.22);
  \node at (4.0,1.92) {\textbf{Quit}};
  % ancora al centro
  \draw[red!70!black, thick] (4.0,2.5) circle (0.09);
  \draw[red!70!black, dashed] (0,2.5) -- (8.0,2.5);
  \draw[red!70!black, dashed] (4.0,0) -- (4.0,5.0);
  \node[red!70!black, font=\scriptsize, anchor=west] at (5.7,2.62) {ancora al centro,};
  \node[red!70!black, font=\scriptsize, anchor=west] at (5.7,2.30) {Position 0, 0};
  \node[font=\scriptsize, anchor=west, align=left] at (5.7,3.35) {Alignment\\0,5 / 1,5};
  \node[font=\scriptsize, anchor=west, align=left] at (5.7,1.65) {Alignment\\0,5 / $-0{,}3$\\(circa)};
  % hierarchy
  \begin{scope}[xshift=9.0cm]
    \draw[thick, fill=gray!8] (0,0.6) rectangle (4.2,4.4);
    \node[anchor=north west, font=\scriptsize\bfseries] at (0.1,4.35) {Hierarchy};
    \node[anchor=west] at (0.2,3.6) {[WB\_Menu]};
    \node[anchor=west] at (0.45,3.1) {Canvas Panel};
    \node[anchor=west, fill=blue!10] at (0.7,2.6) {Image};
    \node[anchor=west] at (0.7,2.1) {Button\_StartGame};
    \node[anchor=west] at (0.95,1.65) {Text ``Play Game''};
    \node[anchor=west] at (0.7,1.2) {Button\_QuitGame};
    \node[anchor=west] at (0.95,0.8) {Text ``Quit''};
    \draw[-{Stealth[length=2mm]}, thick] (4.45,3.0) -- node[right, font=\scriptsize,
      align=left]{disegnati\\in ordine:\\l'ultimo\\sta sopra} (4.45,0.9);
  \end{scope}
\end{tikzpicture}
\caption{Il menu nel Designer. I due bottoni sono ancorati al centro dello schermo con la
posizione azzerata; è l'Alignment (il pivot del bottone) a spostarli sopra e sotto il
centro. L'immagine trasparente è stata trascinata in cima alla Hierarchy, così viene
disegnata per prima, dietro ai bottoni. Ricostruzione dalla lezione: i valori del secondo
bottone sono approssimativi.}
\label{ue12:fig:layout}
\end{figure}

\paragraph{Approfondimento: Alignment, ancore e Size To Content.} Con l'ancora al centro e
la posizione a zero, l'\textbf{Alignment} decide quale punto del bottone coincide con il
centro dello schermo: 0 è il bordo sinistro (o superiore), 1 il bordo destro (o inferiore),
0,5 il centro. Valori fuori dall'intervallo da 0 a 1 sono leciti: con Y = 1,5 il punto
d'aggancio sta mezzo bottone \emph{sotto} il bordo inferiore, quindi il bottone sale; con Y
negativo, scende. È il trucco usato a lezione per impilare due bottoni senza toccare la
posizione. La via più ordinata, per menu con molti bottoni, è un \textbf{Vertical Box}
(sezione Panel della Palette) ancorato al centro, con i bottoni dentro: lo spazio fra gli
elementi si regola con il \emph{Padding} di ciascuno e non servono calcoli. Nel libro di
Romero e Sewell (capitolo 8, \emph{WinMenu}) i bottoni sono ancorati al centro e spostati
con la \emph{Position} (per esempio X $-180$, Y 150, Size 360 $\times$ 100), con l'alignment
di default: è la terza via, equivalente. Funzionamento generale dell'UMG, salvo il
riferimento al libro.

% ---------------------------------------------------------------------
\section{Il Level Blueprint del menu: il widget, la Game Mode e la camera}
\label{ue12:sec:levelbp}

\paragraph{Perché il Level Blueprint.} ``Adesso che ce famo con questa roba?'' Il widget
va mostrato appena si apre la mappa del menu, e solo lì. Per questo si usa il
\textbf{Level Blueprint}, ``che abbiamo usato nella prima lezione'': ogni livello ha il
suo, ed è il posto giusto per la logica che appartiene a \emph{quella} mappa, ``perché
vogliamo che questa mappa sia dedicata apposta al nostro menu''. Si apre dalla toolbar
del Level Editor, dal pulsante dei Blueprint (``c'è il cubetto col più''),
\textbf{Open Level Blueprint}.

\paragraph{Il widget sullo schermo.} ``Ci interessa semplicemente usare il Begin Play,
perché vogliamo che si crei già da subito'':
\begin{enumerate}
  \item \textbf{Event BeginPlay} $\rightarrow$ \textbf{Create Widget}, con
        \textbf{Class} = WB\_Menu, ``come abbiamo visto l'altra volta'';
  \item \textbf{Add to Viewport}, con il \textbf{Return Value} del Create Widget nel
        Target. Il docente pensa per un momento di promuovere il widget a variabile, come
        per l'HUD, poi rinuncia: ``no, beh, non lo facciamo, tanto è il level blueprint''.
\end{enumerate}

\paragraph{Il personaggio di troppo.} In Play il menu compare, ma ``è molto strano:
abbiamo il nostro player dietro''. La mappa usa ancora la Game Mode del progetto, che fa
nascere il bot con i suoi controlli. Nel menu, invece, serve ``una pawn senza controlli,
senza niente'', e Unreal ce l'ha già: nelle \textbf{World Settings} della mappa, alla voce
\textbf{GameMode Override}, si sceglie la \textbf{Game Mode base} di Unreal (nella
trascrizione ``la game mode basica che esiste, la default''), che usa come Pawn la
\textbf{Default Pawn}. In Play, ``già è una cosa diversa: non c'è più il nostro character,
sparito, perché in questa mappa ho overridato la modalità Game Mode classica'', quella
``proprio base di Unreal, in cui non c'è character in terza persona''. ``Ma, come potete
vedere, non ho controllo'' (o meglio, come si vedrà, la Default Pawn si muove ancora:
il controllo vero si toglie nel Graph del widget, sezione~\ref{ue12:sec:menugraph}).

\paragraph{La camera fissa.} ``Andiamo a vedere anche un'altra cosa'': nel livello si
aggiunge un \textbf{Camera Actor} (``scrivete \emph{camera actor}, eccola qua''), che sarà
la vista fissa del menu. Insieme, una \textbf{Directional Light} (``e abbiamo messo il
sole, ragazzi''), regolabile a piacere, e un muretto, ``un po' più carino, così capiamo
meglio''. Per dire al gioco di guardare da quella camera, di nuovo nel Level Blueprint:
\begin{itemize}
  \item la \textbf{referenza} alla camera: selezionato il Camera Actor nel livello, nel
        Level Blueprint tasto destro, \textbf{Create a Reference to CameraActor};
  \item il nodo che cambia la vista. Qui il docente sbaglia strada in diretta: prova da
        \emph{Get Player Camera Manager}, non trova il nodo giusto e rimanda (``adesso non
        mi ricordo benissimo, lo vediamo un attimo dopo''). Dopo aver finito il Graph del
        menu, ritrova la soluzione: ``non è che bisogna prendere il camera manager, ma
        bisogna andare sul controller''. Quindi \textbf{Get Player Controller}
        $\rightarrow$ \textbf{Set View Target with Blend}, con la camera nel pin
        \textbf{New View Target}; il \emph{Blend Time} resta al default, ``non ci
        interessa''.
\end{itemize}
``Questo blocco ci permette di blendare una nuova camera'': ``gli diciamo: guarda, questa
è la tua nuova camera, da d'ora in poi'' (figura~\ref{ue12:fig:menugraph}, in alto).

\paragraph{Il Player Start fuori campo.} Ultimo ritocco, emerso aiutando gli studenti: la
Default Pawn ``è una palla'', e se nasce davanti alla camera si vede, o si vede il suo
riflesso. ``Prendete un \textbf{Player Start} e mettetelo fuori campo'', ``così almeno non vi
si vede la pallina''.

\paragraph{Approfondimento: Game Mode, Default Pawn e Player Start.} La \textbf{Game Mode}
è la classe che stabilisce le regole di una mappa e, fra le altre cose, quale
\emph{Pawn} far nascere per il giocatore (\emph{Default Pawn Class}) e con quale
\emph{Player Controller}. Il progetto ne ha una predefinita (in \emph{Project Settings},
\emph{Maps \& Modes}); il \textbf{GameMode Override} delle World Settings la sostituisce
\emph{solo per quella mappa}, ed è il modo standard di avere un menu senza personaggio.
La classe scelta a lezione, \emph{GameModeBase}, usa come Pawn la \textbf{DefaultPawn}:
una sfera con un movimento ``da spettatore'' (WASD e mouse, senza gravità), pensata per
fare prove. Il Pawn nasce in un \textbf{Player Start}; se la mappa non ne ha, nasce
all'origine o nella posizione della camera dell'editor. Un'alternativa ancora più pulita
per un menu è una Game Mode con la \emph{Default Pawn Class} impostata su \emph{None}: il
Player Controller esiste (e quindi esistono il cursore e gli input dell'interfaccia), ma
nessun Pawn nasce. Funzionamento generale dell'engine.

\paragraph{Approfondimento: Set View Target with Blend nel libro.} Romero e Sewell
descrivono il nodo nel capitolo 15 (\emph{Blueprints Tips}): ``è una funzione della classe
Player Controller'', ``usata per passare la vista di gioco fra camere diverse''; ``il
parametro di input New View Target è l'actor da usare come view target, di solito una
camera''. Il loro esempio è proprio in un Level Blueprint: quando il giocatore entra nella
stanza del tesoro, la vista passa alla camera della stanza. Questo spiega l'errore del
docente: il \emph{Player Camera Manager} gestisce la vista, ma la funzione esposta ai
Blueprint per cambiarla appartiene al \emph{Player Controller}. Il \textbf{Blend Time} a 0
fa uno stacco netto; un valore positivo fa una transizione morbida dalla vista precedente,
con una curva scelta in \emph{Blend Func}.

% ---------------------------------------------------------------------
\section{Il Graph del menu: cursore, input e bottoni}
\label{ue12:sec:menugraph}

\paragraph{Event Construct.} Si torna in WB\_Menu, scheda \textbf{Graph}: ``vediamo le
funzionalità del nostro menu quando clicchiamo le cose''. Dei tre eventi già presenti
``non ci serve il Pre Construct, non ci serve il Tick''. Serve l'\textbf{Event Construct}:
``vuol dire quando questo widget è costruito'', ``è un po' il Begin Play'', ``viene
differenziato dal nome, giusto per non confonderti''.

\paragraph{Il cursore e l'input dell'interfaccia.} In Play, per ora, ``mi posso muovere''
(la Default Pawn) e il mouse non si vede. Nell'Event Construct:
\begin{enumerate}
  \item \textbf{Get Player Controller} $\rightarrow$ \textbf{Set Show Mouse Cursor},
        con la spunta attiva: il cursore resta visibile anche cliccando nella finestra di
        gioco. ``Però, ovviamente, posso ancora muovermi'';
  \item \textbf{Set Input Mode UI Only} (``enable input UI'', dice prima il docente, poi
        trova il nodo), con lo stesso Player Controller: ``questo ci permette di
        utilizzare tranquillamente il nostro widget''. Il pin \emph{Widget to Focus} resta
        vuoto: ``no, tanto non ci serve''.
\end{enumerate}
In Play, ora, ``se noi proviamo a muovere lateralmente, non si muove più''. ``Mi
raccomando: questa cosa, quando noi cliccheremo un pulsante, dovrà essere disattivata e
rimessa in un'altra modalità.''

\paragraph{I colori di hover e pressione.} Con il cursore visibile si possono provare gli
stati dello Style. Nel Designer il docente cambia il colore di \textbf{Hovered}: ``voglio
che quando io ci passo sopra con il mouse diventi tipo arancione, sì, un bell'arancione''
(``lo vediamo in preview, è normale''), e lo copia, con il codice esadecimale, sull'altro
bottone. Il risultato, ammette, è ``molto Ikea''. Il \textbf{Pressed} di default diventa
grigio, ``si vede pochissimo'': lo si mette \textbf{rosso}, su entrambi i bottoni.

\paragraph{I bottoni.} Nel Graph, a sinistra, fra le variabili compaiono
\textbf{Button\_QuitGame} e \textbf{Button\_StartGame} (sono variabili perché i bottoni,
nel Designer, hanno \emph{Is Variable} attivo di default). Ogni bottone ha i suoi eventi,
``l'abbiamo visto anche l'altra volta'': \emph{On Pressed}, \emph{On Released},
\emph{On Hovered}, \emph{On Unhovered} e \textbf{On Clicked}. Si prende l'On Clicked di
entrambi.
\begin{itemize}
  \item \textbf{On Clicked (Button\_QuitGame)} $\rightarrow$ \textbf{Quit Game}: ``è
        piuttosto semplice, letteralmente Quit Game è già presente, di standard
        proprio''. (``Adesso non mi ricordo se ti esce proprio dall'editor'': si veda
        l'approfondimento più sotto.)
  \item \textbf{On Clicked (Button\_StartGame)}: prima \textbf{Get Player Controller}
        $\rightarrow$ \textbf{Set Input Mode Game Only}, ``molto importante'', ``in modo
        tale che ritornino gli input al gioco''; poi \textbf{Open Level (by Name)}, con il
        nome della mappa di gioco. Il docente lo cerca nel Content Browser (``era
        MLevel1? Ah, questa no, era M level 2?'') e trova la mappa del template,
        \textbf{ThirdPersonMap}; prova a copiarne il nome con F2 (``non lo fa, va bene''),
        e lo scrive a mano.
\end{itemize}
Tornati nella mappa del menu e salvato tutto, in Play: Play Game ``ci carica la nostra
bellissima mappina'' (figura~\ref{ue12:fig:menugraph}).

\begin{figure}[htbp]
\centering
\begin{tikzpicture}[font=\footnotesize,
  ue12node/.style={draw, thick, rounded corners=4pt, align=center, fill=gray!8,
                  minimum height=0.9cm, text width=2.3cm},
  ue12ev/.style={ue12node, fill=red!12},
  ue12pure/.style={ue12node, fill=green!10, rounded corners=8pt},
  ue12exec/.style={-{Stealth[length=2.5mm]}, very thick, black!80},
  ue12data/.style={-{Stealth[length=2mm]}, thick, blue!60!black}]
  \node[font=\scriptsize\bfseries, anchor=west] at (-1.3,1.0) {Level Blueprint di M\_Menu};
  \node[ue12ev, text width=1.7cm] (bp) at (0,0) {\textbf{Event BeginPlay}};
  \node[ue12node] (cw) at (2.8,0) {\textbf{Create Widget}\\{\scriptsize Class: WB\_Menu}};
  \node[ue12node, text width=2.5cm] (av) at (5.8,0) {\textbf{Add to Viewport}\\{\scriptsize Target: Return
    Value}};
  \node[ue12node, text width=2.6cm] (vt) at (8.9,0) {\textbf{Set View Target with
    Blend}\\{\scriptsize Blend Time 0}};
  \node[ue12pure] (gpc1) at (7.4,-1.8) {Get Player Controller};
  \node[ue12pure, fill=blue!10] (cam) at (10.2,-1.8) {CameraActor\\{\scriptsize
    (referenza)}};
  \draw[ue12exec] (bp) -- (cw); \draw[ue12exec] (cw) -- (av); \draw[ue12exec] (av) -- (vt);
  \draw[ue12data] ([yshift=-3mm]cw.east) -- ([yshift=-3mm]av.west);
  \draw[ue12data] (gpc1.north) -- node[left, font=\scriptsize]{Target} ([xshift=-6mm]vt.south);
  \draw[ue12data] (cam.north) -- node[right, font=\scriptsize]{New View Target}
    ([xshift=6mm]vt.south);
  % widget graph
  \node[font=\scriptsize\bfseries, anchor=west] at (-1.3,-3.2) {Graph di WB\_Menu};
  \node[ue12ev, text width=1.7cm] (con) at (0,-4.2) {\textbf{Event Construct}};
  \node[ue12node] (smc) at (2.8,-4.2) {\textbf{Set Show Mouse Cursor}\\{\scriptsize
    spunta attiva}};
  \node[ue12node] (ui) at (5.7,-4.2) {\textbf{Set Input Mode UI Only}};
  \node[ue12pure] (gpc2) at (4.2,-5.9) {Get Player Controller};
  \draw[ue12exec] (con) -- (smc); \draw[ue12exec] (smc) -- (ui);
  \draw[ue12data] (gpc2.north) -- (smc.south); \draw[ue12data] (gpc2.north) -- (ui.south);
  \node[ue12ev, text width=2.6cm] (oq) at (0.3,-7.6) {\textbf{On Clicked}\\{\scriptsize
    (Button\_QuitGame)}};
  \node[ue12node] (qg) at (3.6,-7.6) {\textbf{Quit Game}};
  \draw[ue12exec] (oq) -- (qg);
  \node[ue12ev, text width=2.6cm] (os) at (0.3,-9.3) {\textbf{On Clicked}\\{\scriptsize
    (Button\_StartGame)}};
  \node[ue12node, text width=2.7cm] (gm) at (3.8,-9.3) {\textbf{Set Input Mode Game Only}};
  \node[ue12node, text width=2.8cm] (ol) at (7.2,-9.3) {\textbf{Open Level (by
    Name)}\\{\scriptsize ThirdPersonMap}};
  \node[ue12pure] (gpc3) at (3.8,-11.0) {Get Player Controller};
  \draw[ue12exec] (os) -- (gm); \draw[ue12exec] (gm) -- (ol);
  \draw[ue12data] (gpc3) -- (gm);
\end{tikzpicture}
\caption{I grafi del menu. In alto, il Level Blueprint della mappa M\_Menu crea il widget,
lo mostra e passa la vista alla camera fissa. In basso, il widget prende il controllo del
cursore quando nasce, esce dal gioco con Quit e, con Play Game, restituisce l'input al
gioco e apre la mappa di gioco. Ricostruzione dalla lezione.}
\label{ue12:fig:menugraph}
\end{figure}

\paragraph{Approfondimento: il menu nel libro.} Romero e Sewell costruiscono un menu molto
simile nel capitolo 8 (\emph{Creating Constraints and Gameplay Objectives}), il
\emph{WinMenu} che compare a fine partita: due bottoni \emph{Btn restart} e \emph{Btn
quit} ancorati al centro, ciascuno con un Text dentro (``font size 48''), e gli eventi
aggiunti dal pulsante \emph{+} accanto a \emph{On Clicked}, in fondo ai Details del bottone.
Il bottone Restart usa \emph{Open Level (by Object Reference)}, con la mappa scelta da un
menu a tendina, seguito da \emph{Remove from Parent}, ``che rimuove il widget dalla
vista''; il bottone Quit usa \emph{Quit Game}, ``così che il giocatore possa chiudere il
gioco''. Il menu viene mostrato dal personaggio, in un Custom Event \emph{End Game}: Set
Game Paused, \emph{SET Show Mouse Cursor} dal Get Player Controller (``per permettere al
giocatore di riprendere il controllo del mouse''), Create Widget, Add to Viewport. Nel
capitolo 15 il libro elenca anche le tre modalità di input:
\begin{itemize}
  \item \textbf{Set Input Mode Game Only}: ``solo il Player Controller riceve gli eventi
        di input'';
  \item \textbf{Set Input Mode UI Only}: ``solo la UI riceve gli eventi di input'';
  \item \textbf{Set Input Mode Game and UI}: ``la UI ha la priorità nel gestire un evento
        di input, ma se la UI non lo gestisce, lo riceve il Player Controller''; l'esempio
        è un negozio, in cui si scelgono le opzioni con il mouse ma ci si può allontanare
        con le frecce.
\end{itemize}
Una differenza utile da notare: \emph{Open Level (by Object Reference)} non si rompe se la
mappa viene rinominata, mentre \emph{Open Level (by Name)}, usato a lezione, richiede il
nome esatto (si può scrivere il nome semplice, \emph{ThirdPersonMap}, o il percorso
completo, se esistono due mappe con lo stesso nome).

\paragraph{Approfondimento: Construct, Quit Game e il cursore nella mappa di gioco.}
Alcune nozioni generali, non della lezione. \textbf{Event Pre Construct} viene eseguito
anche nell'editor, ogni volta che il widget viene ridisegnato nel Designer: serve per
anteprime e impostazioni visive; \textbf{Event Construct} solo in gioco, quando il widget
viene aggiunto allo schermo (e di nuovo se lo si rimuove e lo si aggiunge). \textbf{Quit
Game}, durante il Play in Editor, \emph{termina la sessione di gioco} e riporta
all'editor, senza chiudere l'editor; in un gioco pacchettizzato chiude l'applicazione. Il
dubbio del docente si scioglie così. Infine, \emph{Open Level} distrugge il mondo corrente
con tutti i suoi actor, widget compresi, e crea un Player Controller nuovo: per questo,
nella mappa di gioco, il menu sparisce da solo e il cursore torna nascosto, anche senza un
\emph{Remove from Parent} o un \emph{Set Show Mouse Cursor} a falso. Il \emph{Set Input
Mode Game Only} prima dell'Open Level, raccomandato dal docente, è comunque una buona
abitudine: il caso in cui non è superfluo è quello di un menu che si chiude \emph{senza}
cambiare livello, come un menu di pausa (nel libro, capitolo 11, il \emph{Resume} fa
Remove from Parent, Set Show Mouse Cursor a falso e Set Game Paused a falso).

% ---------------------------------------------------------------------
\section{Il diorama, e i problemi degli studenti}
\label{ue12:sec:diorama}

\paragraph{Il diorama.} Con il menu funzionante, ``passiamo sopra'' all'aspetto. Sullo
sfondo, ``un piccolo diorama, una piccola box, in cui ci andiamo a mettere tutto quello che
vogliamo''. ``Volendo, potete fare tutto un menu interattivo'': gli ostacoli, i Blueprint
creati nel corso, i bottoni, la torretta. Il docente si limita ai modelli:
\begin{itemize}
  \item il \textbf{bot}: dal Content Browser, o la sua Skeletal Mesh, ``o li metto
        direttamente l'animazione in Idle'': trascinando nel livello un asset di animazione,
        Unreal crea un actor con la mesh che la riproduce;
  \item la \textbf{porta} modulare del compito della volta scorsa, \textbf{SM\_Modular\_\allowbreak Wall\_Door}
        (``wall door, eccola qua''), con le ante \textbf{SM\_Door};
  \item qualche ritocco alla \textbf{camera} e alla \textbf{luce}: l'intensità a 10 è
        ``enorme'', poi 5, poi valori molto più bassi (la trascrizione, qui, è quasi
        incomprensibile: si sentono ``0,01'' e ``1'').
\end{itemize}
``Questo caso è più una cosa artistica, ragazzi, quindi potete fare un po' come volete.''
Le funzioni di base, entrare nel gioco e uscire, ci sono. ``Questo è un menu che io di
solito faccio quando faccio i miei prototipi'', cioè ``quando devo semplicemente mostrare
qualche cosa a livello di gameplay''. Un menu delle \emph{opzioni} ``è già un discorso
molto più'' complesso, ``perché bisogna andare a vedere gli slider'': non si tratta, ``perché
tanto poi andiamo a fare qualcosa di più gameplayoso''.

\paragraph{I problemi dei presenti.} Diversi studenti hanno ancora il personaggio nel
menu.
\begin{itemize}
  \item \textbf{La camera.} Ripasso passo per passo: selezionare la camera nel livello
        (``proprio diventa arancione''), Open Level Blueprint, tasto destro, Create a
        Reference to CameraActor (``devi stare dentro il Level Blueprint''); Get Player
        Controller, Set View Target with Blend, la camera nel New View Target.
  \item \textbf{Il personaggio che nasce ancora.} ``Devi passare alla Game Mode
        default'': \textbf{World Settings}, \textbf{GameMode Override}. Uno studente non
        trova la finestra delle World Settings (si apre dal menu \emph{Window} dell'editor,
        se è chiusa). Poi il Player Start fuori dalla camera.
  \item \textbf{L'errore su Add to Viewport.} ``Mi dà un problema proprio su Add to
        Viewport'': ``perché il target di base è Self, quindi non so se devo cambiare il
        target''. Sì: il Target deve essere il \textbf{Return Value} del Create Widget.
        ``Probabilmente devi dare return value'' al target.
  \item \textbf{Il mouse che non si vede.} Va impostato nell'Event Construct del widget:
        Get Player Controller, Set Show Mouse Cursor, Set Input Mode UI Only. Lo studente
        aveva cancellato l'Event Construct (``basta che fai tasto destro'' e lo si ricrea)
        e non aveva messo la \textbf{spunta} nel Set Show Mouse Cursor (``ah, però non
        avevo messo il tick''): il nodo, senza spunta, imposta il cursore a \emph{falso}.
        ``Funziona? Sì.''
\end{itemize}
Chiuso il menu, si torna alla ThirdPersonMap.

% ---------------------------------------------------------------------
\section{Interfacce e Blueprint Component}
\label{ue12:sec:concetti}

``Andiamo sul nostro amico fidato'': il docente apre il suo bot (``ah no, questo è Boto:
dobbiamo andare su Boto2''), cioè BP\_Bot2. ``Oggi parleremo delle interfacce, ragazzi.''

\paragraph{L'interfaccia, un intermediario.} ``Immaginatevi che un'interfaccia è un
intermediario fra due o più Blueprint.'' Il docente mostra quella già pronta nel suo
progetto, nella cartella dei componenti dei Blueprint: un tipo particolare di Blueprint
che contiene solo tre \textbf{funzioni}, \emph{Interaction Start}, \emph{Interaction End}
e \emph{Interaction}. L'immagine: un \textbf{volo di linea}, ``un'interconnessione tra due
o più poli'', che ``trasporta molto spesso le stesse cose, passeggeri di un tipo,
passeggeri di un altro''. ``Quindi è un qualcosa che va a trasportare varie funzioni da una
parte all'altra.'' Nel caso di oggi l'interfaccia comunica tre messaggi: \emph{sto
iniziando un'interazione} (ti sto guardando), \emph{sto finendo un'interazione} (non ti
guardo più), \emph{sto interagendo} (ho premuto il tasto).

\paragraph{Il componente, una funzionalità isolata.} Il lavoro di \emph{vedere se c'è
qualcosa di interagibile} non lo fa il personaggio, ma un \textbf{Blueprint Component}:
``dei componenti che noi possiamo aggiungere all'interno dei nostri actor per dare
funzionalità in modo isolato''. Per esempio: ``questo actor può tracciare degli
interactable''; ``questo actor ha la vita''. La vita, fatta la volta scorsa \emph{dentro}
il personaggio, si potrebbe spostare in un componente, ``perché magari abbiamo anche i
nemici che hanno la vita'', o ``degli oggetti che si possono rompere''. Lo scopo è
``isolare sempre più possibile le funzionalità, e quindi essere anche modulari''. Infatti
il BP\_Bot2 del docente, a parte un input nuovo, ``non ha niente di nuovo'': si gira verso
l'oggetto e suona un Montage, ma ``non ha niente per tracciare'': il trace è tutto nel
componente.

\paragraph{Actor Component o Scene Component.} Si crea una \textbf{Blueprint Class} e, fra
le classi comuni, non un Actor ma un componente. Ce ne sono due:
\begin{itemize}
  \item l'\textbf{Actor Component} ``mette semplicemente la funzionalità'';
  \item lo \textbf{Scene Component} ``è un Actor Component che ha anche una scena'',
        cioè un \textbf{Transform}: ``non so se vi ricordate la Scene, quando l'abbiamo
        messa: il punto vuoto''. ``Almeno si può spostare nello spazio.''
\end{itemize}
Qui serve il secondo, perché il trace è ``un raggio che da quel punto parte e in un punto
finisce'', e il punto di partenza si deve poter spostare. Il docente lo mostra nel suo bot:
spostando il componente nel Blueprint del personaggio, in Play il raggio parte dal nuovo
punto; lo rimette ``al centro della testa, più o meno''. Il nuovo Blueprint:
\textbf{BPC\_RayInteract} (BPC sta per \emph{Blueprint Component}; il docente, che ne ha già
uno, chiama il suo BPC\_RayInteract1).

\paragraph{Il componente aperto.} ``Come potete vedere, presenta delle similitudini con
tutti i Blueprint.'' ``L'unica cosa: non ha una viewport, perché non ha una forma fisica'',
``tra virgolette, perché è una scena, ma è vuota''. ``E ricordiamoci che è un componente:
quindi sarà qualcosa che deve essere per forza attaccato a qualcos'altro.''

\paragraph{Che cos'è una funzione.} Prima di scrivere la logica, una domanda: chi ha già
sentito parlare di \textbf{funzioni} in informatica? Risponde \textbf{Federico}, che ha
studiato da programmatore: ``la funzione è un insieme di istruzioni: chiami la funzione e
fai quelle istruzioni''. ``Esattamente.'' Il docente indica una catena di nodi nel bot: ``gli
eventi non possono essere messi come funzioni, ma questo qua può essere una funzione''. È
``un modo per ordinare anche il lavoro in varie tipologie di funzionalità, che possono
essere magari trasportabili da un punto all'altro'': ``sempre questo senso di modularità''.
(``Grande Federico, sempre miglior assistente.'')

\begin{figure}[htbp]
\centering
\begin{tikzpicture}[font=\footnotesize,
  ue12box/.style={draw, thick, rounded corners=5pt, align=center, inner sep=5pt},
  ue12msg/.style={-{Stealth[length=2.5mm]}, very thick, violet!70!black},
  ue12own/.style={-{Stealth[length=2.2mm]}, thick, dashed, black!70}]
  \node[ue12box, fill=gray!8, minimum width=4.4cm, minimum height=4.6cm] (bot) at (0,-0.3) {};
  \node[anchor=north, font=\bfseries] at (bot.north) {BP\_Bot2};
  \node[ue12box, fill=blue!10, text width=3.6cm] (cam) at (0,0.8) {Camera Component};
  \node[ue12box, fill=orange!18, text width=3.6cm] (comp) at (0,-1.4)
    {\textbf{BPC\_RayInteract}\\{\scriptsize Scene Component: ogni frame una Sphere Trace
    lungo la camera; variabile Interactable}};
  \draw[ue12own] (comp.north) -- node[right, font=\scriptsize, text width=1.9cm,
    align=left]{Get Owner,\\Get Component by Class} (cam.south);
  \node[ue12box, fill=violet!10, text width=3.2cm] (bpi) at (6.6,-0.3)
    {\textbf{BPI\_InteractionRay}\\{\scriptsize Blueprint Interface}\\[2pt]{\scriptsize
    Interaction Start\\Interaction End\\Interaction}};
  \node[ue12box, fill=green!10, text width=3.3cm] (obj1) at (11.6,1.0)
    {\textbf{BP\_InteractableTest}\\{\scriptsize implementa la BPI: Print String, widget
    ``E'', luce}};
  \node[ue12box, fill=green!5, text width=3.3cm] (obj2) at (11.6,-1.6) {{\scriptsize un
    altro oggetto}\\{\scriptsize (bottone, ascensore\dots) con la stessa BPI e un'altra
    risposta}};
  \draw[ue12msg] (bot.east |- cam) -- node[above, font=\scriptsize, text width=2.4cm,
    align=center]{Interaction\\(tasto E)} (bpi.west |- cam);
  \draw[ue12msg] (comp.east) -- node[below, font=\scriptsize, text width=2.4cm,
    align=center]{Start / End\\(guardare, smettere)} (bpi.west |- comp);
  \draw[ue12msg] (bpi.east) -- (obj1.west);
  \draw[ue12msg] (bpi.east) -- (obj2.west);
\end{tikzpicture}
\caption{I ruoli del sistema di interazione. Il componente, attaccato al personaggio, trova
la camera del suo proprietario, traccia ciò che il giocatore guarda e invia i messaggi di
inizio e fine; il personaggio invia il messaggio Interaction quando si preme E. I messaggi
viaggiano attraverso l'interfaccia: chi li riceve decide che cosa fare, e un actor che non
la implementa li ignora. Schema ricostruito dalla lezione.}
\label{ue12:fig:ruoli}
\end{figure}

\paragraph{Approfondimento: i componenti nel libro.} Romero e Sewell dedicano ai componenti
la seconda metà del capitolo 18 (\emph{Creating Blueprint Libraries and Components}).
``Quando creiamo un Actor Blueprint aggiungiamo spesso componenti con una funzionalità
incapsulata, pronta all'uso, come Projectile Movement, Static Mesh e Collision''; ``possiamo
anche creare i nostri Actor Component con i Blueprint''. Nel pannello \emph{Pick Parent
Class} le due classi sono fra quelle \emph{comuni}: ``lo Scene Component è una classe figlia
dell'Actor Component che ha la struttura Transform (location, rotation e scale). Grazie al
Transform, uno Scene Component può essere attaccato a un altro Scene Component.'' E, come a
lezione: ``quando si scrive la logica di un componente nell'Event Graph, si può ottenere la
referenza all'actor che usa il componente con il nodo \textbf{Get Owner}''. I due esempi del
libro mostrano bene la differenza: \emph{BP\_ExpLevelComp} è un Actor Component che gestisce
punti esperienza e livelli (non ha bisogno di una posizione), mentre \emph{BP\_CircularMovComp}
è uno Scene Component che ruota attorno all'actor e a cui si attacca una Static Mesh, per
simulare uno scudo rotante: ``lo Scene Component è usato come posizione di riferimento per
altri componenti''. Il componente della lezione è dello stesso tipo: la sua posizione è il
punto di partenza del raggio.

% ---------------------------------------------------------------------
\section{Il componente BPC\_RayInteract: la Sphere Trace}
\label{ue12:sec:componente}

\paragraph{La funzione LookForInteractions.} Nel componente si crea una nuova funzione
(``eccolo che vi si crea questo pannellino violaceo''), \textbf{LookForInteractions}.
Dentro, una \textbf{Sphere Trace By Channel}. ``Seguitemi per bene.''
\begin{itemize}
  \item \textbf{Radius}: \textbf{5}, ``perché se no non vi apparirà mai''. Con raggio zero
        la Sphere Trace degenera in un raggio sottilissimo (e il debug quasi non si vede);
        5 unità bastano a renderla visibile e un po' più tollerante.
  \item \textbf{Start}: \textbf{Get World Location} del componente stesso, ``la posizione
        nel mondo del nostro Scene Component: dovrà partire da qui''.
  \item \textbf{End}: da costruire, ed è la parte con il trucco.
\end{itemize}

\paragraph{Prima la camera: Get Owner e Get Component by Class.} Il raggio deve seguire lo
sguardo del giocatore, cioè la \emph{rotazione della camera}. Il componente, però, non sa
nulla della camera del personaggio: deve trovarla. Nell'\textbf{Event Graph} del
componente, all'\textbf{Event BeginPlay}:
\begin{enumerate}
  \item \textbf{Get Owner}: ``la referenza del proprietario di questo Scene Component'';
  \item dal proprietario, \textbf{Get Component by Class} (``non \emph{components},
        \emph{component}, perché se no vi dà un array''), con \textbf{Component Class} =
        \textbf{Camera Component}: ``sarà il nostro punto principale da cui andremo a fare il
        nostro trace'';
  \item il risultato si promuove a variabile, \textbf{CameraComponent}, e si collega tutto
        al Begin Play.
\end{enumerate}
``Perdonate, sono 30.000 componenti: un componente su componente su componente.'' In
sintesi: ``in Begin Play io prendo il mio proprietario, prendo la sua camera, mi salvo la
referenza alla sua camera, e poi la useremo all'interno del nostro componente''.

\paragraph{La camera, se c'è.} Il componente deve funzionare anche su un actor
\emph{senza} camera: ``se il nostro proprietario non ha una camera, allora prendi te stesso
come pivot''. Quindi un \textbf{Select} fra due oggetti:
\begin{itemize}
  \item l'indice è un booleano, \textbf{Is Valid} (``la F verde'', la versione \emph{pura}
        della funzione) sulla variabile CameraComponent: ``se il componente camera esiste,
        quindi il nostro actor ha effettivamente una camera, e quindi questa variabile
        CameraComponent ha qualcosa dentro'';
  \item su \textbf{True}, CameraComponent: ``usami la camera come punto di riferimento'';
  \item su \textbf{False}, \textbf{Self}: ``se no, usi te stesso''.
\end{itemize}
Il Select era già comparso con la torretta, per scegliere fra le due bocche di fuoco: lì
con un \emph{Random Bool}, qui con l'Is Valid. ``So che è un po' complesso.''

\paragraph{La ``mandragata'': il tipo del Select.} Il problema è il \emph{tipo} delle
opzioni. ``Il Select praticamente c'è un piccolo problema'': il nodo è jolly (\emph{wildcard})
e assume il tipo della \emph{prima} cosa che gli si collega. Collegando prima la camera, il
Select diventa un Select di \emph{Camera Component Object Reference}, e Self (il componente
BPC\_RayInteract) non entra più; collegando prima Self, il Select diventa un Select di
riferimenti a BPC\_RayInteract, e la camera non entra. Entrambi però \emph{sono} Scene
Component: la camera è uno Scene Component, e lo è anche il nostro componente. Il trucco
(``qui bisogna fare una mandragata, ragazzi''):
\begin{enumerate}
  \item si crea una variabile qualunque di tipo \textbf{Scene Component Object Reference}
        (il docente la chiama \emph{Scene}: ``e poi non ci servirà a niente'');
  \item se ne mette un Get nel Select: ora le opzioni sono di tipo \emph{Scene Component};
  \item si \textbf{cancella} il Get e si collegano CameraComponent e Self, che sono entrambi
        accettati.
\end{enumerate}
``Ecco fatto che abbiamo praticamente fregato Unreal.''

\paragraph{La domanda di Filippo.} ``Non ho capito perché abbiamo creato la variabile Scene
prima.'' Il docente rifà la prova davanti a tutti: un Select nuovo, il cursore sul pin, la
scritta ``Scene Component Object Reference''; poi Self in un Select vuoto, e il nodo diventa
``Self Object Reference'', ``mi farà fare solamente un select fra due self, non due Scene
Component''. ``Per fregare Unreal'' serve ``una variabile che appartiene generalmente allo
Scene Component''. Uno studente prova a trasformare la variabile in Self: ``non è che te lo
trasforma: lo devi cancellare, e poi tasto destro, scrivi \emph{self}''. Filippo insiste: e
se avessimo collegato subito la camera? Prova del docente: ``Select camera component, camera
component object reference, e io, se provo ad attaccare Self, non me lo attacca''. ``Ok,
capito.''

\paragraph{Il punto finale.} Dal risultato del Select, ``quello che abbiamo sempre fatto'':
\begin{enumerate}
  \item \textbf{Get Forward Vector} (della camera, o del componente stesso);
  \item moltiplicato per una distanza: il secondo pin del Multiply si converte in float
        (``tasto destro, mi raccomando: \emph{To Float}''), si promuove a variabile,
        \textbf{DistanceTrace}, con valore \textbf{300};
  \item il prodotto si \textbf{somma} (Add) allo Start, e il risultato va nell'\textbf{End}.
\end{enumerate}

\paragraph{Perché lo Start è il componente.} Una studentessa chiede perché lo Start sia la
World Location: ``in relazione a dove attacchiamo il componente?'' ``Esatto.'' ``Il punto
iniziale deve essere sempre il punto, il Transform, del nostro componente''; ``il punto che
deve diversificare è la camera''. Il docente lo mostra nel bot: guardando in basso il
raggio va in basso, guardando di lato va di lato, ma parte sempre dal componente, ``così
magari posso posizionarlo avanti, posizionarlo indietro, posizionarlo qui, posizionarlo
là'' (figura~\ref{ue12:fig:look}).

\paragraph{Il componente nel personaggio.} Nel BP\_Bot2 (``a bot2, in questo caso, è il
vostro character''), \textbf{Add} e il nome del componente: BPC\_RayInteract compare fra i
componenti del personaggio. Per vederlo si imposta il \textbf{Draw Debug Type} della Sphere
Trace (``lo metto a \emph{For One Frame}''). In Play, però, non succede niente: ``ah,
certo, ecco perché: non facciamo LookForInteractions''. La funzione va chiamata
nell'\textbf{Event Tick} del componente (``ah, nel tick, infatti''), così il controllo
avviene a ogni frame. Con il debug persistente (\emph{Persistent}) il raggio lascia una scia
``gigantesca''; con \emph{For Duration} resta qualche istante; con \emph{For One Frame}
``rimane un po' più carino'', e si vede solo il raggio attuale. (La trascrizione, qui, è
piena di prove e di parole storpiate: ``la duration vado a activate, troppo corto'', ``Lugano'',
``Adobe''; anche una confusione momentanea fra il bot e il bot2 del docente.)

\paragraph{Approfondimento: Sphere Trace, Tick e costo.} Una \textbf{Sphere Trace} è un
\emph{line trace} ``spesso'': fa scorrere una sfera del raggio indicato dallo Start
all'End e restituisce il primo oggetto che la sfera tocca nel canale scelto. Rispetto a un
raggio sottile perdona di più la mira (con 5 unità di raggio, circa 5 cm, conta anche
sfiorare il bordo di un oggetto), ed è il motivo per cui molti sistemi di interazione la
preferiscono. Farla nel \textbf{Tick}, una volta per frame, costa poco: una singola trace di
300 unità è un'operazione leggera, e il sistema la esegue solo per il personaggio del
giocatore. Se gli actor con il componente fossero molti (per esempio dei nemici), conviene
ridurre la frequenza con un timer (come quelli della stamina, a 0,1 secondi) o disattivare
il Tick del componente quando non serve (\emph{Set Component Tick Enabled}). Una Sphere
Trace può stare senza problemi dentro una funzione, perché non ha uno stato interno da
ricordare fra una chiamata e l'altra; il Do Once e il Flip Flop, invece, sì (si veda più
avanti). Funzionamento generale dell'engine.

% ---------------------------------------------------------------------
\section{Il Trace Channel e la Blueprint Interface}
\label{ue12:sec:bpi}

\paragraph{Un canale per le interazioni.} ``Ha colpito qualcosa? Sì, no.'' Prima di
leggere il risultato serve un canale dedicato, ``che ci permette di identificare gli
oggetti che vogliamo che siano interagibili'': ``se no, praticamente, interagirà con
tutto''. ``Come abbiamo visto anche nella nostra torretta'': \textbf{Project Settings},
\emph{Collision}, \textbf{New Trace Channel}, nome \textbf{InteractionTrace}, risposta di
default \textbf{Ignore}. Il docente, che ce l'ha già, non lo rifà. Nella Sphere Trace, il
\textbf{Trace Channel} diventa InteractionTrace: ``mi raccomando, se non vi appare, fate
compila, perché ha bisogno un attimo di compilare il nuovo Trace Channel''.

\paragraph{Che cosa abbiamo colpito.} Dopo la trace, un \textbf{Branch} sul suo Return
Value (``ha colpito qualcosa?'') e un \textbf{Break Hit Result} sull'Out Hit, da cui, ``vi
ricordate, no?'', escono l'actor colpito, il Physical Material ``e tutto''. Il piano:
``useremo l'interfaccia, la Blueprint Interface, per andare a comunicare con questo actor'':
``ha questo porto in cui può accogliere questa tipologia di interazione? Sì? Benissimo,
allora fai l'interazione''. Dall'\textbf{Hit Actor} si chiede \textbf{Does Object
Implement Interface}: ``ehi, questo actor ce l'ha l'interface?'' Ma l'interfaccia non
esiste ancora: ``ci fermiamo qua per adesso, creiamo l'interface, e ci fermiamo''. (Il
docente parla qui anche di promuovere l'Hit Actor a variabile; la variabile di appoggio
vera e propria, \emph{Interactable}, si crea dopo la pausa, sezione~\ref{ue12:sec:logica}.)

\paragraph{La Blueprint Interface.} Nella cartella \emph{Blueprints}, \emph{Components},
tasto destro, \textbf{Blueprint}, ``non Blueprint Class questo, ma Blueprint questo'': il
sottomenu offre \emph{Blueprint Class}, \emph{Blueprint Function Library} e
\textbf{Blueprint Interface}. Il nome: \textbf{BPI\_InteractionRay}. Aperta:
\begin{itemize}
  \item il grafo è \textbf{in sola lettura}: ``non possiamo far niente'', ``cannot create
        new nodes in read only graph''. Un'interfaccia contiene solo \emph{dichiarazioni};
  \item si possono aggiungere \textbf{funzioni}: la prima c'è già e si rinomina (con F2),
        \textbf{Interaction Start}; con \emph{Add Function} le altre due,
        \textbf{Interaction End} e \textbf{Interaction}.
\end{itemize}
``Abbiamo messo che in questo cargo, questo trasporto, ci stanno Interaction Start,
Interaction End e Interaction: quindi queste funzioni possono essere passate tra di loro.''

\paragraph{Il collegamento.} Nel componente, il pin \textbf{Interface} di Does Object
Implement Interface diventa \mbox{BPI\_}\allowbreak InteractionRay. ``Se quel qualcosa che ho colpito ha questa
interfaccia, allora andremo a fare tutto il procedimento di interazione.'' ``Ci sono domande?
Mi sembra di capire no.'' Sono le 17:26: pausa fino alle 17:55, ``diciotto''.

\paragraph{Approfondimento: la Blueprint Interface nel libro.} Romero e Sewell presentano le
interfacce nel capitolo 16 (\emph{Introduction to VR Development}), nel paragrafo
\emph{Blueprint Communication using interfaces}: ``una Blueprint Interface è un tipo speciale
di Blueprint che contiene solo nomi di funzioni e parametri. Si usa per permettere la
comunicazione fra tipi diversi di Blueprint.'' Si crea dal pulsante \emph{Add} del Content
Browser, sottomenu \emph{Blueprints}, \emph{Blueprint Interface}. L'esempio è il template VR:
l'interfaccia \emph{VRInteraction BPI} è aggiunta alla pistola dalle \emph{Class Settings},
categoria \emph{Interfaces}, e la pistola implementa la funzione \emph{Trigger Pressed}:
``siccome la funzione non ha parametri di output, è implementata come un \emph{evento}''. Il
Pawn VR chiama Trigger Pressed sull'actor che tiene in mano, ``e se l'actor ha implementato
l'interfaccia VRInteraction BPI, Trigger Pressed viene eseguito. Non succede niente se l'actor
non l'ha implementata.'' È esattamente lo schema della lezione: il personaggio non sa
\emph{che cosa} sta guardando (un cubo, un bottone, un ascensore) e non deve fare un
\emph{cast} a nessuna classe; manda il messaggio, e ciascun oggetto risponde a modo suo.
Questo è il vantaggio delle interfacce rispetto al cast e all'ereditarietà viste con
BP\_Activable: classi che non hanno niente in comune (un actor, un personaggio, un oggetto
fisico) possono condividere la stessa ``porta''.

% ---------------------------------------------------------------------
\section{La pausa: la stamina di Tommaso}
\label{ue12:sec:pausa}

\paragraph{Il sintomo.} In pausa, come promesso, il problema rimasto aperto la volta
precedente: \textbf{Tommaso} tiene premuto Shift, la barra della stamina non scende, e
``appena mollo, si vede''. Il docente ripercorre il codice condiviso: l'evento StopRunning,
il Clear and Invalidate, il Set Timer by Event e la variabile TimeStamina (``quanto sta?
0,92''; nella trascrizione il valore è questo, ma va preso con cautela: nella lezione
precedente il valore indicato era 0,05), RegainStamina, il Clamp con CurrentStamina e
MaxStamina, l'InitializeStaminaBar nel Begin Play. ``Sembra tutto a posto.''

\paragraph{Il debug.} Il metodo è quello della volta scorsa: un \textbf{Print String} con la
CurrentStamina (``attaccalo al viola'') dopo l'evento di aggiornamento. I valori scendono
correttamente (``500, 500, 300'' dice la trascrizione, forse leggendo la velocità), poi
risalgono. Ma a schermo ci sono ``troppi Print String'': si cercano e si tolgono quelli
vecchi, uno per volta (quello della torretta, uno della piattaforma, ``forse i mulini?'',
``forse il collectable''; ``nessuno è perfetto''). Il comportamento: la barra
\emph{ricarica} visibilmente, ma \emph{non scende}.

\paragraph{La causa.} ``Ah, ecco perché: non chiami l'evento Call Update Stamina.'' Nella
catena di \textbf{StaminaDeduction}, dopo il Set di CurrentStamina, mancava la
\textbf{Call E\_UpdateStamina}: il valore calava, ma nessuno lo comunicava alla barra. La
Call c'era solo in RegainStamina, che parte al rilascio del tasto; per questo la barra si
aggiornava ``solo quando mollo''. Il rimedio: la Call anche in StaminaDeduction, con
CurrentStamina e MaxStamina collegati. ``Perfetto. Ottimo.''

\paragraph{Il rumore.} Il resto della pausa, nella trascrizione, è il blocco di testo
ripetuto e il frammento di video già segnalati nell'introduzione. Alla ripresa, uno studente
chiede aiuto per BPI\_InteractionRay, che dopo un riavvio di Unreal si apre ``male'': il
grafo in sola lettura non si vede. Il docente prova i menu \emph{Window} (Bookmarks, Find
Results, Compiler Results\dots), carica per sbaglio il layout di default dell'editor
(``mi sono caricato il bellissimo layout di Unreal''), e alla fine trova la via più
semplice: ``fai doppio clic sulla funzione'' nel pannello \emph{My Blueprint}, e il suo
grafo si apre. ``Tanto non ti serve'': nell'interfaccia non si scrive nulla.

% ---------------------------------------------------------------------
\section{La logica: la variabile Interactable}
\label{ue12:sec:logica}

``Vai, allora finiamo questo bellissimo sistema di interazioni.'' Si torna in
BPC\_RayInteract, nella funzione LookForInteractions. Il ramo di partenza è quello in cui
la trace ha colpito qualcosa e l'actor colpito implementa BPI\_InteractionRay.

\paragraph{La variabile d'appoggio.} Serve una variabile di tipo \textbf{Actor}
(\emph{Object Reference}), \textbf{Interactable}: ``sarà la nostra variabile d'appoggio'',
cioè l'oggetto interagibile che il giocatore sta guardando \emph{in questo momento}, o
niente. Il suo scopo è evitare di ripetere il lavoro a ogni frame: ``se l'actor che io ho
intercettato con il mio trace è lo stesso di prima, allora non voglio che mi rifai alcune
sezioni del codice''; ``se è lo stesso, non lo svuoto''.

\paragraph{È un oggetto nuovo?} Un secondo \textbf{Branch}, con la condizione \textbf{Hit
Actor} \textbf{diverso da} Interactable: il nodo \emph{Not Equal} (``diverso si fa con il
punto esclamativo e l'uguale''). ``Il punto esclamativo in informatica vuol dire che questo è
un condizionale che ci dice: guarda, questo è diverso da questo, e ci restituisce un
booleano. Se è diverso ci dà True, se sono uguali ci dà False.'' Sul False non si fa nulla:
è lo stesso oggetto di prima.

\paragraph{I messaggi dell'interfaccia.} ``Adesso, importante'': scrivendo
\emph{interaction} nel menu dei nodi, compaiono già le funzioni di BPI\_InteractionRay,
``senza fare niente'': ``ho già la possibilità di poter comunicare e messaggiare con
l'interfaccia''. Sono i nodi \textbf{(Message)}, con l'icona della busta, che si possono
chiamare su qualunque actor. Servono \textbf{Interaction End} e \textbf{Interaction Start}.

\paragraph{Il ramo True: chiudere il vecchio, aprire il nuovo.} Sul True del diverso:
\begin{enumerate}
  \item un \textbf{Validated Get} di Interactable (si mette il Get, tasto destro,
        \emph{Convert to Validated Get}): ``questa variabile Interactable è valida? C'è
        qualcosa dentro?'';
  \item \textbf{Is Valid}: c'era già un oggetto, quindi prima si chiude la vecchia
        interazione, \textbf{Interaction End} con target il \emph{vecchio} Interactable;
        poi \textbf{Set Interactable} = Hit Actor; poi \textbf{Interaction Start} con
        target il \emph{nuovo} Interactable;
  \item \textbf{Is Not Valid}: ``la variabile è vuota'', ``non è mai iniziata
        un'interazione precedente''; si salta l'End e si fa direttamente il Set
        Interactable = Hit Actor e l'\textbf{Interaction Start}.
\end{enumerate}

\paragraph{Il ramo False: non guardo niente.} Se la trace non colpisce nulla, \emph{oppure}
colpisce un actor senza l'interfaccia (``il muro, ad esempio''), il giocatore ha distolto lo
sguardo: ``allora ti finisco l'interazione''. Entrambi i False (quello del Branch sul Return
Value e quello di Does Object Implement Interface) portano a un altro Validated Get di
Interactable:
\begin{itemize}
  \item \textbf{Is Valid}: \textbf{Interaction End} sul vecchio Interactable, poi
        \textbf{Set Interactable} senza collegare nulla al pin, cioè lo si
        \textbf{svuota};
  \item \textbf{Is Not Valid}: non c'era nessuna interazione da chiudere; per sicurezza
        lo si svuota comunque (``semplicemente dirgli: guarda, svuotati'').
\end{itemize}
``Quindi mettere il Set così è come mettere uguale a null?'' ``È nulla, proprio. La
annulliamo: la mettiamo vuota, e quindi diventerà non valida.'' Il caso del primo ramo,
il passaggio diretto da un oggetto all'altro, serve quando due interagibili sono vicini:
``lui finisce con quello precedente e inizia con il nuovo''; di solito, invece, gli oggetti
sono distanti, la variabile si svuota e poi si riempie di nuovo
(figure~\ref{ue12:fig:look} e~\ref{ue12:fig:stati}).

\paragraph{Il problema di tipo.} Uno studente non riesce a collegare l'Hit Actor al Set
Interactable: ``mi dice che non sono compatibili, perché uno è un riferimento oggetto e
l'altro è un componente di scena''. La variabile Interactable era stata creata con il tipo
sbagliato (ed era cambiata per un clic involontario in \emph{Geometry}, ``che non ho
cliccato''). Rimedio: nei Details della variabile, il tipo \textbf{Actor}, \textbf{Object
Reference} (``non quell'Actor: tipo oggetto, riferimento oggetto''), confermare la
modifica del tipo, poi staccare e riattaccare i fili rossi dei nodi che usavano la
variabile, ``era semplicemente sparito''.

\paragraph{Il Tick.} ``Importante, ragazzi: far sì che questa funzione venga fatta su ogni
frame.'' Nell'Event Graph del componente, \textbf{Event Tick} $\rightarrow$
LookForInteractions (se non lo si è già fatto): ``a ogni frame farà il controllo: c'ho
qualcosa di interagibile davanti a me, dove sto puntando la camera? Sì o no?''.

\begin{figure}[htbp]
\centering
\begin{tikzpicture}[font=\footnotesize,
  ue12node/.style={draw, thick, rounded corners=4pt, align=center, fill=gray!8,
                  minimum height=0.85cm, text width=2.4cm},
  ue12dec/.style={draw, thick, diamond, aspect=2.2, align=center, fill=yellow!15,
                  inner sep=1pt, text width=2.0cm, font=\scriptsize},
  ue12msgn/.style={ue12node, fill=violet!12},
  ue12exec/.style={-{Stealth[length=2.5mm]}, very thick, black!80}]
  \node[ue12node, fill=red!12, text width=3.4cm] (st) at (0,0) {\textbf{Sphere Trace By
    Channel}\\{\scriptsize Start: posizione del componente; End: Start + Forward (camera o
    self) $\times$ 300; Radius 5; InteractionTrace}};
  \node[ue12dec] (hit) at (0,-2.2) {colpito qualcosa?};
  \node[ue12dec] (impl) at (0,-4.4) {Hit Actor implementa la BPI?};
  \node[ue12dec] (diff) at (0,-6.6) {Hit Actor $\neq$ Interactable?};
  \node[ue12dec] (val1) at (0,-8.9) {Interactable valido?};
  \node[ue12msgn] (end1) at (-3.6,-10.8) {\textbf{Interaction End}\\{\scriptsize sul vecchio}};
  \node[ue12node] (set1) at (0,-12.4) {\textbf{Set Interactable}\\{\scriptsize = Hit Actor}};
  \node[ue12msgn] (start) at (0,-14.0) {\textbf{Interaction Start}\\{\scriptsize sul nuovo}};
  \node[ue12node, fill=gray!3, text width=2.0cm] (none) at (3.7,-6.6) {niente\\{\scriptsize
    (stesso oggetto)}};
  \node[ue12dec] (val2) at (6.4,-3.3) {Interactable valido?};
  \node[ue12msgn] (end2) at (6.4,-5.6) {\textbf{Interaction End}\\{\scriptsize sul vecchio}};
  \node[ue12node] (clr) at (6.4,-7.6) {\textbf{Set Interactable}\\{\scriptsize vuoto (null)}};
  \draw[ue12exec] (st) -- (hit);
  \draw[ue12exec] (hit) -- node[left, font=\scriptsize]{True} (impl);
  \draw[ue12exec] (impl) -- node[left, font=\scriptsize]{True} (diff);
  \draw[ue12exec] (diff) -- node[left, font=\scriptsize]{True} (val1);
  \draw[ue12exec] (diff) -- node[above, font=\scriptsize]{False} (none);
  \draw[ue12exec] (val1.west) -| node[pos=0.25, above, font=\scriptsize]{Is Valid} (end1.north);
  \draw[ue12exec] (end1.south) |- (set1.west);
  \draw[ue12exec] (val1) -- node[right, font=\scriptsize]{Is Not Valid} (set1);
  \draw[ue12exec] (set1) -- (start);
  \draw[ue12exec] (hit.east) -| node[pos=0.2, above, font=\scriptsize]{False} (val2.north);
  \draw[ue12exec] (impl.east) -- node[above, font=\scriptsize]{False} ++(2.5,0) |-
    (val2.west);
  \draw[ue12exec] (val2) -- node[right, font=\scriptsize]{Is Valid} (end2);
  \draw[ue12exec] (end2) -- (clr);
  \draw[ue12exec] (val2.east) -- ++(0.6,0) |- node[pos=0.3, right, font=\scriptsize, text
    width=1.0cm, align=left]{Is Not Valid} (clr.east);
\end{tikzpicture}
\caption{La funzione LookForInteractions del componente, chiamata nell'Event Tick. Se il
giocatore guarda un nuovo interagibile, si chiude l'interazione con il vecchio (se c'era)
e si apre con il nuovo; se guarda un oggetto non interagibile, o niente, si chiude
l'interazione in corso e si svuota la variabile. I nodi viola sono messaggi
dell'interfaccia. Ricostruzione dalla lezione.}
\label{ue12:fig:look}
\end{figure}

\begin{figure}[htbp]
\centering
\begin{tikzpicture}[font=\footnotesize,
  ue12state/.style={draw, very thick, rounded corners=10pt, align=center,
                    minimum height=1.3cm, text width=3.4cm},
  ue12tr/.style={-{Stealth[length=2.5mm]}, very thick}]
  \node[ue12state, fill=gray!10] (vuoto) at (0,0) {\textbf{Interactable vuoto}\\{\scriptsize
    non guardo niente di interagibile}};
  \node[ue12state, fill=green!12] (a) at (7.0,0) {\textbf{Interactable = A}\\{\scriptsize
    widget ``E'' di A visibile}};
  \node[ue12state, fill=green!6] (b) at (7.0,-3.4) {\textbf{Interactable = B}\\{\scriptsize
    widget ``E'' di B visibile}};
  \draw[ue12tr] ([yshift=3mm]vuoto.east) -- node[above, font=\scriptsize]{guardo A: Start(A)}
    ([yshift=3mm]a.west);
  \draw[ue12tr] ([yshift=-3mm]a.west) -- node[below, font=\scriptsize, text width=3.2cm,
    align=center]{guardo il muro o niente: End(A), svuota} ([yshift=-3mm]vuoto.east);
  \draw[ue12tr] (a.south) -- node[right, font=\scriptsize, text width=3.4cm]{passo
    direttamente a B: End(A), Set, Start(B)} (b.north);
  \draw[ue12tr] (b.west) -| node[pos=0.3, below, font=\scriptsize]{guardo il muro: End(B),
    svuota} (vuoto.south);
  \draw[ue12tr] (a.east) .. controls +(1.6,0.9) and +(1.6,-0.9) .. node[right,
    font=\scriptsize, text width=1.9cm]{guardo ancora A: nulla} (a.east);
\end{tikzpicture}
\caption{La variabile Interactable come stato. Ogni Interaction Start è sempre seguito, prima
o poi, dal suo Interaction End: l'oggetto può contare sulla coppia per mostrare e nascondere
il suggerimento. Schema ricostruito.}
\label{ue12:fig:stati}
\end{figure}

\paragraph{Approfondimento: messaggi, Does Implement Interface e Validated Get.} Tre
precisazioni generali sui nodi usati. Primo: un nodo \emph{(Message)} di un'interfaccia si
può chiamare su qualunque riferimento a oggetto; se l'oggetto non implementa l'interfaccia,
la chiamata semplicemente non fa nulla, senza errori (la frase del libro citata sopra: ``non
succede niente''). Il controllo \textbf{Does Object Implement Interface}, quindi, non serve
a evitare un errore, ma a \emph{distinguere} gli interagibili dagli altri oggetti, così da
chiudere l'interazione quando si guarda un muro. Secondo: il confronto \emph{Not Equal} fra
due riferimenti confronta l'\emph{identità} degli oggetti (è lo stesso actor?), non la loro
classe; due cubi identici sono comunque diversi. Terzo: il \textbf{Validated Get} è un Get
con l'\emph{Is Valid} incorporato, con le due uscite \emph{Is Valid} e \emph{Is Not Valid};
``validare'' un riferimento prima di usarlo è la raccomandazione con cui il docente chiude
il corso, perché un riferimento vuoto o a un actor distrutto (per esempio un interagibile
raccolto o eliminato) produce, se usato, gli errori \emph{Accessed None} a fine Play.
Funzionamento generale dell'engine.

% ---------------------------------------------------------------------
\section{L'oggetto di prova: BP\_InteractableTest}
\label{ue12:sec:test}

\paragraph{L'actor.} ``Adesso creiamo il nostro actor interagibile.'' Nella cartella dei
Blueprint, una Blueprint Class di tipo \textbf{Actor}, ``veramente semplicissimo'':
\textbf{BP\_InteractableTest}. Due componenti: un \textbf{Cube} e un \textbf{Widget}
(``fate widget, che poi faremo''), messo sopra il cubo.

\paragraph{Implementare l'interfaccia.} ``Non si trova'', se si cerca l'interfaccia fra i
componenti: va aggiunta in \textbf{Class Settings}, nella toolbar. Nei Details compaiono le
categorie \emph{Imports} (\emph{Default Namespace}, \emph{Imported Namespaces}) e
\textbf{Interfaces}, con \emph{Inherited Interfaces} e \textbf{Implemented Interfaces}; in
quest'ultima, \textbf{Add}, BPI\_InteractionRay. ``Per vedere che l'hai aggiunto'': nel
pannello My Blueprint compare la sezione \textbf{Interfaces}, con Interaction, Interaction End
e Interaction Start.

\paragraph{Implementare le funzioni.} Trascinare Interaction Start nel grafo non funziona:
l'editor avverte che la funzione ``was not marked as Blueprint Callable'' e non può essere
piazzata così. ``Dobbiamo overrideare la loro funzione, che di base è nulla'': le funzioni
dell'interfaccia ``sono solo degli appoggi'', vuote, e il loro contenuto ``sarà da
overrideare proprio dal livello di design''. Con un doppio clic (o tasto destro,
\emph{Implement Event}) si ottengono tre eventi rossi: \textbf{Event Interaction Start},
\textbf{Event Interaction End} ed \textbf{Event Interaction}. Ciascuno con un
\textbf{Print String}:
\begin{itemize}
  \item Interaction Start: ``inizio interazione'', \textbf{Text Color} verde;
  \item Interaction End: ``fine interazione'', rosso;
  \item Interaction: ``interagendo'', giallo-arancione (``ancora non l'abbiamo usato,
        quindi non lo vedremo, è normale'').
\end{itemize}
Il colore si trova aprendo la tendina del Print String (la freccia in basso del nodo).

\paragraph{La collisione.} ``È importante però una cosa'': il canale InteractionTrace è in
Ignore di default, quindi il cubo, così com'è, non viene visto dal raggio. Nella
\textbf{Collision} del Cube, preset \emph{Custom} e \textbf{InteractionTrace} su
\textbf{Block}. ``Su Cube, tanto la mesh è quella.''

\paragraph{La prova.} Si mette il cubo in scena. In Play: ``inizio interazione, fine
interazione, inizio interazione, fine interazione''. Con la camera in terza persona colpire
il cubo è un po' difficile, ``ma ci sono delle volte in cui diventa verde''. ``Vi funziona?''
``Sì, funziona.'' ``Quindi adesso abbiamo creato il nostro trace per quanto riguarda le
interazioni.'' Il componente ``fa le veci di \emph{vado a controllarti se ci sono delle
interazioni, e gestisco tutta la cosa}''; l'interazione vera, il tasto, la fa il personaggio,
``in modo classico''.

\paragraph{Approfondimento: eventi o funzioni dell'interfaccia.} Le tre funzioni di
BPI\_InteractionRay non hanno \emph{output}: per questo, nell'actor che le implementa,
diventano \emph{eventi} nell'Event Graph (è la stessa regola citata dal libro per la Trigger
Pressed della pistola VR). Se una funzione dell'interfaccia avesse un valore di ritorno (per
esempio un \emph{GetInteractionText} che restituisce ``Apri'' o ``Accendi''), andrebbe
implementata come \emph{funzione}, con un doppio clic nella sezione Interfaces, e potrebbe
essere chiamata per leggere un dato. Un'interfaccia non ha stato né variabili: ogni classe
che la implementa decide da sé che cosa fare, e può anche lasciare vuote le funzioni che non
le servono. Funzionamento generale dell'engine.

% ---------------------------------------------------------------------
\section{L'input: IA\_Interact e il tasto E}
\label{ue12:sec:input}

Nel BP\_Bot2 serve un nuovo input. ``Come si crea un nuovo input? Con l'IMC e con l'Input
Action.''
\begin{enumerate}
  \item Nel Content Browser si va alle Input Action del progetto (cliccando su IA\_Jump,
        \emph{Show in Folder View});
  \item ``basta che fate una copia di IA\_Run'': \textbf{Duplicate}, nome
        \textbf{IA\_Interact}; ``basta che qua dentro abbia i \textbf{Triggers}
        \emph{Pressed} e \emph{Released}'' (come quelli di IA\_Run), da lasciare così;
  \item nell'\textbf{IMC} di default (\emph{IMC\_\allowbreak Default}), il più accanto alle
        \emph{Mappings} aggiunge una riga per IA\_Interact; poi, cliccando sull'icona
        della tastiera e premendo il tasto, si assegna \textbf{E}.
\end{enumerate}
Il bot2 ``eredita la stessa cosa'': usa lo stesso IMC\_Default del template. Le domande dei
presenti sono tutte sulla tendina dell'IMC (``non mi fa mettere l'azione di input'': ``è la
tendina'') e sui trigger (``il primo è premuto e l'altro qual è?'': ``premuto e rilascio va
benissimo così, lasciato così''). Nel personaggio, cercando \emph{IA\_Interact}, compare il
nuovo evento d'input.

\paragraph{Approfondimento: Pressed e Released nello stesso trigger.} Un'Input Action
digitale senza trigger scatta in \emph{Triggered} a ogni frame finché il tasto è premuto; con
il trigger \emph{Pressed} scatta una sola volta alla pressione, con \emph{Released} una
volta al rilascio. Con \emph{entrambi}, come in IA\_Run e quindi in IA\_Interact, l'uscita
\emph{Triggered} scatta due volte per ogni pressione: quando si preme e quando si rilascia.
Nella lezione questo non dà problemi, perché il grafo del personaggio passa da un \emph{Do
Once} che si resetta solo alla fine del Montage: il secondo impulso, al rilascio, arriva
mentre l'animazione è in corso e viene ignorato. Ma se il Montage fosse molto breve, o il
giocatore tenesse premuto a lungo, il rilascio potrebbe far partire una seconda
interazione: per un'interazione ``a pressione'' basta il solo trigger Pressed, o l'uscita
\emph{Started}. La trascrizione non dice quale uscita dell'evento IA\_Interact usi il docente.
Funzionamento generale dell'Enhanced Input.

% ---------------------------------------------------------------------
\section{L'interazione nel personaggio}
\label{ue12:sec:player}

``Allora adesso partiamo a fare quello che ci servirà per indirizzare il nostro player,
cioè farlo ruotare verso il nostro oggetto, e allo stesso modo far sì che l'interazione
parta.''

\paragraph{Un Do Once.} ``Vi dico che è un'azione che deve essere fatta su un \textbf{Do
Once}'': l'interazione è una sequenza (girarsi, far partire il Montage, aspettarne la fine)
e solo ``a fine del montaggio può reiniziare una prossima interazione''. Quindi IA\_Interact
$\rightarrow$ Do Once.

\paragraph{Chi stiamo guardando.} Dal componente \textbf{BPC\_RayInteract} (trascinato dal
pannello dei componenti nel grafo: ``prendilo e trascinalo'') si prende la sua variabile
\textbf{Interactable} e la si converte in \textbf{Validated Get}: ``questo Interactable è
presente? Ci sta o non ci sta?''. Il personaggio non rifà la trace: legge il risultato del
componente.
\begin{itemize}
  \item \textbf{Is Valid}: il messaggio \textbf{Interaction} di BPI\_InteractionRay
        (``tira fuori e fai \emph{interaction}, cioè Interaction Message''), con target
        Interactable: ``e questa è proprio la base: così interagiamo''.
  \item \textbf{Is Not Valid}: \textbf{Reset} del Do Once (aggiunto dopo, quando il docente
        si accorge che senza il reset, premendo E senza guardare nulla, il Do Once resta
        chiuso per sempre: ``se non è valido, resetta il Do Once comunque, in ogni caso'').
\end{itemize}

\paragraph{Girarsi verso l'oggetto.} Dopo il messaggio, \textbf{Set Actor Rotation} con
target \emph{Self} (``siamo già nel player''). Si splitta il pin \emph{New Rotation}, e si
splitta anche l'uscita di un \textbf{Get Actor Rotation}: ``l'unica asse che ci interessa è
la Z'', perché il personaggio deve girarsi solo attorno all'asse verticale (``se
l'interactable sta qua, deve girare così''). Quindi:
\begin{itemize}
  \item X (Roll) e Y (Pitch) restano quelle attuali, dal Get Actor Rotation;
  \item Z (Yaw) viene da un \textbf{Find Look at Rotation}, ``quello che abbiamo fatto
        l'altra volta per la torretta'', con \textbf{Start} = \textbf{Get Actor Location}
        del personaggio (il docente dice prima ``rotation'', poi si corregge: ``location,
        ragazzi, scusatemi'') e \textbf{Target} = Get Actor Location dell'Interactable. Il
        pin di uscita non si lascia splittare (``non me lo splitta''), quindi un
        \textbf{Break Rotator}, e il suo Z va nella Z della nuova rotazione.
\end{itemize}

\paragraph{La prova provvisoria.} Prima del Montage, una prova: la scritta ``interagendo''
compare, ma il personaggio non si gira; e dopo la prima volta non succede più nulla. ``Ah,
giustamente, ragazzi'': il Do Once non si resetta mai. Il docente collega temporaneamente
il Reset alla fine della catena (``lo faccio solamente per farvelo vedere''): ``mi giro di
qua, e si gira verso di lui''. Poi lo stacca, ``perché ci servirà più avanti''.

\paragraph{Il Montage.} Nel Content Browser si cerca \emph{interact}: fra le animazioni del
bot c'è \textbf{Interact Press Button}. ``Vi ricordate come si fa il montaggio?'' Tasto
destro sull'animazione, \textbf{Create}, \textbf{Create AnimMontage}: si crea un asset che
``è sempre la stessa animazione, ma in questo caso è un montaggio''. Nel personaggio,
\textbf{Play Montage}: nel pin \textbf{In Skeletal Mesh Component} la \textbf{Mesh} del
personaggio (``basta che prendi la Mesh da qui e la metti così''), nel \textbf{Montage to
Play} il nuovo Montage (\emph{Interact Press Button Montage}).

\paragraph{La variabile Interacting.} ``Mettiamo, giusto perché in futuro ci possa
servire, una variabile'' booleana, \textbf{Interacting}: ``dobbiamo verificare se stiamo
interagendo, non stiamo interagendo, eccetera. Sempre buono mettere una variabile che dice:
guarda, io sto interagendo''. Prima del Play Montage, Set Interacting a \textbf{vero}; sull'
\textbf{On Completed}, ``quindi è finita l'animazione'', a \textbf{falso}.

\paragraph{Fermare il giocatore.} Durante l'interazione il giocatore non deve muoversi:
\begin{itemize}
  \item prima del Montage, \textbf{Disable Input}, con il \textbf{Get Player Controller}
        nel pin \emph{Player Controller}: ``chi dobbiamo levare? Gli input al controller del
        nostro player'';
  \item lo \textbf{Stop Movement Immediately}, ``andiamo a fermare qualsiasi tipologia
        di movimento''. Il target non viene detto: il nodo appartiene al componente di
        movimento del personaggio, il \emph{Character Movement};
  \item sull'On Completed, \textbf{Enable Input}, con lo stesso Player Controller;
  \item e infine il \textbf{Reset} del Do Once, ``in modo tale che io possa reinteragire di
        nuovo''.
\end{itemize}
In Play: ``io giro di qua, si gira, fa l'animazione, non mi posso muovere''; ``e adesso mi
riposso muovere appena finito il Montage'' (figura~\ref{ue12:fig:player}). ``Ci siete? Vi
funziona? Perfetto.''

\paragraph{I problemi dei presenti.} Uno studente non trova come prendere il componente
(``l'hai aggiunto come components? Prendilo e trascinalo''); altri si perdono i due nodi
successivi, ripetuti dal docente (Get Interactable, Convert to Validated Get, il messaggio
Interaction). Un altro ha il Montage che non parte: manca la Mesh nel pin del Play Montage.
Uno studente, a cui prima non funzionava, conferma che ora va, ma ha lo
schermo ``pieno di \emph{Hello}'' (forse i Print String rimasti con il testo di default),
``quindi li vedo con la coda dell'occhio''.

\begin{figure}[htbp]
\centering
\begin{tikzpicture}[font=\footnotesize,
  ue12node/.style={draw, thick, rounded corners=4pt, align=center, fill=gray!8,
                  minimum height=0.9cm, text width=2.2cm},
  ue12ev/.style={ue12node, fill=red!12},
  ue12msgn/.style={ue12node, fill=violet!12},
  ue12pure/.style={ue12node, fill=green!10, rounded corners=8pt},
  ue12exec/.style={-{Stealth[length=2.5mm]}, very thick, black!80},
  ue12data/.style={-{Stealth[length=2mm]}, thick, blue!60!black}]
  \node[ue12ev, text width=1.9cm] (ia) at (0,-0.0) {\textbf{IA\_Interact}\\{\scriptsize tasto
    E}};
  \node[ue12node, text width=1.5cm] (do) at (2.3,0) {\textbf{Do Once}};
  \node[ue12node, text width=2.4cm] (vg) at (4.9,0) {\textbf{Validated Get}\\{\scriptsize
    BPC\_RayInteract.\\Interactable}};
  \node[ue12msgn] (msg) at (8.6,0) {\textbf{Interaction}\\{\scriptsize (Message) su
    Interactable}};
  \node[ue12node, text width=2.5cm] (rot) at (11.9,0) {\textbf{Set Actor Rotation}\\{\scriptsize
    X, Y attuali; Z dal Look at}};
  \node[ue12pure, text width=3.0cm] (look) at (11.9,-1.85) {\textbf{Find Look at
    Rotation}\\{\scriptsize mia posizione $\rightarrow$ posizione di Interactable}\\
    $\rightarrow$ \textbf{Break Rotator}: Z};
  \node[ue12node, text width=3.5cm] (pre) at (8.0,-4.4) {\textbf{Set Interacting} = vero\\
    \textbf{Disable Input}\\\textbf{Stop Movement Immediately}};
  \node[ue12node, text width=2.8cm] (pm) at (12.0,-4.4) {\textbf{Play Montage}\\{\scriptsize
    Mesh; Interact Press Button Montage}};
  \node[ue12node, text width=3.5cm] (post) at (12.0,-6.8) {\textbf{Set Interacting} =
    falso\\\textbf{Enable Input}\\\textbf{Reset} del Do Once};
  \node[ue12pure, text width=1.9cm] (gpc) at (8.0,-6.8) {Get Player Controller};
  \draw[ue12exec] (ia) -- (do); \draw[ue12exec] (do) -- (vg);
  \draw[ue12exec] (vg) -- node[below, font=\scriptsize]{Is Valid} (msg);
  \draw[ue12exec] (msg) -- (rot);
  \draw[ue12exec] (rot.east) -- ++(0.35,0) -- ++(0,-3.15) -| (pre.north);
  \draw[ue12exec] (pre) -- (pm);
  \draw[ue12exec] (pm) -- node[right, font=\scriptsize]{On Completed} (post);
  \draw[ue12exec] (vg.south) -- ++(0,-1.0) -| node[pos=0.25, above, font=\scriptsize]{Is
    Not Valid: Reset} ([xshift=-3mm]do.south);
  \draw[ue12data] (look) -- (rot);
  \draw[ue12data] (gpc) -- (pre);
  \draw[ue12data] (gpc) -- (post);
\end{tikzpicture}
\caption{L'interazione nel personaggio. Il personaggio non rifà la trace: legge la variabile
Interactable del suo componente, invia il messaggio, si gira solo sull'asse Z verso
l'oggetto, blocca l'input e suona il Montage; a Montage finito restituisce l'input e riapre
il Do Once. Ricostruzione dalla lezione; la posizione esatta del Set Interacting e del
Disable Input nella catena è dedotta dalle parole del docente.}
\label{ue12:fig:player}
\end{figure}

\paragraph{Approfondimento: Do Once, Disable Input e la rotazione.} Romero e Sewell
descrivono il \textbf{Do Once} nel capitolo 13 (\emph{Data Structures and Flow Control}):
``il nodo Do Once esegue il pin di uscita una sola volta. Se viene attivato di nuovo, il pin
di uscita non viene eseguito. Il pin di input Reset deve essere attivato per permettere al
Do Once di eseguire di nuovo il pin di uscita.'' Il loro esempio è un'arma caricata, che si
ricarica dopo 30 secondi con un timer che resetta il Do Once: lo stesso schema
della lezione, con il Montage al posto del timer. Nel capitolo 15 il libro descrive
\textbf{Enable Input} e \textbf{Disable Input}: ``definiscono se un actor deve rispondere agli
eventi di input'', e ``hanno bisogno di un riferimento al Player Controller in uso''. Due
note generali sull'engine. La prima: \emph{Disable Input} chiamato sul personaggio blocca
gli eventi d'input di quel Blueprint (movimento, camera, salto), ma non ferma un movimento
già in corso, cioè la velocità residua del Character Movement; per questo serve anche
\emph{Stop Movement Immediately}, che azzera la velocità sul posto. La seconda: il
\emph{Set Actor Rotation} è istantaneo, quindi il personaggio ``scatta'' verso l'oggetto. Per
una rotazione morbida si può usare un \emph{RInterp To} nel Tick o una Timeline, come per le
porte; e il template Third Person, che ha l'\emph{Orient Rotation to Movement} attivo, non
contrasta la rotazione perché durante l'interazione il personaggio è fermo.

% ---------------------------------------------------------------------
\section{Il widget ``premi E''}
\label{ue12:sec:widget}

``Andiamo semplicemente a rendere tutto un po' più carino'': un piccolo widget che mostri il
tasto da premere sopra l'oggetto.

\paragraph{Gli asset della UI.} Il docente aveva messo nel materiale didattico una cartella
di immagini per la UI, ``la UI asset''. Uno studente non l'ha trovata (``pensavo
fosse sul Drive''): il docente controlla il Drive condiviso e i canali del Discord, scopre che
non era stata caricata e la carica al momento, ``tanto pesa pochi kilobyte''. ``Estraetela e
schiaffatela dentro Unreal'', nella cartella \textbf{UI}. Il trascinamento nel Content
Browser a qualcuno dà ``fallimento''; il ``metodo abbastanza brusco, però va benissimo'' è
copiare la cartella \emph{estratta} con il gestore dei file direttamente nella cartella
\emph{Content/UI} del progetto. ``Adesso funziona.''

\paragraph{WB\_InteractUI.} Nella cartella UI, un nuovo Widget Blueprint, classe User
Widget: \textbf{WB\_InteractUI}. Nel Designer:
\begin{enumerate}
  \item un \textbf{Border}, con l'alfa del colore a \textbf{0}, ridimensionato a un
        quadrato di circa \textbf{96 $\times$ 96} (``94 va bene anche''), e il
        \textbf{Padding} tutto a \textbf{zero};
  \item un secondo \textbf{Border} \emph{dentro} il primo;
  \item le immagini: nel \textbf{Brush} di ciascun Border (``vedi che c'è il brush, apri il
        brush'') si sceglie un'immagine del pacchetto, ``il panel che volete voi'' (ce ne
        sono diversi, Panel 0, 1, 2, 3, 5\dots). Dopo una correzione (``questo è il
        contrario, scusate''), il Border esterno prende \textbf{Panel Border 5}, con l'alfa a
        \textbf{1} e la tinta \textbf{nera}, e quello interno \textbf{Panel 5};
  \item un \textbf{Text} con la lettera \textbf{E}, con \emph{Horizontal Alignment} e
        \emph{Vertical Alignment} al \textbf{centro}, e la dimensione del \textbf{font} a
        35, poi 40--45 (``una bella E'').
\end{enumerate}
``Piccola cosa, ragazzi, prima non l'ho fatto fare a mano'': la dimensione dei Border si può
impostare anche con precisione, dall'\textbf{Image Size} del Brush (96 $\times$ 96).

\paragraph{Il widget che contiene la barra della vita.} Una studentessa crea un nuovo
widget, ma aprendolo trova le barre della vita e della stamina, e non riesce a modificarle.
Rifacendo la creazione insieme, la causa emerge: nella finestra di scelta della classe aveva
preso un widget esistente come genitore (``anteprima widget praticamente ti dà
un'anteprima dei widget che già ci sono''), cioè il suo HUD; il nuovo widget ne
\emph{ereditava} il contenuto, in sola lettura. ``È buona prassi fare quello tutto in
alto'', \textbf{User Widget}.

\paragraph{L'immagine che non si vede.} Un'altra studentessa ha il bordo esterno che non
compare. Il docente controlla tinta, alfa, Image Size, Padding, \emph{Content Opacity}: ``è
a 1, va bene così''. ``Perdonami, purtroppo non saprei dire'': si cancellano i Border e si
rifanno, ``basta funzionare''. Rifatti, funzionano.

\paragraph{Il widget nell'oggetto.} In BP\_InteractableTest, il componente \textbf{Widget}:
\begin{itemize}
  \item \textbf{Space}: \textbf{Screen}, non World, ``in modo tale che ci sia una specie di
        2D, quindi relativo allo schermo, non al mondo'';
  \item \textbf{Draw Size}: \textbf{100 $\times$ 100}, ``un po' più piccolo'';
  \item \textbf{Widget Class}: WB\_InteractUI.
\end{itemize}
In Play, il quadratino con la E ``scala con lo schermo'', ma, se non lo si nasconde,
``rimane per sempre''.

\paragraph{Mostrare e nascondere.} Nel grafo di BP\_InteractableTest:
\begin{itemize}
  \item \textbf{Event BeginPlay} $\rightarrow$ \textbf{Set Hidden in Game} sul Widget, con
        \emph{New Hidden} \textbf{attivo}: nascosto da subito;
  \item \textbf{Event Interaction Start} $\rightarrow$ Set Hidden in Game,
        \textbf{disattivo}: ``allora la faccio apparire'';
  \item \textbf{Event Interaction End} $\rightarrow$ Set Hidden in Game, \textbf{attivo}.
\end{itemize}
Il docente prova anche un'idea per l'Event Interaction (nascondere il widget, un
\textbf{Delay} di un secondo, mostrarlo di nuovo), ma la toglie: ``questo potrebbe dare
errori, ragazzi'', perché il Delay ``va in contrasto'' con Start ed End (se lo sguardo si
sposta durante il secondo d'attesa, il Delay rimostrerebbe un widget che l'End ha appena
nascosto). ``Lascialo così per adesso.''

\paragraph{La spunta al contrario.} Una studentessa vede la E anche quando non guarda il
cubo: ``allora vuol dire che hai sbagliato l'Hidden in Game''. ``Il tick è positivo sull'End
e negativo sullo Start'', ``e anche positivo sul Begin Play''. Dopo la correzione resta un
dubbio, poi chiarito: all'inizio della prova la camera stava già inquadrando il cubo da una
posizione strana, e il widget era giustamente visibile. ``Non lo vedo più. Perfetto,
ragazzi.''

\paragraph{Approfondimento: il Widget Component.} Il \textbf{Widget Component} mostra un
Widget Blueprint attaccato a un actor. In \textbf{World} space il widget è un piano nel
mondo 3D: ha prospettiva, può essere coperto da altri oggetti, si rimpicciolisce con la
distanza, ed è adatto a schermi diegetici (un terminale, un cartello). In \textbf{Screen}
space il widget è disegnato sopra l'immagine, sempre rivolto alla camera, con una dimensione
costante in pixel, e segue sullo schermo la posizione proiettata del componente: è la scelta
tipica per i suggerimenti ``premi E'' e per le barre della vita sopra i nemici. La
\textbf{Draw Size} è la risoluzione in pixel a cui il widget viene disegnato. Per nascondere
il suggerimento, l'alternativa a \emph{Set Hidden in Game} è \emph{Set Visibility} sul
componente; il risultato è lo stesso. Funzionamento generale dell'engine.

% ---------------------------------------------------------------------
\section{Un esempio: la luce che si accende}
\label{ue12:sec:luce}

\paragraph{L'idea.} ``Sappiate che tutti gli oggetti che create'', se implementano
l'interfaccia, ``avete già un qualcosa che vi identifica se state iniziando
un'interazione'', se la state finendo, ``e la tipologia di interazione''. Le funzioni
costruite per i bottoni nelle lezioni passate ``possiamo traslarle sul nostro
interactable'': ``un diverso modo di interagire e di attivare le cose''. Con Davide si
vedranno ``anche delle piccole cose di modeling'': si può modellare un pulsante, dargli
BPI\_InteractionRay, ``e a quel punto magari potete chiamare l'ascensore''. ``Cosa succede se
interagisco? Ci metto quello che voglio: ci metto un timer, ci metto una cosa, accendo una
luce.''

\paragraph{Velocissimo.} ``Facciamo velocissimo'', in BP\_InteractableTest. Dall'Event
Interaction si toglie il Print String, e:
\begin{enumerate}
  \item una variabile del tipo dell'attivabile luce, la classe figlia di BP\_Activable delle
        lezioni sul pulsante (``la Point Light'', \textbf{BP\_ActivableLight}; la
        trascrizione dice ``BP Activables'' e ``inter light''), con i suoi eventi
        \textbf{Activation} e \textbf{Deactivation};
  \item un \textbf{Flip Flop}: ``questo fa A, poi fa B, poi fa di nuovo A: fondamentalmente
        gioca su due'';
  \item su \textbf{A}, \textbf{Activation} della luce; su \textbf{B},
        \textbf{Deactivation};
  \item la variabile va resa modificabile nell'istanza e collegata alla luce del livello
        (``però gli dobbiamo dare quale light'', ``dammi la light, dammi questa'').
\end{enumerate}
``Abbiamo acceso la luce e la rispegniamo.'' ``Questo è un esempio, ragazzi'': ``qualsiasi
actor può essere separato dagli altri con la sua specifica interazione'', e ``quello che la
identifica già ce l'abbiamo''.

\paragraph{Approfondimento: Flip Flop nel libro.} Romero e Sewell (capitolo 13): ``ogni volta
che un nodo Flip Flop viene eseguito, alterna fra due pin di uscita, A e B''; ha anche un
output booleano \emph{Is A}, vero quando sta eseguendo A. Con una nota importante: ``il Flip
Flop non funziona correttamente dentro una funzione, perché ha una variabile interna che
viene cancellata quando la funzione finisce'', e quindi dentro una funzione esegue sempre A.
Lo stesso vale per il Do Once. È il motivo per cui, nella lezione, il Flip Flop sta
nell'Event Graph dell'oggetto e il Do Once nell'Event Graph del personaggio, mentre la
funzione LookForInteractions, che non deve ricordare nulla fra un frame e l'altro (lo stato
è nella variabile Interactable, che è del componente), può essere una funzione.

% ---------------------------------------------------------------------
\section{La chiusura del corso}
\label{ue12:sec:chiusura}

``Ci sono domande, ragazzi? Domande su tutto il corso, in generale.'' Il docente chiede
anche un \textbf{feedback} sul corso, ``adesso o dopo, tranquillamente sul server''. La
risposta di uno studente, scherzosa: ``tutto chiaro, tutto bellissimo, incredibile''. ``Sono
contento.''

Per il futuro, anche ``per chi guarda in differita'': ``se avete un problema, state facendo il
Big One, eccetera, buttate il vostro problema sul server: sia alunni sia professori
provvederanno a rispondervi e darvi una mano''. ``Se ci sono ultime cose che volete dire,
ragazzi, ditele ora, o effettivamente questa volta tacete per sempre.'' E il congedo: ``vi
ringrazio per la vostra attenzione, vi ringrazio per aver partecipato al corso. Spero che i
vostri progettini possano prendere vita per tutte le cose che poi andrete a fare in Unreal.
Per benino, ricordatevi di \textbf{validare sempre} molte variabili.'' ``Da questo punto in
poi'' le lezioni continuano con Davide. ``È stato un piacere.''

\begin{table}[htbp]
\centering
\small
\begin{tabular}{@{}p{3.7cm}p{3.0cm}p{2.9cm}p{4.4cm}@{}}
\toprule
\textbf{Asset o variabile} & \textbf{Dove} & \textbf{Tipo, valore} & \textbf{A che serve} \\
\midrule
M\_Menu & mappe & Empty Level & la mappa del menu \\
WB\_Menu & UI & Widget Blueprint & il menu con i due bottoni \\
Button\_StartGame, Button\_QuitGame & WB\_Menu & Button, 300 $\times$ 80 & gioca; esci \\
CameraActor & M\_Menu & Camera Actor & la vista fissa del menu \\
GameMode Override & World Settings di M\_Menu & Game Mode base & niente personaggio nel menu \\
ThirdPersonMap & mappe & livello & la mappa aperta da Play Game \\
BPC\_RayInteract & Blueprints, Components & Scene Component & traccia gli interagibili \\
CameraComponent & BPC\_RayInteract & Camera Component (ref.) & la camera del proprietario \\
DistanceTrace & BPC\_RayInteract & Float, 300 & la lunghezza del raggio \\
Interactable & BPC\_RayInteract & Actor (ref.) & l'interagibile guardato ora, o vuoto \\
Scene & BPC\_RayInteract & Scene Component (ref.) & solo per il trucco del Select \\
InteractionTrace & Project Settings & Trace Channel, Ignore & il canale delle interazioni \\
BPI\_InteractionRay & Blueprints, Components & Blueprint Interface & Start, End, Interaction \\
BP\_InteractableTest & Blueprints & Actor (Cube, Widget) & l'oggetto di prova \\
IA\_Interact & Input & Input Action, Pressed e Released & l'interazione, tasto E \\
Interacting & BP\_Bot2 & Boolean & vero durante l'interazione \\
Interact Press Button Montage & animazioni del bot & Anim Montage & l'animazione dell'interazione \\
WB\_InteractUI & UI & Widget Blueprint, 96 $\times$ 96 & il suggerimento ``E'' \\
Widget (componente) & BP\_InteractableTest & Screen, 100 $\times$ 100 & mostra WB\_InteractUI \\
Light & BP\_InteractableTest & BP\_ActivableLight (ref.) & la luce accesa e spenta \\
\bottomrule
\end{tabular}
\caption{Gli asset e le variabili della giornata, con i nomi usati dal docente (BPC\_RayInteract
è, nel progetto del docente, BPC\_RayInteract1; il nome della variabile della luce è
incerto).}
\label{ue12:tab:asset}
\end{table}

\begin{table}[htbp]
\centering
\small
\begin{tabular}{@{}p{5.4cm}p{8.6cm}@{}}
\toprule
\textbf{Sintomo} & \textbf{Causa e rimedio} \\
\midrule
La barra della stamina scende solo al rilascio & manca la Call E\_UpdateStamina in
StaminaDeduction: c'era solo in RegainStamina \\
Il secondo bottone incollato non compare & è sovrapposto al primo (controllare la
Hierarchy); oppure aggiungerne uno nuovo dalla Palette \\
L'immagine copre i bottoni & è in fondo alla Hierarchy, cioè disegnata sopra: trascinarla
prima dei bottoni \\
Il personaggio compare dietro il menu & GameMode Override nelle World Settings della
mappa del menu: Game Mode base \\
Si vede la pallina (Default Pawn) & spostare il Player Start fuori dalla camera \\
La camera del menu non si usa & Set View Target with Blend dal Get Player Controller (non
dal Camera Manager), con la referenza alla camera \\
Errore su Add to Viewport & il Target è Self: collegare il Return Value del Create Widget \\
Il cursore non si vede & Set Show Mouse Cursor nell'Event Construct, con la spunta attiva \\
Nel menu ci si muove ancora & Set Input Mode UI Only nell'Event Construct \\
Il Select non accetta Self (o la camera) & il Select prende il tipo del primo collegamento:
mettere prima una variabile Scene Component, poi cancellarla \\
In Play non si vede il raggio & LookForInteractions non chiamata nell'Event Tick; Draw
Debug Type a None \\
InteractionTrace non compare nella trace & compilare dopo aver creato il canale \\
Il raggio ignora il cubo & la mesh non blocca InteractionTrace (default Ignore): Block \\
L'evento dell'interfaccia non si trascina & non è Blueprint Callable: implementarlo come
evento (doppio clic o Implement Event) \\
Hit Actor non si collega a Interactable & Interactable di tipo sbagliato: Actor, Object
Reference \\
Si interagisce una sola volta & il Do Once non viene resettato (On Completed e Is Not
Valid) \\
Il Montage non parte & manca la Mesh nell'In Skeletal Mesh Component \\
Il nuovo widget contiene le barre dell'HUD & creato con l'HUD come genitore: scegliere User
Widget \\
Il bordo del widget non si vede & causa non trovata; rifare i Border \\
La E si vede quando non si guarda & spunte di Set Hidden in Game invertite: attivo nel
Begin Play e nell'End, disattivo nello Start \\
Trascinare la cartella UI dà ``fallimento'' & copiare la cartella estratta in Content/UI dal
gestore dei file \\
La BPI si apre senza grafo & doppio clic sulla funzione in My Blueprint \\
\bottomrule
\end{tabular}
\caption{I problemi incontrati durante la lezione, con la causa e il rimedio.}
\label{ue12:tab:errori}
\end{table}

\begin{table}[htbp]
\centering
\footnotesize
\begin{tabular}{@{}p{5.8cm}p{8.2cm}@{}}
\toprule
\textbf{Nodo, asset o impostazione} & \textbf{Che cosa fa} \\
\midrule
New Level: Empty / Basic & livello vuoto; con luce e cielo già pronti \\
Widget Blueprint (User Widget) & l'asset dell'interfaccia: Designer e Graph \\
Button: Style (Normal, Hovered, Pressed) & i colori del bottone nei tre stati \\
Hex sRGB / Hex Linear & il codice del colore, da copiare \\
Text dentro un Button & l'etichetta del bottone, figlia del bottone \\
Image: Color and Opacity (alfa) & la trasparenza dell'immagine \\
Hierarchy (ordine) & l'ultimo elemento è disegnato sopra \\
Anchors / Position / Alignment / Size & ancora; scostamento; pivot; dimensioni \\
Open Level Blueprint & il Blueprint del livello corrente \\
Create Widget / Add to Viewport & crea il widget; lo mostra \\
World Settings: GameMode Override & la Game Mode per questa mappa \\
Default Pawn / Player Start & la sfera della Game Mode base; dove nasce il Pawn \\
Camera Actor / Directional Light & una camera nel livello; il sole \\
Create a Reference to (actor) & la referenza a un actor del livello nel Level Blueprint \\
Set View Target with Blend & cambia la camera del Player Controller \\
Event Pre Construct / Construct & alla costruzione nel Designer; alla costruzione in gioco \\
Set Show Mouse Cursor & mostra o nasconde il cursore \\
Set Input Mode UI Only / Game Only & l'input va solo all'interfaccia; solo al gioco \\
On Clicked (Button) & l'evento del clic su un bottone \\
Quit Game & esce dal gioco (in editor, chiude il Play) \\
Open Level (by Name) & apre una mappa dal nome \\
Blueprint Interface & solo dichiarazioni di funzioni, in sola lettura \\
Class Settings: Implemented Interfaces & aggiunge un'interfaccia a un Blueprint \\
Interaction Start / End / Interaction (Message) & i messaggi dell'interfaccia \\
Event Interaction Start / End / Interaction & la loro implementazione nell'oggetto \\
Actor Component / Scene Component & componente senza / con Transform \\
Get Owner & l'actor che possiede il componente \\
Get Component by Class & il primo componente di una classe (non l'array) \\
Sphere Trace By Channel & una sfera lungo un segmento: Radius, Draw Debug Type \\
Get World Location / Get Forward Vector & posizione nel mondo; direzione davanti \\
Select (wildcard) & sceglie fra due valori; il tipo è quello del primo collegato \\
Is Valid (pura) / Validated Get & il riferimento esiste?; Get con le due uscite \\
Not Equal (!=) & vero se i due valori sono diversi \\
New Trace Channel (Ignore / Block) & un canale di collisione personalizzato \\
Break Hit Result: Hit Actor & l'actor colpito \\
Does Object Implement Interface & l'oggetto implementa l'interfaccia? \\
Event Tick & a ogni frame \\
Print String: Text Color & il colore del messaggio \\
Input Action (Pressed, Released) / IMC & l'azione e il tasto (E) \\
Do Once (Reset) & esegue una volta, finché non si resetta \\
Set Actor Rotation / Get Actor Rotation & imposta; legge la rotazione dell'actor \\
Find Look at Rotation / Break Rotator & rotazione per guardare un punto; i suoi tre assi \\
Create AnimMontage / Play Montage & un Montage da un'animazione; lo suona (On Completed) \\
Disable Input / Enable Input & toglie; restituisce l'input all'actor \\
Stop Movement Immediately & azzera il movimento del personaggio \\
Border: Brush, Padding, Image Size & il riquadro con immagine; margine interno; dimensione \\
Widget Component: Space, Draw Size, Widget Class & un widget attaccato a un actor \\
Set Hidden in Game (New Hidden) & nasconde o mostra un componente \\
Delay & aspetta prima di proseguire \\
Flip Flop & alterna fra le uscite A e B \\
\bottomrule
\end{tabular}
\caption{I nodi, gli asset e le impostazioni della dodicesima lezione.}
\label{ue12:tab:nodi}
\end{table}

% ---------------------------------------------------------------------
\section*{Punti chiave del capitolo}
\addcontentsline{toc}{section}{Punti chiave del capitolo}

\begin{itemize}
  \item Il designer usa Unreal per \textbf{prototipare} le meccaniche e verificare che
        funzionino; il progetto d'esame degli anni successivi è un \textbf{pezzo da
        portfolio}. \textbf{Less is more}: il giocatore impiega molto più tempo di chi
        ha progettato il gioco per capirne ogni elemento.
  \item Il menu vive in una \textbf{mappa a parte} (Empty Level), con un \textbf{Widget
        Blueprint} (User Widget) creato e mostrato nel \textbf{Level Blueprint} (Create
        Widget, Add to Viewport con il Return Value).
  \item Nel Designer: Canvas Panel, bottoni con l'\textbf{ancora} al centro e
        l'\textbf{Alignment} per impilarli, un \textbf{Text} dentro ogni bottone, i colori
        di \textbf{Normal}, \textbf{Hovered} e \textbf{Pressed} nello Style (copiabili con
        il codice Hex), e l'ordine della \textbf{Hierarchy} per mettere un'immagine
        dietro (l'ultimo elemento sta sopra).
  \item Niente personaggio nel menu: \textbf{GameMode Override} con la Game Mode base
        (Default Pawn) e Player Start fuori campo. La vista fissa: un Camera Actor e
        \textbf{Set View Target with Blend}, dal \textbf{Player Controller}.
  \item Nell'\textbf{Event Construct} del widget: \textbf{Set Show Mouse Cursor} (con la
        spunta) e \textbf{Set Input Mode UI Only}. Sui bottoni, \textbf{On Clicked}:
        \textbf{Quit Game}; \textbf{Set Input Mode Game Only} e \textbf{Open Level (by
        Name)}.
  \item Una \textbf{Blueprint Interface} è un intermediario (``un volo di linea''): solo
        nomi di funzioni, in sola lettura. Chi la implementa (Class Settings,
        Implemented Interfaces) decide che cosa fare; chi non la implementa ignora i
        messaggi.
  \item Un \textbf{Blueprint Component} isola una funzionalità da aggiungere a qualunque
        actor. L'\textbf{Actor Component} non ha posizione; lo \textbf{Scene Component}
        ha un Transform. Il componente trova il suo actor con \textbf{Get Owner} e la
        camera con \textbf{Get Component by Class}.
  \item \textbf{BPC\_RayInteract}: a ogni frame (\textbf{Event Tick}) una \textbf{Sphere
        Trace} (Radius 5) dalla posizione del componente, lungo il Forward Vector della
        camera (o del componente, se la camera non c'è) per 300 unità, sul canale
        \textbf{InteractionTrace} (Ignore di default, Block sugli interagibili).
  \item Il \textbf{Select} prende il tipo del primo collegamento: per scegliere fra la
        camera e Self, prima una variabile \emph{Scene Component}, poi la si cancella (la
        ``mandragata'').
  \item La variabile \textbf{Interactable} ricorda l'oggetto guardato: se cambia,
        \textbf{Interaction End} sul vecchio, Set, \textbf{Interaction Start} sul nuovo;
        se non si guarda niente di interagibile, End e svuota. Sempre con il
        \textbf{Validated Get}.
  \item Nel personaggio: \textbf{IA\_Interact} (tasto E), \textbf{Do Once}, il messaggio
        \textbf{Interaction}, \textbf{Set Actor Rotation} con il solo \textbf{Yaw} da
        \textbf{Find Look at Rotation}, \textbf{Disable Input} e \textbf{Stop Movement
        Immediately}, \textbf{Play Montage}, poi Enable Input e \textbf{Reset} del Do Once
        (anche quando non c'è nulla da interagire).
  \item Il suggerimento ``E'': un \textbf{Widget Component} in \textbf{Screen} space,
        nascosto nel Begin Play e mostrato e nascosto da Interaction Start ed End con
        \textbf{Set Hidden in Game}.
  \item Ogni oggetto può dare la sua risposta all'interazione: per esempio un
        \textbf{Flip Flop} che chiama Activation e Deactivation di una
        \textbf{BP\_ActivableLight}. Flip Flop e Do Once non vanno dentro le funzioni.
  \item La raccomandazione finale del corso: \textbf{validare sempre} le variabili.
\end{itemize}
